% Euclidean minimality of Markov representations on thirteen Farey rows
% Authors: Dr. Denys Dutykh (Mathematics Department, Khalifa University of Science and Technology,
% Abu Dhabi, UAE) and Prof. Laurent Vuillon (Univ. Savoie Mont Blanc, CNRS, LAMA, Chambery, France)
\documentclass[11pt,reqno]{amsart}

\usepackage[T1]{fontenc}
\usepackage{lmodern}
\usepackage{microtype}
\usepackage[a4paper,margin=32mm]{geometry}
\usepackage{mathtools}
\usepackage{amssymb}
\usepackage{mathrsfs}
\usepackage{booktabs}
\usepackage{array}
\usepackage{float}
\usepackage{enumitem}
\usepackage{xcolor}
\usepackage{tikz}
\usetikzlibrary{arrows.meta,positioning,shapes.geometric,backgrounds}
% Author: Dr. Denys Dutykh (Khalifa University of Science and Technology, Abu Dhabi, UAE)
% House palette shared by the figures of the article and by their standalone builds.
\definecolor{navy}{HTML}{17365D}
\definecolor{wine}{HTML}{8E2A3F}
\definecolor{teal}{HTML}{007F86}
\definecolor{ochre}{HTML}{B66A19}
\definecolor{plum}{HTML}{6E4B9E}
\definecolor{lilac}{HTML}{9C86C8}
\definecolor{ink}{HTML}{1C1C1C}
\definecolor{slate}{HTML}{5C6670}
\definecolor{mist}{HTML}{C7CDD4}
\definecolor{sand}{HTML}{F1EBDD}
\definecolor{paper}{HTML}{FBF9F4}

\usepackage[colorlinks=true,linkcolor=blue!55!black,citecolor=green!40!black,urlcolor=blue!65!black]{hyperref}
\usepackage{aliascnt}

\newtheorem{theorem}{Theorem}[section]
\newaliascnt{proposition}{theorem}
\newtheorem{proposition}[proposition]{Proposition}
\aliascntresetthe{proposition}
\newaliascnt{lemma}{theorem}
\newtheorem{lemma}[lemma]{Lemma}
\aliascntresetthe{lemma}
\newaliascnt{corollary}{theorem}
\newtheorem{corollary}[corollary]{Corollary}
\aliascntresetthe{corollary}
\theoremstyle{definition}
\newaliascnt{definition}{theorem}
\newtheorem{definition}[definition]{Definition}
\aliascntresetthe{definition}
\newaliascnt{remark}{theorem}
\newtheorem{remark}[remark]{Remark}
\aliascntresetthe{remark}
\theoremstyle{plain}
\newaliascnt{conjecture}{theorem}
\newtheorem{conjecture}[conjecture]{Conjecture}
\aliascntresetthe{conjecture}

\usepackage[nameinlink,noabbrev]{cleveref}
\crefname{theorem}{Theorem}{Theorems}
\Crefname{theorem}{Theorem}{Theorems}
\crefname{proposition}{Proposition}{Propositions}
\Crefname{proposition}{Proposition}{Propositions}
\crefname{lemma}{Lemma}{Lemmas}
\Crefname{lemma}{Lemma}{Lemmas}
\crefname{corollary}{Corollary}{Corollaries}
\Crefname{corollary}{Corollary}{Corollaries}
\crefname{definition}{Definition}{Definitions}
\Crefname{definition}{Definition}{Definitions}
\crefname{remark}{Remark}{Remarks}
\Crefname{remark}{Remark}{Remarks}
\crefname{conjecture}{Conjecture}{Conjectures}
\Crefname{conjecture}{Conjecture}{Conjectures}
\crefname{section}{Section}{Sections}
\Crefname{section}{Section}{Sections}
\crefname{subsection}{Section}{Sections}
\Crefname{subsection}{Section}{Sections}
\crefname{figure}{Figure}{Figures}
\Crefname{figure}{Figure}{Figures}
\crefname{table}{Table}{Tables}
\Crefname{table}{Table}{Tables}
\crefname{equation}{}{}
\Crefname{equation}{Equation}{Equations}

\newcommand{\Z}{\mathbb{Z}}
\newcommand{\Q}{\mathbb{Q}}
\newcommand{\R}{\mathbb{R}}
\newcommand{\abs}[1]{\lvert#1\rvert}
\newcommand{\norm}[1]{\left\lVert#1\right\rVert}
\newcommand{\cont}[1]{\operatorname{cont}\!\left(#1\right)}
\newcommand{\height}[1]{h_{\mathrm{abs}}\!\left(#1\right)}
\newcommand{\dist}[1]{\left\lVert#1\right\rVert}


\newaliascnt{example}{theorem}
\theoremstyle{definition}
\newtheorem{example}[example]{Example}
\aliascntresetthe{example}
\crefname{example}{Example}{Examples}
\Crefname{example}{Example}{Examples}
\DeclareMathOperator{\tr}{tr}
\DeclareMathOperator{\SL}{SL}
\DeclareMathOperator{\GL}{GL}
\DeclareMathOperator{\arcosh}{arcosh}
\newcommand{\eps}{\varepsilon}
\newcommand{\mc}[1]{\mathcal{#1}}
\frenchspacing


\title[Euclidean minimality on thirteen Farey rows]{Euclidean minimality of Markov representations\\on thirteen Farey rows}

\author{Denys Dutykh}
\address{Mathematics Department, Khalifa University of Science and Technology, PO Box 127788, Abu Dhabi, United Arab Emirates}
\email{denys.dutykh@ku.ac.ae}

\author{Laurent Vuillon}
\address{Univ. Savoie Mont Blanc, CNRS, LAMA, 73000 Chamb\'ery, France}
\email{laurent.vuillon@univ-smb.fr}

\hypersetup{
 pdftitle={Euclidean minimality of Markov representations on thirteen Farey rows},
 pdfauthor={Denys Dutykh and Laurent Vuillon},
 pdfsubject={Primitive sums of two squares, Euclidean cost and unbounded Farey rows},
 pdfkeywords={Markov numbers, Christoffel words, sums of two squares, Euclidean algorithm, Gaussian integers, palindromic closure}
}
\subjclass[2020]{Primary 11D25, 68R15; Secondary 11A55, 11E25, 11J06, 11Y50}
\keywords{Markov numbers, Christoffel words, sums of two squares, Euclidean algorithm, Gaussian integers, palindromic closure}
\begin{document}
\begin{abstract}
The Markov number attached to a reduced Farey fraction is canonically a sum
of two coprime squares, and the subtractive Euclidean algorithm reduces this
pair in exactly as many steps as the numerator plus the denominator minus
two. We prove that no other primitive representation of the same integer as
a sum of two squares has a shorter Euclidean algorithm when the numerator is
at most thirteen, with no bound on the denominator. A parity-order theorem
for the Fibonacci-normalized norm reduces infinite families of Euclidean
paths to exact comparisons, complete Gaussian factorizations settle finitely
many labels, and the remaining finite assertions are verified by computer in
exact arithmetic. Minimality on every row would imply the Frobenius--Markov
uniqueness conjecture. We conjecture that it holds, and state a second
conjecture, on iterated palindromic closure and the Thue--Morse
morphism, that also implies uniqueness and is verified for every Markov
number up to $10^{12}$.
\end{abstract}
\maketitle
% Authors: Dr. Denys Dutykh (Mathematics Department, Khalifa University of Science and Technology,
% Abu Dhabi, UAE) and Prof. Laurent Vuillon (Univ. Savoie Mont Blanc, CNRS, LAMA, Chambery, France)
\section{A representation problem attached to Markov numbers}

The Markov equation
$$
 x^2+y^2+z^2=3xyz
$$
has a tree of positive integer solutions generated from $(1,1,1)$ by permutations and
the replacement $(x,y,z)\mapsto(x,y,3xy-z)$. The Frobenius--Markov uniqueness
conjecture asks whether the largest entry determines a solution up to
permutation. A useful approach is to attach an arithmetic optimization
problem to each occurrence of that entry in the tree. If two occurrences
are both optimal, a known ordering theorem can distinguish them.

We use the Farey indexing $m_{P/Q}$ of Markov numbers, with
$0<P<Q$ and $\gcd(P,Q)=1$, as in \cite{Aigner2013} and
\cite{LagisquetEtAl2021}. An occurrence includes its reduced fraction;
its label is the integer $m_{P/Q}$. Keeping these notions separate is
essential in a uniqueness problem. The first labels in our convention
are $m_{1/2}=5$, $m_{1/3}=13$, and $m_{2/3}=29$.

Every such occurrence supplies a primitive positive representation
$m_{P/Q}=r^2+s^2$. The subtractive Euclidean algorithm gives this pair
an integer cost $\ell(r,s)$: count single subtractions until $(1,1)$ is
reached. The cost can be read in two equivalent ways, and the paper uses
both. If $s/r=[q_0;\ldots,q_j]$ is the canonical finite continued
fraction, with last quotient at least two, the cost is the sum of the
partial quotients minus one, $q_0+\cdots+q_j-1$. In matrix form,
$$
 \binom sr=P_{q_0}\cdots P_{q_{j-1}}P_{q_j-1}\binom11,\qquad
 P_q=\begin{pmatrix}q&1\\1&0\end{pmatrix},
$$
and the cost is the sum of the indices of the factors;
\cref{lem:cost-descriptions} proves that the two readings agree. The
occurrence pair has cost $P+Q-2$; this follows
from the continued fraction of \c{C}anak\c{c}{\i} and Schiffler
\cite{CanakciSchiffler2020}, as we recall in \cref{sec:occurrences}.
Our question is whether any other primitive representation of the same
integer can have smaller cost. The competitors include all partial
quotients and need not arise from another Markov occurrence.

\begin{theorem}[Thirteen Farey rows]\label{thm:thirteen-rows}
Let $0<P<Q$ be coprime integers with $P\leqslant13$. For every
primitive positive representation
$$
 m_{P/Q}=r^2+s^2,\qquad 0<r<s,
$$
the subtractive Euclidean cost satisfies
$$
 \ell(r,s)\geqslant P+Q-2,
$$
with equality for the occurrence-induced representation.
\end{theorem}

There is no upper bound on $Q$ or on the label. The minimum is weak:
the theorem permits other representations of the same cost. For the
$248$ occurrences of the finite windows of \cref{sec:gaussian}, the
occurrence pair is moreover the only representation of minimum cost.
Since equal labels with equal costs come from the same fraction
(\cref{lem:fixed-cost-injectivity}), weak minimality separates occurrences.

\begin{corollary}\label{cor:thirteen-injective}
If $P/Q$ and $P'/Q'$ are reduced fractions of $(0,1)$ whose numerators are
at most thirteen, then $m_{P/Q}=m_{P'/Q'}$ implies $P/Q=P'/Q'$.
\end{corollary}

\Cref{cor:thirteen-injective} is not new in strength. Lee, Li, Rabideau
and Schiffler show that two distinct fractions with equal labels have a
slope $(P-P')/(Q-Q')$ strictly between $-5/4$ and $-8/7$, and their
numerators then differ by at least six
\cite[Corollary~1.5]{LeeEtAl2023}; this gives the injectivity on the first
six rows, and on the first thirteen rows it leaves sixteen one-parameter
families of candidate pairs, which their fixed-row limits
\cite[Lemma~6.4 and (6.4)]{LeeEtAl2023} (for the first row, the exact value
$m_{1/Q}=F_{2Q+1}$) and a few explicit comparisons of labels exclude. What \cref{cor:thirteen-injective} adds is its
route: it comes from a minimality principle, which makes sense on every row
and would give the uniqueness conjecture if it held on every row. The new
content of the paper is \cref{thm:thirteen-rows} itself, a statement about
every primitive representation of the label, Markov or not, and the method
that proves it.

\subsection{The mechanism of the proof}

Two descriptions of a primitive pair organize the argument. Gaussian
factorization enumerates every representation of one fixed integer.
Inverse Euclidean steps organize every pair of one fixed cost. The
first description treats a finite set of labels completely; the second
controls the unbounded part of each row. \Cref{fig:proof-map} shows how
the two descriptions divide the thirteen rows between them.

At cost $c$, the largest norm is the Fibonacci value $F_{2c+3}$.
Every other path departs from the Fibonacci path at some steps, called
defects. Record the lengths of the Fibonacci runs between successive
defects. The main general tool proved here, \cref{thm:tail-parity}, is
that, on fixing their parities, the full normalized norm $N/F_{2c+3}$ is
strictly monotone in each run length, with an explicit direction. This
holds for every positive number of defects, not only for the strata used
in the thirteen-row theorem. It reduces each infinite box of run lengths to
two exact corner values in $\Q(\sqrt5)$.

An exact expansion of the Cohn trace, in the eigenbasis of the first Cohn
matrix, places the source labels of each row in narrow normalized
intervals around an explicit constant, and we show that every shorter path
misses the interval belonging to its row. The covers used at that step admit unbounded run
lengths: their role is to prove uniform inequalities, not to extrapolate
from a finite search. What remains after them is a finite list of labels,
for which the complete two-square fiber is determined by Gaussian
factorization. The finitely many assertions that these two steps leave to
a computer are collected in one proposition, \cref{prop:finite-verifications};
everything else is proved in the text.

\begin{figure}[tb]
\centering
% Authors: Dr. Denys Dutykh (Mathematics Department, Khalifa University of Science and
% Technology, Abu Dhabi, UAE) and Prof. Laurent Vuillon (Univ. Savoie Mont Blanc, CNRS, LAMA)
% Generated by supplement/code/make_figures.py from exact data; do not edit by hand.
\begin{tikzpicture}[font=\footnotesize, >={Stealth[length=4pt]}]
\draw[teal, line width=1.1pt] (1.049,5.65) -- (1.794,5.65);
\node[ink, right] at (1.831,5.65) {norm envelopes};
\fill[wine] (4.625,5.65) circle (1.1pt);
\node[ink, right] at (4.737,5.65) {complete Gaussian fibers};
\draw[navy, line width=1.7pt] (8.648,5.65) -- (9.393,5.65);
\node[ink, right] at (9.43,5.65) {parity-order covers};
\node[ink, left] at (0.8,0) {$1$};
\draw[teal, line width=1.1pt, ->] (1.049,0) -- (13.267,0);
\node[ink, left] at (0.8,0.4) {$2$};
\draw[teal, line width=1.1pt, ->] (1.123,0.4) -- (13.267,0.4);
\node[ink, left] at (0.8,0.8) {$3$};
\draw[teal, line width=1.1pt, ->] (1.198,0.8) -- (13.267,0.8);
\node[ink, left] at (0.8,1.2) {$4$};
\draw[teal, line width=1.1pt, ->] (1.272,1.2) -- (13.267,1.2);
\node[ink, left] at (0.8,1.6) {$5$};
\draw[teal, line width=1.1pt, ->] (1.347,1.6) -- (13.267,1.6);
\node[ink, left] at (0.8,2) {$6$};
\draw[teal, line width=1.1pt, ->] (1.421,2) -- (13.267,2);
\node[ink, left] at (0.8,2.4) {$7$};
\draw[teal, line width=1.1pt, ->] (1.496,2.4) -- (13.267,2.4);
\node[ink, left] at (0.8,2.8) {$8$};
\draw[mist, line width=0.5pt] (1.571,2.8) -- (3.06,2.8);
\fill[wine] (1.571,2.8) circle (1.1pt);
\fill[wine] (1.72,2.8) circle (1.1pt);
\fill[wine] (1.869,2.8) circle (1.1pt);
\fill[wine] (2.018,2.8) circle (1.1pt);
\fill[wine] (2.167,2.8) circle (1.1pt);
\fill[wine] (2.316,2.8) circle (1.1pt);
\fill[wine] (2.465,2.8) circle (1.1pt);
\fill[wine] (2.614,2.8) circle (1.1pt);
\fill[wine] (2.762,2.8) circle (1.1pt);
\fill[wine] (2.911,2.8) circle (1.1pt);
\draw[navy, line width=1.7pt, ->] (3.06,2.8) -- (13.267,2.8);
\draw[ink, line width=0.6pt] (3.06,2.67) -- (3.06,2.93);
\node[ink, above right, inner sep=1pt, font=\scriptsize] at (3.06,2.82) {$29$};
\node[ink, left] at (0.8,3.2) {$9$};
\draw[mist, line width=0.5pt] (1.645,3.2) -- (3.507,3.2);
\fill[wine] (1.645,3.2) circle (1.1pt);
\fill[wine] (1.72,3.2) circle (1.1pt);
\fill[wine] (1.869,3.2) circle (1.1pt);
\fill[wine] (1.943,3.2) circle (1.1pt);
\fill[wine] (2.092,3.2) circle (1.1pt);
\fill[wine] (2.167,3.2) circle (1.1pt);
\fill[wine] (2.316,3.2) circle (1.1pt);
\fill[wine] (2.39,3.2) circle (1.1pt);
\fill[wine] (2.539,3.2) circle (1.1pt);
\fill[wine] (2.614,3.2) circle (1.1pt);
\fill[wine] (2.762,3.2) circle (1.1pt);
\fill[wine] (2.837,3.2) circle (1.1pt);
\fill[wine] (2.986,3.2) circle (1.1pt);
\fill[wine] (3.06,3.2) circle (1.1pt);
\fill[wine] (3.209,3.2) circle (1.1pt);
\fill[wine] (3.284,3.2) circle (1.1pt);
\fill[wine] (3.433,3.2) circle (1.1pt);
\draw[navy, line width=1.7pt, ->] (3.507,3.2) -- (13.267,3.2);
\draw[ink, line width=0.6pt] (3.507,3.07) -- (3.507,3.33);
\node[ink, above right, inner sep=1pt, font=\scriptsize] at (3.507,3.22) {$35$};
\node[ink, left] at (0.8,3.6) {$10$};
\draw[mist, line width=0.5pt] (1.72,3.6) -- (3.805,3.6);
\fill[wine] (1.72,3.6) circle (1.1pt);
\fill[wine] (1.869,3.6) circle (1.1pt);
\fill[wine] (2.167,3.6) circle (1.1pt);
\fill[wine] (2.316,3.6) circle (1.1pt);
\fill[wine] (2.465,3.6) circle (1.1pt);
\fill[wine] (2.614,3.6) circle (1.1pt);
\fill[wine] (2.911,3.6) circle (1.1pt);
\fill[wine] (3.06,3.6) circle (1.1pt);
\fill[wine] (3.209,3.6) circle (1.1pt);
\fill[wine] (3.358,3.6) circle (1.1pt);
\fill[wine] (3.656,3.6) circle (1.1pt);
\draw[navy, line width=1.7pt, ->] (3.805,3.6) -- (13.267,3.6);
\draw[ink, line width=0.6pt] (3.805,3.47) -- (3.805,3.73);
\node[ink, above right, inner sep=1pt, font=\scriptsize] at (3.805,3.62) {$39$};
\node[ink, left] at (0.8,4) {$11$};
\draw[mist, line width=0.5pt] (1.794,4) -- (6.562,4);
\fill[wine] (1.794,4) circle (1.1pt);
\fill[wine] (1.869,4) circle (1.1pt);
\fill[wine] (1.943,4) circle (1.1pt);
\fill[wine] (2.018,4) circle (1.1pt);
\fill[wine] (2.092,4) circle (1.1pt);
\fill[wine] (2.167,4) circle (1.1pt);
\fill[wine] (2.241,4) circle (1.1pt);
\fill[wine] (2.316,4) circle (1.1pt);
\fill[wine] (2.39,4) circle (1.1pt);
\fill[wine] (2.465,4) circle (1.1pt);
\fill[wine] (2.614,4) circle (1.1pt);
\fill[wine] (2.688,4) circle (1.1pt);
\fill[wine] (2.762,4) circle (1.1pt);
\fill[wine] (2.837,4) circle (1.1pt);
\fill[wine] (2.911,4) circle (1.1pt);
\fill[wine] (2.986,4) circle (1.1pt);
\fill[wine] (3.06,4) circle (1.1pt);
\fill[wine] (3.135,4) circle (1.1pt);
\fill[wine] (3.209,4) circle (1.1pt);
\fill[wine] (3.284,4) circle (1.1pt);
\fill[wine] (3.433,4) circle (1.1pt);
\fill[wine] (3.507,4) circle (1.1pt);
\fill[wine] (3.582,4) circle (1.1pt);
\fill[wine] (3.656,4) circle (1.1pt);
\fill[wine] (3.731,4) circle (1.1pt);
\fill[wine] (3.805,4) circle (1.1pt);
\fill[wine] (3.88,4) circle (1.1pt);
\fill[wine] (3.954,4) circle (1.1pt);
\fill[wine] (4.029,4) circle (1.1pt);
\fill[wine] (4.104,4) circle (1.1pt);
\fill[wine] (4.253,4) circle (1.1pt);
\fill[wine] (4.327,4) circle (1.1pt);
\fill[wine] (4.402,4) circle (1.1pt);
\fill[wine] (4.476,4) circle (1.1pt);
\fill[wine] (4.55,4) circle (1.1pt);
\fill[wine] (4.625,4) circle (1.1pt);
\fill[wine] (4.699,4) circle (1.1pt);
\fill[wine] (4.774,4) circle (1.1pt);
\fill[wine] (4.848,4) circle (1.1pt);
\fill[wine] (4.923,4) circle (1.1pt);
\fill[wine] (5.072,4) circle (1.1pt);
\fill[wine] (5.147,4) circle (1.1pt);
\fill[wine] (5.221,4) circle (1.1pt);
\fill[wine] (5.296,4) circle (1.1pt);
\fill[wine] (5.37,4) circle (1.1pt);
\fill[wine] (5.445,4) circle (1.1pt);
\fill[wine] (5.519,4) circle (1.1pt);
\fill[wine] (5.594,4) circle (1.1pt);
\fill[wine] (5.668,4) circle (1.1pt);
\fill[wine] (5.742,4) circle (1.1pt);
\fill[wine] (5.891,4) circle (1.1pt);
\fill[wine] (5.966,4) circle (1.1pt);
\fill[wine] (6.04,4) circle (1.1pt);
\fill[wine] (6.115,4) circle (1.1pt);
\fill[wine] (6.189,4) circle (1.1pt);
\fill[wine] (6.264,4) circle (1.1pt);
\fill[wine] (6.338,4) circle (1.1pt);
\fill[wine] (6.413,4) circle (1.1pt);
\fill[wine] (6.487,4) circle (1.1pt);
\draw[navy, line width=1.7pt, ->] (6.562,4) -- (13.267,4);
\draw[ink, line width=0.6pt] (6.562,3.87) -- (6.562,4.13);
\node[ink, above right, inner sep=1pt, font=\scriptsize] at (6.562,4.02) {$76$};
\node[ink, left] at (0.8,4.4) {$12$};
\draw[mist, line width=0.5pt] (1.869,4.4) -- (8.425,4.4);
\fill[wine] (1.869,4.4) circle (1.1pt);
\fill[wine] (2.167,4.4) circle (1.1pt);
\fill[wine] (2.316,4.4) circle (1.1pt);
\fill[wine] (2.614,4.4) circle (1.1pt);
\fill[wine] (2.762,4.4) circle (1.1pt);
\fill[wine] (3.06,4.4) circle (1.1pt);
\fill[wine] (3.209,4.4) circle (1.1pt);
\fill[wine] (3.507,4.4) circle (1.1pt);
\fill[wine] (3.656,4.4) circle (1.1pt);
\fill[wine] (3.954,4.4) circle (1.1pt);
\fill[wine] (4.104,4.4) circle (1.1pt);
\fill[wine] (4.402,4.4) circle (1.1pt);
\fill[wine] (4.55,4.4) circle (1.1pt);
\fill[wine] (4.848,4.4) circle (1.1pt);
\fill[wine] (4.998,4.4) circle (1.1pt);
\fill[wine] (5.296,4.4) circle (1.1pt);
\fill[wine] (5.445,4.4) circle (1.1pt);
\fill[wine] (5.742,4.4) circle (1.1pt);
\fill[wine] (5.891,4.4) circle (1.1pt);
\fill[wine] (6.189,4.4) circle (1.1pt);
\fill[wine] (6.338,4.4) circle (1.1pt);
\fill[wine] (6.636,4.4) circle (1.1pt);
\fill[wine] (6.785,4.4) circle (1.1pt);
\fill[wine] (7.083,4.4) circle (1.1pt);
\fill[wine] (7.232,4.4) circle (1.1pt);
\fill[wine] (7.53,4.4) circle (1.1pt);
\fill[wine] (7.679,4.4) circle (1.1pt);
\fill[wine] (7.978,4.4) circle (1.1pt);
\fill[wine] (8.127,4.4) circle (1.1pt);
\draw[navy, line width=1.7pt, ->] (8.425,4.4) -- (13.267,4.4);
\draw[ink, line width=0.6pt] (8.425,4.27) -- (8.425,4.53);
\node[ink, above right, inner sep=1pt, font=\scriptsize] at (8.425,4.42) {$101$};
\node[ink, left] at (0.8,4.8) {$13$};
\draw[mist, line width=0.5pt] (1.943,4.8) -- (11.777,4.8);
\fill[wine] (1.943,4.8) circle (1.1pt);
\fill[wine] (2.018,4.8) circle (1.1pt);
\fill[wine] (2.092,4.8) circle (1.1pt);
\fill[wine] (2.167,4.8) circle (1.1pt);
\fill[wine] (2.241,4.8) circle (1.1pt);
\fill[wine] (2.316,4.8) circle (1.1pt);
\fill[wine] (2.39,4.8) circle (1.1pt);
\fill[wine] (2.465,4.8) circle (1.1pt);
\fill[wine] (2.539,4.8) circle (1.1pt);
\fill[wine] (2.614,4.8) circle (1.1pt);
\fill[wine] (2.688,4.8) circle (1.1pt);
\fill[wine] (2.762,4.8) circle (1.1pt);
\fill[wine] (2.911,4.8) circle (1.1pt);
\fill[wine] (2.986,4.8) circle (1.1pt);
\fill[wine] (3.06,4.8) circle (1.1pt);
\fill[wine] (3.135,4.8) circle (1.1pt);
\fill[wine] (3.209,4.8) circle (1.1pt);
\fill[wine] (3.284,4.8) circle (1.1pt);
\fill[wine] (3.358,4.8) circle (1.1pt);
\fill[wine] (3.433,4.8) circle (1.1pt);
\fill[wine] (3.507,4.8) circle (1.1pt);
\fill[wine] (3.582,4.8) circle (1.1pt);
\fill[wine] (3.656,4.8) circle (1.1pt);
\fill[wine] (3.731,4.8) circle (1.1pt);
\fill[wine] (3.88,4.8) circle (1.1pt);
\fill[wine] (3.954,4.8) circle (1.1pt);
\fill[wine] (4.029,4.8) circle (1.1pt);
\fill[wine] (4.104,4.8) circle (1.1pt);
\fill[wine] (4.178,4.8) circle (1.1pt);
\fill[wine] (4.253,4.8) circle (1.1pt);
\fill[wine] (4.327,4.8) circle (1.1pt);
\fill[wine] (4.402,4.8) circle (1.1pt);
\fill[wine] (4.476,4.8) circle (1.1pt);
\fill[wine] (4.55,4.8) circle (1.1pt);
\fill[wine] (4.625,4.8) circle (1.1pt);
\fill[wine] (4.699,4.8) circle (1.1pt);
\fill[wine] (4.848,4.8) circle (1.1pt);
\fill[wine] (4.923,4.8) circle (1.1pt);
\fill[wine] (4.998,4.8) circle (1.1pt);
\fill[wine] (5.072,4.8) circle (1.1pt);
\fill[wine] (5.147,4.8) circle (1.1pt);
\fill[wine] (5.221,4.8) circle (1.1pt);
\fill[wine] (5.296,4.8) circle (1.1pt);
\fill[wine] (5.37,4.8) circle (1.1pt);
\fill[wine] (5.445,4.8) circle (1.1pt);
\fill[wine] (5.519,4.8) circle (1.1pt);
\fill[wine] (5.594,4.8) circle (1.1pt);
\fill[wine] (5.668,4.8) circle (1.1pt);
\fill[wine] (5.817,4.8) circle (1.1pt);
\fill[wine] (5.891,4.8) circle (1.1pt);
\fill[wine] (5.966,4.8) circle (1.1pt);
\fill[wine] (6.04,4.8) circle (1.1pt);
\fill[wine] (6.115,4.8) circle (1.1pt);
\fill[wine] (6.189,4.8) circle (1.1pt);
\fill[wine] (6.264,4.8) circle (1.1pt);
\fill[wine] (6.338,4.8) circle (1.1pt);
\fill[wine] (6.413,4.8) circle (1.1pt);
\fill[wine] (6.487,4.8) circle (1.1pt);
\fill[wine] (6.562,4.8) circle (1.1pt);
\fill[wine] (6.636,4.8) circle (1.1pt);
\fill[wine] (6.785,4.8) circle (1.1pt);
\fill[wine] (6.86,4.8) circle (1.1pt);
\fill[wine] (6.934,4.8) circle (1.1pt);
\fill[wine] (7.009,4.8) circle (1.1pt);
\fill[wine] (7.083,4.8) circle (1.1pt);
\fill[wine] (7.158,4.8) circle (1.1pt);
\fill[wine] (7.232,4.8) circle (1.1pt);
\fill[wine] (7.307,4.8) circle (1.1pt);
\fill[wine] (7.381,4.8) circle (1.1pt);
\fill[wine] (7.456,4.8) circle (1.1pt);
\fill[wine] (7.53,4.8) circle (1.1pt);
\fill[wine] (7.605,4.8) circle (1.1pt);
\fill[wine] (7.754,4.8) circle (1.1pt);
\fill[wine] (7.829,4.8) circle (1.1pt);
\fill[wine] (7.903,4.8) circle (1.1pt);
\fill[wine] (7.978,4.8) circle (1.1pt);
\fill[wine] (8.052,4.8) circle (1.1pt);
\fill[wine] (8.127,4.8) circle (1.1pt);
\fill[wine] (8.201,4.8) circle (1.1pt);
\fill[wine] (8.275,4.8) circle (1.1pt);
\fill[wine] (8.35,4.8) circle (1.1pt);
\fill[wine] (8.425,4.8) circle (1.1pt);
\fill[wine] (8.499,4.8) circle (1.1pt);
\fill[wine] (8.573,4.8) circle (1.1pt);
\fill[wine] (8.723,4.8) circle (1.1pt);
\fill[wine] (8.797,4.8) circle (1.1pt);
\fill[wine] (8.871,4.8) circle (1.1pt);
\fill[wine] (8.946,4.8) circle (1.1pt);
\fill[wine] (9.021,4.8) circle (1.1pt);
\fill[wine] (9.095,4.8) circle (1.1pt);
\fill[wine] (9.169,4.8) circle (1.1pt);
\fill[wine] (9.244,4.8) circle (1.1pt);
\fill[wine] (9.319,4.8) circle (1.1pt);
\fill[wine] (9.393,4.8) circle (1.1pt);
\fill[wine] (9.467,4.8) circle (1.1pt);
\fill[wine] (9.542,4.8) circle (1.1pt);
\fill[wine] (9.691,4.8) circle (1.1pt);
\fill[wine] (9.765,4.8) circle (1.1pt);
\fill[wine] (9.84,4.8) circle (1.1pt);
\fill[wine] (9.915,4.8) circle (1.1pt);
\fill[wine] (9.989,4.8) circle (1.1pt);
\fill[wine] (10.063,4.8) circle (1.1pt);
\fill[wine] (10.138,4.8) circle (1.1pt);
\fill[wine] (10.213,4.8) circle (1.1pt);
\fill[wine] (10.287,4.8) circle (1.1pt);
\fill[wine] (10.361,4.8) circle (1.1pt);
\fill[wine] (10.436,4.8) circle (1.1pt);
\fill[wine] (10.511,4.8) circle (1.1pt);
\fill[wine] (10.659,4.8) circle (1.1pt);
\fill[wine] (10.734,4.8) circle (1.1pt);
\fill[wine] (10.809,4.8) circle (1.1pt);
\fill[wine] (10.883,4.8) circle (1.1pt);
\fill[wine] (10.957,4.8) circle (1.1pt);
\fill[wine] (11.032,4.8) circle (1.1pt);
\fill[wine] (11.107,4.8) circle (1.1pt);
\fill[wine] (11.181,4.8) circle (1.1pt);
\fill[wine] (11.255,4.8) circle (1.1pt);
\fill[wine] (11.33,4.8) circle (1.1pt);
\fill[wine] (11.405,4.8) circle (1.1pt);
\fill[wine] (11.479,4.8) circle (1.1pt);
\fill[wine] (11.628,4.8) circle (1.1pt);
\fill[wine] (11.703,4.8) circle (1.1pt);
\draw[navy, line width=1.7pt, ->] (11.777,4.8) -- (13.267,4.8);
\draw[ink, line width=0.6pt] (11.777,4.67) -- (11.777,4.93);
\node[ink, above right, inner sep=1pt, font=\scriptsize] at (11.777,4.82) {$146$};
\draw[ink, line width=0.5pt] (0.9,-0.45) -- (13.267,-0.45);
\draw[ink, line width=0.5pt] (0.9,-0.45) -- (0.9,-0.53);
\node[ink, below] at (0.9,-0.53) {$0$};
\draw[ink, line width=0.5pt] (2.762,-0.45) -- (2.762,-0.53);
\node[ink, below] at (2.762,-0.53) {$25$};
\draw[ink, line width=0.5pt] (4.625,-0.45) -- (4.625,-0.53);
\node[ink, below] at (4.625,-0.53) {$50$};
\draw[ink, line width=0.5pt] (6.487,-0.45) -- (6.487,-0.53);
\node[ink, below] at (6.487,-0.53) {$75$};
\draw[ink, line width=0.5pt] (8.35,-0.45) -- (8.35,-0.53);
\node[ink, below] at (8.35,-0.53) {$100$};
\draw[ink, line width=0.5pt] (10.213,-0.45) -- (10.213,-0.53);
\node[ink, below] at (10.213,-0.53) {$125$};
\draw[ink, line width=0.5pt] (12.075,-0.45) -- (12.075,-0.53);
\node[ink, below] at (12.075,-0.53) {$150$};
\node[ink] at (7.083,-1.2) {denominator $Q$};
\node[ink, rotate=90] at (0.15,2.4) {numerator $P$};
\end{tikzpicture}

\caption{How the proof covers the thirteen rows. Each horizontal line is a
Farey row of numerator $P$, drawn against the denominator $Q$. The rows
$P\leqslant7$ are settled by the norm envelopes of \cref{sec:early-rows}
(teal). In the rows $8\leqslant P\leqslant13$ the dots are the $248$ reduced
fractions of the finite windows, settled by complete Gaussian fibers in
\cref{sec:gaussian} (wine); from the junction denominator $Q_P$ on, marked by
a tick, the unbounded tail is settled by the parity-order covers of
\cref{sec:tails} and the comparisons of \cref{sec:assembly} (navy). The arrows
indicate that the rows continue without bound.}
\label{fig:proof-map}
\end{figure}

\Cref{sec:occurrences} gives the two readings of the cost, recalls the
occurrence cost identity and proves the restricted injectivity reduction.
\Cref{sec:gaussian} describes the complete finite fibers and locates the
occurrence pair in them, by one exact Gaussian division per step of the
Farey tree (\cref{prop:vieta-jump}).
\Cref{sec:early-rows} develops the path envelopes and proves the
first seven rows. \Cref{sec:tails} establishes the
general parity order and the construction of the covers.
\Cref{sec:assembly} states the finite verifications, draws their
consequences, joins the domains and proves the theorem. \Cref{sec:outlook}
places the result in the larger question, with three conjectures that imply
the uniqueness conjecture and the computations that support them.

% Authors: Dr. Denys Dutykh (Mathematics Department, Khalifa University of Science and Technology,
% Abu Dhabi, UAE) and Prof. Laurent Vuillon (Univ. Savoie Mont Blanc, CNRS, LAMA, Chambery, France)
\subsection{Relation to earlier work}

The representations studied here belong to an established connection
between Markov numbers, words, and continued fractions. Reutenauer's
Christoffel correspondence parametrizes whole Markov triples
\cite{Reutenauer2009}; the palindromic-closure formula of Reutenauer and
Vuillon gives an occurrence-wise construction of their entries
\cite{ReutenauerVuillon2017}. The underlying central-word and incidence
identities are developed in \cite{BertheDeLucaReutenauer2008}.
These correspondences distinguish occurrences without presupposing that
their largest numerical entries are distinct.

A particularly close arithmetic formulation appears in
\c{C}anak\c{c}{\i} and Schiffler \cite{CanakciSchiffler2020}.
They relate primitive sums of two squares to even palindromic continued
fractions \cite[Theorem~3.10 and Corollary~3.14]{CanakciSchiffler2020}. They
construct the special continued fraction attached to a Markov occurrence
\cite[Theorem~4.1]{CanakciSchiffler2020}, and its right half, through the
proof of \cite[Corollary~4.4]{CanakciSchiffler2020}, is the occurrence pair of
\cref{prop:occurrence-cost} together with its cost $P+Q-2$, as we show after
that proposition.
They conjecture the uniqueness of the representation whose
palindromification is a Markov snake graph
\cite[Conjecture~4.5]{CanakciSchiffler2020} and identify this conjecture
with Markov uniqueness \cite[Theorem~4.8]{CanakciSchiffler2020}. Their
Conjecture~4.6 asks for uniqueness
under the weaker requirement that the half continued fraction use only
the digits $1$ and $2$; they checked both conjectures by computer for all
Markov numbers with denominator below $70$ \cite[Remark~4.7]{CanakciSchiffler2020}.
Our comparison retains every primitive positive representation of the
same integer, with no restriction on its partial quotients. The assertion
is that the occurrence-induced representation minimizes the sum of those
quotients, with the final Euclidean step omitted. It permits ties and
does not assert either of these uniqueness conjectures.

Ordering results compare a different set of objects: the Markov labels
at specified lattice points. Aigner's book records conjectures for rows of
fixed numerator, denominator and sum \cite[Section~10.1]{Aigner2013}; Rabideau and
Schiffler proved the first \cite{RabideauSchiffler2020}, and Lagisquet,
Pelantov\'a, Tavenas, and Vuillon proved all three
\cite{LagisquetEtAl2021}. Their path minimization holds at a fixed lattice
endpoint. Lee, Li, Rabideau, and Schiffler obtain further comparison
inequalities \cite{LeeEtAl2023}, which confine any coincidence of labels to
a narrow cone of slopes; their fixed-row ratio decreases to a constant,
which \cref{sec:row-labels} computes and bounds on both sides from the Cohn
trace, so that each infinite source row lies in a normalized interval. Gaster gives sharp monotonicity ranges using the
stable norm on the modular torus \cite{Gaster2022}. None of these ordering
statements, as stated, minimizes Euclidean cost over the full primitive
two-square fiber of a label. The classical uniqueness theorem for labels
that are prime powers, or twice prime powers, is the case of a fiber with a
single element; \cref{rem:single-representation} derives it from the
results of \cref{sec:occurrences,sec:gaussian}.

Two of our submitted articles treat neighboring questions. The first
studies the same Euclidean cost on the representations produced by a single
Gaussian reflection \cite{DutykhVuillon2026Endpoint}, and contains the cost
identity of \cref{prop:occurrence-cost} in its source dictionary; the second
studies palindromes in the Cohn monoid \cite{DutykhVuillon2026Reflections},
and contains \cref{lem:fixed-cost-injectivity} as the case of equal sums in
the proof of its proposition on the lengths of directing words. Their full texts can be consulted in the repository of
this series, given at the end of the article.

Related extremal problems fix the length and sum of the entries of a
continuant; Kan determines the numerical maximum under these constraints
\cite[Theorem~5]{Kan2012}.\footnote{For sequences of length $t$ and sum $S$,
the published statement assumes $2\leqslant t\leqslant S<nt$, where $n$, a bound
on the entries, belongs to the next theorem of that article; the corrected
version \cite[Theorem~3.5]{Kan2021} states the result for
$2\leqslant t\leqslant S$.}
Our proof derives the required path bounds
directly. In its infinite part, a parity order for the full normalized
ratio reduces whole boxes of gap lengths to exact corner comparisons.
The finitely many remaining labels are settled through their complete
Gaussian factorizations, so that an infinite comparison argument is
joined to a finite piece of exact arithmetic. The resulting thirteen-row
theorem assumes no upper bound on the denominator or on the Markov
value. Its injectivity consequence applies when both occurrences
belong to the proved rows; it does not exclude a collision with an
occurrence outside them.


% Authors: Dr. Denys Dutykh (Mathematics Department, Khalifa University of Science and Technology,
% Abu Dhabi, UAE) and Prof. Laurent Vuillon (Univ. Savoie Mont Blanc, CNRS, LAMA, Chambery, France)
\section{Occurrences, representations, and Euclidean cost}\label{sec:occurrences}

The arithmetic comparison in this paper concerns every primitive
representation of a fixed integer, including representations that do not
come from a Markov occurrence. We define the cost, put the two sets of
objects on the same cost scale, and explain why minimality of the cost is
enough to separate two occurrences.

\subsection{The Euclidean cost and its two readings}\label{sec:cost}

The cost of a pair is read in two ways in this paper: representations are
compared through their continued fractions, while the cost of an occurrence
is computed from a product of matrices. We define the cost by subtractions
and prove that both readings give it. For coprime positive integers $r,s$,
let $\ell(r,s)$ be the number of single subtractions needed to reach
$(1,1)$: subtract the smaller coordinate from the larger until equality. In particular,
$\ell(1,1)=0$, and exchanging the coordinates does not change $\ell$. For a
column $v$ with coprime positive entries we write $\ell(v)$ for the cost of
its two entries. Define
\begin{equation}\label{eq:continued-fraction-matrices}
 P_j=\begin{pmatrix}j&1\\1&0\end{pmatrix}\quad(j\geqslant1).
\end{equation}
For every primitive positive column $v$, exactly $j$ initial
subtractions carry $P_jv=(jv_1+v_2,v_1)^{\mathsf T}$ to
$(v_2,v_1)^{\mathsf T}$. Each of those subtractions is permitted,
including the final step when $v_1=v_2=1$. Consequently
\begin{equation}\label{eq:digit-adds-cost}
 \ell(P_jv)=j+\ell(v).
\end{equation}

\begin{lemma}[The cost by continued fractions and by matrices]\label{lem:cost-descriptions}
Let $q_0,q_1,\ldots,q_j$ be positive integers, with $q_j\geqslant2$ unless
$j=0$ and $q_0=1$. Then the column
\begin{equation}\label{eq:continued-fraction-product}
 \binom sr=P_{q_0}P_{q_1}\cdots P_{q_j}\binom10
\end{equation}
has coprime positive entries, $s/r=[q_0;q_1,\ldots,q_j]$, and
\begin{equation}\label{eq:quotient-cost}
 \ell(r,s)=q_0+q_1+\cdots+q_j-1 .
\end{equation}
Every pair of coprime integers $0<r\leqslant s$ arises in this way, from
the canonical continued fraction of $s/r$. If $(r,s)\neq(1,1)$, then
moreover
\begin{equation}\label{eq:index-product}
 \binom sr=P_{q_0}\cdots P_{q_{j-1}}P_{q_j-1}\binom11 ,
\end{equation}
and $\ell(r,s)$ is the sum of the indices of the matrices in this product.
\end{lemma}

\begin{proof}
\emph{The product.} We argue by induction on $j$. If $j=0$, the column is
$P_{q_0}(1,0)^{\mathsf T}=(q_0,1)^{\mathsf T}$, and $s/r=q_0=[q_0]$. If $j\geqslant1$,
the induction hypothesis applied to $q_1,\ldots,q_j$ gives a column
$(s',r')^{\mathsf T}=P_{q_1}\cdots P_{q_j}(1,0)^{\mathsf T}$ with positive entries and
$s'/r'=[q_1;q_2,\ldots,q_j]$, so that
$$
 \binom sr=P_{q_0}\binom{s'}{r'}=\binom{q_0s'+r'}{s'},\qquad
 \frac sr=q_0+\frac{r'}{s'}=q_0+\frac1{[q_1;q_2,\ldots,q_j]}=[q_0;q_1,\ldots,q_j].
$$
In both cases the entries are positive, and they are coprime because each
matrix $P_q$ has determinant $-1$, so that the column is the image of the
primitive column $(1,0)^{\mathsf T}$ under a matrix invertible over $\Z$.

\emph{Every pair arises.} Let $0<r\leqslant s$ be coprime. The canonical
continued fraction $[q_0;q_1,\ldots,q_j]$ of $s/r$ has positive quotients, since
$s/r\geqslant1$, and its last quotient is at least $2$ unless $s/r=1=[1]$. The
column built from it has coprime positive entries in the ratio $s/r$, and two
primitive positive columns with the same ratio are equal; hence it is
$(s,r)^{\mathsf T}$.

\emph{The cost.} For $0\leqslant i\leqslant j$ put
$v_i=P_{q_i}\cdots P_{q_j}(1,0)^{\mathsf T}$, so that $v_0=(s,r)^{\mathsf T}$,
$v_j=(q_j,1)^{\mathsf T}$ and $v_i=P_{q_i}v_{i+1}$ for $i<j$. By the first step,
every $v_i$ is a primitive positive column, so \eqref{eq:digit-adds-cost},
applied $j$ times, gives
$$
 \ell(r,s)=q_0+q_1+\cdots+q_{j-1}+\ell(q_j,1).
$$
From the pair $(q_j,1)$ every subtraction lowers the first coordinate by one,
and $(1,1)$ is reached after $q_j-1$ of them, so that $\ell(q_j,1)=q_j-1$. This
proves \eqref{eq:quotient-cost}.

\emph{The second product.} If $(r,s)\neq(1,1)$, then $q_j\geqslant2$, so that the
matrix $P_{q_j-1}$ is defined, and
$P_{q_j-1}(1,1)^{\mathsf T}=(q_j,1)^{\mathsf T}=P_{q_j}(1,0)^{\mathsf T}$. Substituting
this in \eqref{eq:continued-fraction-product} gives \eqref{eq:index-product}.
The indices of its factors add up to $q_0+\cdots+q_{j-1}+(q_j-1)$, which is
$\ell(r,s)$ by \eqref{eq:quotient-cost}.
\end{proof}

The introduction described the cost through the continued fraction, as in
\eqref{eq:quotient-cost}. The proof of \cref{prop:occurrence-cost} below
computes it through products of the matrices $P_1$ and $P_2$, by
\eqref{eq:digit-adds-cost}, as the sum of the indices in
\eqref{eq:index-product}. \Cref{lem:cost-descriptions} shows that the two
readings agree, and \cref{lem:euclidean-bijection} adds a third, the
length of a word.

\subsection{The occurrence-induced pair}

For $0<P<Q$ with $\gcd(P,Q)=1$, let $\gamma_{P/Q}$ be the binary
word of length $Q$ whose $i$th letter is $b$ precisely when
$\lfloor iP/Q\rfloor-\lfloor(i-1)P/Q\rfloor=1$.
It is the proper lower Christoffel word with
\begin{equation}\label{eq:farey-counts}
 |\gamma_{P/Q}|_a=Q-P,\qquad |\gamma_{P/Q}|_b=P.
\end{equation}
Thus $Q$ is the total length of the Christoffel word. This convention
distinguishes the denominator from the number of letters $a$.
The word starts with $a$ and ends with $b$, so write
$\gamma_{P/Q}=azb$. Its interior $z$ is a palindrome.
Indeed, for $1\leqslant j\leqslant Q-1$, coprimality gives
$\lfloor(Q-j)P/Q\rfloor=P-1-\lfloor jP/Q\rfloor$.
The increments defining the letters at positions $i$ and $Q+1-i$
therefore agree for $2\leqslant i\leqslant Q-1$.
This is the central-word property of Christoffel words
\cite{BertheDeLucaReutenauer2008}. We write $\widetilde w$ for reversal.

For later exact evaluation, the same occurrence has the Cohn description
\begin{equation}\label{eq:cohn-source}
 \mu(a)=\begin{pmatrix}2&1\\1&1\end{pmatrix},\qquad
 \mu(b)=\begin{pmatrix}5&2\\2&1\end{pmatrix},\qquad
 m_{P/Q}=\frac13\tr\mu(\gamma_{P/Q}).
\end{equation}
The standard Christoffel--Markov correspondence identifies these
occurrences with normalized Markov triples having largest entry greater
than $2$; see, for example,
\cite[Theorem~1 and its proof]{ReutenauerVuillon2017}.
No injectivity of $P/Q\mapsto m_{P/Q}$ is assumed.

The representation attached to an occurrence, and its cost, are known.
\c{C}anak\c{c}{\i} and Schiffler show that $m_{P/Q}$ is the numerator of a
palindromic continued fraction whose entries are $1$ and $2$, with
$2(Q-P-1)$ entries equal to $1$ and $2P$ entries equal to $2$, and that
$m_{P/Q}=r^2+s^2$ where $s/r$ is the continued fraction of its right half
\cite[Section~4.1, Theorem~4.1 and the proof of Corollary~4.4]{CanakciSchiffler2020}.
That half has partial quotients summing to $(Q-P-1)+2P=P+Q-1$ and ending
with $2$, so \cref{lem:cost-descriptions} gives the cost $P+Q-2$. The same statement, in the language of Euclidean paths, is the
source dictionary of our companion article \cite{DutykhVuillon2026Endpoint}.
Alessandro De Luca and the second author have another proof of the
representation, within the combinatorics of words (private communication), in
which the Christoffel word of length $m_{P/Q}$ given by the construction of
Reutenauer and Vuillon \cite{ReutenauerVuillon2017} is the image of a word with
$s$ letters $a$ and $r$ letters $b$ under a substitution that sends $a$ to a
word of length $s$ and $b$ to a word of length $r$, so that its length is
$s^2+r^2$.
We include a proof in the matrix language of the present paper, because it
defines the pair that the computations of \cref{sec:gaussian} use; the
paragraph after the proof shows that it is the pair of \c{C}anak\c{c}{\i} and
Schiffler.

\begin{proposition}[Occurrence pair and its cost]\label{prop:occurrence-cost}
Every occurrence $P/Q$ with $0<P<Q$ induces a primitive positive
representation
$$
 m_{P/Q}=r_{P/Q}^2+s_{P/Q}^2,\qquad 0<r_{P/Q}<s_{P/Q}.
$$
Its subtractive Euclidean cost is $P+Q-2$.
\end{proposition}

\begin{proof}
\emph{The label as a quadratic form.} The matrices $\mu(a)$ and $\mu(b)$
are symmetric, so reversing a word transposes its Cohn matrix, and the
palindrome $z$ has a symmetric matrix
$\mu(z)=\left(\begin{smallmatrix}\alpha&\beta\\\beta&\delta\end{smallmatrix}\right)$.
Multiplying out,
$$
 \mu(a)\,\mu(z)\,\mu(b)=
 \begin{pmatrix}
 10\alpha+9\beta+2\delta & 4\alpha+4\beta+\delta\\
 5\alpha+7\beta+2\delta & 2\alpha+3\beta+\delta
 \end{pmatrix},
$$
whose trace is $3(4\alpha+4\beta+\delta)$. With $e=(2,1)^{\mathsf T}$ one has
$e^{\mathsf T}\mu(z)e=4\alpha+4\beta+\delta$, and therefore
\begin{equation}\label{eq:square-count-identity}
 m_{P/Q}=\tfrac13\tr\mu(\gamma_{P/Q})=e^{\mathsf T}\mu(z)\,e .
\end{equation}

\emph{A square root of the form.} Cut the palindrome at its middle:
$z=\widetilde w\,w$ when $\abs z$ is even, and $z=\widetilde w\,a\,w$ or
$z=\widetilde w\,b\,w$ when $\abs z$ is odd, according to its middle letter.
Since $\mu(\widetilde w)=\mu(w)^{\mathsf T}$, $\mu(a)=P_1^2$ and
$\mu(b)=P_2^2$, with $P_1$ and $P_2$ symmetric, we obtain
$\mu(z)=H^{\mathsf T}H$, where $H$ is $\mu(w)$, $P_1\mu(w)$ or $P_2\mu(w)$ in
the three cases; the empty palindrome gives $H=I$. By
\eqref{eq:square-count-identity}, $m_{P/Q}=(He)^{\mathsf T}(He)$.

\emph{The pair.} The integer matrix $H$ has determinant $\pm1$, so the
column $He$ is primitive. Its entries are positive, because $H$ is a product
of nonnegative matrices with no zero row and $e$ is positive. They are
distinct: two equal coprime entries would both be $1$, giving
$m_{P/Q}=2$, whereas $m_{P/Q}\geqslant5$, because $\mu(z)\geqslant I$ entrywise
(a product of nonnegative matrices each at least $I$) and
\eqref{eq:square-count-identity} then gives $m_{P/Q}\geqslant e^{\mathsf T}e=5$. Sorting the two entries defines
$r_{P/Q}<s_{P/Q}$. The construction uses the occurrence alone and assumes
nothing about other occurrences with the same label.

\emph{Its cost.} The column $e$ reaches $(1,1)$ after one subtraction, so
$\ell(e)=1$. Apply the factors of $H$ to $e$ from the right. By
\eqref{eq:digit-adds-cost}, a factor $\mu(a)=P_1P_1$ coming from a letter
$a$ of $w$ adds $2$ to the cost, a factor $\mu(b)=P_2P_2$ adds $4$, and the
middle factor $P_1$ or $P_2$, when present, adds $1$ or $2$. Each letter of
$w$ occurs twice in $z$, so these increments add up to $\abs z_a+2\abs z_b$,
and \eqref{eq:farey-counts} gives
$$
 \ell(He)=1+|z|_a+2|z|_b
 =1+(Q-P-1)+2(P-1)=P+Q-2.
$$
No restriction on the parity of $|z|$ or of $m_{P/Q}$ is needed.
\end{proof}

The pair of \c{C}anak\c{c}{\i} and Schiffler is this pair. Write
$\gamma_{P/Q}=a^{u_1}b\,a^{u_2}b\cdots a^{u_P}b$. The $k$th letter $b$ stands at
position $\lceil kQ/P\rceil$, and $kQ/P$ is not an integer for $1\leqslant k<P$;
comparison with their crossing numbers $v_k$, which satisfy
$v_1+\cdots+v_k=\lfloor kQ/P\rfloor$ for $k<P$ and $v_1+\cdots+v_P=Q-1$, gives
$u_1=v_1$ and $u_k=v_k-1$ for $2\leqslant k\leqslant P$. Their palindromic
continued fraction is therefore $[2,\sigma(z),2]$, where $\sigma$ replaces each
letter $a$ by the two entries $1,1$ and each letter $b$ by $2,2$. Its right
half is $[\sigma(w),2]$, or $[j,\sigma(w),2]$ with $j=1$ or $2$ for a middle
letter $a$ or $b$, and since $\mu(a)=P_1^2$, $\mu(b)=P_2^2$ and $e=P_2(1,0)^{\mathsf T}$,
this is the continued fraction of the column $He$.

\subsection{Every primitive pair has a unique Euclidean directive}

Introduce the elementary matrices
\begin{equation}\label{eq:directive-incidence}
 L_a=\begin{pmatrix}1&1\\0&1\end{pmatrix},\qquad
 L_b=\begin{pmatrix}1&0\\1&1\end{pmatrix},\qquad
 L_{uv}=L_uL_v,\qquad {\bf1}=\binom11.
\end{equation}

\begin{lemma}[Euclidean bijection]\label{lem:euclidean-bijection}
The map $u\mapsto L_u{\bf1}$ is a bijection from binary words to
ordered coprime positive pairs, and
\begin{equation}\label{eq:euclidean-length}
 \ell(L_u{\bf1})=|u|.
\end{equation}
\end{lemma}

\begin{proof}
If $x>y>0$, the first letter of any directive producing $(x,y)$ must
be $a$, because $L_a$ maps positive pairs to pairs with first coordinate
larger, whereas $L_b$ makes the second coordinate larger. The unique
predecessor is $(x-y,y)$. For $y>x$, the first letter is $b$ and the
predecessor is $(x,y-x)$. Each subtraction preserves positivity and
coprimality and decreases $x+y$. The only equal coprime positive pair
is $(1,1)$, so this process terminates there. Reading the uniquely
determined letters gives both existence and uniqueness, and one letter
is removed at each subtraction. This proves
\eqref{eq:euclidean-length}.
\end{proof}

The word $u$ is the \emph{directive} of the pair $L_u{\bf1}$, and it can be
read on the continued fraction. With
$J=\left(\begin{smallmatrix}0&1\\1&0\end{smallmatrix}\right)$, which exchanges
the two coordinates, one has $P_q=L_a^qJ$, $JL_a^qJ=L_b^q$ and $J{\bf1}={\bf1}$.
Grouping the factors of \eqref{eq:index-product} in consecutive pairs
therefore gives, for coprime $0<r<s$ with $s/r=[q_0;q_1,\ldots,q_j]$,
$$
 \binom sr=L_a^{q_0}L_b^{q_1}L_a^{q_2}L_b^{q_3}\cdots\,{\bf1},
$$
where the exponents are $q_0,q_1,\ldots,q_{j-1},q_j-1$ and the letters
alternate. The directive of $(s,r)$ is thus the word
$a^{q_0}b^{q_1}a^{q_2}\cdots$ whose runs have these lengths, and its length
$q_0+\cdots+q_j-1$ agrees with \eqref{eq:quotient-cost} and
\eqref{eq:euclidean-length}. For instance $21/13=[1;1,1,1,1,2]$, and the
directive of $(21,13)$ is $ababab$.

For unordered representations we always choose $0<r<s$. This loses
neither norms nor costs. The exceptional positive pair $(1,1)$ represents
$2$ and can be treated separately; $1=1^2+0^2$ is outside the positive
pair convention.

\begin{definition}\label{def:weak-minimality}
An occurrence $P/Q$ is \emph{weakly Euclidean-minimal} if
\begin{equation}\label{eq:weak-minimality}
 P+Q-2\leqslant\ell(r,s)
 \quad\text{whenever}\quad
 m_{P/Q}=r^2+s^2,\quad 0<r<s,\quad\gcd(r,s)=1.
\end{equation}
\end{definition}

The competitors in \eqref{eq:weak-minimality} range over the entire
primitive two-square fiber. Their Euclidean directives need not be
directives induced by a central Cohn word as above.
Ties with such other representations are allowed.

\subsection{Why weak minimality is enough within a proved family}

Among the ordering theorems recalled in the introduction, we use the
fixed-sum theorem of Lagisquet, Pelantov\'a, Tavenas and Vuillon
\cite[Theorem~2, item~3]{LagisquetEtAl2021}, which compares two lattice points
differing by a unit step along a line of fixed sum.
Applying it $i$ times gives, for reduced $P/Q$ and
$(P-i)/(Q+i)$ with $0<i<P<Q$,
\begin{equation}\label{eq:fixed-sum-order}
 m_{P/Q}<m_{(P-i)/(Q+i)}.
\end{equation}
The intermediate points of that iteration need not be reduced; this is
harmless because the authors define the label at every lattice point,
not only at Farey fractions.
The inequality compares distinct points on one line of fixed sum $P+Q$.
It does not compare arbitrary two-square representations.

The next lemma is an immediate consequence of \eqref{eq:fixed-sum-order};
it is also the case of equal sums in the proof of the proposition on the
lengths of directing words of Markov words in
\cite{DutykhVuillon2026Reflections}, where the cost $P+Q-2$ appears as the
length of a directing word.

\begin{lemma}[Fixed-cost injectivity]\label{lem:fixed-cost-injectivity}
Two occurrences having the same label and the same occurrence cost
have the same reduced Farey coordinate.
\end{lemma}

\begin{proof}
Their costs are $P+Q-2$ and $P'+Q'-2$. If they are equal and
$P>P'$, set $i=P-P'$. Then $Q'=Q+i$ and $0<i<P$, so
\eqref{eq:fixed-sum-order} contradicts equality of labels. The other
strict ordering of the numerators is excluded in the same way.
Consequently $P=P'$ and $Q=Q'$. A reduced Farey coordinate has a unique
lower Christoffel word, hence a unique normalized occurrence.
\end{proof}

\begin{corollary}\label{cor:restricted-injectivity}
The label map is injective on any family of occurrences all of which
are weakly Euclidean-minimal.
\end{corollary}

\begin{proof}
If two members have equal label, each occurrence-induced pair is a
competitor in the other's two-square fiber. Weak minimality in both
directions makes their costs equal. Apply
Lemma~\ref{lem:fixed-cost-injectivity}.
\end{proof}

Thus weak minimality on every occurrence would imply the
Frobenius--Markov uniqueness conjecture. For a restricted family, both occurrences
must lie in that family. If only one is known to be weakly minimal, a
distinct equal-label occurrence is merely forced to have strictly
larger cost.

% Authors: Dr. Denys Dutykh (Mathematics Department, Khalifa University of Science and Technology,
% Abu Dhabi, UAE) and Prof. Laurent Vuillon (Univ. Savoie Mont Blanc, CNRS, LAMA, Chambery, France)
\section{Complete Gaussian fibers and the finite windows}
\label{sec:gaussian}

The finite part of the proof must compare an occurrence with every
primitive representation of its label. Searching a coordinate rectangle
would obscure completeness. Gaussian factorization gives an exact finite
parameter space whose size depends on the distinct prime divisors, not
on the size of the coordinates, and once the prime factorizations are
known the enumeration is a matter of integer arithmetic alone.

\subsection{The whole primitive representation fiber}

For $M\geqslant5$, define
$$
 \mathcal F(M)=\{(x,y)\in\Z^2:0<x<y,\ \gcd(x,y)=1,\ x^2+y^2=M\}.
$$
The source construction in Section~\ref{sec:occurrences} proves that this
set is nonempty for every occurrence considered here.

\begin{proposition}[Primitive Gaussian allocation]
\label{prop:gaussian-fibre}
If $\mathcal F(M)\ne\varnothing$, then
\begin{equation}\label{eq:gaussian-factorization}
 M=2^\epsilon\prod_{j=1}^t p_j^{e_j},\qquad
 \epsilon\in\{0,1\},\quad t\geqslant1,\quad
 p_j\equiv1\pmod4,
\end{equation}
where the $p_j$ are distinct primes and $e_j\geqslant1$.
Choose $\pi_j=a_j+ib_j$ with $p_j=a_j^2+b_j^2$ and $0<a_j<b_j$.
The elements of $\mathcal F(M)$ are exactly the sorted absolute
coordinate pairs of
\begin{equation}\label{eq:gaussian-sign-cube}
 (1+i)^\epsilon\prod_{j=1}^t
 \begin{cases}
 \pi_j^{e_j},&\sigma_j=1,\\
 \overline{\pi_j}^{\,e_j},&\sigma_j=-1,
 \end{cases}
 \qquad (\sigma_1,\ldots,\sigma_t)\in\{1,-1\}^t.
\end{equation}
Each normalized pair occurs exactly twice. Equivalently, fixing
$\sigma_1=1$ lists the fiber without repetition, and
$|\mathcal F(M)|=2^{t-1}$.
\end{proposition}

\begin{proof}
The Gaussian integers form a Euclidean domain for the squared norm,
hence have unique factorization. If $z=x+iy$ is primitive and a rational
prime $p\equiv3\pmod4$ divides $z\overline z$, its Gaussian primality
forces it to divide $z$ or $\overline z$. Either condition gives
$p\mid x,y$, a contradiction. Likewise $4\nmid M$: a sum of two
squares divisible by four has both coordinates even. This proves the
shape \eqref{eq:gaussian-factorization}; $t=0$ would leave $M=1$ or $2$.

For a split prime $p_j=\pi_j\overline{\pi_j}$, an allocation using both
$\pi_j$ and $\overline{\pi_j}$ in $z$ would make $p_j$ divide both
coordinates. Primitivity therefore puts the whole exponent $e_j$ on
one of the two conjugate factors. Conversely, each allocation in
\eqref{eq:gaussian-sign-cube} has primitive coordinates: a common odd
prime would require both conjugate Gaussian factors, and a common factor
two would require a higher ramified exponent. Unique factorization proves
that these allocations, up to multiplication by a unit, exhaust all
representations. Since $M\geqslant5$, a primitive Gaussian integer of norm $M$
has nonzero coordinates of distinct absolute values, for otherwise $M=1$ or
$M=2$; its sorted absolute pair therefore lies in $\mathcal F(M)$.

It remains to count. Two Gaussian integers $x+iy$ and $x'+iy'$ have the
same normalized pair, that is $\{\abs x,\abs y\}=\{\abs{x'},\abs{y'}\}$,
exactly when $x'+iy'$ is $u(x+iy)$ or $u(x-iy)$ for one of the four units
$u$. Write $z_\sigma$ for the product \eqref{eq:gaussian-sign-cube}. For an
odd prime $p_j=a_j^2+b_j^2$ the factors $\pi_j$ and $\overline{\pi_j}$ are
not associate: $\overline{\pi_j}=u\pi_j$ would force $b_j=0$, $a_j=0$ or
$a_j=\pm b_j$, none of which is possible since $0<a_j<b_j$. Suppose now that
$z_{\sigma'}$ and $z_\sigma$ give the same normalized pair. If
$z_{\sigma'}=uz_\sigma$, unique factorization compares the exponents of the
non-associate primes $\pi_j$ and $\overline{\pi_j}$ and gives $\sigma'=\sigma$.
If $z_{\sigma'}=u\overline{z_\sigma}$, the same comparison gives
$\sigma'=-\sigma$, because conjugation exchanges $\pi_j$ and
$\overline{\pi_j}$ and replaces $1+i$ by the associate $1-i=-i(1+i)$.
Conversely $z_{-\sigma}$ is an associate of $\overline{z_\sigma}$, so $\sigma$
and $-\sigma$ always give the same pair. As $t\geqslant1$, the vectors
$\sigma$ and $-\sigma$ are distinct, so each normalized pair comes from
exactly two sign vectors, which proves the count and the last assertion.
\end{proof}

The exponents in \eqref{eq:gaussian-factorization} do not enlarge the
primitive sign cube. Mixed allocations of a prime power generate
imprimitive representations, so omitting them is a proved reduction,
not a filter based on observed costs.

\begin{remark}[Labels with a single representation]\label{rem:single-representation}
By \cref{prop:gaussian-fibre}, a label $m>2$ has exactly one primitive
positive representation precisely when $t=1$, that is when $m=p^e$ or
$m=2p^e$ for a prime $p\equiv1\pmod4$. Such a label belongs to exactly one
reduced fraction. \emph{Step 1.} Let $m=m_{P/Q}=m_{P'/Q'}$; by
\cref{prop:occurrence-cost} the two occurrence pairs are elements of
$\mathcal F(m)$, of costs $P+Q-2$ and $P'+Q'-2$. \emph{Step 2.} Since
$\abs{\mathcal F(m)}=2^{t-1}=1$, the two pairs coincide, and so do their
costs: $P+Q=P'+Q'$. \emph{Step 3.} \Cref{lem:fixed-cost-injectivity} gives
$P/Q=P'/Q'$. This recovers the theorem of Baragar, Button and Schmutz that a
Markov number which is a prime power or twice a prime power is unique
\cite{Baragar1996,Button1998,Schmutz1996}, in the form stated in
\cite[Theorem~C]{Zhang2007}; see also \cite{LangTan2007}. Here it follows
from the fixed-sum ordering \eqref{eq:fixed-sum-order} and the Gaussian
count, and \cref{thm:thirteen-rows} extends the same fiber argument, on
thirteen rows, to fibers of every size.
\end{remark}

\begin{example}[Two representations, two costs]
The first-row label $m_{1/7}=610=2\cdot5\cdot61$ has exactly two
primitive positive representations:
$$
 610=13^2+21^2=9^2+23^2.
$$
In the normalization of the proposition the two split primes are
$1+2i$ and $5+6i$, and the two allocations give
$(1+i)(1+2i)(5+6i)=-23+9i$ and $(1+i)(1-2i)(5+6i)=21+13i$. Proposition~\ref{prop:gaussian-fibre}
proves that the list is complete. Their continued fractions are
$$
 \frac{21}{13}=[1;1,1,1,1,2],\qquad
 \frac{23}{9}=[2;1,1,4],
$$
so the Euclidean costs are $6$ and $7$, respectively. The occurrence
cost $1+7-2=6$ is therefore the minimum. This illustrates why the
optimization question persists even when the label and its Markov
occurrence are fixed: the full arithmetic fiber contains another
primitive pair whose directive is unrestricted.
\end{example}

\Cref{prop:gaussian-fibre} lists the whole fiber but does not say which of
its elements comes from the occurrence. The two labels next to the
occurrence in the Markov tree do: the occurrence pair is the element of the
sign cube whose prime factors all divide one explicit Gaussian integer.

\begin{proposition}[The occurrence pair in the sign cube]\label{prop:occurrence-in-cube}
Let $0<P<Q$ be coprime, let $L=P_L/Q_L<R=P_R/Q_R$ be its two Farey
parents, the reduced fractions with $P_L+P_R=P$, $Q_L+Q_R=Q$ and
$P_RQ_L-P_LQ_R=1$, and let $m_L$, $m_R$ be their labels, with $m_{0/1}=1$ and
$m_{1/1}=2$. Write $(s,r)^{\mathsf T}=He$ for the
occurrence pair of $P/Q$ before sorting, and $m=m_{P/Q}$. Then the ideal of
$\Z[i]$ generated by $m$ and $m_L+i\,m_R$ is generated by $s+ir$ or by its
conjugate. In particular, in the sign cube \eqref{eq:gaussian-sign-cube} of
$m$, the occurrence pair is the element whose prime factors all divide
$m_L+i\,m_R$, up to the conjugation of the whole product.
\end{proposition}

\begin{proof}
\emph{The root read from the Cohn matrix.} The lower Christoffel word of
$P/Q$ is the product of those of its two Farey parents, the words of $0/1$
and $1/1$ being the letters $a$ and $b$ \cite{Reutenauer2009}. Writing
$\gamma_L$ and $\gamma_R$ for the words of the parents, we have
$\mu(\gamma_{P/Q})=\mu(\gamma_L)\,\mu(\gamma_R)$. Each factor has
determinant one and upper right entry equal to its label: for the letters,
$\mu(a)$ and $\mu(b)$ show it, and for a longer Christoffel word the
computation in the proof of \cref{prop:occurrence-cost} gives an upper right
entry $4\alpha+4\beta+\delta$, which is the label by
\eqref{eq:square-count-identity}. Writing
$\mu(\gamma_{P/Q})=\left(\begin{smallmatrix}a&m\\c&d\end{smallmatrix}\right)$ and
$\mu(\gamma_X)=\left(\begin{smallmatrix}a_X&m_X\\c_X&d_X\end{smallmatrix}\right)$
for $X=L,R$, the upper row of the product gives $m=a_Lm_R+m_Ld_R$ and
$a=a_La_R+m_Lc_R$, so that
$$
 a\,m_R=a_R(m-m_Ld_R)+m_Lc_Rm_R\equiv-m_L(a_Rd_R-m_Rc_R)=-m_L\pmod m .
$$

\emph{The root read from the occurrence pair.} By the proof of
\cref{prop:occurrence-cost}, $\mu(\gamma_{P/Q})=\mu(a)H^{\mathsf T}H\mu(b)$. With
$e_1=(1,0)^{\mathsf T}$, the symmetric matrix $\mu(a)$ satisfies $\mu(a)e_1=e$,
and $\mu(b)e_1=(5,2)^{\mathsf T}=2e+e_1$. Put $v=He=(s,r)^{\mathsf T}$ and
$h=He_1=(p,q)^{\mathsf T}$. Then
$$
 a=e_1^{\mathsf T}\mu(\gamma_{P/Q})e_1=(He)^{\mathsf T}H(2e+e_1)=2m+sp+rq .
$$
The matrix with columns $v$ and $h$ is $HP_2$, of determinant
$\eta=sq-rp=\pm1$, so that $rp=sq-\eta$ and
$$
 a\,r\equiv(sp+rq)\,r=s(sq-\eta)+r^2q=qm-\eta s\equiv-\eta s\pmod m .
$$

\emph{The ideals.} The labels $m_L$, $m_R$, $m$ form a Markov triple, whose
entries are pairwise coprime \cite{Aigner2013}; and $\gcd(s,r)=1$ with
$s^2+r^2=m$. Hence $m_R$ and $r$ are invertible modulo $m$. By the first step,
$m_L+i\,m_R\equiv m_R\,(i-a)\pmod m$, and by the second, $s\equiv-\eta ar$, so that
$s+ir\equiv r\,(i-\eta a)$ and $s-ir\equiv-r\,(i+\eta a)\pmod m$. If $\eta=1$, the
first congruence and the third give equal ideals
$(m,s+ir)=(m,i-a)=(m,m_L+i\,m_R)$; if $\eta=-1$, the fourth does,
$(m,s-ir)=(m,m_L+i\,m_R)$. Finally $(m,s\pm ir)=(s\pm ir)$, because $s\pm ir$
divides $m=(s+ir)(s-ir)$.

\emph{The sign cube.} The generator $s\pm ir$ of the ideal is the greatest
common divisor of $m$ and $m_L+i\,m_R$. A Gaussian prime factor of the label
divides it exactly when it divides $m_L+i\,m_R$, which identifies the
allocation of \eqref{eq:gaussian-sign-cube}, up to the conjugation of the
whole product.
\end{proof}

The congruence $a\,m_R\equiv-m_L\pmod m$ of the first step is classical:
up to sign, the residue of $a$ modulo $m$ is the characteristic number of
the Markov triple $(m,m_L,m_R)$, defined by Frobenius as the solution of
$m_Rx\equiv\pm m_L\pmod m$, a square root of $-1$ modulo $m$
\cite[\S2, p.~462]{Frobenius1913}, see also \cite[Section~2.2]{Aigner2013}.
The proposition translates it into the Gaussian integers and identifies the
corresponding element of the sign cube with the occurrence pair.
The cofactor $(m_L+i\,m_R)/(s\pm ir)$ has norm
$(m_L^2+m_R^2)/m=3m_Lm_R-m$, by the Markov equation, in the form
$mm'=m_L^2+m_R^2$ with $m'=3m_Lm_R-m$ that Frobenius uses
\cite[\S1, p.~460]{Frobenius1913}. For the label $610$ of
the example above, the parents $0/1$ and $1/6$ have labels $1$ and $233$, and
$$
 1+233\,i=(21+13\,i)(5+8\,i):
$$
the first factor is the occurrence pair $(1+i)(1-2i)(5+6i)$, and the second
has norm $89=3\cdot1\cdot233-610$, the label of $1/5$. The second factor is no
accident: $1/5$ is the third vertex of the Farey triangle $(0/1,1/6,1/5)$ of
$1/7$, defined below, and the next subsection shows that the cofactor is
always the occurrence pair of that vertex, up to a unit and a conjugation;
here $5+8i=i\,(8-5i)$, and $(8,5)$ is the occurrence pair of $1/5$.

\subsection{The occurrence pairs along the Farey tree}\label{sec:vieta-jump}

\Cref{prop:occurrence-in-cube} finds the occurrence pair as a greatest
common divisor, a computation that still involves the label itself. The
Markov tree offers a recursion instead. Along the tree the labels obey two
relations of Vieta, a sum and a product, and the product lifts exactly to
the Gaussian integers. The Gaussian integer $m_L+i\,m_R$ of
\cref{prop:occurrence-in-cube} is the product of two Gaussian integers
attached to $P/Q$ and to the third vertex of its Farey triangle, which are
their occurrence pairs up to units and conjugation; once an orientation is
fixed at the root of the tree, the product holds exactly, without any unit.
Every occurrence pair is thus obtained from pairs found earlier by one exact
division in $\Z[i]$, with no factorization, no greatest common divisor and
no continued fraction. This is the Gaussian form of the replacement
$(x,y,z)\mapsto(x,y,3xy-z)$ that generates the Markov tree.

Two fractions $L=P_L/Q_L<R=P_R/Q_R$ with nonnegative terms are \emph{Farey
neighbors} when $P_RQ_L-P_LQ_R=1$. Their mediant is
$M=(P_L+P_R)/(Q_L+Q_R)$, and we call their \emph{difference} the fraction $D$
whose terms are $\abs{P_R-P_L}$ and $\abs{Q_R-Q_L}$. The Farey tree of $[0,1]$
starts from the pair $(0/1,1/1)$, whose difference is $1/0$, and replaces
each pair $(L,R)$ by the two pairs $(L,M)$ and $(M,R)$. Every reduced
fraction of $(0,1)$ is the mediant of exactly one pair of the tree, the pair
of its two Farey parents \cite[Section~3.2, Theorem~3.9 and its proof]{Aigner2013}.
The differences of the two new pairs are $R$ and $L$, because $M-L=R$ and
$R-M=-L$ term by term, so that the difference of a pair of the tree is
either one of the three fractions $0/1$, $1/1$ and $1/0$, or a mediant
formed earlier. We write $\gamma_X$ and $m_X$ for the lower Christoffel word
and the label of a fraction $X$ of $[0,1]$, with $\gamma_{0/1}=a$ and
$\gamma_{1/1}=b$, we set $m_{1/0}=1$, and we call $(L,R,D)$ the \emph{Farey
triangle} of the mediant $M$. Finally, $N(g+ih)=g^2+h^2$ denotes the norm of
$\Z[i]$.

The first lemma records the two relations of Vieta on every triangle. In
this indexing they are the content of Theorem~1 of Reutenauer and of its
proof \cite[Theorem~1 and its proof]{Reutenauer2009}, where the Fricke
identity gives the Markov triple of a standard factorization and two trace
identities give the two cases below; we include the argument in the
notation of the Farey tree.

\begin{lemma}[The labels on a Farey triangle]\label{lem:farey-triangle}
For every pair $(L,R)$ of the tree, with mediant $M$ and difference $D$, the
triple $(m_L,m_R,m_D)$ is a Markov triple with pairwise coprime entries, and
\begin{equation}\label{eq:vieta-relations}
 m_M=3m_Lm_R-m_D,\qquad m_Mm_D=m_L^2+m_R^2 .
\end{equation}
\end{lemma}

\begin{proof}
\emph{The sum relation.} Put $A_X=\mu(\gamma_X)$ for every fraction $X\neq1/0$
of the tree. Since $\gamma_M=\gamma_L\gamma_R$ \cite{Reutenauer2009}, we have
$A_M=A_LA_R$, and each $A_X$ has determinant one and trace $3m_X$ by
\eqref{eq:cohn-source}. For two matrices $A$ and $B$ of determinant one, the
Cayley--Hamilton theorem gives $B+B^{-1}=(\tr B)\,I$; multiplying by $A$ and
taking traces,
$$
 \tr(AB)+\tr(AB^{-1})=\tr A\,\tr B .
$$
With $A=A_L$ and $B=A_R$ this reads $3m_M+\tr(A_LA_R^{-1})=9m_Lm_R$, and it
remains to show that $\tr(A_LA_R^{-1})=3m_D$. For the starting pair,
$$
 A_{0/1}A_{1/1}^{-1}
 =\begin{pmatrix}2&1\\1&1\end{pmatrix}\begin{pmatrix}1&-2\\-2&5\end{pmatrix}
 =\begin{pmatrix}0&1\\-1&3\end{pmatrix},
$$
whose trace is $3=3m_{1/0}$. Every other pair $(L,R)$ was created from a pair
of the tree in one of two ways. Either it is $(L,M')$, where $M'=R$ is the
mediant of a pair $(L,R')$; then $D=R'$ and $A_R=A_LA_D$, so that
$A_LA_R^{-1}=A_LA_D^{-1}A_L^{-1}$, whose trace is $\tr A_D^{-1}=\tr A_D=3m_D$, a
matrix of determinant one and its inverse having the same trace. Or it is
$(M',R)$, where $M'=L$ is the mediant of a pair $(L',R)$; then $D=L'$ and
$A_L=A_DA_R$, so that $A_LA_R^{-1}=A_D$, whose trace is $3m_D$.

\emph{The Markov triples.} The starting triple $(1,2,1)$ satisfies
$x^2+y^2+z^2=3xyz$. If $(m_L,m_R,m_D)$ satisfies it, then $m_D$ is a root of
the polynomial $X^2-3m_Lm_RX+m_L^2+m_R^2$, and by the sum relation
$m_M=3m_Lm_R-m_D$ is its other root. The product of the two roots is
$m_Mm_D=m_L^2+m_R^2$, which is the second relation in
\eqref{eq:vieta-relations}, and $(m_L,m_R,m_M)$ satisfies the Markov equation
as well. The triples of the two new pairs, $(m_L,m_M,m_R)$ and
$(m_M,m_R,m_L)$, are permutations of it, and induction on the depth in the
tree proves that every triple $(m_L,m_R,m_D)$ is a Markov triple.

\emph{Coprimality.} If a prime divides two entries of a Markov triple, the
Markov equation shows that it divides the square of the third, hence the
third. The triple of a new pair has the entries $m_L$, $m_R$ and $m_M$, so a
prime dividing two of them would divide $m_L$ and $m_R$, two entries of the
triple of the parent pair. Since the starting triple $(1,2,1)$ has pairwise
coprime entries, induction on the depth proves the last assertion.
\end{proof}

The product relation is the one to lift. Put $\alpha_{0/1}=1$,
$\alpha_{1/1}=1+i$ and $\alpha_{1/0}=1$, and, for each pair $(L,R)$ of the
tree, with mediant $M$ and difference $D$,
\begin{equation}\label{eq:vieta-jump}
 \overline{\alpha_M}=\frac{m_L+i\,m_R}{\alpha_D}.
\end{equation}
The definition proceeds down the tree: when a pair is reached, $\alpha_D$ is
already known, since $D$ is a starting fraction or a mediant formed earlier,
and it remains to see that the division is exact. Three divisibilities
attached to each triangle ensure it. They are exchanged by the circular
permutation $L\to R\to D\to L$: the first says that the division
\eqref{eq:vieta-jump} is exact, and the other two prepare the divisions of the
next generation. The first lemma below carries them from a pair to its two
new pairs; the second runs the induction down the tree.

\begin{lemma}[One step down the tree]\label{lem:vieta-step}
Let $(L,R)$ be a pair of the tree, with mediant $M$ and difference $D$, and
let $\alpha_L$, $\alpha_R$, $\alpha_D$ be Gaussian integers of norms $m_L$, $m_R$,
$m_D$ such that
\begin{equation}\label{eq:three-divisibilities}
 \alpha_D\mid m_L+i\,m_R,\qquad \alpha_L\mid m_R+i\,m_D,\qquad
 \alpha_R\mid m_D+i\,m_L .
\end{equation}
Then the quotient \eqref{eq:vieta-jump} is a Gaussian integer, $N(\alpha_M)=m_M$,
and
\begin{equation}\label{eq:new-divisibilities}
 \alpha_M\mid m_R+i\,m_L,\qquad \alpha_L\mid m_M+i\,m_R,\qquad
 \alpha_R\mid m_L+i\,m_M .
\end{equation}
\end{lemma}

\begin{proof}
\emph{The norm of $\alpha_M$.} The first divisibility in
\eqref{eq:three-divisibilities} makes the quotient \eqref{eq:vieta-jump} a
Gaussian integer, and by \cref{lem:farey-triangle},
$$
 N(\alpha_M)=\frac{N(m_L+i\,m_R)}{N(\alpha_D)}=\frac{m_L^2+m_R^2}{m_D}
 =\frac{m_Mm_D}{m_D}=m_M .
$$

\emph{Divisibility by $\alpha_M$.} Conjugating \eqref{eq:vieta-jump} gives
$\alpha_M\overline{\alpha_D}=m_L-i\,m_R$, and multiplying by $i$,
$$
 \alpha_M\cdot i\,\overline{\alpha_D}=m_R+i\,m_L .
$$

\emph{Divisibility by $\alpha_L$.} We compute modulo $\alpha_L$, which divides
$m_L=\alpha_L\overline{\alpha_L}$. The definition gives
$\overline{\alpha_M}\alpha_D=m_L+i\,m_R\equiv i\,m_R$, and its conjugate gives
$\alpha_M\overline{\alpha_D}=m_L-i\,m_R$. With $m_M=\alpha_M\overline{\alpha_M}$ and
$m_D=\alpha_D\overline{\alpha_D}$ it follows that, modulo $\alpha_L$,
\begin{align*}
 m_M+i\,m_R&\equiv\alpha_M\overline{\alpha_M}+\overline{\alpha_M}\alpha_D
 =\overline{\alpha_M}\,(\alpha_M+\alpha_D),\\
 \overline{\alpha_D}\,(\alpha_M+\alpha_D)&=m_L-i\,m_R+m_D\equiv m_D-i\,m_R
 =-i\,(m_R+i\,m_D).
\end{align*}
The last member is divisible by $\alpha_L$, by the second divisibility in
\eqref{eq:three-divisibilities}. The Gaussian integers $\alpha_L$ and
$\overline{\alpha_D}$ have no common prime factor, because their norms $m_L$ and
$m_D$ are coprime by \cref{lem:farey-triangle}. Unique factorization in $\Z[i]$
therefore gives $\alpha_L\mid\alpha_M+\alpha_D$, and then $\alpha_L\mid m_M+i\,m_R$
by the first line.

\emph{Divisibility by $\alpha_R$.} We compute modulo $\alpha_R$, which divides
$m_R$. The definition gives $\overline{\alpha_M}\alpha_D=m_L+i\,m_R\equiv m_L$,
and therefore, modulo $\alpha_R$,
\begin{align*}
 m_L+i\,m_M&\equiv\overline{\alpha_M}\alpha_D+i\,\alpha_M\overline{\alpha_M}
 =\overline{\alpha_M}\,(\alpha_D+i\,\alpha_M),\\
 \overline{\alpha_D}\,(\alpha_D+i\,\alpha_M)&=m_D+i\,(m_L-i\,m_R)
 =m_D+m_R+i\,m_L\equiv m_D+i\,m_L .
\end{align*}
The last member is divisible by $\alpha_R$, by the third divisibility in
\eqref{eq:three-divisibilities}. The norms $m_R$ and $m_D$ are coprime, so
$\alpha_R$ and $\overline{\alpha_D}$ have no common prime factor, and unique
factorization gives $\alpha_R\mid\alpha_D+i\,\alpha_M$, and then
$\alpha_R\mid m_L+i\,m_M$ by the first line.
\end{proof}

\begin{lemma}[Three divisibilities]\label{lem:three-divisibilities}
Every quotient \eqref{eq:vieta-jump} is a Gaussian integer, and every pair
$(L,R)$ of the tree, with difference $D$, satisfies
\eqref{eq:three-divisibilities} and
\begin{equation}\label{eq:three-norms}
 N(\alpha_L)=m_L,\qquad N(\alpha_R)=m_R,\qquad N(\alpha_D)=m_D.
\end{equation}
\end{lemma}

\begin{proof}
We argue by induction on the depth of the pair in the tree.

\emph{The starting pair.} For $(L,R,D)=(0/1,1/1,1/0)$, the integers $1$, $1+i$
and $1$ have the norms $1$, $2$ and $1$. The first two divisibilities have
the divisor $1$, and the third says that $1+i$ divides
$m_{1/0}+i\,m_{0/1}=1+i$.

\emph{The step.} Assume \eqref{eq:three-norms} and
\eqref{eq:three-divisibilities} for a pair $(L,R)$. By \cref{lem:vieta-step},
the quotient \eqref{eq:vieta-jump} is a Gaussian integer of norm $m_M$, and
\eqref{eq:new-divisibilities} holds. The pair $(L,M)$ has difference $R$, and
its three divisibilities read $\alpha_R\mid m_L+i\,m_M$, $\alpha_L\mid m_M+i\,m_R$
and $\alpha_M\mid m_R+i\,m_L$; the pair $(M,R)$ has difference $L$, and its
three divisibilities read $\alpha_L\mid m_M+i\,m_R$, $\alpha_M\mid m_R+i\,m_L$ and
$\alpha_R\mid m_L+i\,m_M$. Both are the three divisibilities
\eqref{eq:new-divisibilities}, and the norms of the two new pairs are among
$m_L$, $m_R$ and $m_M=N(\alpha_M)$. Both new pairs therefore satisfy
\eqref{eq:three-norms} and \eqref{eq:three-divisibilities}, which completes the
induction.
\end{proof}

The lemma defines $\alpha_M$ for every reduced fraction $M$ of $(0,1)$, since
$M$ is the mediant of exactly one pair of the tree. It remains to identify
this Gaussian integer.

\begin{proposition}[The Vieta jump in the Gaussian integers]\label{prop:vieta-jump}
Let $M=P/Q$ be a reduced fraction of $(0,1)$ with Farey triangle $(L,R,D)$.
Then
\begin{equation}\label{eq:vieta-product}
 m_L+i\,m_R=\overline{\alpha_M}\,\alpha_D,\qquad N(\alpha_M)=m_M,
\end{equation}
and $\alpha_M$ is the occurrence pair of $P/Q$ up to a unit and a conjugation:
if $(s,r)^{\mathsf T}=He$, then $\alpha_M=\varepsilon(s+ir)$ or
$\alpha_M=\varepsilon(s-ir)$ for a unit $\varepsilon$ of $\Z[i]$. In particular, when
$D$ lies in $(0,1)$, the cofactor $\alpha_D$ of $\overline{\alpha_M}$ in
$m_L+i\,m_R$ is the occurrence pair of $D$, up to a unit and a conjugation.
\end{proposition}

\begin{proof}
\Cref{lem:three-divisibilities} shows that \eqref{eq:vieta-jump} is a Gaussian
integer of norm $m_M$, and multiplying \eqref{eq:vieta-jump} by $\alpha_D$
gives the product in \eqref{eq:vieta-product}. The Gaussian integer
$\overline{\alpha_M}$ divides $m_M=\overline{\alpha_M}\alpha_M$ and
$m_L+i\,m_R=\overline{\alpha_M}\alpha_D$, hence every element of the ideal of
$\Z[i]$ generated by $m_M$ and $m_L+i\,m_R$. By \cref{prop:occurrence-in-cube},
this ideal is generated by $s+ir$ or by $s-ir$, so that $\overline{\alpha_M}$
divides one of them. The two sides of this divisibility have the same norm
$m_M$, so their quotient is a Gaussian integer of norm one, a unit:
$\overline{\alpha_M}=\varepsilon'(s\pm ir)$, and conjugation gives
$\alpha_M=\overline{\varepsilon'}(s\mp ir)$. The last assertion is the first
one applied to the fraction $D$, which has its own Farey triangle.
\end{proof}

Up to units and conjugation, the product form is classical in substance: it
follows from the identity $mm'=m_L^2+m_R^2$ and from the characteristic
numbers of Frobenius, both recalled after \cref{prop:occurrence-in-cube} and
applied to $m$ and to $m'=m_D$, together with unique factorization in $\Z[i]$
\cite[\S\S1--2]{Frobenius1913}. For $m_D$ one more step is needed, since
$(m_D,m_L,m_R)$ is not the triple of the Farey parents of $D$: the Farey
triangles with the vertex $D$ are joined by flips that keep $D$, each of which
replaces a triple $(m_D,x,y)$ by $(m_D,y,3m_Dy-x)$, and this leaves the residue
of $x/y$ modulo $m_D$ unchanged up to sign, because $3m_Dy-x\equiv-x$ and
$(x/y)^2\equiv-1\pmod{m_D}$ by the Markov equation. What
\cref{lem:vieta-step,lem:three-divisibilities} add is the exact orientation,
without units, along the whole tree, and it is this orientation that makes the
recursion possible.

The proposition turns the computation of occurrence pairs into a descent.
To obtain the occurrence pair of $P/Q$, one follows the path of $P/Q$ in the
Farey tree and keeps only the three Gaussian integers of the current
triangle; each step costs one exact division in $\Z[i]$.
\Cref{fig:vieta-tree} shows the first four generations. For the triangle
$L=1/2$, $R=2/3$, $D=1/1$, whose labels are $5$, $29$ and $2$, the rule
\eqref{eq:vieta-jump} gives
$$
 \overline{\alpha_{3/5}}=\frac{5+29\,i}{1+i}=\frac{(5+29\,i)(1-i)}2
 =\frac{34+24\,i}2=17+12\,i ,
$$
and indeed $17^2+12^2=433=m_{3/5}$. Along the fractions $P/(P+1)$, whose right
parent is $1/1$ and whose labels are Pell numbers, the rule gives
successively
\begin{align*}
 29+2i&=(5+12i)(1-2i),\\
 169+2i&=(29+12i)(5-2i),\\
 985+2i&=(29+70i)(5-12i),
\end{align*}
that is $169=5^2+12^2$, $985=29^2+12^2$ and $5741=29^2+70^2$, each second
factor being the Gaussian integer found two generations earlier. Along the
fractions $1/Q$ it recovers the identity $1+233i=(21+13i)(5+8i)$ of the
label $610$, with both factors multiplied by $-1$.

\begin{figure}[tb]
\centering
% Authors: Dr. Denys Dutykh (Mathematics Department, Khalifa University of Science and
% Technology, Abu Dhabi, UAE) and Prof. Laurent Vuillon (Univ. Savoie Mont Blanc, CNRS, LAMA)
% Generated by supplement/code/make_figures.py from exact data; do not edit by hand.
\begin{tikzpicture}[font=\scriptsize]
\node[draw=slate, line width=0.5pt, dotted, fill=white, rounded corners=1.5pt, inner sep=2pt, align=center] (f0x1) at (0,0) {$0/1$\\$1$\\$1$};
\node[draw=wine, line width=0.9pt, fill=white, rounded corners=1.5pt, inner sep=2pt, align=center] (f1x1) at (13.44,0) {$D=1/1$\\$2$\\$1+i$};
\node[draw=slate, line width=0.5pt, dotted, fill=white, rounded corners=1.5pt, inner sep=2pt, align=center] (f1x0) at (6.72,0) {$1/0$\\$1$\\$1$};
\node[draw=wine, line width=0.9pt, fill=white, rounded corners=1.5pt, inner sep=2pt, align=center] (f1x2) at (6.72,-1.42) {$L=1/2$\\$5$\\$1-2i$};
\node[draw=mist, line width=0.5pt, fill=white, rounded corners=1.5pt, inner sep=2pt, align=center] (f1x3) at (3.36,-2.84) {$1/3$\\$13$\\$3-2i$};
\node[draw=wine, line width=0.9pt, fill=white, rounded corners=1.5pt, inner sep=2pt, align=center] (f2x3) at (10.08,-2.84) {$R=2/3$\\$29$\\$5-2i$};
\node[draw=mist, line width=0.5pt, fill=white, rounded corners=1.5pt, inner sep=2pt, align=center] (f1x4) at (1.68,-4.26) {$1/4$\\$34$\\$-5-3i$};
\node[draw=mist, line width=0.5pt, fill=white, rounded corners=1.5pt, inner sep=2pt, align=center] (f2x5) at (5.04,-4.26) {$2/5$\\$194$\\$13-5i$};
\node[draw=navy, line width=1.1pt, fill=sand, rounded corners=1.5pt, inner sep=2pt, align=center] (f3x5) at (8.4,-4.26) {$M=3/5$\\$433$\\$17-12i$};
\node[draw=mist, line width=0.5pt, fill=white, rounded corners=1.5pt, inner sep=2pt, align=center] (f3x4) at (11.76,-4.26) {$3/4$\\$169$\\$5-12i$};
\node[draw=mist, line width=0.5pt, fill=white, rounded corners=1.5pt, inner sep=2pt, align=center] (f1x5) at (0.84,-5.68) {$1/5$\\$89$\\$-5-8i$};
\node[draw=mist, line width=0.5pt, fill=white, rounded corners=1.5pt, inner sep=2pt, align=center] (f2x7) at (2.52,-5.68) {$2/7$\\$1325$\\$34-13i$};
\node[draw=mist, line width=0.5pt, fill=white, rounded corners=1.5pt, inner sep=2pt, align=center] (f3x8) at (4.2,-5.68) {$3/8$\\$7561$\\$-75-44i$};
\node[draw=mist, line width=0.5pt, fill=white, rounded corners=1.5pt, inner sep=2pt, align=center] (f3x7) at (5.88,-5.68) {$3/7$\\$2897$\\$44-31i$};
\node[draw=mist, line width=0.5pt, fill=white, rounded corners=1.5pt, inner sep=2pt, align=center] (f4x7) at (7.56,-5.68) {$4/7$\\$6466$\\$-29-75i$};
\node[draw=mist, line width=0.5pt, fill=white, rounded corners=1.5pt, inner sep=2pt, align=center] (f5x8) at (9.24,-5.68) {$5/8$\\$37666$\\$75-179i$};
\node[draw=mist, line width=0.5pt, fill=white, rounded corners=1.5pt, inner sep=2pt, align=center] (f5x7) at (10.92,-5.68) {$5/7$\\$14701$\\$99-70i$};
\node[draw=mist, line width=0.5pt, fill=white, rounded corners=1.5pt, inner sep=2pt, align=center] (f4x5) at (12.6,-5.68) {$4/5$\\$985$\\$29-12i$};
\begin{scope}[on background layer]
\draw[mist, line width=0.6pt] (f0x1) -- (f1x2);
\draw[mist, line width=0.6pt] (f1x2) -- (f1x3);
\draw[mist, line width=0.6pt] (f1x3) -- (f1x4);
\draw[mist, line width=0.6pt] (f1x3) -- (f2x5);
\draw[mist, line width=0.6pt] (f1x4) -- (f1x5);
\draw[mist, line width=0.6pt] (f1x4) -- (f2x7);
\draw[mist, line width=0.6pt] (f2x3) -- (f3x4);
\draw[mist, line width=0.6pt] (f2x5) -- (f3x7);
\draw[mist, line width=0.6pt] (f2x5) -- (f3x8);
\draw[mist, line width=0.6pt] (f3x4) -- (f4x5);
\draw[mist, line width=0.6pt] (f3x4) -- (f5x7);
\draw[mist, line width=0.6pt] (f3x5) -- (f4x7);
\draw[mist, line width=0.6pt] (f3x5) -- (f5x8);
\draw[wine, line width=0.8pt, dash pattern=on 3pt off 2pt] (f1x2) -- (f2x3);
\draw[wine, line width=0.8pt, dash pattern=on 3pt off 2pt] (f2x3) -- (f1x1);
\draw[wine, line width=0.8pt, dash pattern=on 3pt off 2pt] (f1x1) -- (f1x2);
\draw[navy, line width=1.0pt] (f1x2) -- (f3x5);
\draw[navy, line width=1.0pt] (f2x3) -- (f3x5);
\end{scope}
\end{tikzpicture}

\caption{The first four generations of the Farey tree, with the label $m$
and the Gaussian integer $\alpha$ of each fraction, computed by the rule
\eqref{eq:vieta-jump} alone from the values at the three starting fractions
$0/1$, $1/1$ and $1/0$, drawn at the top. Each fraction is joined to the younger of its
two Farey parents, and $1/2$ to both. In wine and dashed, the Farey triangle
$L=1/2$, $R=2/3$, $D=1/1$; in navy, the edges from its parents to the
mediant $M=3/5$, whose Gaussian integer $17-12i$ is the conjugate of the
quotient of $5+29i$ by~$1+i$.}
\label{fig:vieta-tree}
\end{figure}

\begin{remark}[The pair determines the occurrence]\label{rem:pair-determines}
Two distinct occurrences of one label have distinct occurrence pairs, even
up to units and conjugation: equal pairs have equal costs, and
\cref{lem:fixed-cost-injectivity} applies. \Cref{prop:vieta-jump} computes the
occurrence pair from the Farey triangle, that is, from the position of the
occurrence in the Markov tree, and not from its label. A rule that produced
the occurrence pair from the label alone, correctly for every occurrence,
would therefore prove the uniqueness conjecture; \cref{sec:directives}
states a conjecture of this kind.
\end{remark}

The finite part of the proof, to which we now turn, needs no recursion: it
computes the whole fiber of each label, in which the occurrence pair is the
one given by the Gram factor of \cref{prop:occurrence-cost}.

\subsection{An elementary primality criterion}

The factorizations used below rest on the following classical criterion,
which we state for completeness. Suppose $n>2$, the complete factorization of $n-1$ into proved
primes is supplied, and an integer $a$ satisfies
\begin{equation}\label{eq:full-lucas}
 a^{n-1}\equiv1\pmod n,\qquad
 \gcd\bigl(a^{(n-1)/r}-1,n\bigr)=1
 \quad\text{for every prime }r\mid n-1.
\end{equation}
Then $n$ is prime. Indeed, for any prime divisor $s$ of $n$, these
conditions make the order of $a$ modulo $s$ exactly $n-1$. Since that
order divides $s-1$, one has $s\geqslant n$, hence $s=n$.

Applied recursively, with the primes of $n-1$ themselves established the
same way, the criterion terminates at the prime two and establishes
primality by integer arithmetic alone: no probabilistic test and no search
for factors enters the verification, only the confirmation of factors
already proposed.

Each label is computed twice and the two values are compared: once from
the floor increments of the lower Christoffel word, and once from the Cohn
matrix built by Stern--Brocot mediants from the two boundary matrices. The
occurrence pair is computed from the central Gram factor $H$ of
\cref{prop:occurrence-cost}. For each label one then enumerates the
$2^t$ sign choices in \eqref{eq:gaussian-sign-cube} and passes to the
quotient by conjugation; that every normalized pair is obtained exactly
twice, as \cref{prop:gaussian-fibre} predicts, is itself a check on the
enumeration. The Euclidean algorithm then supplies the quotient sequence
and the cost of every pair.

\subsection{The finite windows}

The finite part of the proof concerns the reduced fractions listed in
\cref{tab:gaussian-windows}: in row $P$, for $8\leqslant P\leqslant13$, every
denominator from $P+1$ up to the last denominator before the junction
$Q_P$ of \cref{sec:assembly}. For each of these $248$ labels, the prime
factorization, the primality of its factors by the recursive use of
\eqref{eq:full-lucas}, the splitting of each prime and the costs of the
resulting pairs are finitely many assertions about integers. They are
recorded as part~\ref{it:fv-windows} of the finite verifications,
\cref{prop:finite-verifications}, and \cref{prop:finite-windows} draws
their consequence: in every finite window the occurrence pair is the unique
representation of minimum cost.

\begin{table}[tb]
\centering
\caption{The complete finite part of the minimality proof. Each row includes
all $Q$ in the displayed interval with $\gcd(P,Q)=1$. The pair count includes
every normalized primitive positive representation of every source, and the
last column gives the largest individual fiber. The later sections cover the
unbounded denominators beyond these windows.}
\label{tab:gaussian-windows}
\small
\setlength{\tabcolsep}{11pt}
\renewcommand{\arraystretch}{1.12}
\begin{tabular}{@{}crrrr@{}}
\toprule
$P$ & Denominators $Q$ & Sources & Pairs & Largest fiber \\
\midrule
8 & $9$--$28$ & 10 & 22 & 4 \\
9 & $10$--$34$ & 17 & 56 & 8 \\
10 & $11$--$38$ & 11 & 33 & 8 \\
11 & $12$--$75$ & 59 & 324 & 32 \\
12 & $13$--$100$ & 29 & 211 & 32 \\
13 & $14$--$145$ & 122 & 992 & 64 \\
\midrule
Total & & 248 & 1638 & 64 \\
\bottomrule
\end{tabular}
\end{table}

The statement on the finite windows is stronger than weak minimality, since
it also asserts uniqueness of the minimizing representation. The unbounded
part requires only weak minimality. It uses the finite windows as exact
boundary data and separately excludes every lower-cost competitor on each
infinite tail.

% Authors: Dr. Denys Dutykh (Mathematics Department, Khalifa University of Science and Technology,
% Abu Dhabi, UAE) and Prof. Laurent Vuillon (Univ. Savoie Mont Blanc, CNRS, LAMA, Chambery, France)
\section{Path envelopes and the first seven rows}\label{sec:early-rows}

The Euclidean bijection turns a cost bound into a finite binary tree.
Enumerating that tree at every denominator would obscure the reason that
an entire row can be treated. Instead, we compare its branches with a
small number of explicit paths. Fibonacci growth controls the cost,
and deviations from the Fibonacci path determine the finer norm bands.

\subsection{Sorted paths and a stable order}

On sorted positive pairs define
\begin{equation}\label{eq:path-maps}
 \mathsf A(x,y)=(x,x+y),\qquad
 \mathsf B(x,y)=(y,x+y).
\end{equation}
The first move from $(1,1)$ is symmetric and gives $(1,2)$. Thereafter
these two moves are distinct, and every sorted primitive pair of cost
$c\geqslant1$ is obtained by a unique sequence of $c-1$ moves.
Indeed, from $0<x<y$ its sorted predecessor is the sorted pair formed
from $x,y-x$. If $y>2x$, the last move was $\mathsf A$; if $y<2x$,
it was $\mathsf B$. The equality $y=2x$ is primitive only at $(1,2)$.
This proves the asserted uniqueness directly, including its initial
boundary.

A move $\mathsf A$ after the fixed first move is called a \emph{defect}.
Write $v_c(D)$ for the endpoint whose defect costs form
$D\subseteq\{2,\ldots,c\}$. Thus a defect at cost $j$ is the move
arriving at depth $j$. Products written as words in the moves are read
chronologically from left to right; an expression $\mathsf B^nz$,
with a vector $z$, means the $n$-fold application of $\mathsf B$ to that vector.

Let $F_0=0,F_1=1,F_{n+1}=F_n+F_{n-1}$, and use
$F_{-n}=(-1)^{n+1}F_n$ when a displayed identity reaches a negative
index. Put
$$
 N(x,y)=x^2+y^2,\qquad
 \Lambda(x,y)=x+\varphi y,\qquad
 \varphi=\frac{1+\sqrt5}{2}.
$$

\begin{lemma}[Maximum, replacement, and propagation]\label{lem:path-orders}
At exact cost $c$, the unique maximum norm is
$F_{2c+3}$, attained at $(F_{c+1},F_{c+2})$.
Replacing any defect by $\mathsf B$ strictly increases the final norm.
Finally, simultaneous weak inequalities in $N$ and $\Lambda$ are
preserved by every common $\mathsf B$-tail.
\end{lemma}

\begin{proof}
Coordinate induction in \eqref{eq:path-maps}, starting from $(1,1)$,
gives $x\leqslant F_{c+1}$ and $y\leqslant F_{c+2}$.
The identity $F_{c+1}^2+F_{c+2}^2=F_{2c+3}$ proves the norm bound.
After the first move, $x<y$, and
$$
 \mathsf B(x,y)-\mathsf A(x,y)=(y-x,0).
$$
This difference is nonnegative and nonzero. Each subsequent matrix
is nonnegative and invertible, so it preserves a nonnegative nonzero
difference. Since both endpoint coordinates are positive, the norm
strictly increases. Replacing every defect therefore proves uniqueness
of the Fibonacci maximum as well.
For propagation, direct calculation gives
\begin{equation}\label{eq:perron-propagation}
 \Lambda(\mathsf Bz)=\varphi\Lambda(z),\qquad
 N(\mathsf Bz)=\varphi^{-2}N(z)+\varphi^{-1}\Lambda(z)^2.
\end{equation}
Both coefficients are positive, and $\Lambda$ is positive on the
states under consideration.
\end{proof}

\begin{lemma}[One-defect minimum]\label{lem:one-defect-minimum}
For $c\geqslant2$, the earliest-defect endpoint
$e_c=\mathsf B^{c-2}(1,3)$ minimizes both $N$ and $\Lambda$ among
one-defect paths, and
\begin{equation}\label{eq:one-defect-minimum}
 N(e_c)=5F_{2c-1}.
\end{equation}
\end{lemma}

\begin{proof}
The assertion starts at $c=2$. An inherited one-defect path is covered
by \eqref{eq:perron-propagation}. The only new path creates its defect
at the last move and has endpoint $b_c=(F_c,F_{c+2})$. Fibonacci
addition gives
$$
 b_c-e_c=F_{c-2}(-1,1),\qquad
 N(b_c)-N(e_c)=4F_{c-2}F_{c-1}.
$$
Its $\Lambda$-difference is $(\varphi-1)F_{c-2}$, so induction closes.
The equality for $N(e_c)$ follows by substitution from
$e_c=(F_c+F_{c-2},F_{c+1}+F_{c-1})$, or from the recurrence
$N(\mathsf B^2z)=3N(\mathsf Bz)-N(z)$ and its first two values.
\end{proof}

\subsection{The second envelope}

\begin{lemma}[The second norm envelope]\label{lem:second-envelope}
At cost $c\geqslant5$, every non-Fibonacci endpoint satisfies
\begin{equation}\label{eq:second-envelope}
 N(x,y)\leqslant E_c:=F_c^2+F_{c+2}^2.
\end{equation}
Equality is attained at $(F_c,F_{c+2})$, and also at $(12,19)$ when
$c=6$; there are no other equality cases.
\end{lemma}

\begin{proof}
Let $j$ be the last defect cost and $r=c-j$. Before that defect,
the coordinate bound in Lemma~\ref{lem:path-orders} gives the upper
state $(F_j,F_{j+1})$. Its final endpoint is
$$
 X_{c,j}=F_c+F_rF_j,\qquad
 Y_{c,j}=F_{c+2}-F_rF_{j-1}.
$$
Any earlier defect makes at least one final coordinate strictly smaller.
Expansion and Cassini's identity give the exact deficit
\begin{align}\label{eq:second-envelope-deficit}
 E_c-X_{c,j}^2-Y_{c,j}^2
 =F_r\bigl[
 2(F_jF_{j-1}+(-1)^j)F_{r+1}
 -F_{j-2}F_{j+1}F_r\bigr].
\end{align}
If $r=0$, this is zero. For $r\geqslant1$ and $j\geqslant4$,
replace $F_{r+1}$ by its lower bound $F_r$; the remaining coefficient is
$$
 2F_{j-1}^2-F_{j-2}^2+2(-1)^j>0.
$$
The case $j=2$ is positive directly. When $j=3$, the bracket in
\eqref{eq:second-envelope-deficit} is $2F_{r+1}-3F_r$.
Here $r\geqslant2$; its values at $r=2,3$ are $1,0$, and for
$r\geqslant4$ it is $F_{r-3}>0$. The zero cases give exactly the
two stated endpoints.
\end{proof}

It follows that an occurrence of cost $L\geqslant6$ and label $m$ is
weakly Euclidean-minimal whenever
\begin{equation}\label{eq:one-hole}
 E_{L-1}<m,\qquad m\ne F_{2L+1}.
\end{equation}
At cost $L-1$ only the Fibonacci maximum can lie above $E_{L-1}$.
Every smaller cost has norm at most
$F_{2L-1}<E_{L-1}$. Thus \eqref{eq:one-hole} excludes every
smaller-cost representation.

\Cref{fig:norm-cloud} shows the situation that this criterion exploits.
Normalized by the Fibonacci maximum, the norms of all pairs of one cost
lie between zero and the second envelope, near $0.829$, except for the
isolated value $1$, so that a hole separates the two. The normalized source labels
of a row tend to the constant $C_P$ computed in \cref{sec:row-labels} below.
For $P\leqslant6$ these limits exceed $1$, and $C_7$ falls in the hole; for
$P\geqslant8$ they fall inside the cloud, where a finer separation is needed,
and that is the business of \cref{sec:tails}.

\begin{figure}[tb]
\centering
% Authors: Dr. Denys Dutykh (Mathematics Department, Khalifa University of Science and
% Technology, Abu Dhabi, UAE) and Prof. Laurent Vuillon (Univ. Savoie Mont Blanc, CNRS, LAMA)
% Generated by supplement/code/make_figures.py from exact data; do not edit by hand.
\begin{tikzpicture}[font=\footnotesize, >={Stealth[length=4pt]}]
\fill[sand] (9.038,0) rectangle (10.9,3.45);
\fill[navy] (0,0) rectangle (0.054,3.277);
\fill[navy] (0.054,0) rectangle (0.109,3.3);
\fill[navy] (0.109,0) rectangle (0.164,3.282);
\fill[navy] (0.164,0) rectangle (0.218,3.259);
\fill[navy] (0.218,0) rectangle (0.273,3.233);
\fill[navy] (0.273,0) rectangle (0.327,3.207);
\fill[navy] (0.327,0) rectangle (0.382,3.186);
\fill[navy] (0.382,0) rectangle (0.436,3.157);
\fill[navy] (0.436,0) rectangle (0.49,3.135);
\fill[navy] (0.49,0) rectangle (0.545,3.11);
\fill[navy] (0.545,0) rectangle (0.6,3.089);
\fill[navy] (0.6,0) rectangle (0.654,3.061);
\fill[navy] (0.654,0) rectangle (0.709,3.04);
\fill[navy] (0.709,0) rectangle (0.763,3.019);
\fill[navy] (0.763,0) rectangle (0.818,2.996);
\fill[navy] (0.818,0) rectangle (0.872,2.972);
\fill[navy] (0.872,0) rectangle (0.927,2.947);
\fill[navy] (0.927,0) rectangle (0.981,2.932);
\fill[navy] (0.981,0) rectangle (1.036,2.903);
\fill[navy] (1.036,0) rectangle (1.09,2.889);
\fill[navy] (1.09,0) rectangle (1.145,2.864);
\fill[navy] (1.145,0) rectangle (1.199,2.847);
\fill[navy] (1.199,0) rectangle (1.254,2.819);
\fill[navy] (1.254,0) rectangle (1.308,2.797);
\fill[navy] (1.308,0) rectangle (1.363,2.779);
\fill[navy] (1.363,0) rectangle (1.417,2.756);
\fill[navy] (1.417,0) rectangle (1.472,2.737);
\fill[navy] (1.472,0) rectangle (1.526,2.714);
\fill[navy] (1.526,0) rectangle (1.581,2.695);
\fill[navy] (1.581,0) rectangle (1.635,2.676);
\fill[navy] (1.635,0) rectangle (1.69,2.651);
\fill[navy] (1.69,0) rectangle (1.744,2.642);
\fill[navy] (1.744,0) rectangle (1.799,2.614);
\fill[navy] (1.799,0) rectangle (1.853,2.59);
\fill[navy] (1.853,0) rectangle (1.907,2.568);
\fill[navy] (1.907,0) rectangle (1.962,2.553);
\fill[navy] (1.962,0) rectangle (2.017,2.536);
\fill[navy] (2.017,0) rectangle (2.071,2.515);
\fill[navy] (2.071,0) rectangle (2.126,2.486);
\fill[navy] (2.126,0) rectangle (2.18,2.466);
\fill[navy] (2.18,0) rectangle (2.235,2.451);
\fill[navy] (2.235,0) rectangle (2.289,2.44);
\fill[navy] (2.289,0) rectangle (2.344,2.399);
\fill[navy] (2.344,0) rectangle (2.398,2.381);
\fill[navy] (2.398,0) rectangle (2.453,2.381);
\fill[navy] (2.453,0) rectangle (2.507,2.353);
\fill[navy] (2.507,0) rectangle (2.562,2.323);
\fill[navy] (2.562,0) rectangle (2.616,2.31);
\fill[navy] (2.616,0) rectangle (2.671,2.289);
\fill[navy] (2.671,0) rectangle (2.725,2.274);
\fill[navy] (2.725,0) rectangle (2.78,2.243);
\fill[navy] (2.78,0) rectangle (2.834,2.236);
\fill[navy] (2.834,0) rectangle (2.889,2.203);
\fill[navy] (2.889,0) rectangle (2.943,2.191);
\fill[navy] (2.943,0) rectangle (2.998,2.178);
\fill[navy] (2.998,0) rectangle (3.052,2.148);
\fill[navy] (3.052,0) rectangle (3.107,2.117);
\fill[navy] (3.107,0) rectangle (3.161,2.109);
\fill[navy] (3.161,0) rectangle (3.216,2.089);
\fill[navy] (3.216,0) rectangle (3.27,2.063);
\fill[navy] (3.27,0) rectangle (3.325,2.047);
\fill[navy] (3.325,0) rectangle (3.379,2.022);
\fill[navy] (3.379,0) rectangle (3.433,2.011);
\fill[navy] (3.433,0) rectangle (3.488,1.969);
\fill[navy] (3.488,0) rectangle (3.543,1.988);
\fill[navy] (3.543,0) rectangle (3.597,1.938);
\fill[navy] (3.597,0) rectangle (3.652,1.943);
\fill[navy] (3.652,0) rectangle (3.706,1.903);
\fill[navy] (3.706,0) rectangle (3.76,1.896);
\fill[navy] (3.76,0) rectangle (3.815,1.859);
\fill[navy] (3.815,0) rectangle (3.869,1.827);
\fill[navy] (3.869,0) rectangle (3.924,1.771);
\fill[navy] (3.924,0) rectangle (3.978,1.8);
\fill[navy] (3.978,0) rectangle (4.033,1.8);
\fill[navy] (4.033,0) rectangle (4.088,1.803);
\fill[navy] (4.088,0) rectangle (4.142,1.766);
\fill[navy] (4.142,0) rectangle (4.197,1.703);
\fill[navy] (4.197,0) rectangle (4.251,1.683);
\fill[navy] (4.251,0) rectangle (4.306,1.688);
\fill[navy] (4.306,0) rectangle (4.36,1.575);
\fill[navy] (4.36,0) rectangle (4.415,1.678);
\fill[navy] (4.415,0) rectangle (4.469,1.531);
\fill[navy] (4.469,0) rectangle (4.524,1.6);
\fill[navy] (4.524,0) rectangle (4.578,1.598);
\fill[navy] (4.578,0) rectangle (4.633,1.523);
\fill[navy] (4.633,0) rectangle (4.687,1.531);
\fill[navy] (4.687,0) rectangle (4.742,1.606);
\fill[navy] (4.742,0) rectangle (4.796,1.464);
\fill[navy] (4.796,0) rectangle (4.851,1.523);
\fill[navy] (4.851,0) rectangle (4.905,1.482);
\fill[navy] (4.905,0) rectangle (4.96,1.421);
\fill[navy] (4.96,0) rectangle (5.014,1.485);
\fill[navy] (5.014,0) rectangle (5.069,1.28);
\fill[navy] (5.069,0) rectangle (5.123,1.491);
\fill[navy] (5.123,0) rectangle (5.178,1.358);
\fill[navy] (5.178,0) rectangle (5.232,1.055);
\fill[navy] (5.232,0) rectangle (5.287,1.219);
\fill[navy] (5.287,0) rectangle (5.341,0.964);
\fill[navy] (5.341,0) rectangle (5.396,0.948);
\fill[navy] (5.396,0) rectangle (5.45,1.244);
\fill[navy] (5.45,0) rectangle (5.505,1.268);
\fill[navy] (5.505,0) rectangle (5.559,1.134);
\fill[navy] (5.559,0) rectangle (5.614,1.225);
\fill[navy] (5.614,0) rectangle (5.668,1.331);
\fill[navy] (5.668,0) rectangle (5.723,1.183);
\fill[navy] (5.723,0) rectangle (5.777,1.076);
\fill[navy] (5.777,0) rectangle (5.832,1.143);
\fill[navy] (5.832,0) rectangle (5.886,1.205);
\fill[navy] (5.886,0) rectangle (5.941,0.978);
\fill[navy] (5.941,0) rectangle (5.995,1.134);
\fill[navy] (5.995,0) rectangle (6.05,1.311);
\fill[navy] (6.05,0) rectangle (6.104,0.812);
\fill[navy] (6.104,0) rectangle (6.158,0.786);
\fill[navy] (6.158,0) rectangle (6.213,1.031);
\fill[navy] (6.213,0) rectangle (6.268,0.729);
\fill[navy] (6.268,0) rectangle (6.322,0.219);
\fill[navy] (6.322,0) rectangle (6.377,0.992);
\fill[navy] (6.377,0) rectangle (6.431,0.509);
\fill[navy] (6.431,0) rectangle (6.486,0.219);
\fill[navy] (6.486,0) rectangle (6.54,0.219);
\fill[navy] (6.54,0) rectangle (6.595,0.348);
\fill[navy] (6.649,0) rectangle (6.704,0.219);
\fill[navy] (6.758,0) rectangle (6.812,0.759);
\fill[navy] (6.812,0) rectangle (6.867,0.439);
\fill[navy] (6.867,0) rectangle (6.921,0.759);
\fill[navy] (6.921,0) rectangle (6.976,0.992);
\fill[navy] (6.976,0) rectangle (7.031,1.125);
\fill[navy] (7.031,0) rectangle (7.085,0.835);
\fill[navy] (7.085,0) rectangle (7.14,0.348);
\fill[navy] (7.14,0) rectangle (7.194,0.219);
\fill[navy] (7.194,0) rectangle (7.249,0.964);
\fill[navy] (7.249,0) rectangle (7.303,0.509);
\fill[navy] (7.412,0) rectangle (7.467,0.219);
\fill[navy] (7.467,0) rectangle (7.521,0.812);
\fill[navy] (7.521,0) rectangle (7.575,0.348);
\fill[navy] (7.63,0) rectangle (7.684,0.219);
\fill[navy] (7.684,0) rectangle (7.739,0.219);
\fill[navy] (7.902,0) rectangle (7.957,0.219);
\fill[navy] (8.556,0) rectangle (8.611,0.219);
\fill[navy] (8.666,0) rectangle (8.72,0.567);
\fill[navy] (8.72,0) rectangle (8.775,0.759);
\fill[navy] (8.992,0) rectangle (9.047,0.348);
\draw[wine, line width=1.4pt] (10.9,0) -- (10.9,0.219);
\fill[wine] (10.9,0.219) circle (1.6pt);
\node[wine, align=center, font=\scriptsize] (fib) at (9.969,1.386) {Fibonacci\\ maximum};
\draw[slate, line width=0.4pt] (fib.south) -- (10.86,0.279);
\node[slate, align=center, font=\scriptsize] at (9.969,2.805) {hole};
\draw[slate, densely dashed, line width=0.6pt] (7.951,0) -- (7.951,3.45);
\node[slate, anchor=south east, align=right, font=\scriptsize] at (7.911,3.47) {one-defect\\ minimum};
\node[slate, anchor=south west, align=left, font=\scriptsize] at (9.078,3.47) {second\\ envelope};
\draw[ink, line width=0.5pt] (0,0) -- (10.9,0);
\draw[ink, line width=0.5pt] (0,0) -- (0,-0.08);
\node[ink, below] at (0,-0.08) {$0.0$};
\draw[ink, line width=0.5pt] (1.09,0) -- (1.09,-0.08);
\node[ink, below] at (1.09,-0.08) {$0.1$};
\draw[ink, line width=0.5pt] (2.18,0) -- (2.18,-0.08);
\node[ink, below] at (2.18,-0.08) {$0.2$};
\draw[ink, line width=0.5pt] (3.27,0) -- (3.27,-0.08);
\node[ink, below] at (3.27,-0.08) {$0.3$};
\draw[ink, line width=0.5pt] (4.36,0) -- (4.36,-0.08);
\node[ink, below] at (4.36,-0.08) {$0.4$};
\draw[ink, line width=0.5pt] (5.45,0) -- (5.45,-0.08);
\node[ink, below] at (5.45,-0.08) {$0.5$};
\draw[ink, line width=0.5pt] (6.54,0) -- (6.54,-0.08);
\node[ink, below] at (6.54,-0.08) {$0.6$};
\draw[ink, line width=0.5pt] (7.63,0) -- (7.63,-0.08);
\node[ink, below] at (7.63,-0.08) {$0.7$};
\draw[ink, line width=0.5pt] (8.72,0) -- (8.72,-0.08);
\node[ink, below] at (8.72,-0.08) {$0.8$};
\draw[ink, line width=0.5pt] (9.81,0) -- (9.81,-0.08);
\node[ink, below] at (9.81,-0.08) {$0.9$};
\draw[ink, line width=0.5pt] (10.9,0) -- (10.9,-0.08);
\node[ink, below] at (10.9,-0.08) {$1.0$};
\node[ink] at (5.45,-0.75) {normalized norm $N/F_{43}$ of the pairs of cost $20$};
\draw[mist, line width=0.5pt] (0,-1.35) -- (10.9,-1.35);
\draw[teal, line width=1.1pt] (9.275,-1.47) -- (9.275,-1.23);
\node[teal, below, font=\scriptsize] at (9.275,-1.47) {$C_{7}$};
\draw[wine, line width=1.1pt] (7.69,-1.47) -- (7.69,-1.23);
\node[wine, below, font=\scriptsize] at (7.69,-1.47) {$C_{8}$};
\draw[wine, line width=1.1pt] (6.377,-1.47) -- (6.377,-1.23);
\node[wine, below, font=\scriptsize] at (6.377,-1.47) {$C_{9}$};
\draw[wine, line width=1.1pt] (5.287,-1.47) -- (5.287,-1.23);
\node[wine, below, font=\scriptsize] at (5.287,-1.47) {$C_{10}$};
\draw[wine, line width=1.1pt] (4.384,-1.47) -- (4.384,-1.23);
\node[wine, below, font=\scriptsize] at (4.384,-1.47) {$C_{11}$};
\draw[wine, line width=1.1pt] (3.635,-1.47) -- (3.635,-1.23);
\node[wine, below, font=\scriptsize] at (3.635,-1.47) {$C_{12}$};
\draw[wine, line width=1.1pt] (3.014,-1.47) -- (3.014,-1.23);
\node[wine, below, font=\scriptsize] at (3.014,-1.47) {$C_{13}$};
\draw[teal, ->, line width=0.7pt] (10.95,-1.35) -- (11.65,-1.35);
\node[teal, right, font=\scriptsize, align=left] at (11.65,-1.35) {$C_2,\ldots,C_6$\\ above $1$};
\node[ink, left, font=\scriptsize] at (-0.05,-1.35) {sources};
\node[ink, rotate=90, font=\scriptsize] at (-0.45,1.65) {pairs (log scale)};
\end{tikzpicture}

\caption{The normalized norms $N/F_{43}$ of all $2^{19}$ sorted primitive pairs
of cost $20$, counted in bins of width $1/200$ on a logarithmic scale (navy).
The Fibonacci maximum is the isolated value $1$; every other pair lies at or
below the second envelope of \cref{lem:second-envelope}, and every
one-defect pair at or above the minimum of \cref{lem:one-defect-minimum}
(dashed). The shaded hole contains no normalized norm. The ticks below mark
the limits $C_P$ of the normalized source labels of the rows: $C_7$ lies in
the hole, $C_8,\ldots,C_{13}$ inside the cloud, and $C_2,\ldots,C_6$ above $1$.}
\label{fig:norm-cloud}
\end{figure}

\subsection{The labels of a row}\label{sec:row-labels}

The source labels of the later rows must be placed in narrow normalized
intervals, uniformly in the denominator. This subsection computes the
constant $C_P$ that they approach and bounds the approach on both sides,
with explicit constants. The idea is to diagonalize the Cohn matrix
$\mu(a)$. In its eigenbasis the second Cohn matrix $\mu(b)$ becomes a
matrix with positive entries, so that the trace formula
\eqref{eq:cohn-source} turns the label into a sum of positive terms; the
term that follows the dominant eigenvector at every letter gives $C_P$,
and the others are corrections whose size is governed by the runs of the
letter $a$.

Put
$$
 \lambda=\frac3{\sqrt5\,\varphi}=\frac{15-3\sqrt5}{10},\qquad
 C_P=\varphi^2\lambda^{P-1},\qquad
 \widehat B=\frac1{\sqrt5}\begin{pmatrix}3\varphi^3&2\\2&3\varphi^{-3}\end{pmatrix},
$$
and, for a reduced fraction $P/Q$, write its lower Christoffel word as
$$
 \gamma_{P/Q}=a^{u_1}b\,a^{u_2}b\cdots a^{u_P}b,\qquad u_k\geqslant0 .
$$
The normalized label of the occurrence is
\begin{equation}\label{eq:row-ratio}
 \rho_{P,Q}=\frac{m_{P/Q}}{F_{2(P+Q-2)+1}} ,
\end{equation}
the label divided by the Fibonacci maximum one cost below the occurrence
cost.

\begin{proposition}[The labels of a row]\label{prop:row-labels}
Let $0<P<Q$ be coprime and $n=\lfloor(Q-P)/P\rfloor$.
\begin{enumerate}[label=\textup{(\roman*)},ref=\textup{(\roman*)},leftmargin=*]
\item\label{it:row-runs} Every run satisfies $n\leqslant u_k\leqslant n+1$, and $u_1+\cdots+u_P=Q-P$.
\item\label{it:row-trace} The label is the trace
 $$
 m_{P/Q}=\frac13\tr\prod_{k=1}^{P}
 \begin{pmatrix}\varphi^{2u_k}&0\\0&\varphi^{-2u_k}\end{pmatrix}\widehat B ,
 $$
 the factors being multiplied in the order $k=1,\dots,P$.
\item\label{it:row-bounds} If $P\geqslant2$, then
 \begin{equation}\label{eq:row-source-limit}
 C_P<\rho_{P,Q}<C_P\bigl(1+\varphi^{-4n-6}\bigr)^P ,
 \end{equation}
 while $\rho_{1,Q}<C_1$ for the first row. In particular
 $\rho_{P,Q}$ tends to $C_P$ when $Q$ grows with $P$ fixed.
\end{enumerate}
\end{proposition}

\begin{proof}
\emph{The runs.} The $k$th letter $b$ of $\gamma_{P/Q}$ stands at position
$\lceil kQ/P\rceil$, as recalled after \cref{prop:occurrence-cost}, so that
$u_1=\lceil Q/P\rceil-1$ and
$u_k=\lceil kQ/P\rceil-\lceil(k-1)Q/P\rceil-1$ for $k\geqslant2$. For every real
$x$ and $y$ the difference $\lceil x+y\rceil-\lceil x\rceil$ is $\lfloor y\rfloor$
or $\lceil y\rceil$; with $y=Q/P$, every $u_k$ is therefore
$\lfloor Q/P\rfloor-1=n$ or $\lceil Q/P\rceil-1\leqslant n+1$. The word has $Q$
letters, of which $P$ are $b$, so the runs add up to $Q-P$.

\emph{The eigenbasis.} Let $V$ be the matrix with columns
$(\varphi,1)^{\mathsf T}$ and $(1,-\varphi)^{\mathsf T}$. Its columns are
orthogonal, each of squared length $\varphi^2+1=\sqrt5\,\varphi$, so that
$O=V/(5^{1/4}\varphi^{1/2})$ is an orthogonal matrix. Since
$\varphi^2=\varphi+1$,
$$
 \mu(a)\binom\varphi1=\binom{2\varphi+1}{\varphi+1}=\varphi^2\binom\varphi1,
 \qquad
 \mu(a)\binom1{-\varphi}=\binom{2-\varphi}{1-\varphi}=\varphi^{-2}\binom1{-\varphi},
$$
because $\varphi^3=2\varphi+1$ and $\varphi^{-2}=2-\varphi$. Thus
$O^{\mathsf T}\mu(a)O=\operatorname{diag}(\varphi^2,\varphi^{-2})$. For the second
matrix,
$$
 V^{\mathsf T}\mu(b)V=
 \begin{pmatrix}5\varphi^2+4\varphi+1&2\varphi\\2\varphi&6-3\varphi\end{pmatrix}
 =\begin{pmatrix}3\varphi^4&2\varphi\\2\varphi&3\varphi^{-2}\end{pmatrix},
$$
using $5\varphi^2+4\varphi+1=9\varphi+6=3\varphi^4$ and $6-3\varphi=3\varphi^{-2}$;
dividing by $\sqrt5\,\varphi$ gives $O^{\mathsf T}\mu(b)O=\widehat B$. All four
entries of $\widehat B$ are positive, $\tr\widehat B=6$ and $\det\widehat B=1$.
Since $\mu(\gamma_{P/Q})=\prod_k\mu(a)^{u_k}\mu(b)$ and
$O^{\mathsf T}\mu(a)^{u}O=\operatorname{diag}(\varphi^{2u},\varphi^{-2u})$, conjugating
by $O$ and taking traces turns \eqref{eq:cohn-source} into \ref{it:row-trace}.

\emph{Normalization.} Write $\beta=\widehat B_{12}/\widehat B_{11}=2/(3\varphi^3)$,
$\delta=\widehat B_{22}/\widehat B_{11}=\varphi^{-6}$ and $t_k=\varphi^{-4u_k}$. Taking
the factor $\varphi^{2u_k}\widehat B_{11}$ out of the $k$th matrix of \ref{it:row-trace} leaves
$$
 M_k=\begin{pmatrix}1&\beta\\t_k\beta&t_k\delta\end{pmatrix},\qquad
 m_{P/Q}=\frac13\,\varphi^{2(Q-P)}\Bigl(\frac{3\varphi^3}{\sqrt5}\Bigr)^P\Theta,
 \qquad \Theta=\tr\prod_{k=1}^PM_k ,
$$
by \ref{it:row-runs}. At the odd index $2(P+Q)-3$, Binet's formula reads
$F_{2(P+Q)-3}=\bigl(\varphi^{2(P+Q)-3}+\varphi^{3-2(P+Q)}\bigr)/\sqrt5$. Dividing,
and collecting the powers of $3$, $\sqrt5$ and $\varphi$,
$$
 \rho_{P,Q}
 =\frac{3^{P-1}\varphi^{3-P}}{5^{(P-1)/2}}\cdot\frac{\Theta}{1+\eta}
 =C_P\,\frac{\Theta}{1+\eta},\qquad \eta=\varphi^{6-4(P+Q)} .
$$

\emph{The lower bound.} Expand the trace of the product of the
nonnegative matrices $M_k$ as a sum over closed walks
$i_1,\dots,i_P,i_{P+1}=i_1$ in $\{1,2\}$, the walk contributing
$\prod_k(M_k)_{i_ki_{k+1}}$. Every contribution is nonnegative. The walk that
stays at $1$ contributes $1$. When $P\geqslant2$, the walk that visits $2$
only at the step $j=1$ contributes $(M_P)_{12}(M_1)_{21}=\beta^2t_1$. Hence
$\Theta\geqslant1+\beta^2t_1\geqslant1+\beta^2\varphi^{-4(Q-P)}$, since $u_1\leqslant Q-P$.
The ratio of $\beta^2\varphi^{-4(Q-P)}$ to $\eta$ is
$$
 \frac{4}{9\varphi^6}\,\varphi^{-4Q+4P}\,\varphi^{4P+4Q-6}
 =\frac49\,\varphi^{8P-12}\geqslant\frac49\,\varphi^4>1 ,
$$
because $P\geqslant2$ and $\varphi^4=(7+3\sqrt5)/2>9/4$. Therefore $\Theta>1+\eta$
and $\rho_{P,Q}>C_P$.

\emph{The upper bound.} By \ref{it:row-runs}, $0\leqslant t_k\leqslant\tau:=\varphi^{-4n}$, so that
each $M_k$ is bounded entrywise by
$M=\left(\begin{smallmatrix}1&\beta\\\tau\beta&\tau\delta\end{smallmatrix}\right)$,
and since products of nonnegative matrices are monotone entrywise,
$\Theta\leqslant\tr M^P$. The matrix $M$ has trace $1+\tau\delta>0$ and determinant
$\tau(\delta-\beta^2)=5\tau/(9\varphi^6)>0$; its discriminant
$(1+\tau\delta)^2-\tfrac{20}{9}\tau\delta=1-\tfrac29\tau\delta+(\tau\delta)^2$ is positive because
$\tau\delta\leqslant\varphi^{-6}<1$.
Its eigenvalues $\mu_1,\mu_2$ are therefore positive, and
$$
 \tr M^P=\mu_1^P+\mu_2^P\leqslant(\mu_1+\mu_2)^P=(1+\tau\delta)^P
 =\bigl(1+\varphi^{-4n-6}\bigr)^P ,
$$
the binomial expansion of the right-hand side containing the left-hand
side among nonnegative terms. As $1+\eta>1$, the upper bound in
\eqref{eq:row-source-limit} follows.

\emph{The first row.} When $P=1$ the product has the single factor $M_1$ and
$u_1=Q-1$, so that $\Theta=1+t_1\delta=1+\varphi^{-4Q-2}$, which is smaller than
$1+\eta=1+\varphi^{2-4Q}$. Hence $\rho_{1,Q}<C_1$, and $\rho_{1,Q}=C_1(1+\varphi^{-4Q-2})/(1+\varphi^{2-4Q})$ tends to
$C_1$. Finally, when $P\geqslant2$, the integer $n\geqslant(Q-2P+1)/P$ grows without bound
with $Q$, and both sides of \eqref{eq:row-source-limit} tend to $C_P$.
\end{proof}

The proposition shows where the constant comes from. In the leading
term of \ref{it:row-trace}, each letter $a$ contributes the factor $\varphi^2$ and each
letter $b$ the factor $\widehat B_{11}=3\varphi^3/\sqrt5=\lambda\varphi^4$; the
Fibonacci normalization grows by $\varphi^2$ for each unit of cost, a letter $a$
adding one unit and a letter $b$ two. Every letter $b$ therefore costs
exactly the factor $\lambda$ against the Fibonacci growth, and
$C_P=(\varphi^2/\lambda)\lambda^P$. The same number appears on the side of the
paths. Normalized by the maximum at its own cost, the second envelope of
\cref{lem:second-envelope} tends to $\lambda$, since
$E_c/F_{2c+3}=(F_c^2+F_{c+2}^2)/(F_{c+1}^2+F_{c+2}^2)$ and $1+\varphi^4=3\varphi^2$;
and the second row lies on that envelope, because for odd $Q=2n+3$ the
word $\gamma_{2/Q}$ has central factor $a^nba^n$, whose Gram factor
$P_2\,\mu(a)^n$ sends $e=(F_3,F_2)^{\mathsf T}$ to $(F_{Q+2},F_Q)^{\mathsf T}$, so
that $m_{2/Q}=E_Q$. The value $\lambda$ is also the limit of the one-defect
ratio of \cref{fig:parity-order}(a).

Lee, Li, Rabideau and Schiffler prove, in their coordinates, that for fixed
$P\geqslant2$ the ratio \eqref{eq:row-ratio} does not increase with $Q$
\cite[Lemma~6.4]{LeeEtAl2023}, and they compute its limit in the proof of
their Lemma~6.6 \cite[equation~(6.4)]{LeeEtAl2023}.\footnote{Two printed slips
of that proof should be corrected. The equality labels the first-row value
$m_{q,1}$ as $F_{2q-1}$ on the line that carries the Binet expression of
$F_{2q+1}$, which is the correct index. And the monotonicity is asserted there
for every positive parameter, whereas at the parameter of the first row the
ratio is $F_{2Q+1}/F_{2Q-1}$, which increases strictly to $\varphi^2$ and stays
below its limit, as \cref{prop:row-labels}\ref{it:row-bounds} also shows.}
Neither result is used below: \cref{prop:row-labels} gives the lower bound
for every denominator directly, and its upper bound, completed by finitely
many exact comparisons, replaces the monotonicity in \cref{sec:assembly}.

\subsection{The first seven rows}

With the second envelope and the labels of a row at hand, the first seven
rows follow from norms alone: their source labels sit above every norm of
smaller cost except possibly the Fibonacci maximum, and that maximum is
avoided.

\begin{proposition}\label{prop:rows-seven}
Every occurrence with $1\leqslant P\leqslant7$ is weakly
Euclidean-minimal.
\end{proposition}

\begin{proof}
The first row has $m_{1/Q}=F_{2Q+1}$ and cost $L=Q-1$.
Every smaller cost has norm at most $F_{2Q-1}$, proving this case.
The row formula follows, for example, by iterating the first Cohn
matrix in \eqref{eq:cohn-source}.

For $L\geqslant6$ put $u=F_{L-1}/F_L$. The interval
$3/5\leqslant u\leqslant5/8$ is preserved by $u\mapsto1/(1+u)$.
The inequality $E_{L-1}<(5/6)F_{2L+1}$ reduces to
$7u^2+2u-4<0$, whose maximum on this interval is $-1/64$.
The constants $C_P$ decrease with $P$, and
$$
 C_7=\frac{729}{125\varphi^4}>\frac56,\qquad
 C_6=\frac{243}{25\sqrt5\varphi^3}>1.
$$
The two radical signs reduce respectively to
$4373>1875\sqrt5$ and $59>25\sqrt5$, with squared margins
$1545004$ and $356$. Hence
\eqref{eq:row-source-limit} proves \eqref{eq:one-hole} for
$2\leqslant P\leqslant6$. For $P=7$, the lower bound of
\eqref{eq:row-source-limit} and $C_7>5/6$ give $m>E_{L-1}$, and it remains to
show that $m\neq F_{2L+1}$, that is $\rho_{7,Q}\neq1$; we show $\rho_{7,Q}<1$. When
$Q\geqslant14$, every run of $\gamma_{7/Q}$ has length at least
$\lfloor(Q-7)/7\rfloor\geqslant1$, and \cref{prop:row-labels} gives
$\rho_{7,Q}<C_7(1+\varphi^{-10})^7$. Since $\varphi^{10}=(123+55\sqrt5)/2>120$ and
$\varphi^4=(7+3\sqrt5)/2>6.85$, one has
$(1+\varphi^{-10})^7<e^{7/120}<e^{0.06}<1.07$ and
$C_7=729/(125\varphi^4)<729/856<0.86$, so that $\rho_{7,Q}<0.86\cdot1.07<0.93$.
For $8\leqslant Q\leqslant13$ the six labels
$195025$, $499393$, $1278818$, $3276509$, $8399329$ and $21531778$ lie below
$F_{27}=196418$, $F_{29}=514229$, $F_{31}=1346269$, $F_{33}=3524578$,
$F_{35}=9227465$ and $F_{37}=24157817$ respectively.
Thus this row also avoids the Fibonacci hole.

The remaining $L\leqslant5$ slopes are
$1/2,1/3,1/4,1/5,1/6,2/3,2/5,3/4$.
Their $(L,m)$ values are respectively
$$
 (1,5),(2,13),(3,34),(4,89),(5,233),(3,29),(5,194),(5,169).
$$
Each satisfies $m>F_{2L+1}$, the maximum at costs at most $L-1$.
\end{proof}

The later rows require a finer separation than the second envelope.
The next section classifies fixed-defect paths by their gap parities, and
Section~\ref{sec:assembly} combines those unbounded covers with the
complete finite Gaussian fibers of Section~\ref{sec:gaussian}.

% Authors: Dr. Denys Dutykh (Mathematics Department, Khalifa University of Science and Technology,
% Abu Dhabi, UAE) and Prof. Laurent Vuillon (Univ. Savoie Mont Blanc, CNRS, LAMA, Chambery, France)
\section{Covering the unbounded tails}
\label{sec:tails}

A finite denominator check cannot establish weak minimality along an entire
Farey row. The missing step is a uniform separation between the source
ratio and every shorter inverse-Euclidean path. We obtain this separation
by covering an infinite lattice of gap lengths with finitely many boxes.
On each box the normalized ratio is controlled by its two values at a
matched pair of corners, and the reason two values suffice is an order
theorem for that ratio, Fibonacci denominator included.

\subsection{A common eigenchannel for every defect count}

Use the inverse-Euclidean matrices
$$
 \mathsf A=\begin{pmatrix}1&0\\1&1\end{pmatrix},\qquad
 \mathsf B=\begin{pmatrix}0&1\\1&1\end{pmatrix},\qquad z_0=\binom12.
$$
For $k\geqslant1$, a path with exactly $k$ defects is described uniquely
by nonnegative gaps $g=(g_0,\ldots,g_k)$: apply $\mathsf B^{g_i}$ and
then $\mathsf A$ for $i<k$, and apply only $\mathsf B^{g_k}$ at the end.
The order in this description is chronological. If $S=\sum_i g_i$, the
cost is $c=S+k+1$.

Put $\varphi=(1+\sqrt5)/2$, $\theta=\varphi^{-2}$, and $q=-\theta=\varphi-2$.
The matrix
$$
 T=\begin{pmatrix}1&\varphi\\\varphi&-1\end{pmatrix}
$$
satisfies $T^2=(1+\varphi^2)I$. Moreover
$$
 \varphi^{-g}T\mathsf B^gT^{-1}=\operatorname{diag}(1,q^g),
 \qquad
 \sqrt5\varphi^{-1}T\mathsf AT^{-1}
 =\begin{pmatrix}2&1\\q&-2q\end{pmatrix}.
$$
Thus we may start with $(X,Y)=(\varphi^2,2-\sqrt5)$, apply
\begin{equation}
 \binom XY\longmapsto
 \begin{pmatrix}2&q^{g_i}\\q&-2q^{g_i+1}\end{pmatrix}\binom XY
 \quad(i<k),
 \qquad (X,Y)\longmapsto(X,q^{g_k}Y)
 \quad(i=k),
 \label{eq:tail-channel}
\end{equation}
and denote the final $X^2+Y^2$ by $\Pi_k(g)$. The initial vector is
$\varphi^{-1}Tz_0$, so the product of all removed positive factors gives
$$
 \Pi_k(g)=5^k\varphi^{-2c}(1+\varphi^2)N(g).
$$
Binet's formula, with $2c+3$ odd, consequently gives
\begin{equation}
 \mathcal R_k(g):=\frac{N(g)}{F_{2c+3}}
 =\frac{\Pi_k(g)}{K_k(1-q^{2k+5+2S})},
 \qquad K_k=5^k\varphi^4.
 \label{eq:tail-normalization}
\end{equation}
The denominator is positive. This formula includes zero gaps and the
special final gap, which has no following defect matrix.

The right-hand side of \eqref{eq:tail-normalization} depends on the gaps
only through the numbers $x_j=q^{g_j}$ for $j<k$, which enter the matrices
\eqref{eq:tail-channel}, and $y=\theta^{2g_k}$, since $q^{2S}=y\prod_{j<k}x_j^2$.
It is a rational function of these variables whose denominator stays
positive when any of them is set to zero. We use it to give a meaning to
infinite gaps: the value at $g_j=\infty$ is the value at $x_j=0$, or at
$y=0$ for the final gap. An even gap has $x_j\in(0,1]$, an odd gap has
$x_j\in[-\theta,0)$, and the final gap has $y\in(0,1]$; with the infinite
values added, the ranges become the closed intervals $[0,1]$, $[-\theta,0]$
and $[0,1]$.

\subsection{Order in each parity coordinate}

\begin{theorem}[Normalized parity order]
\label{thm:tail-parity}
For every $k\geqslant1$, hold all gaps but $g_j$ fixed and increase $g_j$
by two. The normalized ratio $\mathcal R_k$ strictly increases if $j<k$
and $g_j$ is even. It strictly decreases if $j<k$ and $g_j$ is odd, or
if $j=k$ with either parity. More precisely, on the closed ranges just
described, $\mathcal R_k$ is a strictly decreasing function of each $x_j$,
$j<k$, and a strictly increasing function of $y$, whatever the values of the
other variables in their closed ranges; in particular the monotonicity
persists when some of the other gaps are infinite.
\end{theorem}

\begin{proof}
\emph{The state before a gap.} Just before any gap, the physical
column has second-to-first coordinate ratio $\xi\geqslant2$. This is true at
$(1,2)$; every intervening $\mathsf B$ block leaves a ratio greater than one,
and the following $\mathsf A$ adds one. Its channel therefore has $X>0$,
$Y<0$, and
$$
 h=-Y/X=\frac{\xi-\varphi}{1+\varphi\xi},\qquad
 \varphi^{-5}\leqslant h<\varphi^{-1}.
$$
This includes the first gap.

\emph{A nonfinal gap.} Write the two columns of the channel defect matrix as
$\mathsf u=(2,q)^{\mathsf T}$ and $\mathsf v=(1,-2q)^{\mathsf T}$. Their
physical preimages are
$$
 T^{-1}\mathsf u=\binom{\varphi^{-2}}1=:\mathsf a,
 \qquad T^{-1}\mathsf v=\binom{\varphi^{-1}}{\varphi^{-3}}=:\mathsf b,
 \qquad \varphi^{-3}\mathsf a\leqslant\mathsf b\leqslant\varphi\mathsf a.
$$
Let $S$ be the product of the physical matrices that follow the defect, a
matrix with nonnegative entries, and put $U=S\mathsf a$ and $V=S\mathsf b$.
Componentwise $\varphi^{-3}U\leqslant V\leqslant\varphi U$, all entries being
nonnegative, so that
$$
 \Delta:=\langle U,V\rangle\geqslant\varphi^{-3}\lVert U\rVert^2>0,
 \qquad \nu:=\lVert V\rVert^2\leqslant\varphi\Delta ;
$$
here $U\neq0$ because every factor of $S$ maps positive columns to positive
columns. The gap and the defect send the channel vector $(X,-hX)$ to
$X(\mathsf u-xh\mathsf v)$, with $x=x_j=q^{g_j}$, and the channel matrices that
follow are positive multiples of $TST^{-1}$. Since $T$ is scaled orthogonal,
$T^{\mathsf T}T=(1+\varphi^2)I$, the numerator is therefore a positive multiple,
independent of $x$, of
$$
 \Pi(x)=X^2\bigl(\lVert U\rVert^2-2xh\Delta+x^2h^2\nu\bigr).
$$

\emph{An even gap.} Here $0\leqslant x\leqslant1$, and $xh\nu<\Delta$ on this
whole closed range, because $h\nu<\varphi^{-1}\varphi\Delta=\Delta$. The
derivative $2X^2h(xh\nu-\Delta)$ of $\Pi$ is therefore negative, and $\Pi$
decreases strictly with $x$. The normalizing denominator
$1-q^{2k+5+2S}=1+\theta^{2k+5}\prod_ix_i^2\cdot y$ increases with $x^2$, so
$\mathcal R_k$ decreases strictly with $x$ on $[0,1]$. Increasing an even gap
by two decreases $x$, and so increases $\mathcal R_k$. This proves the even
case.

\emph{An odd gap.} Here $x=-\tau$ with $0\leqslant\tau\leqslant\theta$. Set
$\zeta=\theta^{2k+5+2\sum_{i\ne j}g_i}$, so that the normalizing denominator is
$1+\zeta\tau^2$. Up to the fixed positive factor $K_k$, the normalized
function is
$$
 F(\tau)=\frac{\omega_0+2\omega_1\tau+\omega_2\tau^2}{1+\zeta\tau^2},\qquad
 \omega_0=X^2\lVert U\rVert^2,\quad \omega_1=X^2h\Delta,\quad \omega_2=X^2h^2\nu.
$$
From the bounds on $\Delta$ and $h$, $\omega_0/\omega_1\leqslant\varphi^8=\theta^{-4}$.
Differentiation gives
$$
 F'(\tau)=\frac{2(\omega_1+\omega_2\tau-\omega_0\zeta\tau-\omega_1\zeta\tau^2)}{(1+\zeta\tau^2)^2},
$$
and the numerator is positive on the closed range: using $\omega_2\tau\geqslant0$,
$\tau\leqslant\theta$ and $\zeta\leqslant\theta^{2k+5}$,
$$
 \omega_1+\omega_2\tau-\omega_0\zeta\tau-\omega_1\zeta\tau^2
 \geqslant \omega_1\bigl(1-\theta^{-4}\theta^{2k+5}\theta-\theta^{2k+5}\theta^2\bigr)
 =\omega_1\bigl(1-\theta^{2k+2}-\theta^{2k+7}\bigr)>0.
$$
Indeed $\theta<1/2$ and $k\geqslant1$, so even the boundary case has
$\theta^4+\theta^9<1$. Thus $F$ increases strictly with $\tau$, that is,
$\mathcal R_k$ decreases strictly with $x=-\tau$ on $[-\theta,0]$. Increasing the
odd gap by two decreases $\tau$, and therefore decreases $\mathcal R_k$.

\emph{The final gap.} Put $\zeta=\theta^{2k+5+2\sum_{i<k}g_i}$. Up to $K_k$, the
normalized function of $y=\theta^{2g_k}$ is
$$
 \frac{X^2+Y^2y}{1+\zeta y},
$$
whose derivative in $y$ is $(Y^2-X^2\zeta)/(1+\zeta y)^2$, of the sign of
$h^2-\zeta$. This is positive since $h^2\geqslant\theta^5>\theta^{2k+5}\geqslant\zeta$,
including $k=1$. So $\mathcal R_k$ increases strictly with $y$ on $[0,1]$;
increasing the final gap decreases $y$ and proves the last assertion.

\emph{The other variables anywhere in their ranges.} The three derivative
computations used only the bounds on $h$, $\Delta$, $\nu$ and $\zeta$, and these
hold whatever the values of the other variables in their closed ranges, not
only at the values $q^{g}$, $\theta^{2g}$ and $0$. Indeed, for
$-\theta\leqslant x\leqslant1$ the channel action $\operatorname{diag}(1,x)$ of a gap
corresponds in physical coordinates to
$$
 M_x=T^{-1}\operatorname{diag}(1,x)\,T=\frac1{1+\varphi^2}
 \begin{pmatrix}1+\varphi^2x&\varphi(1-x)\\\varphi(1-x)&\varphi^2+x\end{pmatrix},
$$
whose entries are nonnegative because $x\geqslant-\theta=-\varphi^{-2}$ and
$x\leqslant1$; it is $\varphi^{-g}\mathsf B^g$ when $x=q^g$, and for $x=0$ it is the
limit that represents an infinite gap. It maps positive columns to positive
columns. The final gap enters only through $y$, and acts on the channel as
$\operatorname{diag}(1,\sqrt y)$ with $0\leqslant\sqrt y\leqslant1$, a matrix of the same
kind. Every suffix is thus still a product of nonnegative matrices, and the
bounds on $\Delta$ and $\nu$ hold as shown above. Next, applied to a column
$(1,\xi)^{\mathsf T}$ with $\xi\geqslant2$, the matrix $M_x$ gives a column with
positive entries whose second coordinate exceeds the first by
$$
 \frac{\varphi^{-1}+\xi+x\varphi^2(\xi-\varphi)}{1+\varphi^2}
 \geqslant\frac{\varphi^{-1}+\varphi}{1+\varphi^2}>0 ,
$$
since for $x<0$ one has $x\varphi^2(\xi-\varphi)\geqslant-\theta\varphi^2(\xi-\varphi)=-(\xi-\varphi)$;
the next $\mathsf A$ therefore gives a ratio above $2$ again, and the bounds on
$h$ hold before every gap. Finally, every other variable has absolute value at
most $1$, so $\zeta\leqslant\theta^{2k+5}$ in both places where it was used. Each of
the three derivative computations therefore applies at every point of the
closed ranges, and the monotonicity is strict there. With $x_i=0$ or $y=0$
this is the case of infinite gaps, which the covers below use.
\end{proof}

\Cref{fig:parity-order} shows the theorem at work in the simplest stratum,
one defect, where the ratio depends on two gaps only: along a nonfinal gap the
even and the odd values form two monotone sequences that approach the common
limit from opposite sides, and along the final gap the ratio decreases.

\begin{figure}[tb]
\centering
% Authors: Dr. Denys Dutykh (Mathematics Department, Khalifa University of Science and
% Technology, Abu Dhabi, UAE) and Prof. Laurent Vuillon (Univ. Savoie Mont Blanc, CNRS, LAMA)
% Generated by supplement/code/make_figures.py from exact data; do not edit by hand.
\begin{tikzpicture}[font=\footnotesize, >={Stealth[length=4pt]}]
\node[ink, right] at (-0.1,4.45) {(a) a nonfinal gap};
\draw[ink, line width=0.5pt] (0,0) -- (5.2,0);
\draw[ink, line width=0.5pt] (0,0) -- (0,4);
\draw[ink, line width=0.5pt] (0.236,0) -- (0.236,-0.08);
\node[ink, below] at (0.236,-0.08) {$0$};
\draw[ink, line width=0.5pt] (1.024,0) -- (1.024,-0.08);
\node[ink, below] at (1.024,-0.08) {$2$};
\draw[ink, line width=0.5pt] (1.812,0) -- (1.812,-0.08);
\node[ink, below] at (1.812,-0.08) {$4$};
\draw[ink, line width=0.5pt] (2.6,0) -- (2.6,-0.08);
\node[ink, below] at (2.6,-0.08) {$6$};
\draw[ink, line width=0.5pt] (3.388,0) -- (3.388,-0.08);
\node[ink, below] at (3.388,-0.08) {$8$};
\draw[ink, line width=0.5pt] (4.176,0) -- (4.176,-0.08);
\node[ink, below] at (4.176,-0.08) {$10$};
\draw[ink, line width=0.5pt] (4.964,0) -- (4.964,-0.08);
\node[ink, below] at (4.964,-0.08) {$12$};
\draw[ink, line width=0.5pt] (0,0) -- (-0.08,0);
\node[ink, left] at (-0.08,0) {$10^{-7}$};
\draw[ink, line width=0.5pt] (0,1.231) -- (-0.08,1.231);
\node[ink, left] at (-0.08,1.231) {$10^{-5}$};
\draw[ink, line width=0.5pt] (0,2.462) -- (-0.08,2.462);
\node[ink, left] at (-0.08,2.462) {$10^{-3}$};
\draw[ink, line width=0.5pt] (0,3.692) -- (-0.08,3.692);
\node[ink, left] at (-0.08,3.692) {$10^{-1}$};
\node[ink] at (2.6,-0.75) {nonfinal gap $g_0$};
\node[ink, rotate=90] at (-1.15,2) {distance to the limit};
\draw[slate, line width=0.5pt] (0.236,3.556) -- (0.63,3.308) -- (1.024,3.047) -- (1.418,2.792) -- (1.812,2.534) -- (2.206,2.277) -- (2.6,2.02) -- (2.994,1.762) -- (3.388,1.505) -- (3.782,1.248) -- (4.176,0.991) -- (4.57,0.733) -- (4.964,0.476);
\filldraw[navy, fill=white, line width=0.7pt] (0.236,3.556) circle (1.5pt);
\filldraw[navy, fill=white, line width=0.7pt] (1.024,3.047) circle (1.5pt);
\filldraw[navy, fill=white, line width=0.7pt] (1.812,2.534) circle (1.5pt);
\filldraw[navy, fill=white, line width=0.7pt] (2.6,2.02) circle (1.5pt);
\filldraw[navy, fill=white, line width=0.7pt] (3.388,1.505) circle (1.5pt);
\filldraw[navy, fill=white, line width=0.7pt] (4.176,0.991) circle (1.5pt);
\filldraw[navy, fill=white, line width=0.7pt] (4.964,0.476) circle (1.5pt);
\filldraw[wine, fill=white, line width=0.7pt] (0.58,3.258) rectangle (0.68,3.358);
\filldraw[wine, fill=white, line width=0.7pt] (1.368,2.742) rectangle (1.468,2.842);
\filldraw[wine, fill=white, line width=0.7pt] (2.156,2.227) rectangle (2.256,2.327);
\filldraw[wine, fill=white, line width=0.7pt] (2.944,1.712) rectangle (3.044,1.812);
\filldraw[wine, fill=white, line width=0.7pt] (3.732,1.198) rectangle (3.832,1.298);
\filldraw[wine, fill=white, line width=0.7pt] (4.52,0.683) rectangle (4.62,0.783);
\filldraw[navy, fill=white, line width=0.7pt] (2.248,3.75) circle (1.5pt);
\node[navy, right, font=\scriptsize] at (2.498,3.75) {even $g_0$: below the limit};
\filldraw[wine, fill=white, line width=0.7pt] (2.198,3.3) rectangle (2.298,3.4);
\node[wine, right, font=\scriptsize] at (2.498,3.35) {odd $g_0$: above the limit};
\node[ink, right] at (6.85,4.45) {(b) the final gap};
\draw[ink, line width=0.5pt] (6.95,0) -- (12.15,0);
\draw[ink, line width=0.5pt] (6.95,0) -- (6.95,4);
\draw[ink, line width=0.5pt] (7.186,0) -- (7.186,-0.08);
\node[ink, below] at (7.186,-0.08) {$0$};
\draw[ink, line width=0.5pt] (7.974,0) -- (7.974,-0.08);
\node[ink, below] at (7.974,-0.08) {$2$};
\draw[ink, line width=0.5pt] (8.762,0) -- (8.762,-0.08);
\node[ink, below] at (8.762,-0.08) {$4$};
\draw[ink, line width=0.5pt] (9.55,0) -- (9.55,-0.08);
\node[ink, below] at (9.55,-0.08) {$6$};
\draw[ink, line width=0.5pt] (10.338,0) -- (10.338,-0.08);
\node[ink, below] at (10.338,-0.08) {$8$};
\draw[ink, line width=0.5pt] (11.126,0) -- (11.126,-0.08);
\node[ink, below] at (11.126,-0.08) {$10$};
\draw[ink, line width=0.5pt] (11.914,0) -- (11.914,-0.08);
\node[ink, below] at (11.914,-0.08) {$12$};
\draw[ink, line width=0.5pt] (6.95,0.315) -- (6.87,0.315);
\node[ink, left] at (6.87,0.315) {$0.79$};
\draw[ink, line width=0.5pt] (6.95,1.115) -- (6.87,1.115);
\node[ink, left] at (6.87,1.115) {$0.80$};
\draw[ink, line width=0.5pt] (6.95,1.915) -- (6.87,1.915);
\node[ink, left] at (6.87,1.915) {$0.81$};
\draw[ink, line width=0.5pt] (6.95,2.715) -- (6.87,2.715);
\node[ink, left] at (6.87,2.715) {$0.82$};
\draw[ink, line width=0.5pt] (6.95,3.515) -- (6.87,3.515);
\node[ink, left] at (6.87,3.515) {$0.83$};
\node[ink] at (9.55,-0.75) {final gap $g_1$};
\draw[navy, line width=0.8pt, ] (7.186,2.733) -- (7.58,0.634) -- (7.974,0.328) -- (8.368,0.283) -- (8.762,0.277) -- (9.156,0.276) -- (9.55,0.276) -- (9.944,0.276) -- (10.338,0.276) -- (10.732,0.276) -- (11.126,0.276) -- (11.52,0.276) -- (11.914,0.276);
\draw[wine, line width=0.8pt, densely dashed] (7.186,3.724) -- (7.58,1.771) -- (7.974,1.486) -- (8.368,1.444) -- (8.762,1.438) -- (9.156,1.437) -- (9.55,1.437) -- (9.944,1.437) -- (10.338,1.437) -- (10.732,1.437) -- (11.126,1.437) -- (11.52,1.437) -- (11.914,1.437);
\filldraw[navy, fill=white, line width=0.6pt] (7.186,2.733) circle (1.25pt);
\filldraw[navy, fill=white, line width=0.6pt] (7.58,0.634) circle (1.25pt);
\filldraw[navy, fill=white, line width=0.6pt] (7.974,0.328) circle (1.25pt);
\filldraw[navy, fill=white, line width=0.6pt] (8.368,0.283) circle (1.25pt);
\filldraw[navy, fill=white, line width=0.6pt] (8.762,0.277) circle (1.25pt);
\filldraw[navy, fill=white, line width=0.6pt] (9.156,0.276) circle (1.25pt);
\filldraw[navy, fill=white, line width=0.6pt] (9.55,0.276) circle (1.25pt);
\filldraw[navy, fill=white, line width=0.6pt] (9.944,0.276) circle (1.25pt);
\filldraw[navy, fill=white, line width=0.6pt] (10.338,0.276) circle (1.25pt);
\filldraw[navy, fill=white, line width=0.6pt] (10.732,0.276) circle (1.25pt);
\filldraw[navy, fill=white, line width=0.6pt] (11.126,0.276) circle (1.25pt);
\filldraw[navy, fill=white, line width=0.6pt] (11.52,0.276) circle (1.25pt);
\filldraw[navy, fill=white, line width=0.6pt] (11.914,0.276) circle (1.25pt);
\filldraw[wine, fill=white, line width=0.6pt] (7.186,3.724) circle (1.25pt);
\filldraw[wine, fill=white, line width=0.6pt] (7.58,1.771) circle (1.25pt);
\filldraw[wine, fill=white, line width=0.6pt] (7.974,1.486) circle (1.25pt);
\filldraw[wine, fill=white, line width=0.6pt] (8.368,1.444) circle (1.25pt);
\filldraw[wine, fill=white, line width=0.6pt] (8.762,1.438) circle (1.25pt);
\filldraw[wine, fill=white, line width=0.6pt] (9.156,1.437) circle (1.25pt);
\filldraw[wine, fill=white, line width=0.6pt] (9.55,1.437) circle (1.25pt);
\filldraw[wine, fill=white, line width=0.6pt] (9.944,1.437) circle (1.25pt);
\filldraw[wine, fill=white, line width=0.6pt] (10.338,1.437) circle (1.25pt);
\filldraw[wine, fill=white, line width=0.6pt] (10.732,1.437) circle (1.25pt);
\filldraw[wine, fill=white, line width=0.6pt] (11.126,1.437) circle (1.25pt);
\filldraw[wine, fill=white, line width=0.6pt] (11.52,1.437) circle (1.25pt);
\filldraw[wine, fill=white, line width=0.6pt] (11.914,1.437) circle (1.25pt);
\draw[navy, line width=0.8pt, ] (9.55,3.7) -- (9.95,3.7);
\filldraw[navy, fill=white, line width=0.6pt] (9.75,3.7) circle (1.25pt);
\node[navy, right, font=\scriptsize] at (10,3.7) {$g_0=2$};
\draw[wine, line width=0.8pt, densely dashed] (9.55,3.3) -- (9.95,3.3);
\filldraw[wine, fill=white, line width=0.6pt] (9.75,3.3) circle (1.25pt);
\node[wine, right, font=\scriptsize] at (10,3.3) {$g_0=3$};
\node[ink, rotate=90] at (5.8,2) {$N/F_{2c+3}$};
\end{tikzpicture}

\caption{The normalized parity order for one defect, computed exactly. (a)
The one-defect paths with final gap $0$, as a function of the gap $g_0$ before
the defect: the distance of the ratio $N/F_{2c+3}$ from its value
$3/(\sqrt5\,\varphi)=(15-3\sqrt5)/10$ at $g_0=\infty$, on a logarithmic axis.
The even values (navy circles) lie below that limit and increase towards it,
the odd values (wine squares) lie above it and decrease. (b) The ratio itself
as a function of the final gap $g_1$, for $g_0=2$ and $g_0=3$: it decreases
for either parity.}
\label{fig:parity-order}
\end{figure}

\subsection{A finite cover of the infinite gap lattice}

Fix the parity of every gap. A coordinate domain is either a finite
step-two interval or an infinite step-two tail. For a product box, its
lower matched corner uses the lower endpoint in increasing coordinates
and the upper endpoint in decreasing coordinates; the upper matched
corner makes the opposite choices. An infinite endpoint means $x_j=0$, or
$y=0$ for the final gap. The numerator and denominator in
\eqref{eq:tail-normalization} are always evaluated at the same corner.
Moving one coordinate at a time from a point of the box to a corner, and
applying \cref{thm:tail-parity} on the closed ranges at each move, gives
\begin{equation}
 \mathcal R_k(g^{\min})\leqslant\mathcal R_k(g)
 \leqslant\mathcal R_k(g^{\max}).
 \label{eq:tail-corner-bounds}
\end{equation}
A corner need not satisfy a subsequently imposed lower bound on the total
gap sum; it still safely bounds that restricted part of the box.

For a required total $S\geqslant S_0$, set
$\eta=\lceil S_0/(k+1)\rceil$. Every eligible tuple has some gap at least
$\eta$. Partition it by the first such gap: earlier coordinates are below
$\eta$, that coordinate is at least $\eta$, and later coordinates are
unrestricted. Intersecting with a fixed parity gives $k+1$ disjoint root
boxes whenever $\eta\geqslant2$. Empty parity intersections are discarded.
Their union contains every eligible tuple exactly once.

Every nonterminal box is split into two disjoint step-two domains along
one coordinate, namely the non-singleton coordinate whose domain has the
smallest lower endpoint, the first such coordinate in case of a tie: a finite
domain of $n$ points is cut into its first $\lfloor n/2\rfloor$ points and the
rest, and an infinite tail into its first point and the rest.
A finite box may be discarded if its maximum total is
less than $S_0$. Otherwise it closes only when its upper matched corner
is strictly below the lower source threshold, or its lower matched corner
is strictly above the upper source threshold. A singleton inside the
closed source band would be an explicit failure. A subdivision tree all of
whose leaves close therefore settles the whole infinite domain: the root
boxes cover the eligible tuples, each exactly once, each split partitions a
box, and every leaf lies entirely on one side of the band by
\eqref{eq:tail-corner-bounds}.
Whether a given tree closes is a finite question, settled for the trees of
this paper in \cref{prop:finite-verifications}.

All the quantities compared here lie in $\Z[\varphi]$: the channel
entries, every power of $q$, and $K_k$. Writing $a+b\varphi$ for the pair
of integers $(a,b)$ and clearing threshold denominators beforehand, the
sign of $a+b\varphi$ is the sign of $2a+b+b\sqrt5$, which is decided by
the signs of the two integers $2a+b$ and $b$ and, when these oppose each
other, by the sign of $(2a+b)^2-5b^2$. Every comparison below is therefore
an exact comparison of integers; no decimal approximation of $\sqrt5$
enters.

\subsection{The six unbounded source windows}

With the constant $C_P=\varphi^2\lambda^{P-1}$ of \cref{sec:row-labels}, which
the source labels of row $P$ approach, define the rational ceilings and
starting costs by
\begin{equation}
 \begin{array}{c@{\qquad}c@{\qquad}c@{\qquad}c}
 P&c_P&U_P&\text{middle defects}\\[2pt]
 8&34&707/1000&2\\
 9&41&2929/5000&2,3\\
 10&46&4857/10000&2,3,4\\
 11&84&20111/50000&2,3,4,5\\
 12&110&333511/10^6&2,3,4,5,6\\
 13&156&2765404267/10^{10}&2,3,4,5,6,7.
 \end{array}
 \label{eq:tail-source-windows}
\end{equation}
For row $P$, the closed source band is $[C_P,U_P]$. The paths with no
defect and with a single defect are excluded directly in \cref{sec:assembly},
by the Fibonacci maximum and by \cref{lem:one-defect-minimum}; the paths with
at least two defects are the business of the covers. For each pair $(P,k)$
with $2\leqslant k\leqslant P-6$ (the middle defects of the table) the cover
is built with $S_0=c_P-k-1$ and both thresholds; for $k=P-5$ with the lower
threshold alone. \Cref{thm:tail-covers} in \cref{sec:assembly} states what
the closed trees prove.

The covers prove only the source-band exclusions needed at the actual tail
junctions. They do not classify every extremal ray or assert the same bands
at smaller costs. That narrower target makes one uniform order theorem
sufficient for every stratum.

% Authors: Dr. Denys Dutykh (Mathematics Department, Khalifa University of Science and Technology,
% Abu Dhabi, UAE) and Prof. Laurent Vuillon (Univ. Savoie Mont Blanc, CNRS, LAMA, Chambery, France)
\section{Finite verifications and the proof of the theorem}\label{sec:assembly}

We now prove the main theorem. Its computer-assisted part consists of
finitely many assertions about integers and about elements of $\Z[\varphi]$;
we state them together, derive from them the three statements on which the
assembly rests, and then join the finite windows to the infinite tails at
explicit denominators, so that no interval between a finite calculation and
an asymptotic statement is left untreated. Throughout this section, for
$P=8,\ldots,13$ respectively, put
\begin{equation}\label{eq:tail-denominators}
 Q_P=29,35,39,76,101,146 .
\end{equation}
These are one beyond the six finite windows of \cref{tab:gaussian-windows}.
They satisfy $P+Q_P-3=c_P$, with $c_P$ as in
\eqref{eq:tail-source-windows}, and all six fractions $P/Q_P$ are reduced.
Beyond a second denominator the upper bound of \cref{prop:row-labels} alone
keeps a source below its ceiling $U_P$; with
\begin{equation}\label{eq:second-junctions}
 n_P=3,4,4,7,8,11,\qquad Q^*_P=P(n_P+1)=32,45,50,88,108,156,
\end{equation}
every reduced $P/Q$ with $Q\geqslant Q^*_P$ has runs of length at least $n_P$,
and the denominators $Q_P\leqslant Q<Q^*_P$ in between are finitely many.

\subsection{The finite verifications}

\begin{proposition}[Finite verifications]\label{prop:finite-verifications}
The following assertions hold.
\begin{enumerate}[label=\textup{(\roman*)},ref=\textup{(\roman*)},leftmargin=*]
\item\label{it:fv-windows} For each of the $248$ reduced fractions $P/Q$ of
 \cref{tab:gaussian-windows}, the label $m_{P/Q}$ is a product of explicit
 prime powers; each of these primes passes the criterion \eqref{eq:full-lucas}
 with explicit witnesses, recursively down to the prime two; each odd one is
 congruent to $1$ modulo $4$ and has an explicit decomposition as a sum of two
 squares; and among the pairs \eqref{eq:gaussian-sign-cube} built from these
 data, the occurrence pair of \cref{prop:occurrence-cost} has cost $P+Q-2$ and
 every other pair has a larger cost.
\item\label{it:fv-covers} For each $P\in\{8,\ldots,13\}$ and each $k$ with $2\leqslant k\leqslant P-6$,
 the subdivision tree of \cref{sec:tails} built from the root boxes of all
 $2^{k+1}$ parity sectors, with $S_0=c_P-k-1$ and the thresholds $C_P$ and
 $U_P$, is finite, and each of its leaves closes: its upper matched corner is
 strictly below $C_P$ or its lower matched corner is strictly above $U_P$.
 For $k=P-5$ and $S_0=c_P-k-1$, every root box has its upper matched corner
 strictly below $C_P$.
\item\label{it:fv-lower} For $k=2$ and $S_0=152$, every root box of the eight parity sectors has
 its upper matched corner strictly below $71/100$.
\item\label{it:fv-sources} For each $P\in\{8,\ldots,13\}$, one has
 $C_P\bigl(1+\varphi^{-4n_P-6}\bigr)^P<U_P$, and
 $m_{P/Q}<U_P\,F_{2(P+Q-2)+1}$ for each of the reduced $Q$ with
 $Q_P\leqslant Q<Q^*_P$, thirty-eight denominators in all.
\end{enumerate}
\end{proposition}

\begin{proof}
Each assertion is a finite statement about integers and about elements of
$\Z[\varphi]$, whose signs are decided exactly by the integer test of
\cref{sec:tails}. They were verified by computer in exact arithmetic, and the
verification was carried out twice independently, the two computations
resting on different formulas. For part~\ref{it:fv-windows}, the first computation builds the
fiber of each label from the sign cube \eqref{eq:gaussian-sign-cube} and
proves every prime by the criterion \eqref{eq:full-lucas}, applied
recursively down to the prime two; the second finds every representation
afresh, from the square roots of $-1$ modulo the label, combined by the
Chinese remainder theorem and reduced by Cornacchia's descent, and checks
every prime again, by the same criterion with independently written
arithmetic and by a probable-prime test. In both, each label is computed
from its Cohn matrix, the first computation obtaining it in two ways, from the
floor increments of the Christoffel word and by Stern--Brocot mediants. For
parts~\ref{it:fv-covers} and~\ref{it:fv-lower}, the first computation evaluates the normalized ratio
through the channel recursion \eqref{eq:tail-channel} in $\Z[\varphi]$; the
second rebuilds every tree from its root boxes and evaluates the ratio
through the spectral projectors of the matrix $\mathsf B$ in $\Z[\sqrt5]$. For
part~\ref{it:fv-sources}, the first computation takes each label from its Cohn matrix and
decides the constant comparison by the integer sign test in $\Q(\sqrt5)$;
the second takes each label as the continuant of the palindromic continued
fraction $[2,\sigma(z),2]$ of \c{C}anak\c{c}{\i} and Schiffler, the Fibonacci
numbers from the doubling formulas, and bounds the constants with rational
enclosures of $\sqrt5$ to sixty decimal places.
\end{proof}

For instance, the first integer comparison of part~\ref{it:fv-sources} is the positive
integer identity
$$
 707F_{71}-1000\,m_{8/29}=25054682255203.
$$

\begin{remark}\label{rem:computer-contribution}
The computer contributes only the assertions of
\cref{prop:finite-verifications}: the factorizations and the costs of
$1638$ explicit pairs, the closing of the boxes of finitely many trees,
thirty-eight comparisons of integers and six comparisons of constants.
Everything that turns them into statements about infinitely many
denominators is proved in the text: that the sign cube lists the whole
fiber (\cref{prop:gaussian-fibre}), that two matched corners bound a whole
box (\cref{thm:tail-parity} and \eqref{eq:tail-corner-bounds}), that the
root boxes cover every admissible gap tuple exactly once and each split
partitions a box (\cref{sec:tails}), that a source row stays in its band
beyond the junction (\cref{prop:row-labels} and \eqref{eq:second-junctions}),
and the comparisons at the remaining costs, below.
\end{remark}

\subsection{Three consequences}

\begin{proposition}[Complete finite windows]\label{prop:finite-windows}
For every reduced $P/Q$ in \cref{tab:gaussian-windows}, the occurrence
pair is the unique element of $\mathcal F(m_{P/Q})$ with minimum
Euclidean cost. That cost is $P+Q-2$.
\end{proposition}

\begin{proof}
By \cref{prop:finite-verifications}\ref{it:fv-windows} the label is a product of proved
primes, the odd ones split as recorded. By \cref{prop:gaussian-fibre}, the
pairs \eqref{eq:gaussian-sign-cube} built from this factorization are the
whole fiber $\mathcal F(m_{P/Q})$. By \cref{prop:finite-verifications}\ref{it:fv-windows}
again, the occurrence pair is the only one of them with cost $P+Q-2$, and
all others cost more. The table contains every reduced denominator of each
window, so no occurrence of the windows is missing.
\end{proof}

\begin{theorem}[Uniform tail exclusion]
\label{thm:tail-covers}
For each $P\in\{8,9,10,11,12,13\}$, each $2\leqslant k\leqslant P-6$, and
every nonnegative gap tuple with cost $c\geqslant c_P$, one has
$$
 \frac{N(g)}{F_{2c+3}}<C_P
 \quad\text{or}\quad
 \frac{N(g)}{F_{2c+3}}>U_P.
$$
The inequalities are strict, so no endpoint in these strata belongs to
the closed source band $[C_P,U_P]$. Every path with $k\geqslant P-5$
defects and cost $c\geqslant c_P$ satisfies the stronger conclusion
$N(g)/F_{2c+3}<C_P$.
\end{theorem}

\begin{proof}
\emph{Middle defect counts.} Fix $P$ and $k$ with $2\leqslant k\leqslant P-6$.
A tuple with $k$ defects has cost $c=S+k+1$, so $c\geqslant c_P$ means
$S\geqslant S_0=c_P-k-1$. By \cref{sec:tails}, the root boxes of all parity
sectors cover the tuples with $S\geqslant S_0$, each exactly once, and each
split partitions a box into two. By \cref{prop:finite-verifications}\ref{it:fv-covers} the resulting tree is
finite and every leaf closes, so each admissible tuple $g$ lies in exactly one
closed leaf. If the upper matched corner of that leaf is below $C_P$, then
\eqref{eq:tail-corner-bounds} gives
$\mathcal R_k(g)\leqslant\mathcal R_k(g^{\max})<C_P$; otherwise its lower
matched corner is above $U_P$ and
$\mathcal R_k(g)\geqslant\mathcal R_k(g^{\min})>U_P$.

\emph{The defect count $P-5$.} By \cref{prop:finite-verifications}\ref{it:fv-covers} every
root box has its upper matched corner below $C_P$, and
\eqref{eq:tail-corner-bounds} gives $\mathcal R_{P-5}(g)<C_P$ on the whole
root box, hence for every admissible tuple.

\emph{More defects.} If a path has more than $P-5$ defects, replace surplus
$\mathsf A$ moves by $\mathsf B$ as in \cref{lem:path-orders}. At an
ordered positive state $(u,v)$ the replacement increases the column by
$(v-u,0)^{\mathsf T}$, and every subsequent matrix is nonnegative and
invertible. The final norm strictly increases, while the cost stays fixed.
Replacing until exactly $P-5$ defects remain gives a path of the same cost
whose norm exceeds the original one and is below $C_PF_{2c+3}$.
\end{proof}

\begin{corollary}[A lower-cost bound]
\label{cor:tail-lower-cost}
Every path of cost $c\geqslant155$ with at least two defects satisfies
$$
 \frac{N(g)}{F_{2c+3}}<\frac{71}{100}.
$$
\end{corollary}

\begin{proof}
A path with exactly two defects and cost $c\geqslant155$ has $S\geqslant152$,
so it lies in one root box of the cover with $S_0=152$. By
\cref{prop:finite-verifications}\ref{it:fv-lower} and \eqref{eq:tail-corner-bounds}, its
normalized ratio is below the value at the upper matched corner of that box,
hence below $71/100$. A path with more defects is handled by replacing
surplus defects, as in the proof of \cref{thm:tail-covers}.
\end{proof}

The corollary is used only at the cost-$(L-2)$ interface in row thirteen.

\subsection{Where the source intervals meet the covers}

Put $L=P+Q-2$ and $c=L-1$. The occurrence pair has cost $L$, and we must
exclude every representation of the same norm at every smaller cost. We
first place the source in its band: for every reduced $P/Q$ with
$Q\geqslant Q_P$,
\begin{equation}\label{eq:joined-source-band}
 C_P<\frac{m_{P/Q}}{F_{2c+3}}<U_P.
\end{equation}
The lower bound is \cref{prop:row-labels}\ref{it:row-bounds}, since $F_{2c+3}=F_{2(P+Q-2)+1}$.
For the upper bound there are two cases. If $Q\geqslant Q^*_P$, every run of
$\gamma_{P/Q}$ has length at least $n=\lfloor(Q-P)/P\rfloor\geqslant n_P$, and since the
upper bound of \cref{prop:row-labels}\ref{it:row-bounds} decreases with $n$,
\begin{equation}\label{eq:junction-source}
 \frac{m_{P/Q}}{F_{2c+3}}<C_P\bigl(1+\varphi^{-4n_P-6}\bigr)^P<U_P ,
\end{equation}
the second inequality being part~\ref{it:fv-sources} of \cref{prop:finite-verifications}.
If $Q_P\leqslant Q<Q^*_P$, the upper bound is one of the thirty-eight integer
comparisons of the same part. \Cref{fig:junction} shows the band, the source
labels and the bound of \cref{prop:row-labels} for the row of numerator
eight.

\begin{figure}[tb]
\centering
% Authors: Dr. Denys Dutykh (Mathematics Department, Khalifa University of Science and
% Technology, Abu Dhabi, UAE) and Prof. Laurent Vuillon (Univ. Savoie Mont Blanc, CNRS, LAMA)
% Generated by supplement/code/make_figures.py from exact data; do not edit by hand.
\begin{tikzpicture}[font=\footnotesize, >={Stealth[length=4pt]}]
\fill[sand] (0,0) rectangle (12,3.649);
\draw[ochre, line width=0.8pt] (0,3.649) -- (12,3.649);
\node[ochre, above left, font=\scriptsize] at (12,3.649) {$U_8-C_8$};
\node[slate, below left, font=\scriptsize, align=right] at (11.95,3.599) {source band};
\draw[slate, densely dashed, line width=0.6pt] (3.243,0) -- (3.243,5.2);
\node[wine, anchor=south east, font=\scriptsize] at (3.193,5.2) {finite window};
\node[navy, anchor=south west, font=\scriptsize] at (3.293,5.2) {tail from $Q_8=29$};
\draw[teal, line width=0.7pt, dash pattern=on 5pt off 2pt] (0.081,4.972) -- (1.216,4.972) -- (1.216,4.474) -- (2.514,4.474) -- (2.514,4.011) -- (3.811,4.011) -- (3.811,3.552) -- (5.108,3.552) -- (5.108,3.095) -- (6.405,3.095) -- (6.405,2.637) -- (7.703,2.637) -- (7.703,2.18) -- (9,2.18) -- (9,1.722) -- (10.297,1.722) -- (10.297,1.264) -- (11.595,1.264) -- (11.595,0.807) -- (11.919,0.807);
\draw[teal, line width=0.5pt] (3.892,3.432) -- (3.892,3.672);
\node[teal, font=\scriptsize, inner sep=1pt] at (9.892,2.86) {bound of \cref*{prop:row-labels}};
\filldraw[wine, fill=white, line width=0.6pt] (0.162,4.731) circle (1.3pt);
\filldraw[wine, fill=white, line width=0.6pt] (0.486,4.658) circle (1.3pt);
\filldraw[wine, fill=white, line width=0.6pt] (0.811,4.561) circle (1.3pt);
\filldraw[wine, fill=white, line width=0.6pt] (1.135,4.411) circle (1.3pt);
\filldraw[wine, fill=white, line width=0.6pt] (1.459,4.251) circle (1.3pt);
\filldraw[wine, fill=white, line width=0.6pt] (1.784,4.185) circle (1.3pt);
\filldraw[wine, fill=white, line width=0.6pt] (2.108,4.095) circle (1.3pt);
\filldraw[wine, fill=white, line width=0.6pt] (2.432,3.949) circle (1.3pt);
\filldraw[wine, fill=white, line width=0.6pt] (2.757,3.791) circle (1.3pt);
\filldraw[wine, fill=white, line width=0.6pt] (3.081,3.726) circle (1.3pt);
\filldraw[navy, fill=white, line width=0.6pt] (3.405,3.636) circle (1.3pt);
\filldraw[navy, fill=white, line width=0.6pt] (3.73,3.49) circle (1.3pt);
\filldraw[navy, fill=white, line width=0.6pt] (4.054,3.333) circle (1.3pt);
\filldraw[navy, fill=white, line width=0.6pt] (4.378,3.268) circle (1.3pt);
\filldraw[navy, fill=white, line width=0.6pt] (4.703,3.178) circle (1.3pt);
\filldraw[navy, fill=white, line width=0.6pt] (5.027,3.033) circle (1.3pt);
\filldraw[navy, fill=white, line width=0.6pt] (5.351,2.875) circle (1.3pt);
\filldraw[navy, fill=white, line width=0.6pt] (5.676,2.81) circle (1.3pt);
\filldraw[navy, fill=white, line width=0.6pt] (6,2.721) circle (1.3pt);
\filldraw[navy, fill=white, line width=0.6pt] (6.324,2.575) circle (1.3pt);
\filldraw[navy, fill=white, line width=0.6pt] (6.649,2.418) circle (1.3pt);
\filldraw[navy, fill=white, line width=0.6pt] (6.973,2.353) circle (1.3pt);
\filldraw[navy, fill=white, line width=0.6pt] (7.297,2.263) circle (1.3pt);
\filldraw[navy, fill=white, line width=0.6pt] (7.622,2.117) circle (1.3pt);
\filldraw[navy, fill=white, line width=0.6pt] (7.946,1.96) circle (1.3pt);
\filldraw[navy, fill=white, line width=0.6pt] (8.27,1.895) circle (1.3pt);
\filldraw[navy, fill=white, line width=0.6pt] (8.595,1.805) circle (1.3pt);
\filldraw[navy, fill=white, line width=0.6pt] (8.919,1.66) circle (1.3pt);
\filldraw[navy, fill=white, line width=0.6pt] (9.243,1.502) circle (1.3pt);
\filldraw[navy, fill=white, line width=0.6pt] (9.568,1.437) circle (1.3pt);
\filldraw[navy, fill=white, line width=0.6pt] (9.892,1.348) circle (1.3pt);
\filldraw[navy, fill=white, line width=0.6pt] (10.216,1.202) circle (1.3pt);
\filldraw[navy, fill=white, line width=0.6pt] (10.541,1.045) circle (1.3pt);
\filldraw[navy, fill=white, line width=0.6pt] (10.865,0.98) circle (1.3pt);
\filldraw[navy, fill=white, line width=0.6pt] (11.189,0.89) circle (1.3pt);
\filldraw[navy, fill=white, line width=0.6pt] (11.514,0.745) circle (1.3pt);
\filldraw[navy, fill=white, line width=0.6pt] (11.838,0.587) circle (1.3pt);
\draw[ink, line width=0.5pt] (0,0) -- (12,0);
\draw[ink, line width=0.5pt] (0,0) -- (0,5.2);
\draw[ink, line width=0.5pt] (0.324,0) -- (0.324,-0.08);
\node[ink, below] at (0.324,-0.08) {$10$};
\draw[ink, line width=0.5pt] (1.946,0) -- (1.946,-0.08);
\node[ink, below] at (1.946,-0.08) {$20$};
\draw[ink, line width=0.5pt] (3.568,0) -- (3.568,-0.08);
\node[ink, below] at (3.568,-0.08) {$30$};
\draw[ink, line width=0.5pt] (5.189,0) -- (5.189,-0.08);
\node[ink, below] at (5.189,-0.08) {$40$};
\draw[ink, line width=0.5pt] (6.811,0) -- (6.811,-0.08);
\node[ink, below] at (6.811,-0.08) {$50$};
\draw[ink, line width=0.5pt] (8.432,0) -- (8.432,-0.08);
\node[ink, below] at (8.432,-0.08) {$60$};
\draw[ink, line width=0.5pt] (10.054,0) -- (10.054,-0.08);
\node[ink, below] at (10.054,-0.08) {$70$};
\draw[ink, line width=0.5pt] (11.676,0) -- (11.676,-0.08);
\node[ink, below] at (11.676,-0.08) {$80$};
\node[ink] at (6,-0.75) {denominator $Q$};
\draw[ink, line width=0.5pt] (0,0.274) -- (-0.08,0.274);
\node[ink, left] at (-0.08,0.274) {$10^{-9}$};
\draw[ink, line width=0.5pt] (0,1.368) -- (-0.08,1.368);
\node[ink, left] at (-0.08,1.368) {$10^{-7}$};
\draw[ink, line width=0.5pt] (0,2.463) -- (-0.08,2.463);
\node[ink, left] at (-0.08,2.463) {$10^{-5}$};
\draw[ink, line width=0.5pt] (0,3.558) -- (-0.08,3.558);
\node[ink, left] at (-0.08,3.558) {$10^{-3}$};
\draw[ink, line width=0.5pt] (0,4.653) -- (-0.08,4.653);
\node[ink, left] at (-0.08,4.653) {$10^{-1}$};
\node[ink, rotate=90] at (-1.25,2.6) {$m_{8/Q}/F_{2Q+13}-C_8$};
\fill[paper] (0.18,0.15) rectangle (3.05,3.45);
\draw[slate, line width=0.4pt] (0.18,0.15) rectangle (3.05,3.45);
\node[ink, font=\scriptsize, anchor=north west] at (0.22,3.43) {linear scale};
\draw[ochre, line width=0.6pt] (1.04,1.933) -- (2.87,1.933);
\node[ochre, above left, font=\tiny, inner sep=1.5pt] at (2.87,1.933) {$U_8-C_8$};
\draw[slate, densely dashed, line width=0.5pt] (1.65,0.48) -- (1.65,3.05);
\filldraw[wine, fill=white, line width=0.5pt] (1.345,2.918) circle (0.9pt);
\filldraw[navy, fill=white, line width=0.5pt] (1.955,1.789) circle (0.9pt);
\filldraw[navy, fill=white, line width=0.5pt] (2.565,0.661) circle (0.9pt);
\draw[ink, line width=0.4pt] (1.04,0.48) -- (2.87,0.48);
\draw[ink, line width=0.4pt] (1.04,0.48) -- (1.04,3.05);
\draw[ink, line width=0.4pt] (1.345,0.48) -- (1.345,0.42);
\node[ink, below, font=\tiny, inner sep=1.5pt] at (1.345,0.42) {$27$};
\draw[ink, line width=0.4pt] (1.955,0.48) -- (1.955,0.42);
\node[ink, below, font=\tiny, inner sep=1.5pt] at (1.955,0.42) {$29$};
\draw[ink, line width=0.4pt] (2.565,0.48) -- (2.565,0.42);
\node[ink, below, font=\tiny, inner sep=1.5pt] at (2.565,0.42) {$31$};
\draw[ink, line width=0.4pt] (1.04,1.1) -- (0.98,1.1);
\node[ink, left, font=\tiny, inner sep=1.5pt] at (0.98,1.1) {$0.001$};
\draw[ink, line width=0.4pt] (1.04,2.873) -- (0.98,2.873);
\node[ink, left, font=\tiny, inner sep=1.5pt] at (0.98,2.873) {$0.002$};
\end{tikzpicture}

\caption{The source band of the row $P=8$. The normalized labels
$m_{8/Q}/F_{2(Q+6)+1}$ (navy circles) decrease with the denominator and tend
to $C_8$; from the junction $Q_8=29$ on they lie in the band between $C_8$ and
$U_8=707/1000$ (shaded), which contains no normalized norm of a path one
step cheaper: \cref{thm:tail-covers} excludes the paths with at least two
defects, and the paths with fewer are excluded at the start of
\cref{sec:one-step-below}. Before the junction, the wine
circles are the labels of the finite window, settled by their complete
Gaussian fibers. The dashed staircase is the upper bound of
\cref{prop:row-labels}, $C_8\bigl((1+\varphi^{-4n-6})^8-1\bigr)$ with
$n=\lfloor(Q-8)/8\rfloor$; it lies above every label and enters the band at
$Q^*_8=32$, so that only the denominators $29$ and $31$ need a direct
comparison. The vertical scale is $(m_{8/Q}/F_{2(Q+6)+1}-C_8)$, on a
logarithmic axis. The inset repeats the denominators $27$, $29$ and $31$ on a
linear scale: at the junction the label lies inside the band by
$8.1\times10^{-5}$.}
\label{fig:junction}
\end{figure}

\subsection{Excluding cost one step below the occurrence}\label{sec:one-step-below}

The path with no defect has norm $F_{2c+3}$, above every source ceiling.
A one-defect path satisfies
$$
 \frac{N}{F_{2c+3}}\geqslant
 \frac{5F_{2c-1}}{F_{2c+3}}
 >5\varphi^{-4}>\frac{29}{40}>U_8\geqslant U_P.
$$
The first inequality is \cref{lem:one-defect-minimum}; the strict
Fibonacci comparison follows from Binet's formula at odd indices.
\Cref{thm:tail-covers} excludes the closed source band for every
remaining defect count. Thus no representation at cost $L-1$ has norm
$m_{P/Q}$.

\subsection{Excluding all smaller costs}

It remains necessary to check smaller costs: exclusion at one depth
alone would not imply minimality. We use the following elementary
consequences of Binet's formula. For positive odd $n$, the ratios
$F_n/F_{n+2}$ decrease strictly to $\varphi^{-2}$, and
$F_n/F_{n+4}$ decrease strictly to $\varphi^{-4}$. In particular,
for odd $n\geqslant5$,
\begin{equation}\label{eq:coarse-fibonacci-ratios}
 \varphi^{-2}<\frac{F_n}{F_{n+2}}\leqslant\frac5{13}<\frac25,
 \qquad
 \frac{F_n}{F_{n+4}}\leqslant\frac5{34}<\frac3{20}.
\end{equation}
All uses below occur strictly beyond these initial indices.

For $8\leqslant P\leqslant11$, every cost $d\leqslant L-2$ has
$$
 N\leqslant F_{2L-1}
 <\frac25F_{2L+1}
 <C_{11}F_{2L+1}\leqslant C_PF_{2L+1}\leqslant m_{P/Q}.
$$
For $P=12,13$, the same Fibonacci bound gives, for every cost $d\leqslant L-3$,
$$
 N\leqslant F_{2L-3}
 <\frac3{20}F_{2L+1}
 <C_{13}F_{2L+1}\leqslant m_{P/Q}.
$$
Only the cost $d=L-2$ in the last two rows remains. Its Fibonacci
maximum lies above the source, since
$$
 \frac{F_{2d+3}}{F_{2d+5}}>\varphi^{-2}>\frac38>U_{12}>U_{13}.
$$
For row twelve, every non-Fibonacci path satisfies
$N\leqslant E_d<F_{2d+5}/3<C_{12}F_{2d+5}\leqslant m_{12/Q}$.
Here the strict envelope comparison has a short exact proof:
$$
 F_{2d+5}-3E_d=2(F_{d+1}^2-2F_d^2)>0.
$$
Indeed $F_{d+1}/F_d>3/2$ for $d\geqslant5$, and the present range
starts at $d=109$.

For row thirteen, a one-defect path is above the source:
$$
 \frac{N}{F_{2d+5}}\geqslant\frac{5F_{2d-1}}{F_{2d+5}}
 >5\varphi^{-6}>\frac5{18}>U_{13}.
$$
The radical inequality uses $\varphi^6=9+4\sqrt5<18$.
Every path with at least two defects is below the source. Since
$d\geqslant155$, \cref{cor:tail-lower-cost} gives
$$
 \frac{N}{F_{2d+5}}
 =\frac{N}{F_{2d+3}}\frac{F_{2d+3}}{F_{2d+5}}
 <\frac{71}{100}\frac5{13}
 =\frac{71}{260}<C_{13}.
$$
All the comparisons made in this section between $C_P$, the rational
constants and $U_P$ are inequalities in $\Q(\sqrt5)$, decided exactly by
the integer sign test of \cref{sec:tails}; no decimal rounding enters the
proof.

\begin{proof}[Proof of \cref{thm:thirteen-rows}]
\Cref{prop:occurrence-cost} supplies a primitive representation
of cost $L=P+Q-2$. \Cref{prop:rows-seven} proves its weak
minimality for $P\leqslant7$. For each $8\leqslant P\leqslant13$,
\cref{prop:finite-windows} covers every reduced denominator
$P<Q<Q_P$. For $Q\geqslant Q_P$, the preceding three subsections
exclude every smaller cost. This proves weak minimality without a
denominator bound.
\end{proof}

\begin{proof}[Proof of \cref{cor:thirteen-injective}]
By \cref{thm:thirteen-rows}, every occurrence with numerator at most
thirteen is weakly Euclidean-minimal, and \cref{cor:restricted-injectivity}
applies to this family.
\end{proof}

The proof uses complete Gaussian fibers only for the finite initial
windows. In the infinite tails it excludes all shorter primitive paths
directly, so no factorization of an unbounded family of Markov integers
is required. These two forms of completeness are complementary.

% Authors: Dr. Denys Dutykh (Mathematics Department, Khalifa University of Science and Technology,
% Abu Dhabi, UAE) and Prof. Laurent Vuillon (Univ. Savoie Mont Blanc, CNRS, LAMA, Chambery, France)
\section{What the comparison isolates}\label{sec:outlook}

The optimization problem separates two sources of difficulty in
Markov uniqueness. The occurrence cost is explicit, but a fixed label
may have many primitive two-square representations. Gaussian
factorization describes that entire fiber without identifying its
minimum. The path comparison supplies the missing order for the thirteen
rows treated here. Its content is therefore stronger than a comparison
between selected Markov labels, while its conclusion remains restricted
to the family actually proved.

The parity-order theorem has a scope independent of the row cutoff. It
turns the normalized norm for any fixed defect count into a function
whose extrema on parity boxes occur at prescribed corners, even when
some gaps are infinite. The useful feature is the normalization itself:
bounding the numerator and Fibonacci denominator separately loses the
correlation that makes narrow source intervals accessible.

Extending the result to another row requires more than another finite
list of labels. One must first separate its source interval from every
relevant shorter-cost stratum, including the smaller costs in the final
assembly. The cover procedure has no proved termination theorem for
arbitrary source intervals. A box that remains inside such an interval
is an unresolved comparison, not evidence of a Markov collision.
This is the precise additional work required beyond row thirteen.

The broader question is whether the conclusion of
\cref{thm:thirteen-rows} holds on every row.

\begin{conjecture}[Euclidean minimality of Markov representations]
\label{conj:euclidean-minimality}
Every Markov occurrence weakly minimizes the Euclidean cost over the whole
primitive positive two-square fiber of its label.
\end{conjecture}

By \cref{cor:restricted-injectivity}, the conjecture implies the
Frobenius--Markov uniqueness conjecture. Beyond the thirteen rows the
evidence is computational. By complete factorization of the labels, every
one of the $4145$ occurrences with denominator at most $100$, or with
$P+Q\leqslant160$, is strictly minimal: its pair is the only representation of
minimum cost. \Cref{fig:margins} shows by how much, up to the denominator
$100$. The margin can be a single unit: the label
$m_{4/5}=985=12^2+29^2=16^2+27^2$ has the occurrence pair of cost $7$ and a
second pair of cost $8$. \c{C}anak\c{c}{\i} and Schiffler checked their
uniqueness conjectures by computer for all denominators below $70$
\cite[Remark~4.7]{CanakciSchiffler2020}.

\begin{figure}[tb]
\centering
% Authors: Dr. Denys Dutykh (Mathematics Department, Khalifa University of Science and
% Technology, Abu Dhabi, UAE) and Prof. Laurent Vuillon (Univ. Savoie Mont Blanc, CNRS, LAMA)
% Generated by supplement/code/make_figures.py from exact data; do not edit by hand.
\begin{tikzpicture}[font=\footnotesize, >={Stealth[length=4pt]}]
\fill[sand] (0.122,0.036) rectangle (12.322,0.965);
\draw[slate, line width=0.4pt] (0.244,0.071) circle (1.0pt);
\draw[slate, line width=0.4pt] (0.366,0.071) circle (1.0pt);
\draw[slate, line width=0.4pt] (0.366,0.143) circle (1.0pt);
\draw[slate, line width=0.4pt] (0.488,0.071) circle (1.0pt);
\draw[slate, line width=0.4pt] (0.488,0.214) circle (1.0pt);
\draw[slate, line width=0.4pt] (0.61,0.071) circle (1.0pt);
\draw[slate, line width=0.4pt] (0.61,0.143) circle (1.0pt);
\draw[slate, line width=0.4pt] (0.61,0.214) circle (1.0pt);
\draw[slate, line width=0.4pt] (0.732,0.071) circle (1.0pt);
\draw[slate, line width=0.4pt] (0.732,0.357) circle (1.0pt);
\draw[slate, line width=0.4pt] (0.854,0.214) circle (1.0pt);
\draw[slate, line width=0.4pt] (0.854,0.429) circle (1.0pt);
\draw[slate, line width=0.4pt] (0.976,0.071) circle (1.0pt);
\draw[slate, line width=0.4pt] (0.976,0.214) circle (1.0pt);
\draw[slate, line width=0.4pt] (1.098,0.286) circle (1.0pt);
\draw[slate, line width=0.4pt] (1.098,0.357) circle (1.0pt);
\draw[slate, line width=0.4pt] (1.342,0.071) circle (1.0pt);
\draw[slate, line width=0.4pt] (1.342,0.357) circle (1.0pt);
\draw[slate, line width=0.4pt] (1.342,0.5) circle (1.0pt);
\draw[slate, line width=0.4pt] (1.464,0.357) circle (1.0pt);
\draw[slate, line width=0.4pt] (1.586,0.143) circle (1.0pt);
\draw[slate, line width=0.4pt] (1.586,0.286) circle (1.0pt);
\draw[slate, line width=0.4pt] (1.586,0.5) circle (1.0pt);
\draw[slate, line width=0.4pt] (1.586,0.715) circle (1.0pt);
\draw[slate, line width=0.4pt] (1.708,0.071) circle (1.0pt);
\draw[slate, line width=0.4pt] (1.708,0.357) circle (1.0pt);
\draw[slate, line width=0.4pt] (1.708,0.786) circle (1.0pt);
\draw[slate, line width=0.4pt] (1.83,0.143) circle (1.0pt);
\draw[slate, line width=0.4pt] (1.83,0.572) circle (1.0pt);
\draw[slate, line width=0.4pt] (1.83,0.929) circle (1.0pt);
\draw[slate, line width=0.4pt] (1.83,1.001) circle (1.0pt);
\draw[slate, line width=0.4pt] (2.074,0.286) circle (1.0pt);
\draw[slate, line width=0.4pt] (2.074,0.929) circle (1.0pt);
\draw[slate, line width=0.4pt] (2.318,0.357) circle (1.0pt);
\draw[slate, line width=0.4pt] (2.318,0.5) circle (1.0pt);
\draw[slate, line width=0.4pt] (2.318,0.715) circle (1.0pt);
\draw[slate, line width=0.4pt] (2.318,1.073) circle (1.0pt);
\draw[slate, line width=0.4pt] (2.318,1.144) circle (1.0pt);
\draw[slate, line width=0.4pt] (2.44,0.214) circle (1.0pt);
\draw[slate, line width=0.4pt] (2.44,0.786) circle (1.0pt);
\draw[slate, line width=0.4pt] (2.44,0.929) circle (1.0pt);
\draw[slate, line width=0.4pt] (2.44,1.215) circle (1.0pt);
\draw[slate, line width=0.4pt] (2.562,0.071) circle (1.0pt);
\draw[slate, line width=0.4pt] (2.562,1.144) circle (1.0pt);
\draw[slate, line width=0.4pt] (2.562,1.43) circle (1.0pt);
\draw[slate, line width=0.4pt] (2.684,0.357) circle (1.0pt);
\draw[slate, line width=0.4pt] (2.806,0.071) circle (1.0pt);
\draw[slate, line width=0.4pt] (2.806,0.357) circle (1.0pt);
\draw[slate, line width=0.4pt] (2.806,0.429) circle (1.0pt);
\draw[slate, line width=0.4pt] (2.806,0.5) circle (1.0pt);
\draw[slate, line width=0.4pt] (2.928,0.357) circle (1.0pt);
\draw[slate, line width=0.4pt] (3.05,0.143) circle (1.0pt);
\draw[slate, line width=0.4pt] (3.05,1.287) circle (1.0pt);
\draw[slate, line width=0.4pt] (3.05,1.644) circle (1.0pt);
\draw[slate, line width=0.4pt] (3.294,0.572) circle (1.0pt);
\draw[slate, line width=0.4pt] (3.294,1.859) circle (1.0pt);
\draw[slate, line width=0.4pt] (3.538,0.143) circle (1.0pt);
\draw[slate, line width=0.4pt] (3.538,1.43) circle (1.0pt);
\draw[slate, line width=0.4pt] (3.538,1.859) circle (1.0pt);
\draw[slate, line width=0.4pt] (3.66,0.929) circle (1.0pt);
\draw[slate, line width=0.4pt] (3.66,1.644) circle (1.0pt);
\draw[slate, line width=0.4pt] (3.66,2.073) circle (1.0pt);
\draw[slate, line width=0.4pt] (3.782,0.429) circle (1.0pt);
\draw[slate, line width=0.4pt] (3.904,0.5) circle (1.0pt);
\draw[slate, line width=0.4pt] (4.026,0.143) circle (1.0pt);
\draw[slate, line width=0.4pt] (4.148,1.787) circle (1.0pt);
\draw[slate, line width=0.4pt] (4.148,2.216) circle (1.0pt);
\draw[slate, line width=0.4pt] (4.27,0.572) circle (1.0pt);
\draw[slate, line width=0.4pt] (4.514,0.715) circle (1.0pt);
\draw[slate, line width=0.4pt] (4.636,0.357) circle (1.0pt);
\draw[slate, line width=0.4pt] (4.636,0.5) circle (1.0pt);
\draw[slate, line width=0.4pt] (4.88,0.214) circle (1.0pt);
\draw[slate, line width=0.4pt] (5.002,0.071) circle (1.0pt);
\draw[slate, line width=0.4pt] (5.002,1.001) circle (1.0pt);
\draw[slate, line width=0.4pt] (5.002,2.073) circle (1.0pt);
\draw[slate, line width=0.4pt] (5.002,2.216) circle (1.0pt);
\draw[slate, line width=0.4pt] (5.124,0.357) circle (1.0pt);
\draw[slate, line width=0.4pt] (5.124,1.787) circle (1.0pt);
\draw[slate, line width=0.4pt] (5.246,0.572) circle (1.0pt);
\draw[slate, line width=0.4pt] (5.246,0.858) circle (1.0pt);
\draw[slate, line width=0.4pt] (5.246,1.573) circle (1.0pt);
\draw[slate, line width=0.4pt] (5.49,0.786) circle (1.0pt);
\draw[slate, line width=0.4pt] (5.49,2.216) circle (1.0pt);
\draw[slate, line width=0.4pt] (5.49,2.931) circle (1.0pt);
\draw[slate, line width=0.4pt] (5.49,3.146) circle (1.0pt);
\draw[slate, line width=0.4pt] (5.612,0.643) circle (1.0pt);
\draw[slate, line width=0.4pt] (5.612,2.216) circle (1.0pt);
\draw[slate, line width=0.4pt] (5.734,1.287) circle (1.0pt);
\draw[slate, line width=0.4pt] (5.734,2.145) circle (1.0pt);
\draw[slate, line width=0.4pt] (5.734,2.645) circle (1.0pt);
\draw[slate, line width=0.4pt] (5.978,3.432) circle (1.0pt);
\draw[slate, line width=0.4pt] (6.1,1.215) circle (1.0pt);
\draw[slate, line width=0.4pt] (6.1,3.074) circle (1.0pt);
\draw[slate, line width=0.4pt] (6.222,0.715) circle (1.0pt);
\draw[slate, line width=0.4pt] (6.222,1.43) circle (1.0pt);
\draw[slate, line width=0.4pt] (6.222,3.575) circle (1.0pt);
\draw[slate, line width=0.4pt] (6.466,0.357) circle (1.0pt);
\draw[slate, line width=0.4pt] (6.466,1.073) circle (1.0pt);
\draw[slate, line width=0.4pt] (6.466,1.358) circle (1.0pt);
\draw[slate, line width=0.4pt] (6.588,1.358) circle (1.0pt);
\draw[slate, line width=0.4pt] (6.588,2.502) circle (1.0pt);
\draw[slate, line width=0.4pt] (6.588,3.074) circle (1.0pt);
\draw[slate, line width=0.4pt] (6.71,2.288) circle (1.0pt);
\draw[slate, line width=0.4pt] (6.71,2.431) circle (1.0pt);
\draw[slate, line width=0.4pt] (6.832,0.643) circle (1.0pt);
\draw[slate, line width=0.4pt] (6.832,2.931) circle (1.0pt);
\draw[slate, line width=0.4pt] (6.954,1.001) circle (1.0pt);
\draw[slate, line width=0.4pt] (6.954,1.644) circle (1.0pt);
\draw[slate, line width=0.4pt] (6.954,2.073) circle (1.0pt);
\draw[slate, line width=0.4pt] (7.076,0.929) circle (1.0pt);
\draw[slate, line width=0.4pt] (7.076,1.358) circle (1.0pt);
\draw[slate, line width=0.4pt] (7.076,1.787) circle (1.0pt);
\draw[slate, line width=0.4pt] (7.076,2.931) circle (1.0pt);
\draw[slate, line width=0.4pt] (7.198,2.073) circle (1.0pt);
\draw[slate, line width=0.4pt] (7.32,2.216) circle (1.0pt);
\draw[slate, line width=0.4pt] (7.32,2.645) circle (1.0pt);
\draw[slate, line width=0.4pt] (7.442,0.286) circle (1.0pt);
\draw[slate, line width=0.4pt] (7.442,0.643) circle (1.0pt);
\draw[slate, line width=0.4pt] (7.442,1.287) circle (1.0pt);
\draw[slate, line width=0.4pt] (7.442,1.644) circle (1.0pt);
\draw[slate, line width=0.4pt] (7.442,2.145) circle (1.0pt);
\draw[slate, line width=0.4pt] (7.442,3.861) circle (1.0pt);
\draw[slate, line width=0.4pt] (7.564,0.643) circle (1.0pt);
\draw[slate, line width=0.4pt] (7.564,0.929) circle (1.0pt);
\draw[slate, line width=0.4pt] (7.686,1.644) circle (1.0pt);
\draw[slate, line width=0.4pt] (7.808,0.5) circle (1.0pt);
\draw[slate, line width=0.4pt] (7.808,3.074) circle (1.0pt);
\draw[slate, line width=0.4pt] (7.93,0.071) circle (1.0pt);
\draw[slate, line width=0.4pt] (8.052,1.644) circle (1.0pt);
\draw[slate, line width=0.4pt] (8.052,2.931) circle (1.0pt);
\draw[slate, line width=0.4pt] (8.052,3.789) circle (1.0pt);
\draw[slate, line width=0.4pt] (8.174,1.573) circle (1.0pt);
\draw[slate, line width=0.4pt] (8.174,2.931) circle (1.0pt);
\draw[slate, line width=0.4pt] (8.174,4.504) circle (1.0pt);
\draw[slate, line width=0.4pt] (8.296,0.071) circle (1.0pt);
\draw[slate, line width=0.4pt] (8.54,1.215) circle (1.0pt);
\draw[slate, line width=0.4pt] (8.662,0.214) circle (1.0pt);
\draw[slate, line width=0.4pt] (8.662,0.357) circle (1.0pt);
\draw[slate, line width=0.4pt] (8.662,0.786) circle (1.0pt);
\draw[slate, line width=0.4pt] (8.662,1.43) circle (1.0pt);
\draw[slate, line width=0.4pt] (8.784,3.932) circle (1.0pt);
\draw[slate, line width=0.4pt] (8.906,0.357) circle (1.0pt);
\draw[slate, line width=0.4pt] (8.906,1.144) circle (1.0pt);
\draw[slate, line width=0.4pt] (8.906,4.433) circle (1.0pt);
\draw[slate, line width=0.4pt] (9.028,0.357) circle (1.0pt);
\draw[slate, line width=0.4pt] (9.15,0.143) circle (1.0pt);
\draw[slate, line width=0.4pt] (9.15,1.859) circle (1.0pt);
\draw[slate, line width=0.4pt] (9.15,3.36) circle (1.0pt);
\draw[slate, line width=0.4pt] (9.394,1.073) circle (1.0pt);
\draw[slate, line width=0.4pt] (9.394,1.93) circle (1.0pt);
\draw[slate, line width=0.4pt] (9.394,2.288) circle (1.0pt);
\draw[slate, line width=0.4pt] (9.394,2.431) circle (1.0pt);
\draw[slate, line width=0.4pt] (9.516,1.787) circle (1.0pt);
\draw[slate, line width=0.4pt] (9.638,2.288) circle (1.0pt);
\draw[slate, line width=0.4pt] (9.638,2.502) circle (1.0pt);
\draw[slate, line width=0.4pt] (9.638,4.147) circle (1.0pt);
\draw[slate, line width=0.4pt] (9.638,4.361) circle (1.0pt);
\draw[slate, line width=0.4pt] (9.882,2.002) circle (1.0pt);
\draw[slate, line width=0.4pt] (9.882,3.074) circle (1.0pt);
\draw[slate, line width=0.4pt] (9.882,3.718) circle (1.0pt);
\draw[slate, line width=0.4pt] (10.004,0.5) circle (1.0pt);
\draw[slate, line width=0.4pt] (10.004,2.073) circle (1.0pt);
\draw[slate, line width=0.4pt] (10.126,2.216) circle (1.0pt);
\draw[slate, line width=0.4pt] (10.126,2.502) circle (1.0pt);
\draw[slate, line width=0.4pt] (10.248,1.358) circle (1.0pt);
\draw[slate, line width=0.4pt] (10.248,5.934) circle (1.0pt);
\draw[slate, line width=0.4pt] (10.37,4.719) circle (1.0pt);
\draw[slate, line width=0.4pt] (10.492,0.929) circle (1.0pt);
\draw[slate, line width=0.4pt] (10.492,3.789) circle (1.0pt);
\draw[slate, line width=0.4pt] (10.492,3.932) circle (1.0pt);
\draw[slate, line width=0.4pt] (10.614,5.72) circle (1.0pt);
\draw[slate, line width=0.4pt] (10.736,2.645) circle (1.0pt);
\draw[slate, line width=0.4pt] (10.858,5.076) circle (1.0pt);
\draw[slate, line width=0.4pt] (10.98,5.219) circle (1.0pt);
\draw[slate, line width=0.4pt] (11.102,5.219) circle (1.0pt);
\draw[slate, line width=0.4pt] (11.102,6.435) circle (1.0pt);
\draw[slate, line width=0.4pt] (11.224,0.643) circle (1.0pt);
\draw[slate, line width=0.4pt] (11.224,6.077) circle (1.0pt);
\draw[slate, line width=0.4pt] (11.346,3.789) circle (1.0pt);
\draw[slate, line width=0.4pt] (11.468,0.929) circle (1.0pt);
\draw[slate, line width=0.4pt] (11.468,1.073) circle (1.0pt);
\draw[slate, line width=0.4pt] (11.468,1.787) circle (1.0pt);
\draw[slate, line width=0.4pt] (11.468,1.93) circle (1.0pt);
\draw[slate, line width=0.4pt] (11.59,1.644) circle (1.0pt);
\draw[slate, line width=0.4pt] (11.59,2.431) circle (1.0pt);
\draw[slate, line width=0.4pt] (11.59,6.363) circle (1.0pt);
\draw[slate, line width=0.4pt] (11.712,6.792) circle (1.0pt);
\draw[slate, line width=0.4pt] (11.834,1.787) circle (1.0pt);
\draw[slate, line width=0.4pt] (11.834,6.292) circle (1.0pt);
\draw[slate, line width=0.4pt] (11.956,1.358) circle (1.0pt);
\draw[slate, line width=0.4pt] (11.956,1.644) circle (1.0pt);
\draw[slate, line width=0.4pt] (11.956,6.649) circle (1.0pt);
\draw[slate, line width=0.4pt] (12.078,1.787) circle (1.0pt);
\draw[slate, line width=0.4pt] (12.078,2.073) circle (1.0pt);
\draw[slate, line width=0.4pt] (12.078,5.005) circle (1.0pt);
\fill[navy] (1.22,0.643) circle (0.75pt);
\fill[navy] (1.342,0.214) circle (0.75pt);
\fill[navy] (2.562,1.215) circle (0.75pt);
\fill[navy] (2.806,0.286) circle (0.75pt);
\fill[navy] (2.928,0.929) circle (0.75pt);
\fill[navy] (3.05,0.572) circle (0.75pt);
\fill[navy] (3.172,1.501) circle (0.75pt);
\fill[navy] (3.172,1.644) circle (0.75pt);
\fill[navy] (3.66,1.215) circle (0.75pt);
\fill[navy] (3.782,0.643) circle (0.75pt);
\fill[navy] (3.782,0.786) circle (0.75pt);
\fill[navy] (3.904,0.214) circle (0.75pt);
\fill[navy] (4.026,1.144) circle (0.75pt);
\fill[navy] (4.026,1.43) circle (0.75pt);
\fill[navy] (4.026,2.216) circle (0.75pt);
\fill[navy] (4.148,0.357) circle (0.75pt);
\fill[navy] (4.148,2.073) circle (0.75pt);
\fill[navy] (4.148,2.359) circle (0.75pt);
\fill[navy] (4.27,0.643) circle (0.75pt);
\fill[navy] (4.392,1.358) circle (0.75pt);
\fill[navy] (4.392,2.216) circle (0.75pt);
\fill[navy] (4.514,0.357) circle (0.75pt);
\fill[navy] (4.514,1.001) circle (0.75pt);
\fill[navy] (4.514,1.644) circle (0.75pt);
\fill[navy] (4.514,1.859) circle (0.75pt);
\fill[navy] (4.514,2.502) circle (0.75pt);
\fill[navy] (4.514,2.574) circle (0.75pt);
\fill[navy] (4.636,1.644) circle (0.75pt);
\fill[navy] (4.636,2.502) circle (0.75pt);
\fill[navy] (4.758,1.644) circle (0.75pt);
\fill[navy] (4.758,1.787) circle (0.75pt);
\fill[navy] (4.758,2.002) circle (0.75pt);
\fill[navy] (4.758,2.216) circle (0.75pt);
\fill[navy] (4.88,2.645) circle (0.75pt);
\fill[navy] (5.002,0.214) circle (0.75pt);
\fill[navy] (5.002,0.429) circle (0.75pt);
\fill[navy] (5.002,1.859) circle (0.75pt);
\fill[navy] (5.002,1.93) circle (0.75pt);
\fill[navy] (5.002,2.288) circle (0.75pt);
\fill[navy] (5.124,0.786) circle (0.75pt);
\fill[navy] (5.124,2.216) circle (0.75pt);
\fill[navy] (5.124,2.931) circle (0.75pt);
\fill[navy] (5.246,0.357) circle (0.75pt);
\fill[navy] (5.246,0.715) circle (0.75pt);
\fill[navy] (5.246,1.215) circle (0.75pt);
\fill[navy] (5.246,1.287) circle (0.75pt);
\fill[navy] (5.246,2.288) circle (0.75pt);
\fill[navy] (5.246,2.717) circle (0.75pt);
\fill[navy] (5.368,0.357) circle (0.75pt);
\fill[navy] (5.368,2.502) circle (0.75pt);
\fill[navy] (5.368,2.788) circle (0.75pt);
\fill[navy] (5.49,1.144) circle (0.75pt);
\fill[navy] (5.49,2.717) circle (0.75pt);
\fill[navy] (5.612,2.502) circle (0.75pt);
\fill[navy] (5.612,3.074) circle (0.75pt);
\fill[navy] (5.734,0.214) circle (0.75pt);
\fill[navy] (5.734,0.5) circle (0.75pt);
\fill[navy] (5.734,0.786) circle (0.75pt);
\fill[navy] (5.734,2.717) circle (0.75pt);
\fill[navy] (5.734,2.788) circle (0.75pt);
\fill[navy] (5.734,2.86) circle (0.75pt);
\fill[navy] (5.734,3.217) circle (0.75pt);
\fill[navy] (5.856,1.787) circle (0.75pt);
\fill[navy] (5.978,2.788) circle (0.75pt);
\fill[navy] (5.978,2.931) circle (0.75pt);
\fill[navy] (5.978,3.146) circle (0.75pt);
\fill[navy] (5.978,3.217) circle (0.75pt);
\fill[navy] (5.978,3.289) circle (0.75pt);
\fill[navy] (5.978,3.36) circle (0.75pt);
\fill[navy] (6.1,0.214) circle (0.75pt);
\fill[navy] (6.1,2.073) circle (0.75pt);
\fill[navy] (6.222,1.787) circle (0.75pt);
\fill[navy] (6.222,2.073) circle (0.75pt);
\fill[navy] (6.222,2.931) circle (0.75pt);
\fill[navy] (6.344,0.357) circle (0.75pt);
\fill[navy] (6.344,1.93) circle (0.75pt);
\fill[navy] (6.344,2.645) circle (0.75pt);
\fill[navy] (6.344,3.646) circle (0.75pt);
\fill[navy] (6.466,0.643) circle (0.75pt);
\fill[navy] (6.466,0.858) circle (0.75pt);
\fill[navy] (6.466,0.929) circle (0.75pt);
\fill[navy] (6.466,1.144) circle (0.75pt);
\fill[navy] (6.466,1.787) circle (0.75pt);
\fill[navy] (6.466,1.93) circle (0.75pt);
\fill[navy] (6.466,2.431) circle (0.75pt);
\fill[navy] (6.466,3.003) circle (0.75pt);
\fill[navy] (6.466,3.503) circle (0.75pt);
\fill[navy] (6.588,0.071) circle (0.75pt);
\fill[navy] (6.71,0.572) circle (0.75pt);
\fill[navy] (6.71,0.929) circle (0.75pt);
\fill[navy] (6.71,1.501) circle (0.75pt);
\fill[navy] (6.71,2.216) circle (0.75pt);
\fill[navy] (6.71,3.003) circle (0.75pt);
\fill[navy] (6.71,3.074) circle (0.75pt);
\fill[navy] (6.71,3.503) circle (0.75pt);
\fill[navy] (6.71,3.646) circle (0.75pt);
\fill[navy] (6.832,0.071) circle (0.75pt);
\fill[navy] (6.832,1.787) circle (0.75pt);
\fill[navy] (6.832,3.074) circle (0.75pt);
\fill[navy] (6.832,3.789) circle (0.75pt);
\fill[navy] (6.954,0.786) circle (0.75pt);
\fill[navy] (6.954,2.645) circle (0.75pt);
\fill[navy] (6.954,3.146) circle (0.75pt);
\fill[navy] (6.954,3.718) circle (0.75pt);
\fill[navy] (7.076,0.643) circle (0.75pt);
\fill[navy] (7.076,3.789) circle (0.75pt);
\fill[navy] (7.076,3.932) circle (0.75pt);
\fill[navy] (7.198,0.929) circle (0.75pt);
\fill[navy] (7.198,1.144) circle (0.75pt);
\fill[navy] (7.198,1.358) circle (0.75pt);
\fill[navy] (7.198,1.43) circle (0.75pt);
\fill[navy] (7.198,1.787) circle (0.75pt);
\fill[navy] (7.198,2.288) circle (0.75pt);
\fill[navy] (7.198,2.717) circle (0.75pt);
\fill[navy] (7.198,2.788) circle (0.75pt);
\fill[navy] (7.198,3.146) circle (0.75pt);
\fill[navy] (7.198,3.217) circle (0.75pt);
\fill[navy] (7.198,3.36) circle (0.75pt);
\fill[navy] (7.198,3.432) circle (0.75pt);
\fill[navy] (7.198,3.646) circle (0.75pt);
\fill[navy] (7.198,3.932) circle (0.75pt);
\fill[navy] (7.198,4.004) circle (0.75pt);
\fill[navy] (7.32,3.36) circle (0.75pt);
\fill[navy] (7.32,3.789) circle (0.75pt);
\fill[navy] (7.442,0.357) circle (0.75pt);
\fill[navy] (7.442,1.144) circle (0.75pt);
\fill[navy] (7.442,1.787) circle (0.75pt);
\fill[navy] (7.442,1.93) circle (0.75pt);
\fill[navy] (7.442,2.073) circle (0.75pt);
\fill[navy] (7.442,2.216) circle (0.75pt);
\fill[navy] (7.442,2.359) circle (0.75pt);
\fill[navy] (7.442,2.574) circle (0.75pt);
\fill[navy] (7.442,2.645) circle (0.75pt);
\fill[navy] (7.442,2.86) circle (0.75pt);
\fill[navy] (7.442,3.003) circle (0.75pt);
\fill[navy] (7.442,3.36) circle (0.75pt);
\fill[navy] (7.442,3.432) circle (0.75pt);
\fill[navy] (7.442,3.503) circle (0.75pt);
\fill[navy] (7.442,3.646) circle (0.75pt);
\fill[navy] (7.442,3.789) circle (0.75pt);
\fill[navy] (7.564,0.786) circle (0.75pt);
\fill[navy] (7.564,1.644) circle (0.75pt);
\fill[navy] (7.564,2.359) circle (0.75pt);
\fill[navy] (7.564,2.931) circle (0.75pt);
\fill[navy] (7.564,3.217) circle (0.75pt);
\fill[navy] (7.564,3.503) circle (0.75pt);
\fill[navy] (7.564,3.646) circle (0.75pt);
\fill[navy] (7.686,0.071) circle (0.75pt);
\fill[navy] (7.686,0.715) circle (0.75pt);
\fill[navy] (7.686,1.144) circle (0.75pt);
\fill[navy] (7.686,1.787) circle (0.75pt);
\fill[navy] (7.686,1.859) circle (0.75pt);
\fill[navy] (7.686,2.431) circle (0.75pt);
\fill[navy] (7.686,3.074) circle (0.75pt);
\fill[navy] (7.686,3.718) circle (0.75pt);
\fill[navy] (7.686,4.218) circle (0.75pt);
\fill[navy] (7.686,4.361) circle (0.75pt);
\fill[navy] (7.808,0.929) circle (0.75pt);
\fill[navy] (7.808,1.073) circle (0.75pt);
\fill[navy] (7.808,1.358) circle (0.75pt);
\fill[navy] (7.808,1.787) circle (0.75pt);
\fill[navy] (7.808,2.073) circle (0.75pt);
\fill[navy] (7.808,2.216) circle (0.75pt);
\fill[navy] (7.808,2.645) circle (0.75pt);
\fill[navy] (7.808,2.931) circle (0.75pt);
\fill[navy] (7.808,3.217) circle (0.75pt);
\fill[navy] (7.808,3.36) circle (0.75pt);
\fill[navy] (7.808,3.503) circle (0.75pt);
\fill[navy] (7.808,3.646) circle (0.75pt);
\fill[navy] (7.808,4.504) circle (0.75pt);
\fill[navy] (7.93,0.286) circle (0.75pt);
\fill[navy] (7.93,0.858) circle (0.75pt);
\fill[navy] (7.93,1.144) circle (0.75pt);
\fill[navy] (7.93,1.215) circle (0.75pt);
\fill[navy] (7.93,1.358) circle (0.75pt);
\fill[navy] (7.93,1.573) circle (0.75pt);
\fill[navy] (7.93,1.644) circle (0.75pt);
\fill[navy] (7.93,2.359) circle (0.75pt);
\fill[navy] (7.93,2.574) circle (0.75pt);
\fill[navy] (7.93,2.645) circle (0.75pt);
\fill[navy] (7.93,3.289) circle (0.75pt);
\fill[navy] (7.93,3.36) circle (0.75pt);
\fill[navy] (7.93,3.432) circle (0.75pt);
\fill[navy] (7.93,4.004) circle (0.75pt);
\fill[navy] (7.93,4.075) circle (0.75pt);
\fill[navy] (7.93,4.218) circle (0.75pt);
\fill[navy] (7.93,4.361) circle (0.75pt);
\fill[navy] (7.93,4.433) circle (0.75pt);
\fill[navy] (7.93,4.504) circle (0.75pt);
\fill[navy] (8.052,0.357) circle (0.75pt);
\fill[navy] (8.052,3.503) circle (0.75pt);
\fill[navy] (8.052,4.218) circle (0.75pt);
\fill[navy] (8.052,4.647) circle (0.75pt);
\fill[navy] (8.174,0.429) circle (0.75pt);
\fill[navy] (8.174,1.001) circle (0.75pt);
\fill[navy] (8.174,1.215) circle (0.75pt);
\fill[navy] (8.174,1.287) circle (0.75pt);
\fill[navy] (8.174,2.288) circle (0.75pt);
\fill[navy] (8.174,2.359) circle (0.75pt);
\fill[navy] (8.174,2.788) circle (0.75pt);
\fill[navy] (8.174,3.074) circle (0.75pt);
\fill[navy] (8.174,3.646) circle (0.75pt);
\fill[navy] (8.174,3.932) circle (0.75pt);
\fill[navy] (8.174,4.361) circle (0.75pt);
\fill[navy] (8.174,4.647) circle (0.75pt);
\fill[navy] (8.296,0.5) circle (0.75pt);
\fill[navy] (8.296,1.073) circle (0.75pt);
\fill[navy] (8.296,1.644) circle (0.75pt);
\fill[navy] (8.296,1.787) circle (0.75pt);
\fill[navy] (8.296,1.93) circle (0.75pt);
\fill[navy] (8.296,2.645) circle (0.75pt);
\fill[navy] (8.296,2.788) circle (0.75pt);
\fill[navy] (8.296,3.074) circle (0.75pt);
\fill[navy] (8.296,4.218) circle (0.75pt);
\fill[navy] (8.296,4.647) circle (0.75pt);
\fill[navy] (8.418,0.572) circle (0.75pt);
\fill[navy] (8.418,1.358) circle (0.75pt);
\fill[navy] (8.418,1.43) circle (0.75pt);
\fill[navy] (8.418,1.573) circle (0.75pt);
\fill[navy] (8.418,1.787) circle (0.75pt);
\fill[navy] (8.418,2.073) circle (0.75pt);
\fill[navy] (8.418,2.717) circle (0.75pt);
\fill[navy] (8.418,2.931) circle (0.75pt);
\fill[navy] (8.418,3.503) circle (0.75pt);
\fill[navy] (8.418,4.147) circle (0.75pt);
\fill[navy] (8.418,4.361) circle (0.75pt);
\fill[navy] (8.418,4.647) circle (0.75pt);
\fill[navy] (8.418,4.862) circle (0.75pt);
\fill[navy] (8.54,1.644) circle (0.75pt);
\fill[navy] (8.54,1.93) circle (0.75pt);
\fill[navy] (8.54,2.073) circle (0.75pt);
\fill[navy] (8.54,2.645) circle (0.75pt);
\fill[navy] (8.54,2.931) circle (0.75pt);
\fill[navy] (8.54,3.36) circle (0.75pt);
\fill[navy] (8.54,4.075) circle (0.75pt);
\fill[navy] (8.54,4.218) circle (0.75pt);
\fill[navy] (8.54,4.361) circle (0.75pt);
\fill[navy] (8.54,4.933) circle (0.75pt);
\fill[navy] (8.662,0.143) circle (0.75pt);
\fill[navy] (8.662,0.286) circle (0.75pt);
\fill[navy] (8.662,0.929) circle (0.75pt);
\fill[navy] (8.662,1.501) circle (0.75pt);
\fill[navy] (8.662,2.431) circle (0.75pt);
\fill[navy] (8.662,2.931) circle (0.75pt);
\fill[navy] (8.662,3.718) circle (0.75pt);
\fill[navy] (8.662,3.789) circle (0.75pt);
\fill[navy] (8.662,3.861) circle (0.75pt);
\fill[navy] (8.662,3.932) circle (0.75pt);
\fill[navy] (8.662,4.004) circle (0.75pt);
\fill[navy] (8.662,4.433) circle (0.75pt);
\fill[navy] (8.662,4.79) circle (0.75pt);
\fill[navy] (8.662,4.862) circle (0.75pt);
\fill[navy] (8.662,4.933) circle (0.75pt);
\fill[navy] (8.784,0.929) circle (0.75pt);
\fill[navy] (8.784,1.215) circle (0.75pt);
\fill[navy] (8.784,1.358) circle (0.75pt);
\fill[navy] (8.784,1.644) circle (0.75pt);
\fill[navy] (8.784,2.073) circle (0.75pt);
\fill[navy] (8.784,2.931) circle (0.75pt);
\fill[navy] (8.784,3.503) circle (0.75pt);
\fill[navy] (8.784,4.218) circle (0.75pt);
\fill[navy] (8.784,4.361) circle (0.75pt);
\fill[navy] (8.784,4.647) circle (0.75pt);
\fill[navy] (8.784,4.79) circle (0.75pt);
\fill[navy] (8.906,0.643) circle (0.75pt);
\fill[navy] (8.906,1.215) circle (0.75pt);
\fill[navy] (8.906,1.501) circle (0.75pt);
\fill[navy] (8.906,1.644) circle (0.75pt);
\fill[navy] (8.906,1.716) circle (0.75pt);
\fill[navy] (8.906,1.787) circle (0.75pt);
\fill[navy] (8.906,1.859) circle (0.75pt);
\fill[navy] (8.906,2.931) circle (0.75pt);
\fill[navy] (8.906,3.074) circle (0.75pt);
\fill[navy] (8.906,3.217) circle (0.75pt);
\fill[navy] (8.906,3.289) circle (0.75pt);
\fill[navy] (8.906,3.36) circle (0.75pt);
\fill[navy] (8.906,3.575) circle (0.75pt);
\fill[navy] (8.906,3.646) circle (0.75pt);
\fill[navy] (8.906,3.718) circle (0.75pt);
\fill[navy] (8.906,3.789) circle (0.75pt);
\fill[navy] (8.906,3.861) circle (0.75pt);
\fill[navy] (8.906,4.004) circle (0.75pt);
\fill[navy] (8.906,4.147) circle (0.75pt);
\fill[navy] (8.906,4.29) circle (0.75pt);
\fill[navy] (8.906,4.361) circle (0.75pt);
\fill[navy] (8.906,4.576) circle (0.75pt);
\fill[navy] (8.906,4.647) circle (0.75pt);
\fill[navy] (8.906,4.79) circle (0.75pt);
\fill[navy] (8.906,4.933) circle (0.75pt);
\fill[navy] (8.906,5.148) circle (0.75pt);
\fill[navy] (9.028,0.643) circle (0.75pt);
\fill[navy] (9.028,1.358) circle (0.75pt);
\fill[navy] (9.028,2.359) circle (0.75pt);
\fill[navy] (9.028,2.502) circle (0.75pt);
\fill[navy] (9.028,2.931) circle (0.75pt);
\fill[navy] (9.028,3.503) circle (0.75pt);
\fill[navy] (9.028,4.361) circle (0.75pt);
\fill[navy] (9.028,4.933) circle (0.75pt);
\fill[navy] (9.15,0.572) circle (0.75pt);
\fill[navy] (9.15,2.073) circle (0.75pt);
\fill[navy] (9.15,2.216) circle (0.75pt);
\fill[navy] (9.15,2.288) circle (0.75pt);
\fill[navy] (9.15,2.431) circle (0.75pt);
\fill[navy] (9.15,3.074) circle (0.75pt);
\fill[navy] (9.15,3.289) circle (0.75pt);
\fill[navy] (9.15,3.718) circle (0.75pt);
\fill[navy] (9.15,4.218) circle (0.75pt);
\fill[navy] (9.15,4.433) circle (0.75pt);
\fill[navy] (9.15,4.79) circle (0.75pt);
\fill[navy] (9.15,5.291) circle (0.75pt);
\fill[navy] (9.272,0.929) circle (0.75pt);
\fill[navy] (9.272,1.215) circle (0.75pt);
\fill[navy] (9.272,1.644) circle (0.75pt);
\fill[navy] (9.272,1.787) circle (0.75pt);
\fill[navy] (9.272,2.359) circle (0.75pt);
\fill[navy] (9.272,2.502) circle (0.75pt);
\fill[navy] (9.272,2.645) circle (0.75pt);
\fill[navy] (9.272,3.503) circle (0.75pt);
\fill[navy] (9.272,3.932) circle (0.75pt);
\fill[navy] (9.272,4.218) circle (0.75pt);
\fill[navy] (9.272,4.361) circle (0.75pt);
\fill[navy] (9.272,4.647) circle (0.75pt);
\fill[navy] (9.272,4.79) circle (0.75pt);
\fill[navy] (9.272,5.076) circle (0.75pt);
\fill[navy] (9.272,5.219) circle (0.75pt);
\fill[navy] (9.272,5.362) circle (0.75pt);
\fill[navy] (9.394,0.429) circle (0.75pt);
\fill[navy] (9.394,1.43) circle (0.75pt);
\fill[navy] (9.394,2.073) circle (0.75pt);
\fill[navy] (9.394,2.216) circle (0.75pt);
\fill[navy] (9.394,2.717) circle (0.75pt);
\fill[navy] (9.394,2.788) circle (0.75pt);
\fill[navy] (9.394,3.432) circle (0.75pt);
\fill[navy] (9.394,3.718) circle (0.75pt);
\fill[navy] (9.394,3.789) circle (0.75pt);
\fill[navy] (9.394,4.433) circle (0.75pt);
\fill[navy] (9.394,4.647) circle (0.75pt);
\fill[navy] (9.394,4.862) circle (0.75pt);
\fill[navy] (9.394,5.148) circle (0.75pt);
\fill[navy] (9.394,5.219) circle (0.75pt);
\fill[navy] (9.394,5.362) circle (0.75pt);
\fill[navy] (9.516,0.357) circle (0.75pt);
\fill[navy] (9.516,0.786) circle (0.75pt);
\fill[navy] (9.516,1.644) circle (0.75pt);
\fill[navy] (9.516,3.074) circle (0.75pt);
\fill[navy] (9.516,3.789) circle (0.75pt);
\fill[navy] (9.516,3.932) circle (0.75pt);
\fill[navy] (9.516,4.361) circle (0.75pt);
\fill[navy] (9.516,4.79) circle (0.75pt);
\fill[navy] (9.516,5.076) circle (0.75pt);
\fill[navy] (9.638,0.643) circle (0.75pt);
\fill[navy] (9.638,1.001) circle (0.75pt);
\fill[navy] (9.638,1.287) circle (0.75pt);
\fill[navy] (9.638,1.358) circle (0.75pt);
\fill[navy] (9.638,1.573) circle (0.75pt);
\fill[navy] (9.638,1.859) circle (0.75pt);
\fill[navy] (9.638,2.073) circle (0.75pt);
\fill[navy] (9.638,2.216) circle (0.75pt);
\fill[navy] (9.638,2.431) circle (0.75pt);
\fill[navy] (9.638,2.717) circle (0.75pt);
\fill[navy] (9.638,2.931) circle (0.75pt);
\fill[navy] (9.638,3.074) circle (0.75pt);
\fill[navy] (9.638,3.432) circle (0.75pt);
\fill[navy] (9.638,3.575) circle (0.75pt);
\fill[navy] (9.638,3.646) circle (0.75pt);
\fill[navy] (9.638,3.861) circle (0.75pt);
\fill[navy] (9.638,4.218) circle (0.75pt);
\fill[navy] (9.638,4.29) circle (0.75pt);
\fill[navy] (9.638,4.433) circle (0.75pt);
\fill[navy] (9.638,4.504) circle (0.75pt);
\fill[navy] (9.638,4.576) circle (0.75pt);
\fill[navy] (9.638,4.719) circle (0.75pt);
\fill[navy] (9.638,4.862) circle (0.75pt);
\fill[navy] (9.638,5.005) circle (0.75pt);
\fill[navy] (9.638,5.076) circle (0.75pt);
\fill[navy] (9.638,5.219) circle (0.75pt);
\fill[navy] (9.638,5.434) circle (0.75pt);
\fill[navy] (9.638,5.505) circle (0.75pt);
\fill[navy] (9.638,5.577) circle (0.75pt);
\fill[navy] (9.76,0.929) circle (0.75pt);
\fill[navy] (9.76,1.358) circle (0.75pt);
\fill[navy] (9.76,1.501) circle (0.75pt);
\fill[navy] (9.76,1.644) circle (0.75pt);
\fill[navy] (9.76,2.359) circle (0.75pt);
\fill[navy] (9.76,2.645) circle (0.75pt);
\fill[navy] (9.76,3.36) circle (0.75pt);
\fill[navy] (9.76,3.646) circle (0.75pt);
\fill[navy] (9.76,3.789) circle (0.75pt);
\fill[navy] (9.76,4.218) circle (0.75pt);
\fill[navy] (9.76,4.504) circle (0.75pt);
\fill[navy] (9.76,4.79) circle (0.75pt);
\fill[navy] (9.76,5.076) circle (0.75pt);
\fill[navy] (9.76,5.219) circle (0.75pt);
\fill[navy] (9.76,5.505) circle (0.75pt);
\fill[navy] (9.882,0.071) circle (0.75pt);
\fill[navy] (9.882,0.143) circle (0.75pt);
\fill[navy] (9.882,0.929) circle (0.75pt);
\fill[navy] (9.882,1.144) circle (0.75pt);
\fill[navy] (9.882,1.215) circle (0.75pt);
\fill[navy] (9.882,1.358) circle (0.75pt);
\fill[navy] (9.882,1.43) circle (0.75pt);
\fill[navy] (9.882,1.644) circle (0.75pt);
\fill[navy] (9.882,1.859) circle (0.75pt);
\fill[navy] (9.882,2.216) circle (0.75pt);
\fill[navy] (9.882,2.288) circle (0.75pt);
\fill[navy] (9.882,2.502) circle (0.75pt);
\fill[navy] (9.882,2.645) circle (0.75pt);
\fill[navy] (9.882,2.717) circle (0.75pt);
\fill[navy] (9.882,2.86) circle (0.75pt);
\fill[navy] (9.882,3.146) circle (0.75pt);
\fill[navy] (9.882,3.289) circle (0.75pt);
\fill[navy] (9.882,4.004) circle (0.75pt);
\fill[navy] (9.882,4.433) circle (0.75pt);
\fill[navy] (9.882,4.576) circle (0.75pt);
\fill[navy] (9.882,4.647) circle (0.75pt);
\fill[navy] (9.882,5.434) circle (0.75pt);
\fill[navy] (9.882,5.505) circle (0.75pt);
\fill[navy] (9.882,5.648) circle (0.75pt);
\fill[navy] (10.004,2.359) circle (0.75pt);
\fill[navy] (10.004,3.074) circle (0.75pt);
\fill[navy] (10.004,3.36) circle (0.75pt);
\fill[navy] (10.004,3.932) circle (0.75pt);
\fill[navy] (10.004,4.361) circle (0.75pt);
\fill[navy] (10.004,4.504) circle (0.75pt);
\fill[navy] (10.004,4.647) circle (0.75pt);
\fill[navy] (10.004,4.79) circle (0.75pt);
\fill[navy] (10.004,5.076) circle (0.75pt);
\fill[navy] (10.004,5.505) circle (0.75pt);
\fill[navy] (10.004,5.648) circle (0.75pt);
\fill[navy] (10.004,5.791) circle (0.75pt);
\fill[navy] (10.126,0.643) circle (0.75pt);
\fill[navy] (10.126,1.144) circle (0.75pt);
\fill[navy] (10.126,1.287) circle (0.75pt);
\fill[navy] (10.126,1.644) circle (0.75pt);
\fill[navy] (10.126,1.716) circle (0.75pt);
\fill[navy] (10.126,2.073) circle (0.75pt);
\fill[navy] (10.126,2.288) circle (0.75pt);
\fill[navy] (10.126,2.359) circle (0.75pt);
\fill[navy] (10.126,3.074) circle (0.75pt);
\fill[navy] (10.126,3.217) circle (0.75pt);
\fill[navy] (10.126,3.432) circle (0.75pt);
\fill[navy] (10.126,3.503) circle (0.75pt);
\fill[navy] (10.126,3.575) circle (0.75pt);
\fill[navy] (10.126,3.718) circle (0.75pt);
\fill[navy] (10.126,4.004) circle (0.75pt);
\fill[navy] (10.126,4.147) circle (0.75pt);
\fill[navy] (10.126,4.29) circle (0.75pt);
\fill[navy] (10.126,4.433) circle (0.75pt);
\fill[navy] (10.126,4.504) circle (0.75pt);
\fill[navy] (10.126,4.647) circle (0.75pt);
\fill[navy] (10.126,4.719) circle (0.75pt);
\fill[navy] (10.126,4.79) circle (0.75pt);
\fill[navy] (10.126,4.862) circle (0.75pt);
\fill[navy] (10.126,4.933) circle (0.75pt);
\fill[navy] (10.126,5.005) circle (0.75pt);
\fill[navy] (10.126,5.076) circle (0.75pt);
\fill[navy] (10.126,5.219) circle (0.75pt);
\fill[navy] (10.126,5.434) circle (0.75pt);
\fill[navy] (10.126,5.577) circle (0.75pt);
\fill[navy] (10.126,5.648) circle (0.75pt);
\fill[navy] (10.126,5.72) circle (0.75pt);
\fill[navy] (10.248,0.357) circle (0.75pt);
\fill[navy] (10.248,2.216) circle (0.75pt);
\fill[navy] (10.248,2.931) circle (0.75pt);
\fill[navy] (10.248,3.932) circle (0.75pt);
\fill[navy] (10.248,4.361) circle (0.75pt);
\fill[navy] (10.248,4.647) circle (0.75pt);
\fill[navy] (10.248,4.79) circle (0.75pt);
\fill[navy] (10.248,5.076) circle (0.75pt);
\fill[navy] (10.248,5.219) circle (0.75pt);
\fill[navy] (10.248,5.648) circle (0.75pt);
\fill[navy] (10.37,0.143) circle (0.75pt);
\fill[navy] (10.37,0.286) circle (0.75pt);
\fill[navy] (10.37,0.572) circle (0.75pt);
\fill[navy] (10.37,0.643) circle (0.75pt);
\fill[navy] (10.37,0.858) circle (0.75pt);
\fill[navy] (10.37,1.287) circle (0.75pt);
\fill[navy] (10.37,1.358) circle (0.75pt);
\fill[navy] (10.37,1.501) circle (0.75pt);
\fill[navy] (10.37,1.573) circle (0.75pt);
\fill[navy] (10.37,1.644) circle (0.75pt);
\fill[navy] (10.37,1.716) circle (0.75pt);
\fill[navy] (10.37,2.073) circle (0.75pt);
\fill[navy] (10.37,2.288) circle (0.75pt);
\fill[navy] (10.37,2.359) circle (0.75pt);
\fill[navy] (10.37,2.574) circle (0.75pt);
\fill[navy] (10.37,3.003) circle (0.75pt);
\fill[navy] (10.37,3.289) circle (0.75pt);
\fill[navy] (10.37,3.718) circle (0.75pt);
\fill[navy] (10.37,3.789) circle (0.75pt);
\fill[navy] (10.37,3.861) circle (0.75pt);
\fill[navy] (10.37,4.075) circle (0.75pt);
\fill[navy] (10.37,4.147) circle (0.75pt);
\fill[navy] (10.37,4.218) circle (0.75pt);
\fill[navy] (10.37,4.433) circle (0.75pt);
\fill[navy] (10.37,4.576) circle (0.75pt);
\fill[navy] (10.37,4.79) circle (0.75pt);
\fill[navy] (10.37,5.076) circle (0.75pt);
\fill[navy] (10.37,5.148) circle (0.75pt);
\fill[navy] (10.37,5.291) circle (0.75pt);
\fill[navy] (10.37,5.434) circle (0.75pt);
\fill[navy] (10.37,5.505) circle (0.75pt);
\fill[navy] (10.37,5.577) circle (0.75pt);
\fill[navy] (10.37,5.648) circle (0.75pt);
\fill[navy] (10.37,5.791) circle (0.75pt);
\fill[navy] (10.37,5.863) circle (0.75pt);
\fill[navy] (10.37,5.934) circle (0.75pt);
\fill[navy] (10.37,6.006) circle (0.75pt);
\fill[navy] (10.492,0.214) circle (0.75pt);
\fill[navy] (10.492,1.501) circle (0.75pt);
\fill[navy] (10.492,1.644) circle (0.75pt);
\fill[navy] (10.492,2.073) circle (0.75pt);
\fill[navy] (10.492,2.359) circle (0.75pt);
\fill[navy] (10.492,2.502) circle (0.75pt);
\fill[navy] (10.492,2.788) circle (0.75pt);
\fill[navy] (10.492,2.931) circle (0.75pt);
\fill[navy] (10.492,4.075) circle (0.75pt);
\fill[navy] (10.492,4.218) circle (0.75pt);
\fill[navy] (10.492,4.361) circle (0.75pt);
\fill[navy] (10.492,4.504) circle (0.75pt);
\fill[navy] (10.492,4.647) circle (0.75pt);
\fill[navy] (10.492,5.076) circle (0.75pt);
\fill[navy] (10.492,5.219) circle (0.75pt);
\fill[navy] (10.492,5.362) circle (0.75pt);
\fill[navy] (10.492,5.648) circle (0.75pt);
\fill[navy] (10.492,5.791) circle (0.75pt);
\fill[navy] (10.492,5.934) circle (0.75pt);
\fill[navy] (10.614,0.357) circle (0.75pt);
\fill[navy] (10.614,0.5) circle (0.75pt);
\fill[navy] (10.614,0.929) circle (0.75pt);
\fill[navy] (10.614,1.001) circle (0.75pt);
\fill[navy] (10.614,1.358) circle (0.75pt);
\fill[navy] (10.614,1.43) circle (0.75pt);
\fill[navy] (10.614,1.787) circle (0.75pt);
\fill[navy] (10.614,2.002) circle (0.75pt);
\fill[navy] (10.614,2.216) circle (0.75pt);
\fill[navy] (10.614,2.288) circle (0.75pt);
\fill[navy] (10.614,2.431) circle (0.75pt);
\fill[navy] (10.614,2.502) circle (0.75pt);
\fill[navy] (10.614,2.645) circle (0.75pt);
\fill[navy] (10.614,2.717) circle (0.75pt);
\fill[navy] (10.614,2.931) circle (0.75pt);
\fill[navy] (10.614,3.146) circle (0.75pt);
\fill[navy] (10.614,3.36) circle (0.75pt);
\fill[navy] (10.614,3.575) circle (0.75pt);
\fill[navy] (10.614,4.361) circle (0.75pt);
\fill[navy] (10.614,4.433) circle (0.75pt);
\fill[navy] (10.614,4.576) circle (0.75pt);
\fill[navy] (10.614,4.647) circle (0.75pt);
\fill[navy] (10.614,4.862) circle (0.75pt);
\fill[navy] (10.614,5.005) circle (0.75pt);
\fill[navy] (10.614,5.076) circle (0.75pt);
\fill[navy] (10.614,5.219) circle (0.75pt);
\fill[navy] (10.614,5.434) circle (0.75pt);
\fill[navy] (10.614,5.648) circle (0.75pt);
\fill[navy] (10.614,5.863) circle (0.75pt);
\fill[navy] (10.614,5.934) circle (0.75pt);
\fill[navy] (10.614,6.149) circle (0.75pt);
\fill[navy] (10.736,1.215) circle (0.75pt);
\fill[navy] (10.736,1.644) circle (0.75pt);
\fill[navy] (10.736,1.787) circle (0.75pt);
\fill[navy] (10.736,2.502) circle (0.75pt);
\fill[navy] (10.736,3.36) circle (0.75pt);
\fill[navy] (10.736,3.503) circle (0.75pt);
\fill[navy] (10.736,4.218) circle (0.75pt);
\fill[navy] (10.736,4.647) circle (0.75pt);
\fill[navy] (10.736,4.79) circle (0.75pt);
\fill[navy] (10.736,4.933) circle (0.75pt);
\fill[navy] (10.736,5.076) circle (0.75pt);
\fill[navy] (10.736,5.219) circle (0.75pt);
\fill[navy] (10.736,5.362) circle (0.75pt);
\fill[navy] (10.736,5.648) circle (0.75pt);
\fill[navy] (10.736,5.791) circle (0.75pt);
\fill[navy] (10.736,6.077) circle (0.75pt);
\fill[navy] (10.858,0.071) circle (0.75pt);
\fill[navy] (10.858,0.143) circle (0.75pt);
\fill[navy] (10.858,0.214) circle (0.75pt);
\fill[navy] (10.858,0.286) circle (0.75pt);
\fill[navy] (10.858,0.572) circle (0.75pt);
\fill[navy] (10.858,0.643) circle (0.75pt);
\fill[navy] (10.858,1.073) circle (0.75pt);
\fill[navy] (10.858,1.215) circle (0.75pt);
\fill[navy] (10.858,1.358) circle (0.75pt);
\fill[navy] (10.858,1.43) circle (0.75pt);
\fill[navy] (10.858,1.501) circle (0.75pt);
\fill[navy] (10.858,1.644) circle (0.75pt);
\fill[navy] (10.858,1.93) circle (0.75pt);
\fill[navy] (10.858,2.002) circle (0.75pt);
\fill[navy] (10.858,2.073) circle (0.75pt);
\fill[navy] (10.858,2.216) circle (0.75pt);
\fill[navy] (10.858,2.288) circle (0.75pt);
\fill[navy] (10.858,2.359) circle (0.75pt);
\fill[navy] (10.858,2.431) circle (0.75pt);
\fill[navy] (10.858,2.788) circle (0.75pt);
\fill[navy] (10.858,2.86) circle (0.75pt);
\fill[navy] (10.858,2.931) circle (0.75pt);
\fill[navy] (10.858,3.074) circle (0.75pt);
\fill[navy] (10.858,3.146) circle (0.75pt);
\fill[navy] (10.858,3.432) circle (0.75pt);
\fill[navy] (10.858,3.575) circle (0.75pt);
\fill[navy] (10.858,3.646) circle (0.75pt);
\fill[navy] (10.858,3.718) circle (0.75pt);
\fill[navy] (10.858,3.789) circle (0.75pt);
\fill[navy] (10.858,3.932) circle (0.75pt);
\fill[navy] (10.858,4.004) circle (0.75pt);
\fill[navy] (10.858,4.075) circle (0.75pt);
\fill[navy] (10.858,4.147) circle (0.75pt);
\fill[navy] (10.858,4.218) circle (0.75pt);
\fill[navy] (10.858,4.361) circle (0.75pt);
\fill[navy] (10.858,4.647) circle (0.75pt);
\fill[navy] (10.858,4.719) circle (0.75pt);
\fill[navy] (10.858,4.862) circle (0.75pt);
\fill[navy] (10.858,5.005) circle (0.75pt);
\fill[navy] (10.858,5.148) circle (0.75pt);
\fill[navy] (10.858,5.219) circle (0.75pt);
\fill[navy] (10.858,5.291) circle (0.75pt);
\fill[navy] (10.858,5.434) circle (0.75pt);
\fill[navy] (10.858,5.505) circle (0.75pt);
\fill[navy] (10.858,5.577) circle (0.75pt);
\fill[navy] (10.858,5.648) circle (0.75pt);
\fill[navy] (10.858,5.72) circle (0.75pt);
\fill[navy] (10.858,5.863) circle (0.75pt);
\fill[navy] (10.858,5.934) circle (0.75pt);
\fill[navy] (10.858,6.006) circle (0.75pt);
\fill[navy] (10.858,6.077) circle (0.75pt);
\fill[navy] (10.858,6.149) circle (0.75pt);
\fill[navy] (10.858,6.22) circle (0.75pt);
\fill[navy] (10.98,0.071) circle (0.75pt);
\fill[navy] (10.98,0.786) circle (0.75pt);
\fill[navy] (10.98,1.644) circle (0.75pt);
\fill[navy] (10.98,2.645) circle (0.75pt);
\fill[navy] (10.98,2.931) circle (0.75pt);
\fill[navy] (10.98,3.36) circle (0.75pt);
\fill[navy] (10.98,4.218) circle (0.75pt);
\fill[navy] (10.98,4.79) circle (0.75pt);
\fill[navy] (10.98,5.076) circle (0.75pt);
\fill[navy] (10.98,5.648) circle (0.75pt);
\fill[navy] (10.98,5.934) circle (0.75pt);
\fill[navy] (10.98,6.363) circle (0.75pt);
\fill[navy] (11.102,0.286) circle (0.75pt);
\fill[navy] (11.102,0.643) circle (0.75pt);
\fill[navy] (11.102,0.715) circle (0.75pt);
\fill[navy] (11.102,0.786) circle (0.75pt);
\fill[navy] (11.102,0.858) circle (0.75pt);
\fill[navy] (11.102,1.073) circle (0.75pt);
\fill[navy] (11.102,1.215) circle (0.75pt);
\fill[navy] (11.102,1.573) circle (0.75pt);
\fill[navy] (11.102,1.644) circle (0.75pt);
\fill[navy] (11.102,1.716) circle (0.75pt);
\fill[navy] (11.102,1.787) circle (0.75pt);
\fill[navy] (11.102,2.073) circle (0.75pt);
\fill[navy] (11.102,2.145) circle (0.75pt);
\fill[navy] (11.102,2.288) circle (0.75pt);
\fill[navy] (11.102,2.431) circle (0.75pt);
\fill[navy] (11.102,3.217) circle (0.75pt);
\fill[navy] (11.102,3.289) circle (0.75pt);
\fill[navy] (11.102,3.36) circle (0.75pt);
\fill[navy] (11.102,3.432) circle (0.75pt);
\fill[navy] (11.102,3.646) circle (0.75pt);
\fill[navy] (11.102,3.789) circle (0.75pt);
\fill[navy] (11.102,3.861) circle (0.75pt);
\fill[navy] (11.102,3.932) circle (0.75pt);
\fill[navy] (11.102,4.075) circle (0.75pt);
\fill[navy] (11.102,4.433) circle (0.75pt);
\fill[navy] (11.102,4.576) circle (0.75pt);
\fill[navy] (11.102,4.719) circle (0.75pt);
\fill[navy] (11.102,4.79) circle (0.75pt);
\fill[navy] (11.102,4.862) circle (0.75pt);
\fill[navy] (11.102,4.933) circle (0.75pt);
\fill[navy] (11.102,5.148) circle (0.75pt);
\fill[navy] (11.102,5.362) circle (0.75pt);
\fill[navy] (11.102,5.72) circle (0.75pt);
\fill[navy] (11.102,5.791) circle (0.75pt);
\fill[navy] (11.102,5.863) circle (0.75pt);
\fill[navy] (11.102,5.934) circle (0.75pt);
\fill[navy] (11.102,6.077) circle (0.75pt);
\fill[navy] (11.102,6.149) circle (0.75pt);
\fill[navy] (11.102,6.22) circle (0.75pt);
\fill[navy] (11.102,6.363) circle (0.75pt);
\fill[navy] (11.224,0.214) circle (0.75pt);
\fill[navy] (11.224,0.786) circle (0.75pt);
\fill[navy] (11.224,1.073) circle (0.75pt);
\fill[navy] (11.224,1.501) circle (0.75pt);
\fill[navy] (11.224,2.073) circle (0.75pt);
\fill[navy] (11.224,2.359) circle (0.75pt);
\fill[navy] (11.224,2.788) circle (0.75pt);
\fill[navy] (11.224,3.36) circle (0.75pt);
\fill[navy] (11.224,3.503) circle (0.75pt);
\fill[navy] (11.224,3.789) circle (0.75pt);
\fill[navy] (11.224,4.218) circle (0.75pt);
\fill[navy] (11.224,4.361) circle (0.75pt);
\fill[navy] (11.224,4.504) circle (0.75pt);
\fill[navy] (11.224,4.79) circle (0.75pt);
\fill[navy] (11.224,5.076) circle (0.75pt);
\fill[navy] (11.224,5.219) circle (0.75pt);
\fill[navy] (11.224,5.362) circle (0.75pt);
\fill[navy] (11.224,5.505) circle (0.75pt);
\fill[navy] (11.224,5.648) circle (0.75pt);
\fill[navy] (11.224,5.791) circle (0.75pt);
\fill[navy] (11.224,5.934) circle (0.75pt);
\fill[navy] (11.224,6.363) circle (0.75pt);
\fill[navy] (11.346,0.357) circle (0.75pt);
\fill[navy] (11.346,0.715) circle (0.75pt);
\fill[navy] (11.346,0.786) circle (0.75pt);
\fill[navy] (11.346,1.001) circle (0.75pt);
\fill[navy] (11.346,1.144) circle (0.75pt);
\fill[navy] (11.346,1.43) circle (0.75pt);
\fill[navy] (11.346,1.573) circle (0.75pt);
\fill[navy] (11.346,1.787) circle (0.75pt);
\fill[navy] (11.346,1.859) circle (0.75pt);
\fill[navy] (11.346,2.073) circle (0.75pt);
\fill[navy] (11.346,2.288) circle (0.75pt);
\fill[navy] (11.346,2.431) circle (0.75pt);
\fill[navy] (11.346,2.645) circle (0.75pt);
\fill[navy] (11.346,2.717) circle (0.75pt);
\fill[navy] (11.346,2.86) circle (0.75pt);
\fill[navy] (11.346,3.074) circle (0.75pt);
\fill[navy] (11.346,3.289) circle (0.75pt);
\fill[navy] (11.346,3.36) circle (0.75pt);
\fill[navy] (11.346,3.503) circle (0.75pt);
\fill[navy] (11.346,3.575) circle (0.75pt);
\fill[navy] (11.346,3.718) circle (0.75pt);
\fill[navy] (11.346,4.147) circle (0.75pt);
\fill[navy] (11.346,4.218) circle (0.75pt);
\fill[navy] (11.346,4.647) circle (0.75pt);
\fill[navy] (11.346,5.005) circle (0.75pt);
\fill[navy] (11.346,5.291) circle (0.75pt);
\fill[navy] (11.346,5.434) circle (0.75pt);
\fill[navy] (11.346,5.505) circle (0.75pt);
\fill[navy] (11.346,5.863) circle (0.75pt);
\fill[navy] (11.346,5.934) circle (0.75pt);
\fill[navy] (11.346,6.149) circle (0.75pt);
\fill[navy] (11.346,6.292) circle (0.75pt);
\fill[navy] (11.346,6.363) circle (0.75pt);
\fill[navy] (11.468,0.786) circle (0.75pt);
\fill[navy] (11.468,1.215) circle (0.75pt);
\fill[navy] (11.468,1.358) circle (0.75pt);
\fill[navy] (11.468,2.359) circle (0.75pt);
\fill[navy] (11.468,2.502) circle (0.75pt);
\fill[navy] (11.468,2.645) circle (0.75pt);
\fill[navy] (11.468,2.788) circle (0.75pt);
\fill[navy] (11.468,3.074) circle (0.75pt);
\fill[navy] (11.468,3.217) circle (0.75pt);
\fill[navy] (11.468,3.646) circle (0.75pt);
\fill[navy] (11.468,3.932) circle (0.75pt);
\fill[navy] (11.468,4.075) circle (0.75pt);
\fill[navy] (11.468,4.218) circle (0.75pt);
\fill[navy] (11.468,4.361) circle (0.75pt);
\fill[navy] (11.468,4.933) circle (0.75pt);
\fill[navy] (11.468,5.076) circle (0.75pt);
\fill[navy] (11.468,5.219) circle (0.75pt);
\fill[navy] (11.468,5.362) circle (0.75pt);
\fill[navy] (11.468,5.648) circle (0.75pt);
\fill[navy] (11.468,5.934) circle (0.75pt);
\fill[navy] (11.468,6.077) circle (0.75pt);
\fill[navy] (11.468,6.22) circle (0.75pt);
\fill[navy] (11.468,6.363) circle (0.75pt);
\fill[navy] (11.468,6.506) circle (0.75pt);
\fill[navy] (11.59,0.214) circle (0.75pt);
\fill[navy] (11.59,0.429) circle (0.75pt);
\fill[navy] (11.59,0.5) circle (0.75pt);
\fill[navy] (11.59,0.929) circle (0.75pt);
\fill[navy] (11.59,1.001) circle (0.75pt);
\fill[navy] (11.59,1.501) circle (0.75pt);
\fill[navy] (11.59,1.573) circle (0.75pt);
\fill[navy] (11.59,1.859) circle (0.75pt);
\fill[navy] (11.59,1.93) circle (0.75pt);
\fill[navy] (11.59,2.216) circle (0.75pt);
\fill[navy] (11.59,2.288) circle (0.75pt);
\fill[navy] (11.59,2.359) circle (0.75pt);
\fill[navy] (11.59,2.788) circle (0.75pt);
\fill[navy] (11.59,2.931) circle (0.75pt);
\fill[navy] (11.59,3.074) circle (0.75pt);
\fill[navy] (11.59,3.432) circle (0.75pt);
\fill[navy] (11.59,3.503) circle (0.75pt);
\fill[navy] (11.59,3.646) circle (0.75pt);
\fill[navy] (11.59,4.004) circle (0.75pt);
\fill[navy] (11.59,4.504) circle (0.75pt);
\fill[navy] (11.59,4.576) circle (0.75pt);
\fill[navy] (11.59,4.79) circle (0.75pt);
\fill[navy] (11.59,5.148) circle (0.75pt);
\fill[navy] (11.59,5.219) circle (0.75pt);
\fill[navy] (11.59,5.291) circle (0.75pt);
\fill[navy] (11.59,5.505) circle (0.75pt);
\fill[navy] (11.59,5.577) circle (0.75pt);
\fill[navy] (11.59,5.648) circle (0.75pt);
\fill[navy] (11.59,5.863) circle (0.75pt);
\fill[navy] (11.59,6.006) circle (0.75pt);
\fill[navy] (11.59,6.149) circle (0.75pt);
\fill[navy] (11.59,6.22) circle (0.75pt);
\fill[navy] (11.59,6.292) circle (0.75pt);
\fill[navy] (11.59,6.506) circle (0.75pt);
\fill[navy] (11.59,6.578) circle (0.75pt);
\fill[navy] (11.59,6.649) circle (0.75pt);
\fill[navy] (11.712,0.786) circle (0.75pt);
\fill[navy] (11.712,1.215) circle (0.75pt);
\fill[navy] (11.712,2.073) circle (0.75pt);
\fill[navy] (11.712,3.503) circle (0.75pt);
\fill[navy] (11.712,3.932) circle (0.75pt);
\fill[navy] (11.712,4.218) circle (0.75pt);
\fill[navy] (11.712,4.361) circle (0.75pt);
\fill[navy] (11.712,4.647) circle (0.75pt);
\fill[navy] (11.712,5.076) circle (0.75pt);
\fill[navy] (11.712,5.219) circle (0.75pt);
\fill[navy] (11.712,5.505) circle (0.75pt);
\fill[navy] (11.712,5.648) circle (0.75pt);
\fill[navy] (11.712,5.934) circle (0.75pt);
\fill[navy] (11.712,6.077) circle (0.75pt);
\fill[navy] (11.712,6.363) circle (0.75pt);
\fill[navy] (11.712,6.506) circle (0.75pt);
\fill[navy] (11.834,0.214) circle (0.75pt);
\fill[navy] (11.834,0.643) circle (0.75pt);
\fill[navy] (11.834,0.858) circle (0.75pt);
\fill[navy] (11.834,1.93) circle (0.75pt);
\fill[navy] (11.834,2.145) circle (0.75pt);
\fill[navy] (11.834,2.288) circle (0.75pt);
\fill[navy] (11.834,2.502) circle (0.75pt);
\fill[navy] (11.834,2.717) circle (0.75pt);
\fill[navy] (11.834,2.788) circle (0.75pt);
\fill[navy] (11.834,2.86) circle (0.75pt);
\fill[navy] (11.834,3.003) circle (0.75pt);
\fill[navy] (11.834,3.074) circle (0.75pt);
\fill[navy] (11.834,3.146) circle (0.75pt);
\fill[navy] (11.834,3.289) circle (0.75pt);
\fill[navy] (11.834,3.503) circle (0.75pt);
\fill[navy] (11.834,3.646) circle (0.75pt);
\fill[navy] (11.834,3.718) circle (0.75pt);
\fill[navy] (11.834,3.789) circle (0.75pt);
\fill[navy] (11.834,3.932) circle (0.75pt);
\fill[navy] (11.834,4.004) circle (0.75pt);
\fill[navy] (11.834,4.075) circle (0.75pt);
\fill[navy] (11.834,4.147) circle (0.75pt);
\fill[navy] (11.834,4.218) circle (0.75pt);
\fill[navy] (11.834,4.29) circle (0.75pt);
\fill[navy] (11.834,4.361) circle (0.75pt);
\fill[navy] (11.834,4.433) circle (0.75pt);
\fill[navy] (11.834,4.504) circle (0.75pt);
\fill[navy] (11.834,4.576) circle (0.75pt);
\fill[navy] (11.834,4.647) circle (0.75pt);
\fill[navy] (11.834,4.79) circle (0.75pt);
\fill[navy] (11.834,4.862) circle (0.75pt);
\fill[navy] (11.834,4.933) circle (0.75pt);
\fill[navy] (11.834,5.005) circle (0.75pt);
\fill[navy] (11.834,5.076) circle (0.75pt);
\fill[navy] (11.834,5.148) circle (0.75pt);
\fill[navy] (11.834,5.291) circle (0.75pt);
\fill[navy] (11.834,5.577) circle (0.75pt);
\fill[navy] (11.834,5.648) circle (0.75pt);
\fill[navy] (11.834,5.72) circle (0.75pt);
\fill[navy] (11.834,5.791) circle (0.75pt);
\fill[navy] (11.834,5.863) circle (0.75pt);
\fill[navy] (11.834,6.006) circle (0.75pt);
\fill[navy] (11.834,6.077) circle (0.75pt);
\fill[navy] (11.834,6.149) circle (0.75pt);
\fill[navy] (11.834,6.22) circle (0.75pt);
\fill[navy] (11.834,6.363) circle (0.75pt);
\fill[navy] (11.834,6.435) circle (0.75pt);
\fill[navy] (11.834,6.578) circle (0.75pt);
\fill[navy] (11.834,6.721) circle (0.75pt);
\fill[navy] (11.834,6.864) circle (0.75pt);
\fill[navy] (11.956,0.357) circle (0.75pt);
\fill[navy] (11.956,0.643) circle (0.75pt);
\fill[navy] (11.956,0.786) circle (0.75pt);
\fill[navy] (11.956,0.929) circle (0.75pt);
\fill[navy] (11.956,1.787) circle (0.75pt);
\fill[navy] (11.956,1.93) circle (0.75pt);
\fill[navy] (11.956,2.359) circle (0.75pt);
\fill[navy] (11.956,2.645) circle (0.75pt);
\fill[navy] (11.956,2.931) circle (0.75pt);
\fill[navy] (11.956,3.217) circle (0.75pt);
\fill[navy] (11.956,3.646) circle (0.75pt);
\fill[navy] (11.956,3.789) circle (0.75pt);
\fill[navy] (11.956,4.075) circle (0.75pt);
\fill[navy] (11.956,4.79) circle (0.75pt);
\fill[navy] (11.956,4.933) circle (0.75pt);
\fill[navy] (11.956,5.076) circle (0.75pt);
\fill[navy] (11.956,5.219) circle (0.75pt);
\fill[navy] (11.956,5.362) circle (0.75pt);
\fill[navy] (11.956,5.648) circle (0.75pt);
\fill[navy] (11.956,5.791) circle (0.75pt);
\fill[navy] (11.956,5.934) circle (0.75pt);
\fill[navy] (11.956,6.077) circle (0.75pt);
\fill[navy] (11.956,6.22) circle (0.75pt);
\fill[navy] (11.956,6.363) circle (0.75pt);
\fill[navy] (11.956,6.792) circle (0.75pt);
\fill[navy] (12.078,0.071) circle (0.75pt);
\fill[navy] (12.078,0.143) circle (0.75pt);
\fill[navy] (12.078,0.286) circle (0.75pt);
\fill[navy] (12.078,0.357) circle (0.75pt);
\fill[navy] (12.078,0.5) circle (0.75pt);
\fill[navy] (12.078,0.572) circle (0.75pt);
\fill[navy] (12.078,0.715) circle (0.75pt);
\fill[navy] (12.078,0.929) circle (0.75pt);
\fill[navy] (12.078,1.001) circle (0.75pt);
\fill[navy] (12.078,1.144) circle (0.75pt);
\fill[navy] (12.078,1.43) circle (0.75pt);
\fill[navy] (12.078,1.644) circle (0.75pt);
\fill[navy] (12.078,1.859) circle (0.75pt);
\fill[navy] (12.078,2.002) circle (0.75pt);
\fill[navy] (12.078,2.645) circle (0.75pt);
\fill[navy] (12.078,2.86) circle (0.75pt);
\fill[navy] (12.078,3.074) circle (0.75pt);
\fill[navy] (12.078,3.575) circle (0.75pt);
\fill[navy] (12.078,3.718) circle (0.75pt);
\fill[navy] (12.078,4.147) circle (0.75pt);
\fill[navy] (12.078,4.218) circle (0.75pt);
\fill[navy] (12.078,4.361) circle (0.75pt);
\fill[navy] (12.078,4.433) circle (0.75pt);
\fill[navy] (12.078,4.576) circle (0.75pt);
\fill[navy] (12.078,4.647) circle (0.75pt);
\fill[navy] (12.078,4.79) circle (0.75pt);
\fill[navy] (12.078,4.862) circle (0.75pt);
\fill[navy] (12.078,5.076) circle (0.75pt);
\fill[navy] (12.078,5.219) circle (0.75pt);
\fill[navy] (12.078,5.434) circle (0.75pt);
\fill[navy] (12.078,5.648) circle (0.75pt);
\fill[navy] (12.078,5.72) circle (0.75pt);
\fill[navy] (12.078,5.863) circle (0.75pt);
\fill[navy] (12.078,5.934) circle (0.75pt);
\fill[navy] (12.078,6.077) circle (0.75pt);
\fill[navy] (12.078,6.149) circle (0.75pt);
\fill[navy] (12.078,6.363) circle (0.75pt);
\fill[navy] (12.078,6.506) circle (0.75pt);
\fill[navy] (12.078,6.578) circle (0.75pt);
\fill[navy] (12.078,6.721) circle (0.75pt);
\fill[navy] (12.078,6.935) circle (0.75pt);
\fill[navy] (12.078,7.007) circle (0.75pt);
\fill[navy] (12.2,0.643) circle (0.75pt);
\fill[navy] (12.2,0.786) circle (0.75pt);
\fill[navy] (12.2,1.215) circle (0.75pt);
\fill[navy] (12.2,1.644) circle (0.75pt);
\fill[navy] (12.2,1.93) circle (0.75pt);
\fill[navy] (12.2,2.216) circle (0.75pt);
\fill[navy] (12.2,2.359) circle (0.75pt);
\fill[navy] (12.2,2.645) circle (0.75pt);
\fill[navy] (12.2,2.931) circle (0.75pt);
\fill[navy] (12.2,3.074) circle (0.75pt);
\fill[navy] (12.2,3.36) circle (0.75pt);
\fill[navy] (12.2,3.646) circle (0.75pt);
\fill[navy] (12.2,3.789) circle (0.75pt);
\fill[navy] (12.2,4.075) circle (0.75pt);
\fill[navy] (12.2,4.218) circle (0.75pt);
\fill[navy] (12.2,4.504) circle (0.75pt);
\fill[navy] (12.2,4.79) circle (0.75pt);
\fill[navy] (12.2,4.933) circle (0.75pt);
\fill[navy] (12.2,5.219) circle (0.75pt);
\fill[navy] (12.2,5.505) circle (0.75pt);
\fill[navy] (12.2,5.648) circle (0.75pt);
\fill[navy] (12.2,5.791) circle (0.75pt);
\fill[navy] (12.2,5.934) circle (0.75pt);
\fill[navy] (12.2,6.22) circle (0.75pt);
\fill[navy] (12.2,6.363) circle (0.75pt);
\fill[navy] (12.2,6.506) circle (0.75pt);
\fill[navy] (12.2,6.649) circle (0.75pt);
\fill[navy] (12.2,6.935) circle (0.75pt);
\fill[navy] (12.2,7.078) circle (0.75pt);
\fill[teal] (0.976,0.357) circle (1.05pt);
\fill[teal] (1.22,0.5) circle (1.05pt);
\fill[teal] (1.342,0.286) circle (1.05pt);
\fill[teal] (1.342,0.643) circle (1.05pt);
\fill[teal] (1.586,0.429) circle (1.05pt);
\fill[teal] (1.586,0.643) circle (1.05pt);
\fill[teal] (1.83,0.5) circle (1.05pt);
\fill[teal] (1.952,0.643) circle (1.05pt);
\fill[teal] (1.952,0.786) circle (1.05pt);
\fill[teal] (1.952,1.073) circle (1.05pt);
\fill[teal] (2.074,0.357) circle (1.05pt);
\fill[teal] (2.074,0.5) circle (1.05pt);
\fill[teal] (2.074,0.572) circle (1.05pt);
\fill[teal] (2.074,0.643) circle (1.05pt);
\fill[teal] (2.074,0.786) circle (1.05pt);
\fill[teal] (2.074,0.858) circle (1.05pt);
\fill[teal] (2.074,1.001) circle (1.05pt);
\fill[teal] (2.196,0.5) circle (1.05pt);
\fill[teal] (2.196,0.786) circle (1.05pt);
\fill[teal] (2.196,0.929) circle (1.05pt);
\fill[teal] (2.318,0.286) circle (1.05pt);
\fill[teal] (2.318,0.643) circle (1.05pt);
\fill[teal] (2.318,0.786) circle (1.05pt);
\fill[teal] (2.318,0.858) circle (1.05pt);
\fill[teal] (2.318,0.929) circle (1.05pt);
\fill[teal] (2.318,1.001) circle (1.05pt);
\fill[teal] (2.318,1.215) circle (1.05pt);
\fill[teal] (2.318,1.287) circle (1.05pt);
\fill[teal] (2.44,0.071) circle (1.05pt);
\fill[teal] (2.44,0.643) circle (1.05pt);
\fill[teal] (2.44,1.358) circle (1.05pt);
\fill[teal] (2.562,0.143) circle (1.05pt);
\fill[teal] (2.562,0.286) circle (1.05pt);
\fill[teal] (2.562,1.358) circle (1.05pt);
\fill[teal] (2.684,0.929) circle (1.05pt);
\fill[teal] (2.684,1.215) circle (1.05pt);
\fill[teal] (2.684,1.358) circle (1.05pt);
\fill[teal] (2.684,1.501) circle (1.05pt);
\fill[teal] (2.806,0.143) circle (1.05pt);
\fill[teal] (2.806,0.572) circle (1.05pt);
\fill[teal] (2.806,0.643) circle (1.05pt);
\fill[teal] (2.806,0.715) circle (1.05pt);
\fill[teal] (2.806,0.786) circle (1.05pt);
\fill[teal] (2.806,0.858) circle (1.05pt);
\fill[teal] (2.806,0.929) circle (1.05pt);
\fill[teal] (2.806,1.001) circle (1.05pt);
\fill[teal] (2.806,1.073) circle (1.05pt);
\fill[teal] (2.806,1.144) circle (1.05pt);
\fill[teal] (2.806,1.287) circle (1.05pt);
\fill[teal] (2.806,1.358) circle (1.05pt);
\fill[teal] (2.806,1.43) circle (1.05pt);
\fill[teal] (2.806,1.501) circle (1.05pt);
\fill[teal] (2.928,0.5) circle (1.05pt);
\fill[teal] (2.928,0.786) circle (1.05pt);
\fill[teal] (2.928,1.215) circle (1.05pt);
\fill[teal] (2.928,1.358) circle (1.05pt);
\fill[teal] (2.928,1.644) circle (1.05pt);
\fill[teal] (3.05,0.429) circle (1.05pt);
\fill[teal] (3.05,0.643) circle (1.05pt);
\fill[teal] (3.05,0.786) circle (1.05pt);
\fill[teal] (3.05,0.858) circle (1.05pt);
\fill[teal] (3.05,0.929) circle (1.05pt);
\fill[teal] (3.05,1.001) circle (1.05pt);
\fill[teal] (3.05,1.144) circle (1.05pt);
\fill[teal] (3.05,1.215) circle (1.05pt);
\fill[teal] (3.05,1.358) circle (1.05pt);
\fill[teal] (3.05,1.573) circle (1.05pt);
\fill[teal] (3.05,1.716) circle (1.05pt);
\fill[teal] (3.172,0.071) circle (1.05pt);
\fill[teal] (3.172,0.357) circle (1.05pt);
\fill[teal] (3.172,0.5) circle (1.05pt);
\fill[teal] (3.172,0.643) circle (1.05pt);
\fill[teal] (3.172,0.786) circle (1.05pt);
\fill[teal] (3.172,1.073) circle (1.05pt);
\fill[teal] (3.172,1.215) circle (1.05pt);
\fill[teal] (3.172,1.358) circle (1.05pt);
\fill[teal] (3.172,1.787) circle (1.05pt);
\fill[teal] (3.294,0.143) circle (1.05pt);
\fill[teal] (3.294,0.286) circle (1.05pt);
\fill[teal] (3.294,0.357) circle (1.05pt);
\fill[teal] (3.294,0.5) circle (1.05pt);
\fill[teal] (3.294,0.715) circle (1.05pt);
\fill[teal] (3.294,0.786) circle (1.05pt);
\fill[teal] (3.294,1.001) circle (1.05pt);
\fill[teal] (3.294,1.144) circle (1.05pt);
\fill[teal] (3.294,1.215) circle (1.05pt);
\fill[teal] (3.294,1.43) circle (1.05pt);
\fill[teal] (3.294,1.573) circle (1.05pt);
\fill[teal] (3.294,1.644) circle (1.05pt);
\fill[teal] (3.416,0.214) circle (1.05pt);
\fill[teal] (3.416,0.357) circle (1.05pt);
\fill[teal] (3.416,0.643) circle (1.05pt);
\fill[teal] (3.416,0.786) circle (1.05pt);
\fill[teal] (3.416,0.929) circle (1.05pt);
\fill[teal] (3.416,1.073) circle (1.05pt);
\fill[teal] (3.416,1.358) circle (1.05pt);
\fill[teal] (3.416,1.644) circle (1.05pt);
\fill[teal] (3.416,1.787) circle (1.05pt);
\fill[teal] (3.416,1.93) circle (1.05pt);
\fill[teal] (3.538,0.071) circle (1.05pt);
\fill[teal] (3.538,0.214) circle (1.05pt);
\fill[teal] (3.538,0.286) circle (1.05pt);
\fill[teal] (3.538,0.429) circle (1.05pt);
\fill[teal] (3.538,0.5) circle (1.05pt);
\fill[teal] (3.538,0.572) circle (1.05pt);
\fill[teal] (3.538,0.643) circle (1.05pt);
\fill[teal] (3.538,0.715) circle (1.05pt);
\fill[teal] (3.538,0.786) circle (1.05pt);
\fill[teal] (3.538,0.858) circle (1.05pt);
\fill[teal] (3.538,0.929) circle (1.05pt);
\fill[teal] (3.538,1.001) circle (1.05pt);
\fill[teal] (3.538,1.073) circle (1.05pt);
\fill[teal] (3.538,1.144) circle (1.05pt);
\fill[teal] (3.538,1.215) circle (1.05pt);
\fill[teal] (3.538,1.287) circle (1.05pt);
\fill[teal] (3.538,1.358) circle (1.05pt);
\fill[teal] (3.538,1.501) circle (1.05pt);
\fill[teal] (3.538,1.573) circle (1.05pt);
\fill[teal] (3.538,1.644) circle (1.05pt);
\fill[teal] (3.538,1.716) circle (1.05pt);
\fill[teal] (3.538,1.787) circle (1.05pt);
\fill[teal] (3.538,1.93) circle (1.05pt);
\fill[teal] (3.538,2.002) circle (1.05pt);
\fill[teal] (3.66,0.071) circle (1.05pt);
\fill[teal] (3.66,0.5) circle (1.05pt);
\fill[teal] (3.66,0.786) circle (1.05pt);
\fill[teal] (3.66,1.358) circle (1.05pt);
\fill[teal] (3.782,0.143) circle (1.05pt);
\fill[teal] (3.782,0.214) circle (1.05pt);
\fill[teal] (3.782,0.286) circle (1.05pt);
\fill[teal] (3.782,0.357) circle (1.05pt);
\fill[teal] (3.782,0.5) circle (1.05pt);
\fill[teal] (3.782,0.572) circle (1.05pt);
\fill[teal] (3.782,0.715) circle (1.05pt);
\fill[teal] (3.782,0.858) circle (1.05pt);
\fill[teal] (3.782,0.929) circle (1.05pt);
\fill[teal] (3.782,1.001) circle (1.05pt);
\fill[teal] (3.782,1.073) circle (1.05pt);
\fill[teal] (3.782,1.144) circle (1.05pt);
\fill[teal] (3.782,1.215) circle (1.05pt);
\fill[teal] (3.782,1.287) circle (1.05pt);
\fill[teal] (3.782,1.358) circle (1.05pt);
\fill[teal] (3.782,1.43) circle (1.05pt);
\fill[teal] (3.782,1.501) circle (1.05pt);
\fill[teal] (3.782,1.573) circle (1.05pt);
\fill[teal] (3.782,1.644) circle (1.05pt);
\fill[teal] (3.782,1.716) circle (1.05pt);
\fill[teal] (3.782,1.787) circle (1.05pt);
\fill[teal] (3.782,1.859) circle (1.05pt);
\fill[teal] (3.782,1.93) circle (1.05pt);
\fill[teal] (3.782,2.002) circle (1.05pt);
\fill[teal] (3.782,2.145) circle (1.05pt);
\fill[teal] (3.904,0.071) circle (1.05pt);
\fill[teal] (3.904,0.643) circle (1.05pt);
\fill[teal] (3.904,0.786) circle (1.05pt);
\fill[teal] (3.904,0.929) circle (1.05pt);
\fill[teal] (3.904,1.073) circle (1.05pt);
\fill[teal] (3.904,1.501) circle (1.05pt);
\fill[teal] (3.904,1.787) circle (1.05pt);
\fill[teal] (3.904,1.93) circle (1.05pt);
\fill[teal] (3.904,2.073) circle (1.05pt);
\fill[teal] (3.904,2.216) circle (1.05pt);
\fill[teal] (4.026,0.071) circle (1.05pt);
\fill[teal] (4.026,0.286) circle (1.05pt);
\fill[teal] (4.026,0.357) circle (1.05pt);
\fill[teal] (4.026,0.5) circle (1.05pt);
\fill[teal] (4.026,0.572) circle (1.05pt);
\fill[teal] (4.026,0.715) circle (1.05pt);
\fill[teal] (4.026,0.929) circle (1.05pt);
\fill[teal] (4.026,1.001) circle (1.05pt);
\fill[teal] (4.026,1.215) circle (1.05pt);
\fill[teal] (4.026,1.358) circle (1.05pt);
\fill[teal] (4.026,1.644) circle (1.05pt);
\fill[teal] (4.026,1.787) circle (1.05pt);
\fill[teal] (4.026,1.859) circle (1.05pt);
\fill[teal] (4.026,2.002) circle (1.05pt);
\fill[teal] (4.026,2.073) circle (1.05pt);
\fill[teal] (4.026,2.288) circle (1.05pt);
\fill[teal] (4.148,0.071) circle (1.05pt);
\fill[teal] (4.148,0.214) circle (1.05pt);
\fill[teal] (4.148,0.643) circle (1.05pt);
\fill[teal] (4.148,0.786) circle (1.05pt);
\fill[teal] (4.148,0.929) circle (1.05pt);
\fill[teal] (4.148,1.073) circle (1.05pt);
\fill[teal] (4.148,1.358) circle (1.05pt);
\fill[teal] (4.148,1.501) circle (1.05pt);
\fill[teal] (4.148,1.644) circle (1.05pt);
\fill[teal] (4.148,1.93) circle (1.05pt);
\fill[teal] (4.27,0.071) circle (1.05pt);
\fill[teal] (4.27,0.143) circle (1.05pt);
\fill[teal] (4.27,0.214) circle (1.05pt);
\fill[teal] (4.27,0.286) circle (1.05pt);
\fill[teal] (4.27,0.429) circle (1.05pt);
\fill[teal] (4.27,0.786) circle (1.05pt);
\fill[teal] (4.27,0.858) circle (1.05pt);
\fill[teal] (4.27,0.929) circle (1.05pt);
\fill[teal] (4.27,1.144) circle (1.05pt);
\fill[teal] (4.27,1.215) circle (1.05pt);
\fill[teal] (4.27,1.287) circle (1.05pt);
\fill[teal] (4.27,1.358) circle (1.05pt);
\fill[teal] (4.27,1.573) circle (1.05pt);
\fill[teal] (4.27,1.644) circle (1.05pt);
\fill[teal] (4.27,1.716) circle (1.05pt);
\fill[teal] (4.27,1.859) circle (1.05pt);
\fill[teal] (4.27,1.93) circle (1.05pt);
\fill[teal] (4.27,2.073) circle (1.05pt);
\fill[teal] (4.27,2.216) circle (1.05pt);
\fill[teal] (4.27,2.288) circle (1.05pt);
\fill[teal] (4.27,2.359) circle (1.05pt);
\fill[teal] (4.27,2.431) circle (1.05pt);
\fill[teal] (4.392,0.071) circle (1.05pt);
\fill[teal] (4.392,0.357) circle (1.05pt);
\fill[teal] (4.392,0.5) circle (1.05pt);
\fill[teal] (4.392,0.786) circle (1.05pt);
\fill[teal] (4.392,1.215) circle (1.05pt);
\fill[teal] (4.392,1.644) circle (1.05pt);
\fill[teal] (4.392,1.787) circle (1.05pt);
\fill[teal] (4.392,2.073) circle (1.05pt);
\fill[teal] (4.392,2.502) circle (1.05pt);
\fill[teal] (4.514,0.143) circle (1.05pt);
\fill[teal] (4.514,0.214) circle (1.05pt);
\fill[teal] (4.514,0.286) circle (1.05pt);
\fill[teal] (4.514,0.572) circle (1.05pt);
\fill[teal] (4.514,0.643) circle (1.05pt);
\fill[teal] (4.514,0.786) circle (1.05pt);
\fill[teal] (4.514,0.858) circle (1.05pt);
\fill[teal] (4.514,0.929) circle (1.05pt);
\fill[teal] (4.514,1.073) circle (1.05pt);
\fill[teal] (4.514,1.144) circle (1.05pt);
\fill[teal] (4.514,1.215) circle (1.05pt);
\fill[teal] (4.514,1.287) circle (1.05pt);
\fill[teal] (4.514,1.358) circle (1.05pt);
\fill[teal] (4.514,1.43) circle (1.05pt);
\fill[teal] (4.514,1.501) circle (1.05pt);
\fill[teal] (4.514,1.573) circle (1.05pt);
\fill[teal] (4.514,1.716) circle (1.05pt);
\fill[teal] (4.514,1.787) circle (1.05pt);
\fill[teal] (4.514,1.93) circle (1.05pt);
\fill[teal] (4.514,2.002) circle (1.05pt);
\fill[teal] (4.514,2.073) circle (1.05pt);
\fill[teal] (4.514,2.145) circle (1.05pt);
\fill[teal] (4.514,2.216) circle (1.05pt);
\fill[teal] (4.514,2.288) circle (1.05pt);
\fill[teal] (4.514,2.359) circle (1.05pt);
\fill[teal] (4.514,2.431) circle (1.05pt);
\fill[teal] (4.636,0.214) circle (1.05pt);
\fill[teal] (4.636,0.643) circle (1.05pt);
\fill[teal] (4.636,0.786) circle (1.05pt);
\fill[teal] (4.636,0.929) circle (1.05pt);
\fill[teal] (4.636,1.073) circle (1.05pt);
\fill[teal] (4.636,1.215) circle (1.05pt);
\fill[teal] (4.636,1.501) circle (1.05pt);
\fill[teal] (4.636,1.787) circle (1.05pt);
\fill[teal] (4.636,2.073) circle (1.05pt);
\fill[teal] (4.636,2.216) circle (1.05pt);
\fill[teal] (4.636,2.359) circle (1.05pt);
\fill[teal] (4.636,2.645) circle (1.05pt);
\fill[teal] (4.758,0.071) circle (1.05pt);
\fill[teal] (4.758,0.143) circle (1.05pt);
\fill[teal] (4.758,0.286) circle (1.05pt);
\fill[teal] (4.758,0.357) circle (1.05pt);
\fill[teal] (4.758,0.5) circle (1.05pt);
\fill[teal] (4.758,0.572) circle (1.05pt);
\fill[teal] (4.758,0.715) circle (1.05pt);
\fill[teal] (4.758,0.786) circle (1.05pt);
\fill[teal] (4.758,1.001) circle (1.05pt);
\fill[teal] (4.758,1.144) circle (1.05pt);
\fill[teal] (4.758,1.215) circle (1.05pt);
\fill[teal] (4.758,1.358) circle (1.05pt);
\fill[teal] (4.758,1.43) circle (1.05pt);
\fill[teal] (4.758,1.573) circle (1.05pt);
\fill[teal] (4.758,2.073) circle (1.05pt);
\fill[teal] (4.758,2.288) circle (1.05pt);
\fill[teal] (4.758,2.431) circle (1.05pt);
\fill[teal] (4.758,2.502) circle (1.05pt);
\fill[teal] (4.758,2.645) circle (1.05pt);
\fill[teal] (4.758,2.717) circle (1.05pt);
\fill[teal] (4.88,0.071) circle (1.05pt);
\fill[teal] (4.88,0.5) circle (1.05pt);
\fill[teal] (4.88,0.643) circle (1.05pt);
\fill[teal] (4.88,0.786) circle (1.05pt);
\fill[teal] (4.88,0.929) circle (1.05pt);
\fill[teal] (4.88,1.215) circle (1.05pt);
\fill[teal] (4.88,1.358) circle (1.05pt);
\fill[teal] (4.88,1.501) circle (1.05pt);
\fill[teal] (4.88,1.644) circle (1.05pt);
\fill[teal] (4.88,1.93) circle (1.05pt);
\fill[teal] (4.88,2.073) circle (1.05pt);
\fill[teal] (4.88,2.216) circle (1.05pt);
\fill[teal] (4.88,2.359) circle (1.05pt);
\fill[teal] (4.88,2.788) circle (1.05pt);
\fill[teal] (5.002,0.143) circle (1.05pt);
\fill[teal] (5.002,0.286) circle (1.05pt);
\fill[teal] (5.002,0.357) circle (1.05pt);
\fill[teal] (5.002,0.5) circle (1.05pt);
\fill[teal] (5.002,0.572) circle (1.05pt);
\fill[teal] (5.002,0.643) circle (1.05pt);
\fill[teal] (5.002,0.715) circle (1.05pt);
\fill[teal] (5.002,0.786) circle (1.05pt);
\fill[teal] (5.002,0.858) circle (1.05pt);
\fill[teal] (5.002,0.929) circle (1.05pt);
\fill[teal] (5.002,1.073) circle (1.05pt);
\fill[teal] (5.002,1.144) circle (1.05pt);
\fill[teal] (5.002,1.215) circle (1.05pt);
\fill[teal] (5.002,1.287) circle (1.05pt);
\fill[teal] (5.002,1.358) circle (1.05pt);
\fill[teal] (5.002,1.43) circle (1.05pt);
\fill[teal] (5.002,1.501) circle (1.05pt);
\fill[teal] (5.002,1.573) circle (1.05pt);
\fill[teal] (5.002,1.644) circle (1.05pt);
\fill[teal] (5.002,1.716) circle (1.05pt);
\fill[teal] (5.002,1.787) circle (1.05pt);
\fill[teal] (5.002,2.002) circle (1.05pt);
\fill[teal] (5.002,2.145) circle (1.05pt);
\fill[teal] (5.002,2.359) circle (1.05pt);
\fill[teal] (5.002,2.431) circle (1.05pt);
\fill[teal] (5.002,2.502) circle (1.05pt);
\fill[teal] (5.002,2.574) circle (1.05pt);
\fill[teal] (5.002,2.645) circle (1.05pt);
\fill[teal] (5.002,2.717) circle (1.05pt);
\fill[teal] (5.002,2.788) circle (1.05pt);
\fill[teal] (5.002,2.86) circle (1.05pt);
\fill[teal] (5.124,0.071) circle (1.05pt);
\fill[teal] (5.124,0.929) circle (1.05pt);
\fill[teal] (5.124,1.215) circle (1.05pt);
\fill[teal] (5.124,1.358) circle (1.05pt);
\fill[teal] (5.124,1.644) circle (1.05pt);
\fill[teal] (5.124,2.073) circle (1.05pt);
\fill[teal] (5.124,2.645) circle (1.05pt);
\fill[teal] (5.246,0.071) circle (1.05pt);
\fill[teal] (5.246,0.143) circle (1.05pt);
\fill[teal] (5.246,0.214) circle (1.05pt);
\fill[teal] (5.246,0.286) circle (1.05pt);
\fill[teal] (5.246,0.429) circle (1.05pt);
\fill[teal] (5.246,0.643) circle (1.05pt);
\fill[teal] (5.246,0.786) circle (1.05pt);
\fill[teal] (5.246,0.929) circle (1.05pt);
\fill[teal] (5.246,1.001) circle (1.05pt);
\fill[teal] (5.246,1.073) circle (1.05pt);
\fill[teal] (5.246,1.144) circle (1.05pt);
\fill[teal] (5.246,1.358) circle (1.05pt);
\fill[teal] (5.246,1.43) circle (1.05pt);
\fill[teal] (5.246,1.501) circle (1.05pt);
\fill[teal] (5.246,1.644) circle (1.05pt);
\fill[teal] (5.246,1.716) circle (1.05pt);
\fill[teal] (5.246,1.787) circle (1.05pt);
\fill[teal] (5.246,1.93) circle (1.05pt);
\fill[teal] (5.246,2.002) circle (1.05pt);
\fill[teal] (5.246,2.073) circle (1.05pt);
\fill[teal] (5.246,2.145) circle (1.05pt);
\fill[teal] (5.246,2.216) circle (1.05pt);
\fill[teal] (5.246,2.359) circle (1.05pt);
\fill[teal] (5.246,2.431) circle (1.05pt);
\fill[teal] (5.246,2.502) circle (1.05pt);
\fill[teal] (5.246,2.574) circle (1.05pt);
\fill[teal] (5.246,2.645) circle (1.05pt);
\fill[teal] (5.246,2.788) circle (1.05pt);
\fill[teal] (5.246,2.86) circle (1.05pt);
\fill[teal] (5.246,2.931) circle (1.05pt);
\fill[teal] (5.246,3.003) circle (1.05pt);
\fill[teal] (5.368,0.071) circle (1.05pt);
\fill[teal] (5.368,0.214) circle (1.05pt);
\fill[teal] (5.368,0.5) circle (1.05pt);
\fill[teal] (5.368,0.643) circle (1.05pt);
\fill[teal] (5.368,0.929) circle (1.05pt);
\fill[teal] (5.368,1.073) circle (1.05pt);
\fill[teal] (5.368,1.215) circle (1.05pt);
\fill[teal] (5.368,1.358) circle (1.05pt);
\fill[teal] (5.368,1.501) circle (1.05pt);
\fill[teal] (5.368,1.644) circle (1.05pt);
\fill[teal] (5.368,1.787) circle (1.05pt);
\fill[teal] (5.368,1.93) circle (1.05pt);
\fill[teal] (5.368,2.073) circle (1.05pt);
\fill[teal] (5.368,2.216) circle (1.05pt);
\fill[teal] (5.368,2.645) circle (1.05pt);
\fill[teal] (5.368,2.931) circle (1.05pt);
\fill[teal] (5.368,3.074) circle (1.05pt);
\fill[teal] (5.49,0.071) circle (1.05pt);
\fill[teal] (5.49,0.143) circle (1.05pt);
\fill[teal] (5.49,0.286) circle (1.05pt);
\fill[teal] (5.49,0.5) circle (1.05pt);
\fill[teal] (5.49,0.572) circle (1.05pt);
\fill[teal] (5.49,0.929) circle (1.05pt);
\fill[teal] (5.49,1.001) circle (1.05pt);
\fill[teal] (5.49,1.215) circle (1.05pt);
\fill[teal] (5.49,1.358) circle (1.05pt);
\fill[teal] (5.49,1.573) circle (1.05pt);
\fill[teal] (5.49,1.644) circle (1.05pt);
\fill[teal] (5.49,1.859) circle (1.05pt);
\fill[teal] (5.49,2.002) circle (1.05pt);
\fill[teal] (5.49,2.073) circle (1.05pt);
\fill[teal] (5.49,2.288) circle (1.05pt);
\fill[teal] (5.49,2.431) circle (1.05pt);
\fill[teal] (5.49,2.645) circle (1.05pt);
\fill[teal] (5.49,3.074) circle (1.05pt);
\fill[teal] (5.612,0.071) circle (1.05pt);
\fill[teal] (5.612,0.214) circle (1.05pt);
\fill[teal] (5.612,0.357) circle (1.05pt);
\fill[teal] (5.612,0.5) circle (1.05pt);
\fill[teal] (5.612,0.786) circle (1.05pt);
\fill[teal] (5.612,0.929) circle (1.05pt);
\fill[teal] (5.612,1.073) circle (1.05pt);
\fill[teal] (5.612,1.215) circle (1.05pt);
\fill[teal] (5.612,1.358) circle (1.05pt);
\fill[teal] (5.612,1.501) circle (1.05pt);
\fill[teal] (5.612,1.787) circle (1.05pt);
\fill[teal] (5.612,1.93) circle (1.05pt);
\fill[teal] (5.612,2.073) circle (1.05pt);
\fill[teal] (5.612,2.359) circle (1.05pt);
\fill[teal] (5.612,2.645) circle (1.05pt);
\fill[teal] (5.612,2.788) circle (1.05pt);
\fill[teal] (5.612,2.931) circle (1.05pt);
\fill[teal] (5.612,3.217) circle (1.05pt);
\fill[teal] (5.734,0.071) circle (1.05pt);
\fill[teal] (5.734,0.143) circle (1.05pt);
\fill[teal] (5.734,0.286) circle (1.05pt);
\fill[teal] (5.734,0.357) circle (1.05pt);
\fill[teal] (5.734,0.429) circle (1.05pt);
\fill[teal] (5.734,0.643) circle (1.05pt);
\fill[teal] (5.734,0.715) circle (1.05pt);
\fill[teal] (5.734,0.858) circle (1.05pt);
\fill[teal] (5.734,0.929) circle (1.05pt);
\fill[teal] (5.734,1.001) circle (1.05pt);
\fill[teal] (5.734,1.073) circle (1.05pt);
\fill[teal] (5.734,1.144) circle (1.05pt);
\fill[teal] (5.734,1.215) circle (1.05pt);
\fill[teal] (5.734,1.358) circle (1.05pt);
\fill[teal] (5.734,1.43) circle (1.05pt);
\fill[teal] (5.734,1.501) circle (1.05pt);
\fill[teal] (5.734,1.573) circle (1.05pt);
\fill[teal] (5.734,1.644) circle (1.05pt);
\fill[teal] (5.734,1.716) circle (1.05pt);
\fill[teal] (5.734,1.787) circle (1.05pt);
\fill[teal] (5.734,1.859) circle (1.05pt);
\fill[teal] (5.734,1.93) circle (1.05pt);
\fill[teal] (5.734,2.073) circle (1.05pt);
\fill[teal] (5.734,2.216) circle (1.05pt);
\fill[teal] (5.734,2.288) circle (1.05pt);
\fill[teal] (5.734,2.359) circle (1.05pt);
\fill[teal] (5.734,2.431) circle (1.05pt);
\fill[teal] (5.734,2.502) circle (1.05pt);
\fill[teal] (5.734,2.574) circle (1.05pt);
\fill[teal] (5.734,2.931) circle (1.05pt);
\fill[teal] (5.734,3.003) circle (1.05pt);
\fill[teal] (5.734,3.074) circle (1.05pt);
\fill[teal] (5.734,3.146) circle (1.05pt);
\fill[teal] (5.734,3.289) circle (1.05pt);
\fill[teal] (5.856,0.071) circle (1.05pt);
\fill[teal] (5.856,0.357) circle (1.05pt);
\fill[teal] (5.856,0.5) circle (1.05pt);
\fill[teal] (5.856,0.786) circle (1.05pt);
\fill[teal] (5.856,0.929) circle (1.05pt);
\fill[teal] (5.856,1.215) circle (1.05pt);
\fill[teal] (5.856,1.358) circle (1.05pt);
\fill[teal] (5.856,1.644) circle (1.05pt);
\fill[teal] (5.856,2.073) circle (1.05pt);
\fill[teal] (5.856,2.216) circle (1.05pt);
\fill[teal] (5.856,2.502) circle (1.05pt);
\fill[teal] (5.856,2.645) circle (1.05pt);
\fill[teal] (5.856,2.931) circle (1.05pt);
\fill[teal] (5.856,3.074) circle (1.05pt);
\fill[teal] (5.856,3.36) circle (1.05pt);
\fill[teal] (5.978,0.071) circle (1.05pt);
\fill[teal] (5.978,0.143) circle (1.05pt);
\fill[teal] (5.978,0.214) circle (1.05pt);
\fill[teal] (5.978,0.286) circle (1.05pt);
\fill[teal] (5.978,0.357) circle (1.05pt);
\fill[teal] (5.978,0.429) circle (1.05pt);
\fill[teal] (5.978,0.572) circle (1.05pt);
\fill[teal] (5.978,0.643) circle (1.05pt);
\fill[teal] (5.978,0.715) circle (1.05pt);
\fill[teal] (5.978,0.786) circle (1.05pt);
\fill[teal] (5.978,0.858) circle (1.05pt);
\fill[teal] (5.978,0.929) circle (1.05pt);
\fill[teal] (5.978,1.073) circle (1.05pt);
\fill[teal] (5.978,1.144) circle (1.05pt);
\fill[teal] (5.978,1.215) circle (1.05pt);
\fill[teal] (5.978,1.287) circle (1.05pt);
\fill[teal] (5.978,1.358) circle (1.05pt);
\fill[teal] (5.978,1.43) circle (1.05pt);
\fill[teal] (5.978,1.573) circle (1.05pt);
\fill[teal] (5.978,1.644) circle (1.05pt);
\fill[teal] (5.978,1.716) circle (1.05pt);
\fill[teal] (5.978,1.787) circle (1.05pt);
\fill[teal] (5.978,1.859) circle (1.05pt);
\fill[teal] (5.978,1.93) circle (1.05pt);
\fill[teal] (5.978,2.073) circle (1.05pt);
\fill[teal] (5.978,2.145) circle (1.05pt);
\fill[teal] (5.978,2.216) circle (1.05pt);
\fill[teal] (5.978,2.288) circle (1.05pt);
\fill[teal] (5.978,2.359) circle (1.05pt);
\fill[teal] (5.978,2.431) circle (1.05pt);
\fill[teal] (5.978,2.574) circle (1.05pt);
\fill[teal] (5.978,2.645) circle (1.05pt);
\fill[teal] (5.978,2.717) circle (1.05pt);
\fill[teal] (5.978,2.86) circle (1.05pt);
\fill[teal] (5.978,3.074) circle (1.05pt);
\fill[teal] (6.1,0.071) circle (1.05pt);
\fill[teal] (6.1,0.5) circle (1.05pt);
\fill[teal] (6.1,0.643) circle (1.05pt);
\fill[teal] (6.1,0.786) circle (1.05pt);
\fill[teal] (6.1,0.929) circle (1.05pt);
\fill[teal] (6.1,1.358) circle (1.05pt);
\fill[teal] (6.1,1.501) circle (1.05pt);
\fill[teal] (6.1,1.644) circle (1.05pt);
\fill[teal] (6.1,1.93) circle (1.05pt);
\fill[teal] (6.1,2.216) circle (1.05pt);
\fill[teal] (6.1,2.359) circle (1.05pt);
\fill[teal] (6.1,2.645) circle (1.05pt);
\fill[teal] (6.1,2.788) circle (1.05pt);
\fill[teal] (6.1,2.931) circle (1.05pt);
\fill[teal] (6.1,3.36) circle (1.05pt);
\fill[teal] (6.1,3.503) circle (1.05pt);
\fill[teal] (6.222,0.071) circle (1.05pt);
\fill[teal] (6.222,0.143) circle (1.05pt);
\fill[teal] (6.222,0.286) circle (1.05pt);
\fill[teal] (6.222,0.357) circle (1.05pt);
\fill[teal] (6.222,0.5) circle (1.05pt);
\fill[teal] (6.222,0.572) circle (1.05pt);
\fill[teal] (6.222,0.786) circle (1.05pt);
\fill[teal] (6.222,0.929) circle (1.05pt);
\fill[teal] (6.222,1.001) circle (1.05pt);
\fill[teal] (6.222,1.144) circle (1.05pt);
\fill[teal] (6.222,1.358) circle (1.05pt);
\fill[teal] (6.222,1.573) circle (1.05pt);
\fill[teal] (6.222,1.644) circle (1.05pt);
\fill[teal] (6.222,1.859) circle (1.05pt);
\fill[teal] (6.222,2.002) circle (1.05pt);
\fill[teal] (6.222,2.216) circle (1.05pt);
\fill[teal] (6.222,2.288) circle (1.05pt);
\fill[teal] (6.222,2.502) circle (1.05pt);
\fill[teal] (6.222,2.645) circle (1.05pt);
\fill[teal] (6.222,2.717) circle (1.05pt);
\fill[teal] (6.222,2.86) circle (1.05pt);
\fill[teal] (6.222,3.074) circle (1.05pt);
\fill[teal] (6.222,3.146) circle (1.05pt);
\fill[teal] (6.222,3.289) circle (1.05pt);
\fill[teal] (6.222,3.36) circle (1.05pt);
\fill[teal] (6.222,3.503) circle (1.05pt);
\fill[teal] (6.344,0.071) circle (1.05pt);
\fill[teal] (6.344,0.214) circle (1.05pt);
\fill[teal] (6.344,0.5) circle (1.05pt);
\fill[teal] (6.344,0.643) circle (1.05pt);
\fill[teal] (6.344,0.786) circle (1.05pt);
\fill[teal] (6.344,1.073) circle (1.05pt);
\fill[teal] (6.344,1.215) circle (1.05pt);
\fill[teal] (6.344,1.358) circle (1.05pt);
\fill[teal] (6.344,1.501) circle (1.05pt);
\fill[teal] (6.344,1.644) circle (1.05pt);
\fill[teal] (6.344,1.787) circle (1.05pt);
\fill[teal] (6.344,2.073) circle (1.05pt);
\fill[teal] (6.344,2.216) circle (1.05pt);
\fill[teal] (6.344,2.359) circle (1.05pt);
\fill[teal] (6.344,2.502) circle (1.05pt);
\fill[teal] (6.344,2.931) circle (1.05pt);
\fill[teal] (6.344,3.074) circle (1.05pt);
\fill[teal] (6.344,3.217) circle (1.05pt);
\fill[teal] (6.344,3.36) circle (1.05pt);
\fill[teal] (6.344,3.503) circle (1.05pt);
\fill[teal] (6.466,0.071) circle (1.05pt);
\fill[teal] (6.466,0.143) circle (1.05pt);
\fill[teal] (6.466,0.214) circle (1.05pt);
\fill[teal] (6.466,0.286) circle (1.05pt);
\fill[teal] (6.466,0.429) circle (1.05pt);
\fill[teal] (6.466,0.5) circle (1.05pt);
\fill[teal] (6.466,0.572) circle (1.05pt);
\fill[teal] (6.466,0.715) circle (1.05pt);
\fill[teal] (6.466,0.786) circle (1.05pt);
\fill[teal] (6.466,1.001) circle (1.05pt);
\fill[teal] (6.466,1.215) circle (1.05pt);
\fill[teal] (6.466,1.287) circle (1.05pt);
\fill[teal] (6.466,1.43) circle (1.05pt);
\fill[teal] (6.466,1.501) circle (1.05pt);
\fill[teal] (6.466,1.573) circle (1.05pt);
\fill[teal] (6.466,1.644) circle (1.05pt);
\fill[teal] (6.466,1.716) circle (1.05pt);
\fill[teal] (6.466,1.859) circle (1.05pt);
\fill[teal] (6.466,2.002) circle (1.05pt);
\fill[teal] (6.466,2.073) circle (1.05pt);
\fill[teal] (6.466,2.145) circle (1.05pt);
\fill[teal] (6.466,2.216) circle (1.05pt);
\fill[teal] (6.466,2.288) circle (1.05pt);
\fill[teal] (6.466,2.359) circle (1.05pt);
\fill[teal] (6.466,2.502) circle (1.05pt);
\fill[teal] (6.466,2.574) circle (1.05pt);
\fill[teal] (6.466,2.645) circle (1.05pt);
\fill[teal] (6.466,2.717) circle (1.05pt);
\fill[teal] (6.466,2.788) circle (1.05pt);
\fill[teal] (6.466,2.86) circle (1.05pt);
\fill[teal] (6.466,2.931) circle (1.05pt);
\fill[teal] (6.466,3.074) circle (1.05pt);
\fill[teal] (6.466,3.146) circle (1.05pt);
\fill[teal] (6.466,3.217) circle (1.05pt);
\fill[teal] (6.466,3.289) circle (1.05pt);
\fill[teal] (6.466,3.36) circle (1.05pt);
\fill[teal] (6.466,3.432) circle (1.05pt);
\fill[teal] (6.466,3.575) circle (1.05pt);
\fill[teal] (6.466,3.646) circle (1.05pt);
\fill[teal] (6.466,3.718) circle (1.05pt);
\fill[teal] (6.588,0.357) circle (1.05pt);
\fill[teal] (6.588,0.5) circle (1.05pt);
\fill[teal] (6.588,0.786) circle (1.05pt);
\fill[teal] (6.588,0.929) circle (1.05pt);
\fill[teal] (6.588,1.215) circle (1.05pt);
\fill[teal] (6.588,1.644) circle (1.05pt);
\fill[teal] (6.588,1.787) circle (1.05pt);
\fill[teal] (6.588,2.073) circle (1.05pt);
\fill[teal] (6.588,2.216) circle (1.05pt);
\fill[teal] (6.588,2.645) circle (1.05pt);
\fill[teal] (6.588,2.931) circle (1.05pt);
\fill[teal] (6.588,3.36) circle (1.05pt);
\fill[teal] (6.588,3.503) circle (1.05pt);
\fill[teal] (6.588,3.789) circle (1.05pt);
\fill[teal] (6.71,0.143) circle (1.05pt);
\fill[teal] (6.71,0.214) circle (1.05pt);
\fill[teal] (6.71,0.286) circle (1.05pt);
\fill[teal] (6.71,0.429) circle (1.05pt);
\fill[teal] (6.71,0.5) circle (1.05pt);
\fill[teal] (6.71,0.643) circle (1.05pt);
\fill[teal] (6.71,0.858) circle (1.05pt);
\fill[teal] (6.71,1.001) circle (1.05pt);
\fill[teal] (6.71,1.144) circle (1.05pt);
\fill[teal] (6.71,1.215) circle (1.05pt);
\fill[teal] (6.71,1.287) circle (1.05pt);
\fill[teal] (6.71,1.358) circle (1.05pt);
\fill[teal] (6.71,1.644) circle (1.05pt);
\fill[teal] (6.71,1.716) circle (1.05pt);
\fill[teal] (6.71,1.859) circle (1.05pt);
\fill[teal] (6.71,1.93) circle (1.05pt);
\fill[teal] (6.71,2.002) circle (1.05pt);
\fill[teal] (6.71,2.073) circle (1.05pt);
\fill[teal] (6.71,2.574) circle (1.05pt);
\fill[teal] (6.71,2.645) circle (1.05pt);
\fill[teal] (6.71,2.717) circle (1.05pt);
\fill[teal] (6.71,2.788) circle (1.05pt);
\fill[teal] (6.71,2.931) circle (1.05pt);
\fill[teal] (6.71,3.289) circle (1.05pt);
\fill[teal] (6.71,3.36) circle (1.05pt);
\fill[teal] (6.71,3.432) circle (1.05pt);
\fill[teal] (6.71,3.718) circle (1.05pt);
\fill[teal] (6.71,3.789) circle (1.05pt);
\fill[teal] (6.71,3.861) circle (1.05pt);
\fill[teal] (6.832,0.214) circle (1.05pt);
\fill[teal] (6.832,0.357) circle (1.05pt);
\fill[teal] (6.832,0.786) circle (1.05pt);
\fill[teal] (6.832,0.929) circle (1.05pt);
\fill[teal] (6.832,1.073) circle (1.05pt);
\fill[teal] (6.832,1.215) circle (1.05pt);
\fill[teal] (6.832,1.358) circle (1.05pt);
\fill[teal] (6.832,1.644) circle (1.05pt);
\fill[teal] (6.832,1.93) circle (1.05pt);
\fill[teal] (6.832,2.073) circle (1.05pt);
\fill[teal] (6.832,2.216) circle (1.05pt);
\fill[teal] (6.832,2.359) circle (1.05pt);
\fill[teal] (6.832,2.645) circle (1.05pt);
\fill[teal] (6.832,2.788) circle (1.05pt);
\fill[teal] (6.832,3.217) circle (1.05pt);
\fill[teal] (6.832,3.36) circle (1.05pt);
\fill[teal] (6.832,3.646) circle (1.05pt);
\fill[teal] (6.832,3.932) circle (1.05pt);
\fill[teal] (6.954,0.071) circle (1.05pt);
\fill[teal] (6.954,0.143) circle (1.05pt);
\fill[teal] (6.954,0.286) circle (1.05pt);
\fill[teal] (6.954,0.357) circle (1.05pt);
\fill[teal] (6.954,0.5) circle (1.05pt);
\fill[teal] (6.954,0.572) circle (1.05pt);
\fill[teal] (6.954,0.715) circle (1.05pt);
\fill[teal] (6.954,0.929) circle (1.05pt);
\fill[teal] (6.954,1.144) circle (1.05pt);
\fill[teal] (6.954,1.215) circle (1.05pt);
\fill[teal] (6.954,1.43) circle (1.05pt);
\fill[teal] (6.954,1.573) circle (1.05pt);
\fill[teal] (6.954,1.787) circle (1.05pt);
\fill[teal] (6.954,1.859) circle (1.05pt);
\fill[teal] (6.954,2.002) circle (1.05pt);
\fill[teal] (6.954,2.216) circle (1.05pt);
\fill[teal] (6.954,2.288) circle (1.05pt);
\fill[teal] (6.954,2.431) circle (1.05pt);
\fill[teal] (6.954,2.502) circle (1.05pt);
\fill[teal] (6.954,2.86) circle (1.05pt);
\fill[teal] (6.954,2.931) circle (1.05pt);
\fill[teal] (6.954,3.074) circle (1.05pt);
\fill[teal] (6.954,3.289) circle (1.05pt);
\fill[teal] (6.954,3.36) circle (1.05pt);
\fill[teal] (6.954,3.503) circle (1.05pt);
\fill[teal] (6.954,3.575) circle (1.05pt);
\fill[teal] (6.954,3.789) circle (1.05pt);
\fill[teal] (6.954,3.932) circle (1.05pt);
\fill[teal] (6.954,4.004) circle (1.05pt);
\fill[teal] (7.076,0.071) circle (1.05pt);
\fill[teal] (7.076,0.214) circle (1.05pt);
\fill[teal] (7.076,0.357) circle (1.05pt);
\fill[teal] (7.076,0.5) circle (1.05pt);
\fill[teal] (7.076,0.786) circle (1.05pt);
\fill[teal] (7.076,1.073) circle (1.05pt);
\fill[teal] (7.076,1.215) circle (1.05pt);
\fill[teal] (7.076,1.501) circle (1.05pt);
\fill[teal] (7.076,1.644) circle (1.05pt);
\fill[teal] (7.076,1.93) circle (1.05pt);
\fill[teal] (7.076,2.216) circle (1.05pt);
\fill[teal] (7.076,2.359) circle (1.05pt);
\fill[teal] (7.076,2.645) circle (1.05pt);
\fill[teal] (7.076,2.788) circle (1.05pt);
\fill[teal] (7.076,3.074) circle (1.05pt);
\fill[teal] (7.076,3.217) circle (1.05pt);
\fill[teal] (7.076,3.36) circle (1.05pt);
\fill[teal] (7.076,3.503) circle (1.05pt);
\fill[teal] (7.076,3.646) circle (1.05pt);
\fill[teal] (7.076,4.075) circle (1.05pt);
\fill[teal] (7.198,0.071) circle (1.05pt);
\fill[teal] (7.198,0.143) circle (1.05pt);
\fill[teal] (7.198,0.214) circle (1.05pt);
\fill[teal] (7.198,0.286) circle (1.05pt);
\fill[teal] (7.198,0.357) circle (1.05pt);
\fill[teal] (7.198,0.429) circle (1.05pt);
\fill[teal] (7.198,0.5) circle (1.05pt);
\fill[teal] (7.198,0.572) circle (1.05pt);
\fill[teal] (7.198,0.643) circle (1.05pt);
\fill[teal] (7.198,0.715) circle (1.05pt);
\fill[teal] (7.198,0.786) circle (1.05pt);
\fill[teal] (7.198,0.858) circle (1.05pt);
\fill[teal] (7.198,1.001) circle (1.05pt);
\fill[teal] (7.198,1.073) circle (1.05pt);
\fill[teal] (7.198,1.215) circle (1.05pt);
\fill[teal] (7.198,1.287) circle (1.05pt);
\fill[teal] (7.198,1.501) circle (1.05pt);
\fill[teal] (7.198,1.573) circle (1.05pt);
\fill[teal] (7.198,1.644) circle (1.05pt);
\fill[teal] (7.198,1.716) circle (1.05pt);
\fill[teal] (7.198,1.859) circle (1.05pt);
\fill[teal] (7.198,1.93) circle (1.05pt);
\fill[teal] (7.198,2.002) circle (1.05pt);
\fill[teal] (7.198,2.145) circle (1.05pt);
\fill[teal] (7.198,2.216) circle (1.05pt);
\fill[teal] (7.198,2.359) circle (1.05pt);
\fill[teal] (7.198,2.431) circle (1.05pt);
\fill[teal] (7.198,2.502) circle (1.05pt);
\fill[teal] (7.198,2.574) circle (1.05pt);
\fill[teal] (7.198,2.645) circle (1.05pt);
\fill[teal] (7.198,2.86) circle (1.05pt);
\fill[teal] (7.198,2.931) circle (1.05pt);
\fill[teal] (7.198,3.003) circle (1.05pt);
\fill[teal] (7.198,3.074) circle (1.05pt);
\fill[teal] (7.198,3.289) circle (1.05pt);
\fill[teal] (7.198,3.503) circle (1.05pt);
\fill[teal] (7.198,3.575) circle (1.05pt);
\fill[teal] (7.198,3.718) circle (1.05pt);
\fill[teal] (7.198,3.789) circle (1.05pt);
\fill[teal] (7.198,3.861) circle (1.05pt);
\fill[teal] (7.198,4.075) circle (1.05pt);
\fill[teal] (7.198,4.147) circle (1.05pt);
\fill[teal] (7.32,0.071) circle (1.05pt);
\fill[teal] (7.32,0.5) circle (1.05pt);
\fill[teal] (7.32,0.786) circle (1.05pt);
\fill[teal] (7.32,0.929) circle (1.05pt);
\fill[teal] (7.32,1.215) circle (1.05pt);
\fill[teal] (7.32,1.358) circle (1.05pt);
\fill[teal] (7.32,1.644) circle (1.05pt);
\fill[teal] (7.32,2.073) circle (1.05pt);
\fill[teal] (7.32,2.931) circle (1.05pt);
\fill[teal] (7.32,3.074) circle (1.05pt);
\fill[teal] (7.32,3.503) circle (1.05pt);
\fill[teal] (7.32,4.218) circle (1.05pt);
\fill[teal] (7.442,0.071) circle (1.05pt);
\fill[teal] (7.442,0.143) circle (1.05pt);
\fill[teal] (7.442,0.214) circle (1.05pt);
\fill[teal] (7.442,0.429) circle (1.05pt);
\fill[teal] (7.442,0.5) circle (1.05pt);
\fill[teal] (7.442,0.572) circle (1.05pt);
\fill[teal] (7.442,0.786) circle (1.05pt);
\fill[teal] (7.442,0.858) circle (1.05pt);
\fill[teal] (7.442,0.929) circle (1.05pt);
\fill[teal] (7.442,1.001) circle (1.05pt);
\fill[teal] (7.442,1.073) circle (1.05pt);
\fill[teal] (7.442,1.215) circle (1.05pt);
\fill[teal] (7.442,1.358) circle (1.05pt);
\fill[teal] (7.442,1.43) circle (1.05pt);
\fill[teal] (7.442,1.501) circle (1.05pt);
\fill[teal] (7.442,1.573) circle (1.05pt);
\fill[teal] (7.442,1.716) circle (1.05pt);
\fill[teal] (7.442,1.859) circle (1.05pt);
\fill[teal] (7.442,2.002) circle (1.05pt);
\fill[teal] (7.442,2.288) circle (1.05pt);
\fill[teal] (7.442,2.431) circle (1.05pt);
\fill[teal] (7.442,2.502) circle (1.05pt);
\fill[teal] (7.442,2.717) circle (1.05pt);
\fill[teal] (7.442,2.788) circle (1.05pt);
\fill[teal] (7.442,2.931) circle (1.05pt);
\fill[teal] (7.442,3.074) circle (1.05pt);
\fill[teal] (7.442,3.146) circle (1.05pt);
\fill[teal] (7.442,3.217) circle (1.05pt);
\fill[teal] (7.442,3.289) circle (1.05pt);
\fill[teal] (7.442,3.575) circle (1.05pt);
\fill[teal] (7.442,3.718) circle (1.05pt);
\fill[teal] (7.442,3.932) circle (1.05pt);
\fill[teal] (7.442,4.004) circle (1.05pt);
\fill[teal] (7.442,4.075) circle (1.05pt);
\fill[teal] (7.442,4.147) circle (1.05pt);
\fill[teal] (7.442,4.218) circle (1.05pt);
\fill[teal] (7.442,4.29) circle (1.05pt);
\fill[teal] (7.564,0.071) circle (1.05pt);
\fill[teal] (7.564,0.214) circle (1.05pt);
\fill[teal] (7.564,0.357) circle (1.05pt);
\fill[teal] (7.564,0.5) circle (1.05pt);
\fill[teal] (7.564,1.073) circle (1.05pt);
\fill[teal] (7.564,1.215) circle (1.05pt);
\fill[teal] (7.564,1.358) circle (1.05pt);
\fill[teal] (7.564,1.501) circle (1.05pt);
\fill[teal] (7.564,1.787) circle (1.05pt);
\fill[teal] (7.564,1.93) circle (1.05pt);
\fill[teal] (7.564,2.073) circle (1.05pt);
\fill[teal] (7.564,2.502) circle (1.05pt);
\fill[teal] (7.564,2.788) circle (1.05pt);
\fill[teal] (7.564,3.074) circle (1.05pt);
\fill[teal] (7.564,3.36) circle (1.05pt);
\fill[teal] (7.564,3.789) circle (1.05pt);
\fill[teal] (7.564,3.932) circle (1.05pt);
\fill[teal] (7.564,4.075) circle (1.05pt);
\fill[teal] (7.564,4.218) circle (1.05pt);
\fill[teal] (7.564,4.361) circle (1.05pt);
\fill[teal] (7.686,0.143) circle (1.05pt);
\fill[teal] (7.686,0.286) circle (1.05pt);
\fill[teal] (7.686,0.357) circle (1.05pt);
\fill[teal] (7.686,0.572) circle (1.05pt);
\fill[teal] (7.686,0.786) circle (1.05pt);
\fill[teal] (7.686,0.929) circle (1.05pt);
\fill[teal] (7.686,1.215) circle (1.05pt);
\fill[teal] (7.686,1.358) circle (1.05pt);
\fill[teal] (7.686,1.43) circle (1.05pt);
\fill[teal] (7.686,1.573) circle (1.05pt);
\fill[teal] (7.686,2.073) circle (1.05pt);
\fill[teal] (7.686,2.216) circle (1.05pt);
\fill[teal] (7.686,2.288) circle (1.05pt);
\fill[teal] (7.686,2.645) circle (1.05pt);
\fill[teal] (7.686,2.717) circle (1.05pt);
\fill[teal] (7.686,2.86) circle (1.05pt);
\fill[teal] (7.686,2.931) circle (1.05pt);
\fill[teal] (7.686,3.146) circle (1.05pt);
\fill[teal] (7.686,3.289) circle (1.05pt);
\fill[teal] (7.686,3.36) circle (1.05pt);
\fill[teal] (7.686,3.575) circle (1.05pt);
\fill[teal] (7.686,3.789) circle (1.05pt);
\fill[teal] (7.686,3.932) circle (1.05pt);
\fill[teal] (7.686,4.147) circle (1.05pt);
\fill[teal] (7.686,4.433) circle (1.05pt);
\fill[teal] (7.808,0.071) circle (1.05pt);
\fill[teal] (7.808,0.214) circle (1.05pt);
\fill[teal] (7.808,0.357) circle (1.05pt);
\fill[teal] (7.808,0.643) circle (1.05pt);
\fill[teal] (7.808,0.786) circle (1.05pt);
\fill[teal] (7.808,1.215) circle (1.05pt);
\fill[teal] (7.808,1.501) circle (1.05pt);
\fill[teal] (7.808,1.644) circle (1.05pt);
\fill[teal] (7.808,1.93) circle (1.05pt);
\fill[teal] (7.808,2.359) circle (1.05pt);
\fill[teal] (7.808,2.502) circle (1.05pt);
\fill[teal] (7.808,2.788) circle (1.05pt);
\fill[teal] (7.808,3.789) circle (1.05pt);
\fill[teal] (7.808,3.932) circle (1.05pt);
\fill[teal] (7.808,4.075) circle (1.05pt);
\fill[teal] (7.808,4.218) circle (1.05pt);
\fill[teal] (7.808,4.361) circle (1.05pt);
\fill[teal] (7.93,0.143) circle (1.05pt);
\fill[teal] (7.93,0.214) circle (1.05pt);
\fill[teal] (7.93,0.429) circle (1.05pt);
\fill[teal] (7.93,0.5) circle (1.05pt);
\fill[teal] (7.93,0.572) circle (1.05pt);
\fill[teal] (7.93,0.643) circle (1.05pt);
\fill[teal] (7.93,1.001) circle (1.05pt);
\fill[teal] (7.93,1.287) circle (1.05pt);
\fill[teal] (7.93,1.501) circle (1.05pt);
\fill[teal] (7.93,1.716) circle (1.05pt);
\fill[teal] (7.93,1.93) circle (1.05pt);
\fill[teal] (7.93,2.002) circle (1.05pt);
\fill[teal] (7.93,2.073) circle (1.05pt);
\fill[teal] (7.93,2.216) circle (1.05pt);
\fill[teal] (7.93,2.288) circle (1.05pt);
\fill[teal] (7.93,2.431) circle (1.05pt);
\fill[teal] (7.93,2.717) circle (1.05pt);
\fill[teal] (7.93,2.931) circle (1.05pt);
\fill[teal] (7.93,3.003) circle (1.05pt);
\fill[teal] (7.93,3.074) circle (1.05pt);
\fill[teal] (7.93,3.146) circle (1.05pt);
\fill[teal] (7.93,3.503) circle (1.05pt);
\fill[teal] (7.93,3.646) circle (1.05pt);
\fill[teal] (7.93,3.789) circle (1.05pt);
\fill[teal] (7.93,3.861) circle (1.05pt);
\fill[teal] (7.93,4.147) circle (1.05pt);
\fill[teal] (7.93,4.576) circle (1.05pt);
\fill[teal] (8.052,0.071) circle (1.05pt);
\fill[teal] (8.052,0.5) circle (1.05pt);
\fill[teal] (8.052,0.929) circle (1.05pt);
\fill[teal] (8.052,1.215) circle (1.05pt);
\fill[teal] (8.052,1.358) circle (1.05pt);
\fill[teal] (8.052,1.787) circle (1.05pt);
\fill[teal] (8.052,2.073) circle (1.05pt);
\fill[teal] (8.052,2.216) circle (1.05pt);
\fill[teal] (8.052,2.502) circle (1.05pt);
\fill[teal] (8.052,2.645) circle (1.05pt);
\fill[teal] (8.052,3.074) circle (1.05pt);
\fill[teal] (8.052,3.36) circle (1.05pt);
\fill[teal] (8.052,4.361) circle (1.05pt);
\fill[teal] (8.174,0.071) circle (1.05pt);
\fill[teal] (8.174,0.143) circle (1.05pt);
\fill[teal] (8.174,0.214) circle (1.05pt);
\fill[teal] (8.174,0.286) circle (1.05pt);
\fill[teal] (8.174,0.357) circle (1.05pt);
\fill[teal] (8.174,0.5) circle (1.05pt);
\fill[teal] (8.174,0.572) circle (1.05pt);
\fill[teal] (8.174,0.643) circle (1.05pt);
\fill[teal] (8.174,0.715) circle (1.05pt);
\fill[teal] (8.174,0.786) circle (1.05pt);
\fill[teal] (8.174,0.858) circle (1.05pt);
\fill[teal] (8.174,0.929) circle (1.05pt);
\fill[teal] (8.174,1.073) circle (1.05pt);
\fill[teal] (8.174,1.144) circle (1.05pt);
\fill[teal] (8.174,1.358) circle (1.05pt);
\fill[teal] (8.174,1.43) circle (1.05pt);
\fill[teal] (8.174,1.501) circle (1.05pt);
\fill[teal] (8.174,1.644) circle (1.05pt);
\fill[teal] (8.174,1.716) circle (1.05pt);
\fill[teal] (8.174,1.787) circle (1.05pt);
\fill[teal] (8.174,1.859) circle (1.05pt);
\fill[teal] (8.174,1.93) circle (1.05pt);
\fill[teal] (8.174,2.002) circle (1.05pt);
\fill[teal] (8.174,2.073) circle (1.05pt);
\fill[teal] (8.174,2.145) circle (1.05pt);
\fill[teal] (8.174,2.216) circle (1.05pt);
\fill[teal] (8.174,2.431) circle (1.05pt);
\fill[teal] (8.174,2.502) circle (1.05pt);
\fill[teal] (8.174,2.574) circle (1.05pt);
\fill[teal] (8.174,2.645) circle (1.05pt);
\fill[teal] (8.174,2.717) circle (1.05pt);
\fill[teal] (8.174,2.86) circle (1.05pt);
\fill[teal] (8.174,3.003) circle (1.05pt);
\fill[teal] (8.174,3.146) circle (1.05pt);
\fill[teal] (8.174,3.217) circle (1.05pt);
\fill[teal] (8.174,3.289) circle (1.05pt);
\fill[teal] (8.174,3.36) circle (1.05pt);
\fill[teal] (8.174,3.432) circle (1.05pt);
\fill[teal] (8.174,3.503) circle (1.05pt);
\fill[teal] (8.174,3.575) circle (1.05pt);
\fill[teal] (8.174,3.718) circle (1.05pt);
\fill[teal] (8.174,3.789) circle (1.05pt);
\fill[teal] (8.174,3.861) circle (1.05pt);
\fill[teal] (8.174,4.004) circle (1.05pt);
\fill[teal] (8.174,4.075) circle (1.05pt);
\fill[teal] (8.174,4.147) circle (1.05pt);
\fill[teal] (8.174,4.218) circle (1.05pt);
\fill[teal] (8.174,4.29) circle (1.05pt);
\fill[teal] (8.174,4.433) circle (1.05pt);
\fill[teal] (8.174,4.576) circle (1.05pt);
\fill[teal] (8.174,4.719) circle (1.05pt);
\fill[teal] (8.296,0.214) circle (1.05pt);
\fill[teal] (8.296,0.357) circle (1.05pt);
\fill[teal] (8.296,0.643) circle (1.05pt);
\fill[teal] (8.296,0.786) circle (1.05pt);
\fill[teal] (8.296,0.929) circle (1.05pt);
\fill[teal] (8.296,1.358) circle (1.05pt);
\fill[teal] (8.296,1.501) circle (1.05pt);
\fill[teal] (8.296,2.073) circle (1.05pt);
\fill[teal] (8.296,2.216) circle (1.05pt);
\fill[teal] (8.296,2.359) circle (1.05pt);
\fill[teal] (8.296,2.502) circle (1.05pt);
\fill[teal] (8.296,2.931) circle (1.05pt);
\fill[teal] (8.296,3.217) circle (1.05pt);
\fill[teal] (8.296,3.36) circle (1.05pt);
\fill[teal] (8.296,3.503) circle (1.05pt);
\fill[teal] (8.296,3.789) circle (1.05pt);
\fill[teal] (8.296,3.932) circle (1.05pt);
\fill[teal] (8.296,4.075) circle (1.05pt);
\fill[teal] (8.296,4.361) circle (1.05pt);
\fill[teal] (8.296,4.504) circle (1.05pt);
\fill[teal] (8.296,4.79) circle (1.05pt);
\fill[teal] (8.418,0.071) circle (1.05pt);
\fill[teal] (8.418,0.143) circle (1.05pt);
\fill[teal] (8.418,0.286) circle (1.05pt);
\fill[teal] (8.418,0.357) circle (1.05pt);
\fill[teal] (8.418,0.5) circle (1.05pt);
\fill[teal] (8.418,0.715) circle (1.05pt);
\fill[teal] (8.418,0.786) circle (1.05pt);
\fill[teal] (8.418,0.929) circle (1.05pt);
\fill[teal] (8.418,1.001) circle (1.05pt);
\fill[teal] (8.418,1.144) circle (1.05pt);
\fill[teal] (8.418,1.215) circle (1.05pt);
\fill[teal] (8.418,1.859) circle (1.05pt);
\fill[teal] (8.418,2.002) circle (1.05pt);
\fill[teal] (8.418,2.216) circle (1.05pt);
\fill[teal] (8.418,2.288) circle (1.05pt);
\fill[teal] (8.418,2.431) circle (1.05pt);
\fill[teal] (8.418,2.502) circle (1.05pt);
\fill[teal] (8.418,2.645) circle (1.05pt);
\fill[teal] (8.418,2.86) circle (1.05pt);
\fill[teal] (8.418,3.074) circle (1.05pt);
\fill[teal] (8.418,3.146) circle (1.05pt);
\fill[teal] (8.418,3.36) circle (1.05pt);
\fill[teal] (8.418,3.575) circle (1.05pt);
\fill[teal] (8.418,3.718) circle (1.05pt);
\fill[teal] (8.418,3.789) circle (1.05pt);
\fill[teal] (8.418,3.932) circle (1.05pt);
\fill[teal] (8.418,4.004) circle (1.05pt);
\fill[teal] (8.418,4.218) circle (1.05pt);
\fill[teal] (8.418,4.433) circle (1.05pt);
\fill[teal] (8.418,4.576) circle (1.05pt);
\fill[teal] (8.418,4.79) circle (1.05pt);
\fill[teal] (8.54,0.071) circle (1.05pt);
\fill[teal] (8.54,0.214) circle (1.05pt);
\fill[teal] (8.54,0.643) circle (1.05pt);
\fill[teal] (8.54,0.786) circle (1.05pt);
\fill[teal] (8.54,0.929) circle (1.05pt);
\fill[teal] (8.54,1.358) circle (1.05pt);
\fill[teal] (8.54,2.216) circle (1.05pt);
\fill[teal] (8.54,2.359) circle (1.05pt);
\fill[teal] (8.54,2.788) circle (1.05pt);
\fill[teal] (8.54,3.074) circle (1.05pt);
\fill[teal] (8.54,3.646) circle (1.05pt);
\fill[teal] (8.54,3.789) circle (1.05pt);
\fill[teal] (8.54,4.79) circle (1.05pt);
\fill[teal] (8.662,0.071) circle (1.05pt);
\fill[teal] (8.662,0.429) circle (1.05pt);
\fill[teal] (8.662,0.5) circle (1.05pt);
\fill[teal] (8.662,0.572) circle (1.05pt);
\fill[teal] (8.662,0.643) circle (1.05pt);
\fill[teal] (8.662,0.715) circle (1.05pt);
\fill[teal] (8.662,0.858) circle (1.05pt);
\fill[teal] (8.662,1.001) circle (1.05pt);
\fill[teal] (8.662,1.073) circle (1.05pt);
\fill[teal] (8.662,1.144) circle (1.05pt);
\fill[teal] (8.662,1.215) circle (1.05pt);
\fill[teal] (8.662,1.287) circle (1.05pt);
\fill[teal] (8.662,1.358) circle (1.05pt);
\fill[teal] (8.662,1.573) circle (1.05pt);
\fill[teal] (8.662,1.644) circle (1.05pt);
\fill[teal] (8.662,1.716) circle (1.05pt);
\fill[teal] (8.662,1.787) circle (1.05pt);
\fill[teal] (8.662,1.859) circle (1.05pt);
\fill[teal] (8.662,1.93) circle (1.05pt);
\fill[teal] (8.662,2.002) circle (1.05pt);
\fill[teal] (8.662,2.073) circle (1.05pt);
\fill[teal] (8.662,2.145) circle (1.05pt);
\fill[teal] (8.662,2.216) circle (1.05pt);
\fill[teal] (8.662,2.288) circle (1.05pt);
\fill[teal] (8.662,2.359) circle (1.05pt);
\fill[teal] (8.662,2.502) circle (1.05pt);
\fill[teal] (8.662,2.574) circle (1.05pt);
\fill[teal] (8.662,2.645) circle (1.05pt);
\fill[teal] (8.662,2.717) circle (1.05pt);
\fill[teal] (8.662,2.788) circle (1.05pt);
\fill[teal] (8.662,2.86) circle (1.05pt);
\fill[teal] (8.662,3.003) circle (1.05pt);
\fill[teal] (8.662,3.074) circle (1.05pt);
\fill[teal] (8.662,3.146) circle (1.05pt);
\fill[teal] (8.662,3.217) circle (1.05pt);
\fill[teal] (8.662,3.289) circle (1.05pt);
\fill[teal] (8.662,3.36) circle (1.05pt);
\fill[teal] (8.662,3.432) circle (1.05pt);
\fill[teal] (8.662,3.503) circle (1.05pt);
\fill[teal] (8.662,3.575) circle (1.05pt);
\fill[teal] (8.662,3.646) circle (1.05pt);
\fill[teal] (8.662,4.075) circle (1.05pt);
\fill[teal] (8.662,4.147) circle (1.05pt);
\fill[teal] (8.662,4.218) circle (1.05pt);
\fill[teal] (8.662,4.29) circle (1.05pt);
\fill[teal] (8.662,4.361) circle (1.05pt);
\fill[teal] (8.662,4.504) circle (1.05pt);
\fill[teal] (8.662,4.576) circle (1.05pt);
\fill[teal] (8.662,4.647) circle (1.05pt);
\fill[teal] (8.662,4.719) circle (1.05pt);
\fill[teal] (8.662,5.005) circle (1.05pt);
\fill[teal] (8.784,0.071) circle (1.05pt);
\fill[teal] (8.784,0.357) circle (1.05pt);
\fill[teal] (8.784,0.5) circle (1.05pt);
\fill[teal] (8.784,0.786) circle (1.05pt);
\fill[teal] (8.784,1.787) circle (1.05pt);
\fill[teal] (8.784,2.216) circle (1.05pt);
\fill[teal] (8.784,2.502) circle (1.05pt);
\fill[teal] (8.784,2.645) circle (1.05pt);
\fill[teal] (8.784,3.074) circle (1.05pt);
\fill[teal] (8.784,3.36) circle (1.05pt);
\fill[teal] (8.784,3.789) circle (1.05pt);
\fill[teal] (8.784,5.076) circle (1.05pt);
\fill[teal] (8.906,0.071) circle (1.05pt);
\fill[teal] (8.906,0.143) circle (1.05pt);
\fill[teal] (8.906,0.214) circle (1.05pt);
\fill[teal] (8.906,0.286) circle (1.05pt);
\fill[teal] (8.906,0.429) circle (1.05pt);
\fill[teal] (8.906,0.5) circle (1.05pt);
\fill[teal] (8.906,0.572) circle (1.05pt);
\fill[teal] (8.906,0.715) circle (1.05pt);
\fill[teal] (8.906,0.786) circle (1.05pt);
\fill[teal] (8.906,0.858) circle (1.05pt);
\fill[teal] (8.906,0.929) circle (1.05pt);
\fill[teal] (8.906,1.001) circle (1.05pt);
\fill[teal] (8.906,1.073) circle (1.05pt);
\fill[teal] (8.906,1.287) circle (1.05pt);
\fill[teal] (8.906,1.358) circle (1.05pt);
\fill[teal] (8.906,1.43) circle (1.05pt);
\fill[teal] (8.906,1.573) circle (1.05pt);
\fill[teal] (8.906,1.93) circle (1.05pt);
\fill[teal] (8.906,2.002) circle (1.05pt);
\fill[teal] (8.906,2.073) circle (1.05pt);
\fill[teal] (8.906,2.145) circle (1.05pt);
\fill[teal] (8.906,2.216) circle (1.05pt);
\fill[teal] (8.906,2.288) circle (1.05pt);
\fill[teal] (8.906,2.359) circle (1.05pt);
\fill[teal] (8.906,2.431) circle (1.05pt);
\fill[teal] (8.906,2.502) circle (1.05pt);
\fill[teal] (8.906,2.574) circle (1.05pt);
\fill[teal] (8.906,2.645) circle (1.05pt);
\fill[teal] (8.906,2.717) circle (1.05pt);
\fill[teal] (8.906,2.788) circle (1.05pt);
\fill[teal] (8.906,2.86) circle (1.05pt);
\fill[teal] (8.906,3.003) circle (1.05pt);
\fill[teal] (8.906,3.432) circle (1.05pt);
\fill[teal] (8.906,3.503) circle (1.05pt);
\fill[teal] (8.906,3.932) circle (1.05pt);
\fill[teal] (8.906,4.075) circle (1.05pt);
\fill[teal] (8.906,4.218) circle (1.05pt);
\fill[teal] (8.906,4.504) circle (1.05pt);
\fill[teal] (8.906,4.719) circle (1.05pt);
\fill[teal] (8.906,4.862) circle (1.05pt);
\fill[teal] (8.906,5.005) circle (1.05pt);
\fill[teal] (8.906,5.076) circle (1.05pt);
\fill[teal] (9.028,0.071) circle (1.05pt);
\fill[teal] (9.028,0.214) circle (1.05pt);
\fill[teal] (9.028,0.5) circle (1.05pt);
\fill[teal] (9.028,0.786) circle (1.05pt);
\fill[teal] (9.028,0.929) circle (1.05pt);
\fill[teal] (9.028,1.073) circle (1.05pt);
\fill[teal] (9.028,1.215) circle (1.05pt);
\fill[teal] (9.028,1.501) circle (1.05pt);
\fill[teal] (9.028,1.644) circle (1.05pt);
\fill[teal] (9.028,1.787) circle (1.05pt);
\fill[teal] (9.028,1.93) circle (1.05pt);
\fill[teal] (9.028,2.073) circle (1.05pt);
\fill[teal] (9.028,2.216) circle (1.05pt);
\fill[teal] (9.028,2.788) circle (1.05pt);
\fill[teal] (9.028,3.074) circle (1.05pt);
\fill[teal] (9.028,3.217) circle (1.05pt);
\fill[teal] (9.028,3.36) circle (1.05pt);
\fill[teal] (9.028,3.646) circle (1.05pt);
\fill[teal] (9.028,3.789) circle (1.05pt);
\fill[teal] (9.028,3.932) circle (1.05pt);
\fill[teal] (9.028,4.075) circle (1.05pt);
\fill[teal] (9.028,4.218) circle (1.05pt);
\fill[teal] (9.028,4.504) circle (1.05pt);
\fill[teal] (9.028,4.647) circle (1.05pt);
\fill[teal] (9.028,4.79) circle (1.05pt);
\fill[teal] (9.028,5.076) circle (1.05pt);
\fill[teal] (9.028,5.219) circle (1.05pt);
\fill[teal] (9.15,0.071) circle (1.05pt);
\fill[teal] (9.15,0.286) circle (1.05pt);
\fill[teal] (9.15,0.5) circle (1.05pt);
\fill[teal] (9.15,0.786) circle (1.05pt);
\fill[teal] (9.15,0.929) circle (1.05pt);
\fill[teal] (9.15,1.001) circle (1.05pt);
\fill[teal] (9.15,1.144) circle (1.05pt);
\fill[teal] (9.15,1.215) circle (1.05pt);
\fill[teal] (9.15,1.358) circle (1.05pt);
\fill[teal] (9.15,1.573) circle (1.05pt);
\fill[teal] (9.15,1.644) circle (1.05pt);
\fill[teal] (9.15,2.002) circle (1.05pt);
\fill[teal] (9.15,2.645) circle (1.05pt);
\fill[teal] (9.15,2.717) circle (1.05pt);
\fill[teal] (9.15,2.931) circle (1.05pt);
\fill[teal] (9.15,3.146) circle (1.05pt);
\fill[teal] (9.15,3.503) circle (1.05pt);
\fill[teal] (9.15,3.789) circle (1.05pt);
\fill[teal] (9.15,4.004) circle (1.05pt);
\fill[teal] (9.15,4.147) circle (1.05pt);
\fill[teal] (9.15,4.361) circle (1.05pt);
\fill[teal] (9.15,4.576) circle (1.05pt);
\fill[teal] (9.15,4.862) circle (1.05pt);
\fill[teal] (9.15,5.076) circle (1.05pt);
\fill[teal] (9.15,5.219) circle (1.05pt);
\fill[teal] (9.272,0.071) circle (1.05pt);
\fill[teal] (9.272,0.214) circle (1.05pt);
\fill[teal] (9.272,0.357) circle (1.05pt);
\fill[teal] (9.272,0.5) circle (1.05pt);
\fill[teal] (9.272,0.643) circle (1.05pt);
\fill[teal] (9.272,0.786) circle (1.05pt);
\fill[teal] (9.272,1.073) circle (1.05pt);
\fill[teal] (9.272,1.501) circle (1.05pt);
\fill[teal] (9.272,1.93) circle (1.05pt);
\fill[teal] (9.272,2.073) circle (1.05pt);
\fill[teal] (9.272,2.216) circle (1.05pt);
\fill[teal] (9.272,2.788) circle (1.05pt);
\fill[teal] (9.272,2.931) circle (1.05pt);
\fill[teal] (9.272,3.074) circle (1.05pt);
\fill[teal] (9.272,3.217) circle (1.05pt);
\fill[teal] (9.272,3.36) circle (1.05pt);
\fill[teal] (9.272,3.646) circle (1.05pt);
\fill[teal] (9.272,3.789) circle (1.05pt);
\fill[teal] (9.272,4.504) circle (1.05pt);
\fill[teal] (9.272,4.933) circle (1.05pt);
\fill[teal] (9.394,0.071) circle (1.05pt);
\fill[teal] (9.394,0.143) circle (1.05pt);
\fill[teal] (9.394,0.214) circle (1.05pt);
\fill[teal] (9.394,0.286) circle (1.05pt);
\fill[teal] (9.394,0.357) circle (1.05pt);
\fill[teal] (9.394,0.572) circle (1.05pt);
\fill[teal] (9.394,0.643) circle (1.05pt);
\fill[teal] (9.394,0.715) circle (1.05pt);
\fill[teal] (9.394,0.858) circle (1.05pt);
\fill[teal] (9.394,0.929) circle (1.05pt);
\fill[teal] (9.394,1.144) circle (1.05pt);
\fill[teal] (9.394,1.215) circle (1.05pt);
\fill[teal] (9.394,1.287) circle (1.05pt);
\fill[teal] (9.394,1.358) circle (1.05pt);
\fill[teal] (9.394,1.644) circle (1.05pt);
\fill[teal] (9.394,1.716) circle (1.05pt);
\fill[teal] (9.394,1.787) circle (1.05pt);
\fill[teal] (9.394,1.859) circle (1.05pt);
\fill[teal] (9.394,2.145) circle (1.05pt);
\fill[teal] (9.394,2.574) circle (1.05pt);
\fill[teal] (9.394,2.645) circle (1.05pt);
\fill[teal] (9.394,2.86) circle (1.05pt);
\fill[teal] (9.394,2.931) circle (1.05pt);
\fill[teal] (9.394,3.074) circle (1.05pt);
\fill[teal] (9.394,3.217) circle (1.05pt);
\fill[teal] (9.394,3.36) circle (1.05pt);
\fill[teal] (9.394,3.575) circle (1.05pt);
\fill[teal] (9.394,3.646) circle (1.05pt);
\fill[teal] (9.394,3.861) circle (1.05pt);
\fill[teal] (9.394,4.075) circle (1.05pt);
\fill[teal] (9.394,4.147) circle (1.05pt);
\fill[teal] (9.394,4.218) circle (1.05pt);
\fill[teal] (9.394,4.29) circle (1.05pt);
\fill[teal] (9.394,4.361) circle (1.05pt);
\fill[teal] (9.394,4.576) circle (1.05pt);
\fill[teal] (9.394,4.79) circle (1.05pt);
\fill[teal] (9.394,4.933) circle (1.05pt);
\fill[teal] (9.394,5.076) circle (1.05pt);
\fill[teal] (9.394,5.291) circle (1.05pt);
\fill[teal] (9.394,5.434) circle (1.05pt);
\fill[teal] (9.516,0.071) circle (1.05pt);
\fill[teal] (9.516,0.5) circle (1.05pt);
\fill[teal] (9.516,1.215) circle (1.05pt);
\fill[teal] (9.516,1.358) circle (1.05pt);
\fill[teal] (9.516,2.073) circle (1.05pt);
\fill[teal] (9.516,2.216) circle (1.05pt);
\fill[teal] (9.516,2.502) circle (1.05pt);
\fill[teal] (9.516,2.645) circle (1.05pt);
\fill[teal] (9.516,2.931) circle (1.05pt);
\fill[teal] (9.516,3.36) circle (1.05pt);
\fill[teal] (9.516,3.503) circle (1.05pt);
\fill[teal] (9.516,4.218) circle (1.05pt);
\fill[teal] (9.516,5.219) circle (1.05pt);
\fill[teal] (9.516,5.505) circle (1.05pt);
\fill[teal] (9.638,0.071) circle (1.05pt);
\fill[teal] (9.638,0.143) circle (1.05pt);
\fill[teal] (9.638,0.214) circle (1.05pt);
\fill[teal] (9.638,0.286) circle (1.05pt);
\fill[teal] (9.638,0.357) circle (1.05pt);
\fill[teal] (9.638,0.429) circle (1.05pt);
\fill[teal] (9.638,0.5) circle (1.05pt);
\fill[teal] (9.638,0.572) circle (1.05pt);
\fill[teal] (9.638,0.715) circle (1.05pt);
\fill[teal] (9.638,0.786) circle (1.05pt);
\fill[teal] (9.638,0.858) circle (1.05pt);
\fill[teal] (9.638,1.073) circle (1.05pt);
\fill[teal] (9.638,1.144) circle (1.05pt);
\fill[teal] (9.638,1.215) circle (1.05pt);
\fill[teal] (9.638,1.43) circle (1.05pt);
\fill[teal] (9.638,1.501) circle (1.05pt);
\fill[teal] (9.638,1.644) circle (1.05pt);
\fill[teal] (9.638,1.716) circle (1.05pt);
\fill[teal] (9.638,1.787) circle (1.05pt);
\fill[teal] (9.638,1.93) circle (1.05pt);
\fill[teal] (9.638,2.002) circle (1.05pt);
\fill[teal] (9.638,2.145) circle (1.05pt);
\fill[teal] (9.638,2.359) circle (1.05pt);
\fill[teal] (9.638,2.574) circle (1.05pt);
\fill[teal] (9.638,2.645) circle (1.05pt);
\fill[teal] (9.638,2.788) circle (1.05pt);
\fill[teal] (9.638,2.86) circle (1.05pt);
\fill[teal] (9.638,3.003) circle (1.05pt);
\fill[teal] (9.638,3.146) circle (1.05pt);
\fill[teal] (9.638,3.217) circle (1.05pt);
\fill[teal] (9.638,3.289) circle (1.05pt);
\fill[teal] (9.638,3.36) circle (1.05pt);
\fill[teal] (9.638,3.503) circle (1.05pt);
\fill[teal] (9.638,3.718) circle (1.05pt);
\fill[teal] (9.638,3.789) circle (1.05pt);
\fill[teal] (9.638,3.932) circle (1.05pt);
\fill[teal] (9.638,4.004) circle (1.05pt);
\fill[teal] (9.638,4.075) circle (1.05pt);
\fill[teal] (9.638,4.647) circle (1.05pt);
\fill[teal] (9.638,4.79) circle (1.05pt);
\fill[teal] (9.638,4.933) circle (1.05pt);
\fill[teal] (9.638,5.148) circle (1.05pt);
\fill[teal] (9.638,5.291) circle (1.05pt);
\fill[teal] (9.638,5.362) circle (1.05pt);
\fill[teal] (9.76,0.071) circle (1.05pt);
\fill[teal] (9.76,0.214) circle (1.05pt);
\fill[teal] (9.76,0.5) circle (1.05pt);
\fill[teal] (9.76,0.643) circle (1.05pt);
\fill[teal] (9.76,0.786) circle (1.05pt);
\fill[teal] (9.76,1.215) circle (1.05pt);
\fill[teal] (9.76,1.93) circle (1.05pt);
\fill[teal] (9.76,2.073) circle (1.05pt);
\fill[teal] (9.76,2.216) circle (1.05pt);
\fill[teal] (9.76,2.788) circle (1.05pt);
\fill[teal] (9.76,2.931) circle (1.05pt);
\fill[teal] (9.76,3.074) circle (1.05pt);
\fill[teal] (9.76,3.503) circle (1.05pt);
\fill[teal] (9.76,4.075) circle (1.05pt);
\fill[teal] (9.76,4.361) circle (1.05pt);
\fill[teal] (9.76,4.933) circle (1.05pt);
\fill[teal] (9.76,5.648) circle (1.05pt);
\fill[teal] (9.882,0.286) circle (1.05pt);
\fill[teal] (9.882,0.357) circle (1.05pt);
\fill[teal] (9.882,0.5) circle (1.05pt);
\fill[teal] (9.882,0.572) circle (1.05pt);
\fill[teal] (9.882,0.715) circle (1.05pt);
\fill[teal] (9.882,0.786) circle (1.05pt);
\fill[teal] (9.882,1.001) circle (1.05pt);
\fill[teal] (9.882,1.573) circle (1.05pt);
\fill[teal] (9.882,1.787) circle (1.05pt);
\fill[teal] (9.882,2.073) circle (1.05pt);
\fill[teal] (9.882,2.431) circle (1.05pt);
\fill[teal] (9.882,2.931) circle (1.05pt);
\fill[teal] (9.882,3.36) circle (1.05pt);
\fill[teal] (9.882,3.503) circle (1.05pt);
\fill[teal] (9.882,3.575) circle (1.05pt);
\fill[teal] (9.882,3.789) circle (1.05pt);
\fill[teal] (9.882,3.932) circle (1.05pt);
\fill[teal] (9.882,4.147) circle (1.05pt);
\fill[teal] (9.882,4.218) circle (1.05pt);
\fill[teal] (9.882,4.361) circle (1.05pt);
\fill[teal] (9.882,4.79) circle (1.05pt);
\fill[teal] (9.882,4.862) circle (1.05pt);
\fill[teal] (9.882,5.005) circle (1.05pt);
\fill[teal] (9.882,5.076) circle (1.05pt);
\fill[teal] (9.882,5.219) circle (1.05pt);
\fill[teal] (9.882,5.291) circle (1.05pt);
\fill[teal] (9.882,5.72) circle (1.05pt);
\fill[teal] (10.004,0.071) circle (1.05pt);
\fill[teal] (10.004,0.214) circle (1.05pt);
\fill[teal] (10.004,0.357) circle (1.05pt);
\fill[teal] (10.004,0.643) circle (1.05pt);
\fill[teal] (10.004,0.786) circle (1.05pt);
\fill[teal] (10.004,0.929) circle (1.05pt);
\fill[teal] (10.004,1.073) circle (1.05pt);
\fill[teal] (10.004,1.215) circle (1.05pt);
\fill[teal] (10.004,1.358) circle (1.05pt);
\fill[teal] (10.004,1.501) circle (1.05pt);
\fill[teal] (10.004,1.644) circle (1.05pt);
\fill[teal] (10.004,1.787) circle (1.05pt);
\fill[teal] (10.004,1.93) circle (1.05pt);
\fill[teal] (10.004,2.216) circle (1.05pt);
\fill[teal] (10.004,2.502) circle (1.05pt);
\fill[teal] (10.004,2.645) circle (1.05pt);
\fill[teal] (10.004,2.788) circle (1.05pt);
\fill[teal] (10.004,3.217) circle (1.05pt);
\fill[teal] (10.004,3.503) circle (1.05pt);
\fill[teal] (10.004,3.646) circle (1.05pt);
\fill[teal] (10.004,3.789) circle (1.05pt);
\fill[teal] (10.004,4.075) circle (1.05pt);
\fill[teal] (10.004,4.218) circle (1.05pt);
\fill[teal] (10.004,4.933) circle (1.05pt);
\fill[teal] (10.004,5.219) circle (1.05pt);
\fill[teal] (10.004,5.362) circle (1.05pt);
\fill[teal] (10.126,0.071) circle (1.05pt);
\fill[teal] (10.126,0.143) circle (1.05pt);
\fill[teal] (10.126,0.214) circle (1.05pt);
\fill[teal] (10.126,0.286) circle (1.05pt);
\fill[teal] (10.126,0.357) circle (1.05pt);
\fill[teal] (10.126,0.429) circle (1.05pt);
\fill[teal] (10.126,0.5) circle (1.05pt);
\fill[teal] (10.126,0.572) circle (1.05pt);
\fill[teal] (10.126,0.715) circle (1.05pt);
\fill[teal] (10.126,0.786) circle (1.05pt);
\fill[teal] (10.126,0.858) circle (1.05pt);
\fill[teal] (10.126,0.929) circle (1.05pt);
\fill[teal] (10.126,1.073) circle (1.05pt);
\fill[teal] (10.126,1.215) circle (1.05pt);
\fill[teal] (10.126,1.358) circle (1.05pt);
\fill[teal] (10.126,1.43) circle (1.05pt);
\fill[teal] (10.126,1.501) circle (1.05pt);
\fill[teal] (10.126,1.573) circle (1.05pt);
\fill[teal] (10.126,1.787) circle (1.05pt);
\fill[teal] (10.126,1.859) circle (1.05pt);
\fill[teal] (10.126,1.93) circle (1.05pt);
\fill[teal] (10.126,2.002) circle (1.05pt);
\fill[teal] (10.126,2.145) circle (1.05pt);
\fill[teal] (10.126,2.431) circle (1.05pt);
\fill[teal] (10.126,2.574) circle (1.05pt);
\fill[teal] (10.126,2.645) circle (1.05pt);
\fill[teal] (10.126,2.717) circle (1.05pt);
\fill[teal] (10.126,2.788) circle (1.05pt);
\fill[teal] (10.126,2.86) circle (1.05pt);
\fill[teal] (10.126,2.931) circle (1.05pt);
\fill[teal] (10.126,3.003) circle (1.05pt);
\fill[teal] (10.126,3.146) circle (1.05pt);
\fill[teal] (10.126,3.289) circle (1.05pt);
\fill[teal] (10.126,3.36) circle (1.05pt);
\fill[teal] (10.126,3.646) circle (1.05pt);
\fill[teal] (10.126,3.789) circle (1.05pt);
\fill[teal] (10.126,3.861) circle (1.05pt);
\fill[teal] (10.126,3.932) circle (1.05pt);
\fill[teal] (10.126,4.075) circle (1.05pt);
\fill[teal] (10.126,4.218) circle (1.05pt);
\fill[teal] (10.126,4.361) circle (1.05pt);
\fill[teal] (10.126,4.576) circle (1.05pt);
\fill[teal] (10.126,5.148) circle (1.05pt);
\fill[teal] (10.126,5.291) circle (1.05pt);
\fill[teal] (10.126,5.362) circle (1.05pt);
\fill[teal] (10.126,5.505) circle (1.05pt);
\fill[teal] (10.126,5.791) circle (1.05pt);
\fill[teal] (10.126,5.863) circle (1.05pt);
\fill[teal] (10.248,0.071) circle (1.05pt);
\fill[teal] (10.248,0.786) circle (1.05pt);
\fill[teal] (10.248,0.929) circle (1.05pt);
\fill[teal] (10.248,1.215) circle (1.05pt);
\fill[teal] (10.248,1.644) circle (1.05pt);
\fill[teal] (10.248,1.787) circle (1.05pt);
\fill[teal] (10.248,2.073) circle (1.05pt);
\fill[teal] (10.248,2.645) circle (1.05pt);
\fill[teal] (10.248,3.074) circle (1.05pt);
\fill[teal] (10.248,3.36) circle (1.05pt);
\fill[teal] (10.248,3.789) circle (1.05pt);
\fill[teal] (10.248,4.218) circle (1.05pt);
\fill[teal] (10.37,0.071) circle (1.05pt);
\fill[teal] (10.37,0.214) circle (1.05pt);
\fill[teal] (10.37,0.429) circle (1.05pt);
\fill[teal] (10.37,0.5) circle (1.05pt);
\fill[teal] (10.37,0.786) circle (1.05pt);
\fill[teal] (10.37,0.929) circle (1.05pt);
\fill[teal] (10.37,1.001) circle (1.05pt);
\fill[teal] (10.37,1.144) circle (1.05pt);
\fill[teal] (10.37,1.859) circle (1.05pt);
\fill[teal] (10.37,1.93) circle (1.05pt);
\fill[teal] (10.37,2.002) circle (1.05pt);
\fill[teal] (10.37,2.216) circle (1.05pt);
\fill[teal] (10.37,2.645) circle (1.05pt);
\fill[teal] (10.37,2.717) circle (1.05pt);
\fill[teal] (10.37,2.788) circle (1.05pt);
\fill[teal] (10.37,2.931) circle (1.05pt);
\fill[teal] (10.37,3.074) circle (1.05pt);
\fill[teal] (10.37,3.146) circle (1.05pt);
\fill[teal] (10.37,3.36) circle (1.05pt);
\fill[teal] (10.37,3.432) circle (1.05pt);
\fill[teal] (10.37,3.503) circle (1.05pt);
\fill[teal] (10.37,4.004) circle (1.05pt);
\fill[teal] (10.37,4.361) circle (1.05pt);
\fill[teal] (10.37,4.504) circle (1.05pt);
\fill[teal] (10.37,4.933) circle (1.05pt);
\fill[teal] (10.37,5.219) circle (1.05pt);
\fill[teal] (10.492,0.071) circle (1.05pt);
\fill[teal] (10.492,0.357) circle (1.05pt);
\fill[teal] (10.492,0.5) circle (1.05pt);
\fill[teal] (10.492,0.643) circle (1.05pt);
\fill[teal] (10.492,0.786) circle (1.05pt);
\fill[teal] (10.492,1.073) circle (1.05pt);
\fill[teal] (10.492,1.215) circle (1.05pt);
\fill[teal] (10.492,1.358) circle (1.05pt);
\fill[teal] (10.492,1.787) circle (1.05pt);
\fill[teal] (10.492,1.93) circle (1.05pt);
\fill[teal] (10.492,2.216) circle (1.05pt);
\fill[teal] (10.492,2.645) circle (1.05pt);
\fill[teal] (10.492,3.217) circle (1.05pt);
\fill[teal] (10.492,3.36) circle (1.05pt);
\fill[teal] (10.492,3.503) circle (1.05pt);
\fill[teal] (10.492,3.646) circle (1.05pt);
\fill[teal] (10.492,4.79) circle (1.05pt);
\fill[teal] (10.492,4.933) circle (1.05pt);
\fill[teal] (10.492,5.505) circle (1.05pt);
\fill[teal] (10.492,6.077) circle (1.05pt);
\fill[teal] (10.614,0.071) circle (1.05pt);
\fill[teal] (10.614,0.143) circle (1.05pt);
\fill[teal] (10.614,0.286) circle (1.05pt);
\fill[teal] (10.614,0.572) circle (1.05pt);
\fill[teal] (10.614,0.715) circle (1.05pt);
\fill[teal] (10.614,0.786) circle (1.05pt);
\fill[teal] (10.614,1.144) circle (1.05pt);
\fill[teal] (10.614,1.215) circle (1.05pt);
\fill[teal] (10.614,1.573) circle (1.05pt);
\fill[teal] (10.614,1.644) circle (1.05pt);
\fill[teal] (10.614,1.859) circle (1.05pt);
\fill[teal] (10.614,2.86) circle (1.05pt);
\fill[teal] (10.614,3.074) circle (1.05pt);
\fill[teal] (10.614,3.289) circle (1.05pt);
\fill[teal] (10.614,3.503) circle (1.05pt);
\fill[teal] (10.614,3.718) circle (1.05pt);
\fill[teal] (10.614,3.789) circle (1.05pt);
\fill[teal] (10.614,3.932) circle (1.05pt);
\fill[teal] (10.614,4.004) circle (1.05pt);
\fill[teal] (10.614,4.218) circle (1.05pt);
\fill[teal] (10.614,4.79) circle (1.05pt);
\fill[teal] (10.614,5.291) circle (1.05pt);
\fill[teal] (10.614,5.505) circle (1.05pt);
\fill[teal] (10.614,6.077) circle (1.05pt);
\fill[teal] (10.736,0.071) circle (1.05pt);
\fill[teal] (10.736,0.214) circle (1.05pt);
\fill[teal] (10.736,0.357) circle (1.05pt);
\fill[teal] (10.736,0.5) circle (1.05pt);
\fill[teal] (10.736,0.643) circle (1.05pt);
\fill[teal] (10.736,0.929) circle (1.05pt);
\fill[teal] (10.736,1.073) circle (1.05pt);
\fill[teal] (10.736,1.358) circle (1.05pt);
\fill[teal] (10.736,1.501) circle (1.05pt);
\fill[teal] (10.736,1.93) circle (1.05pt);
\fill[teal] (10.736,2.073) circle (1.05pt);
\fill[teal] (10.736,2.216) circle (1.05pt);
\fill[teal] (10.736,2.788) circle (1.05pt);
\fill[teal] (10.736,2.931) circle (1.05pt);
\fill[teal] (10.736,3.074) circle (1.05pt);
\fill[teal] (10.736,3.217) circle (1.05pt);
\fill[teal] (10.736,3.646) circle (1.05pt);
\fill[teal] (10.736,4.075) circle (1.05pt);
\fill[teal] (10.736,4.361) circle (1.05pt);
\fill[teal] (10.736,4.504) circle (1.05pt);
\fill[teal] (10.736,5.934) circle (1.05pt);
\fill[teal] (10.736,6.22) circle (1.05pt);
\fill[teal] (10.858,0.357) circle (1.05pt);
\fill[teal] (10.858,0.429) circle (1.05pt);
\fill[teal] (10.858,0.5) circle (1.05pt);
\fill[teal] (10.858,0.715) circle (1.05pt);
\fill[teal] (10.858,0.786) circle (1.05pt);
\fill[teal] (10.858,0.858) circle (1.05pt);
\fill[teal] (10.858,0.929) circle (1.05pt);
\fill[teal] (10.858,1.001) circle (1.05pt);
\fill[teal] (10.858,1.144) circle (1.05pt);
\fill[teal] (10.858,1.287) circle (1.05pt);
\fill[teal] (10.858,1.573) circle (1.05pt);
\fill[teal] (10.858,1.716) circle (1.05pt);
\fill[teal] (10.858,1.787) circle (1.05pt);
\fill[teal] (10.858,1.859) circle (1.05pt);
\fill[teal] (10.858,2.145) circle (1.05pt);
\fill[teal] (10.858,2.502) circle (1.05pt);
\fill[teal] (10.858,2.574) circle (1.05pt);
\fill[teal] (10.858,2.645) circle (1.05pt);
\fill[teal] (10.858,2.717) circle (1.05pt);
\fill[teal] (10.858,3.003) circle (1.05pt);
\fill[teal] (10.858,3.217) circle (1.05pt);
\fill[teal] (10.858,3.289) circle (1.05pt);
\fill[teal] (10.858,3.36) circle (1.05pt);
\fill[teal] (10.858,3.503) circle (1.05pt);
\fill[teal] (10.858,3.861) circle (1.05pt);
\fill[teal] (10.858,4.29) circle (1.05pt);
\fill[teal] (10.858,4.433) circle (1.05pt);
\fill[teal] (10.858,4.504) circle (1.05pt);
\fill[teal] (10.858,4.576) circle (1.05pt);
\fill[teal] (10.858,4.79) circle (1.05pt);
\fill[teal] (10.858,4.933) circle (1.05pt);
\fill[teal] (10.858,5.362) circle (1.05pt);
\fill[teal] (10.858,5.791) circle (1.05pt);
\fill[teal] (10.858,6.292) circle (1.05pt);
\fill[teal] (10.98,0.5) circle (1.05pt);
\fill[teal] (10.98,0.929) circle (1.05pt);
\fill[teal] (10.98,1.215) circle (1.05pt);
\fill[teal] (10.98,1.358) circle (1.05pt);
\fill[teal] (10.98,2.073) circle (1.05pt);
\fill[teal] (10.98,2.216) circle (1.05pt);
\fill[teal] (10.98,3.074) circle (1.05pt);
\fill[teal] (10.98,3.503) circle (1.05pt);
\fill[teal] (10.98,3.789) circle (1.05pt);
\fill[teal] (10.98,4.361) circle (1.05pt);
\fill[teal] (10.98,5.505) circle (1.05pt);
\fill[teal] (11.102,0.071) circle (1.05pt);
\fill[teal] (11.102,0.143) circle (1.05pt);
\fill[teal] (11.102,0.214) circle (1.05pt);
\fill[teal] (11.102,0.357) circle (1.05pt);
\fill[teal] (11.102,0.429) circle (1.05pt);
\fill[teal] (11.102,0.572) circle (1.05pt);
\fill[teal] (11.102,1.144) circle (1.05pt);
\fill[teal] (11.102,1.287) circle (1.05pt);
\fill[teal] (11.102,1.358) circle (1.05pt);
\fill[teal] (11.102,1.43) circle (1.05pt);
\fill[teal] (11.102,1.93) circle (1.05pt);
\fill[teal] (11.102,2.216) circle (1.05pt);
\fill[teal] (11.102,2.359) circle (1.05pt);
\fill[teal] (11.102,2.574) circle (1.05pt);
\fill[teal] (11.102,2.645) circle (1.05pt);
\fill[teal] (11.102,2.717) circle (1.05pt);
\fill[teal] (11.102,2.86) circle (1.05pt);
\fill[teal] (11.102,2.931) circle (1.05pt);
\fill[teal] (11.102,3.074) circle (1.05pt);
\fill[teal] (11.102,3.146) circle (1.05pt);
\fill[teal] (11.102,3.575) circle (1.05pt);
\fill[teal] (11.102,4.147) circle (1.05pt);
\fill[teal] (11.102,4.218) circle (1.05pt);
\fill[teal] (11.102,4.29) circle (1.05pt);
\fill[teal] (11.102,4.361) circle (1.05pt);
\fill[teal] (11.102,5.076) circle (1.05pt);
\fill[teal] (11.102,5.291) circle (1.05pt);
\fill[teal] (11.102,5.434) circle (1.05pt);
\fill[teal] (11.102,5.648) circle (1.05pt);
\fill[teal] (11.102,6.292) circle (1.05pt);
\fill[teal] (11.224,0.071) circle (1.05pt);
\fill[teal] (11.224,0.357) circle (1.05pt);
\fill[teal] (11.224,0.5) circle (1.05pt);
\fill[teal] (11.224,0.929) circle (1.05pt);
\fill[teal] (11.224,1.215) circle (1.05pt);
\fill[teal] (11.224,1.358) circle (1.05pt);
\fill[teal] (11.224,1.787) circle (1.05pt);
\fill[teal] (11.224,1.93) circle (1.05pt);
\fill[teal] (11.224,2.216) circle (1.05pt);
\fill[teal] (11.224,2.502) circle (1.05pt);
\fill[teal] (11.224,2.645) circle (1.05pt);
\fill[teal] (11.224,2.931) circle (1.05pt);
\fill[teal] (11.224,3.074) circle (1.05pt);
\fill[teal] (11.224,3.217) circle (1.05pt);
\fill[teal] (11.224,3.646) circle (1.05pt);
\fill[teal] (11.224,4.075) circle (1.05pt);
\fill[teal] (11.224,4.647) circle (1.05pt);
\fill[teal] (11.224,6.22) circle (1.05pt);
\fill[teal] (11.224,6.506) circle (1.05pt);
\fill[teal] (11.346,0.071) circle (1.05pt);
\fill[teal] (11.346,0.143) circle (1.05pt);
\fill[teal] (11.346,0.286) circle (1.05pt);
\fill[teal] (11.346,0.5) circle (1.05pt);
\fill[teal] (11.346,0.572) circle (1.05pt);
\fill[teal] (11.346,0.929) circle (1.05pt);
\fill[teal] (11.346,1.215) circle (1.05pt);
\fill[teal] (11.346,1.358) circle (1.05pt);
\fill[teal] (11.346,1.644) circle (1.05pt);
\fill[teal] (11.346,2.002) circle (1.05pt);
\fill[teal] (11.346,2.502) circle (1.05pt);
\fill[teal] (11.346,2.931) circle (1.05pt);
\fill[teal] (11.346,3.146) circle (1.05pt);
\fill[teal] (11.346,3.932) circle (1.05pt);
\fill[teal] (11.346,4.004) circle (1.05pt);
\fill[teal] (11.346,4.361) circle (1.05pt);
\fill[teal] (11.346,4.576) circle (1.05pt);
\fill[teal] (11.346,4.79) circle (1.05pt);
\fill[teal] (11.346,4.862) circle (1.05pt);
\fill[teal] (11.346,5.076) circle (1.05pt);
\fill[teal] (11.346,5.219) circle (1.05pt);
\fill[teal] (11.346,5.648) circle (1.05pt);
\fill[teal] (11.346,5.72) circle (1.05pt);
\fill[teal] (11.346,6.077) circle (1.05pt);
\fill[teal] (11.346,6.506) circle (1.05pt);
\fill[teal] (11.346,6.578) circle (1.05pt);
\fill[teal] (11.468,0.071) circle (1.05pt);
\fill[teal] (11.468,0.214) circle (1.05pt);
\fill[teal] (11.468,0.357) circle (1.05pt);
\fill[teal] (11.468,0.5) circle (1.05pt);
\fill[teal] (11.468,0.643) circle (1.05pt);
\fill[teal] (11.468,1.501) circle (1.05pt);
\fill[teal] (11.468,1.644) circle (1.05pt);
\fill[teal] (11.468,2.073) circle (1.05pt);
\fill[teal] (11.468,2.216) circle (1.05pt);
\fill[teal] (11.468,2.931) circle (1.05pt);
\fill[teal] (11.468,3.503) circle (1.05pt);
\fill[teal] (11.468,3.789) circle (1.05pt);
\fill[teal] (11.468,4.504) circle (1.05pt);
\fill[teal] (11.468,4.647) circle (1.05pt);
\fill[teal] (11.468,4.79) circle (1.05pt);
\fill[teal] (11.468,5.505) circle (1.05pt);
\fill[teal] (11.468,5.791) circle (1.05pt);
\fill[teal] (11.468,6.649) circle (1.05pt);
\fill[teal] (11.59,0.071) circle (1.05pt);
\fill[teal] (11.59,0.143) circle (1.05pt);
\fill[teal] (11.59,0.286) circle (1.05pt);
\fill[teal] (11.59,0.572) circle (1.05pt);
\fill[teal] (11.59,0.643) circle (1.05pt);
\fill[teal] (11.59,0.786) circle (1.05pt);
\fill[teal] (11.59,0.858) circle (1.05pt);
\fill[teal] (11.59,1.144) circle (1.05pt);
\fill[teal] (11.59,1.215) circle (1.05pt);
\fill[teal] (11.59,1.287) circle (1.05pt);
\fill[teal] (11.59,1.716) circle (1.05pt);
\fill[teal] (11.59,2.002) circle (1.05pt);
\fill[teal] (11.59,2.073) circle (1.05pt);
\fill[teal] (11.59,2.574) circle (1.05pt);
\fill[teal] (11.59,2.645) circle (1.05pt);
\fill[teal] (11.59,3.003) circle (1.05pt);
\fill[teal] (11.59,3.146) circle (1.05pt);
\fill[teal] (11.59,3.289) circle (1.05pt);
\fill[teal] (11.59,3.36) circle (1.05pt);
\fill[teal] (11.59,3.718) circle (1.05pt);
\fill[teal] (11.59,3.789) circle (1.05pt);
\fill[teal] (11.59,3.861) circle (1.05pt);
\fill[teal] (11.59,4.147) circle (1.05pt);
\fill[teal] (11.59,4.218) circle (1.05pt);
\fill[teal] (11.59,4.361) circle (1.05pt);
\fill[teal] (11.59,4.433) circle (1.05pt);
\fill[teal] (11.59,4.719) circle (1.05pt);
\fill[teal] (11.59,4.862) circle (1.05pt);
\fill[teal] (11.59,4.933) circle (1.05pt);
\fill[teal] (11.59,5.076) circle (1.05pt);
\fill[teal] (11.59,5.791) circle (1.05pt);
\fill[teal] (11.59,5.934) circle (1.05pt);
\fill[teal] (11.59,6.721) circle (1.05pt);
\fill[teal] (11.712,0.071) circle (1.05pt);
\fill[teal] (11.712,0.357) circle (1.05pt);
\fill[teal] (11.712,0.5) circle (1.05pt);
\fill[teal] (11.712,0.929) circle (1.05pt);
\fill[teal] (11.712,1.358) circle (1.05pt);
\fill[teal] (11.712,1.644) circle (1.05pt);
\fill[teal] (11.712,1.787) circle (1.05pt);
\fill[teal] (11.712,2.216) circle (1.05pt);
\fill[teal] (11.712,2.502) circle (1.05pt);
\fill[teal] (11.712,2.645) circle (1.05pt);
\fill[teal] (11.712,2.931) circle (1.05pt);
\fill[teal] (11.712,3.074) circle (1.05pt);
\fill[teal] (11.712,3.36) circle (1.05pt);
\fill[teal] (11.712,3.789) circle (1.05pt);
\fill[teal] (11.712,4.79) circle (1.05pt);
\fill[teal] (11.834,0.071) circle (1.05pt);
\fill[teal] (11.834,0.143) circle (1.05pt);
\fill[teal] (11.834,0.286) circle (1.05pt);
\fill[teal] (11.834,0.357) circle (1.05pt);
\fill[teal] (11.834,0.429) circle (1.05pt);
\fill[teal] (11.834,0.5) circle (1.05pt);
\fill[teal] (11.834,0.572) circle (1.05pt);
\fill[teal] (11.834,0.715) circle (1.05pt);
\fill[teal] (11.834,0.786) circle (1.05pt);
\fill[teal] (11.834,0.929) circle (1.05pt);
\fill[teal] (11.834,1.001) circle (1.05pt);
\fill[teal] (11.834,1.073) circle (1.05pt);
\fill[teal] (11.834,1.144) circle (1.05pt);
\fill[teal] (11.834,1.215) circle (1.05pt);
\fill[teal] (11.834,1.287) circle (1.05pt);
\fill[teal] (11.834,1.358) circle (1.05pt);
\fill[teal] (11.834,1.43) circle (1.05pt);
\fill[teal] (11.834,1.501) circle (1.05pt);
\fill[teal] (11.834,1.573) circle (1.05pt);
\fill[teal] (11.834,1.644) circle (1.05pt);
\fill[teal] (11.834,1.716) circle (1.05pt);
\fill[teal] (11.834,1.859) circle (1.05pt);
\fill[teal] (11.834,2.002) circle (1.05pt);
\fill[teal] (11.834,2.073) circle (1.05pt);
\fill[teal] (11.834,2.216) circle (1.05pt);
\fill[teal] (11.834,2.359) circle (1.05pt);
\fill[teal] (11.834,2.431) circle (1.05pt);
\fill[teal] (11.834,2.574) circle (1.05pt);
\fill[teal] (11.834,2.645) circle (1.05pt);
\fill[teal] (11.834,2.931) circle (1.05pt);
\fill[teal] (11.834,3.217) circle (1.05pt);
\fill[teal] (11.834,3.36) circle (1.05pt);
\fill[teal] (11.834,3.432) circle (1.05pt);
\fill[teal] (11.834,3.575) circle (1.05pt);
\fill[teal] (11.834,3.861) circle (1.05pt);
\fill[teal] (11.834,4.719) circle (1.05pt);
\fill[teal] (11.834,5.219) circle (1.05pt);
\fill[teal] (11.834,5.362) circle (1.05pt);
\fill[teal] (11.834,5.434) circle (1.05pt);
\fill[teal] (11.834,5.505) circle (1.05pt);
\fill[teal] (11.834,5.934) circle (1.05pt);
\fill[teal] (11.834,6.506) circle (1.05pt);
\fill[teal] (11.834,6.649) circle (1.05pt);
\fill[teal] (11.834,6.792) circle (1.05pt);
\fill[teal] (11.956,0.071) circle (1.05pt);
\fill[teal] (11.956,0.214) circle (1.05pt);
\fill[teal] (11.956,1.073) circle (1.05pt);
\fill[teal] (11.956,1.215) circle (1.05pt);
\fill[teal] (11.956,2.073) circle (1.05pt);
\fill[teal] (11.956,2.216) circle (1.05pt);
\fill[teal] (11.956,2.788) circle (1.05pt);
\fill[teal] (11.956,3.074) circle (1.05pt);
\fill[teal] (11.956,3.36) circle (1.05pt);
\fill[teal] (11.956,3.932) circle (1.05pt);
\fill[teal] (11.956,4.218) circle (1.05pt);
\fill[teal] (11.956,4.361) circle (1.05pt);
\fill[teal] (11.956,4.647) circle (1.05pt);
\fill[teal] (11.956,6.935) circle (1.05pt);
\fill[teal] (12.078,1.215) circle (1.05pt);
\fill[teal] (12.078,1.358) circle (1.05pt);
\fill[teal] (12.078,2.216) circle (1.05pt);
\fill[teal] (12.078,2.288) circle (1.05pt);
\fill[teal] (12.078,2.431) circle (1.05pt);
\fill[teal] (12.078,2.502) circle (1.05pt);
\fill[teal] (12.078,2.717) circle (1.05pt);
\fill[teal] (12.078,2.931) circle (1.05pt);
\fill[teal] (12.078,3.289) circle (1.05pt);
\fill[teal] (12.078,3.36) circle (1.05pt);
\fill[teal] (12.078,3.503) circle (1.05pt);
\fill[teal] (12.078,3.789) circle (1.05pt);
\fill[teal] (12.078,4.004) circle (1.05pt);
\fill[teal] (12.078,5.291) circle (1.05pt);
\fill[teal] (12.078,6.792) circle (1.05pt);
\fill[teal] (12.2,0.071) circle (1.05pt);
\fill[teal] (12.2,0.214) circle (1.05pt);
\fill[teal] (12.2,0.5) circle (1.05pt);
\fill[teal] (12.2,0.929) circle (1.05pt);
\fill[teal] (12.2,1.358) circle (1.05pt);
\fill[teal] (12.2,1.501) circle (1.05pt);
\fill[teal] (12.2,2.073) circle (1.05pt);
\fill[teal] (12.2,2.788) circle (1.05pt);
\fill[teal] (12.2,3.503) circle (1.05pt);
\fill[teal] (12.2,4.361) circle (1.05pt);
\fill[teal] (12.2,5.076) circle (1.05pt);
\fill[ochre] (0.854,0.143) ++(-1.15pt,-1.15pt) rectangle ++(2.3pt,2.3pt);
\fill[ochre] (1.098,0.143) ++(-1.15pt,-1.15pt) rectangle ++(2.3pt,2.3pt);
\fill[ochre] (1.098,0.5) ++(-1.15pt,-1.15pt) rectangle ++(2.3pt,2.3pt);
\fill[ochre] (1.342,0.429) ++(-1.15pt,-1.15pt) rectangle ++(2.3pt,2.3pt);
\fill[ochre] (1.342,0.572) ++(-1.15pt,-1.15pt) rectangle ++(2.3pt,2.3pt);
\fill[ochre] (1.342,0.715) ++(-1.15pt,-1.15pt) rectangle ++(2.3pt,2.3pt);
\fill[ochre] (1.464,0.786) ++(-1.15pt,-1.15pt) rectangle ++(2.3pt,2.3pt);
\fill[ochre] (1.586,0.071) ++(-1.15pt,-1.15pt) rectangle ++(2.3pt,2.3pt);
\fill[ochre] (1.586,0.214) ++(-1.15pt,-1.15pt) rectangle ++(2.3pt,2.3pt);
\fill[ochre] (1.586,0.357) ++(-1.15pt,-1.15pt) rectangle ++(2.3pt,2.3pt);
\fill[ochre] (1.586,0.786) ++(-1.15pt,-1.15pt) rectangle ++(2.3pt,2.3pt);
\fill[ochre] (1.586,0.858) ++(-1.15pt,-1.15pt) rectangle ++(2.3pt,2.3pt);
\fill[ochre] (1.708,0.214) ++(-1.15pt,-1.15pt) rectangle ++(2.3pt,2.3pt);
\fill[ochre] (1.708,0.643) ++(-1.15pt,-1.15pt) rectangle ++(2.3pt,2.3pt);
\fill[ochre] (1.708,0.929) ++(-1.15pt,-1.15pt) rectangle ++(2.3pt,2.3pt);
\fill[ochre] (1.83,0.071) ++(-1.15pt,-1.15pt) rectangle ++(2.3pt,2.3pt);
\fill[ochre] (1.83,0.286) ++(-1.15pt,-1.15pt) rectangle ++(2.3pt,2.3pt);
\fill[ochre] (1.83,0.786) ++(-1.15pt,-1.15pt) rectangle ++(2.3pt,2.3pt);
\fill[ochre] (1.952,0.071) ++(-1.15pt,-1.15pt) rectangle ++(2.3pt,2.3pt);
\fill[ochre] (1.952,0.214) ++(-1.15pt,-1.15pt) rectangle ++(2.3pt,2.3pt);
\fill[ochre] (1.952,0.357) ++(-1.15pt,-1.15pt) rectangle ++(2.3pt,2.3pt);
\fill[ochre] (1.952,0.5) ++(-1.15pt,-1.15pt) rectangle ++(2.3pt,2.3pt);
\fill[ochre] (1.952,0.929) ++(-1.15pt,-1.15pt) rectangle ++(2.3pt,2.3pt);
\fill[ochre] (2.074,0.071) ++(-1.15pt,-1.15pt) rectangle ++(2.3pt,2.3pt);
\fill[ochre] (2.074,0.214) ++(-1.15pt,-1.15pt) rectangle ++(2.3pt,2.3pt);
\fill[ochre] (2.074,0.429) ++(-1.15pt,-1.15pt) rectangle ++(2.3pt,2.3pt);
\fill[ochre] (2.074,1.073) ++(-1.15pt,-1.15pt) rectangle ++(2.3pt,2.3pt);
\fill[ochre] (2.074,1.144) ++(-1.15pt,-1.15pt) rectangle ++(2.3pt,2.3pt);
\fill[ochre] (2.196,0.071) ++(-1.15pt,-1.15pt) rectangle ++(2.3pt,2.3pt);
\fill[ochre] (2.196,0.357) ++(-1.15pt,-1.15pt) rectangle ++(2.3pt,2.3pt);
\fill[ochre] (2.196,1.215) ++(-1.15pt,-1.15pt) rectangle ++(2.3pt,2.3pt);
\fill[ochre] (2.318,0.071) ++(-1.15pt,-1.15pt) rectangle ++(2.3pt,2.3pt);
\fill[ochre] (2.318,0.143) ++(-1.15pt,-1.15pt) rectangle ++(2.3pt,2.3pt);
\fill[ochre] (2.318,0.214) ++(-1.15pt,-1.15pt) rectangle ++(2.3pt,2.3pt);
\fill[ochre] (2.318,0.429) ++(-1.15pt,-1.15pt) rectangle ++(2.3pt,2.3pt);
\fill[ochre] (2.318,0.572) ++(-1.15pt,-1.15pt) rectangle ++(2.3pt,2.3pt);
\fill[ochre] (2.44,0.5) ++(-1.15pt,-1.15pt) rectangle ++(2.3pt,2.3pt);
\fill[ochre] (2.562,0.357) ++(-1.15pt,-1.15pt) rectangle ++(2.3pt,2.3pt);
\fill[ochre] (2.562,0.572) ++(-1.15pt,-1.15pt) rectangle ++(2.3pt,2.3pt);
\fill[ochre] (2.562,0.715) ++(-1.15pt,-1.15pt) rectangle ++(2.3pt,2.3pt);
\fill[ochre] (2.562,0.786) ++(-1.15pt,-1.15pt) rectangle ++(2.3pt,2.3pt);
\fill[ochre] (2.562,0.929) ++(-1.15pt,-1.15pt) rectangle ++(2.3pt,2.3pt);
\fill[ochre] (2.684,0.071) ++(-1.15pt,-1.15pt) rectangle ++(2.3pt,2.3pt);
\fill[ochre] (2.684,0.214) ++(-1.15pt,-1.15pt) rectangle ++(2.3pt,2.3pt);
\fill[ochre] (2.684,0.5) ++(-1.15pt,-1.15pt) rectangle ++(2.3pt,2.3pt);
\fill[ochre] (2.684,0.643) ++(-1.15pt,-1.15pt) rectangle ++(2.3pt,2.3pt);
\fill[ochre] (2.684,1.073) ++(-1.15pt,-1.15pt) rectangle ++(2.3pt,2.3pt);
\fill[ochre] (2.806,0.214) ++(-1.15pt,-1.15pt) rectangle ++(2.3pt,2.3pt);
\fill[ochre] (2.806,1.215) ++(-1.15pt,-1.15pt) rectangle ++(2.3pt,2.3pt);
\fill[ochre] (2.806,1.573) ++(-1.15pt,-1.15pt) rectangle ++(2.3pt,2.3pt);
\fill[ochre] (2.928,0.071) ++(-1.15pt,-1.15pt) rectangle ++(2.3pt,2.3pt);
\fill[ochre] (3.05,0.071) ++(-1.15pt,-1.15pt) rectangle ++(2.3pt,2.3pt);
\fill[ochre] (3.05,0.214) ++(-1.15pt,-1.15pt) rectangle ++(2.3pt,2.3pt);
\fill[ochre] (3.05,0.286) ++(-1.15pt,-1.15pt) rectangle ++(2.3pt,2.3pt);
\fill[ochre] (3.05,0.5) ++(-1.15pt,-1.15pt) rectangle ++(2.3pt,2.3pt);
\fill[ochre] (3.05,1.501) ++(-1.15pt,-1.15pt) rectangle ++(2.3pt,2.3pt);
\fill[ochre] (3.172,0.214) ++(-1.15pt,-1.15pt) rectangle ++(2.3pt,2.3pt);
\fill[ochre] (3.294,0.071) ++(-1.15pt,-1.15pt) rectangle ++(2.3pt,2.3pt);
\fill[ochre] (3.294,0.929) ++(-1.15pt,-1.15pt) rectangle ++(2.3pt,2.3pt);
\fill[ochre] (3.294,1.358) ++(-1.15pt,-1.15pt) rectangle ++(2.3pt,2.3pt);
\fill[ochre] (3.294,1.787) ++(-1.15pt,-1.15pt) rectangle ++(2.3pt,2.3pt);
\fill[ochre] (3.416,0.071) ++(-1.15pt,-1.15pt) rectangle ++(2.3pt,2.3pt);
\fill[ochre] (3.538,0.357) ++(-1.15pt,-1.15pt) rectangle ++(2.3pt,2.3pt);
\fill[ochre] (3.782,0.071) ++(-1.15pt,-1.15pt) rectangle ++(2.3pt,2.3pt);
\fill[ochre] (3.782,2.073) ++(-1.15pt,-1.15pt) rectangle ++(2.3pt,2.3pt);
\fill[ochre] (3.904,0.357) ++(-1.15pt,-1.15pt) rectangle ++(2.3pt,2.3pt);
\fill[ochre] (3.904,1.215) ++(-1.15pt,-1.15pt) rectangle ++(2.3pt,2.3pt);
\fill[ochre] (3.904,1.644) ++(-1.15pt,-1.15pt) rectangle ++(2.3pt,2.3pt);
\fill[ochre] (4.148,0.5) ++(-1.15pt,-1.15pt) rectangle ++(2.3pt,2.3pt);
\fill[ochre] (4.392,0.929) ++(-1.15pt,-1.15pt) rectangle ++(2.3pt,2.3pt);
\fill[ochre] (4.514,0.071) ++(-1.15pt,-1.15pt) rectangle ++(2.3pt,2.3pt);
\fill[ochre] (4.514,0.429) ++(-1.15pt,-1.15pt) rectangle ++(2.3pt,2.3pt);
\fill[ochre] (4.514,0.5) ++(-1.15pt,-1.15pt) rectangle ++(2.3pt,2.3pt);
\fill[ochre] (4.636,0.071) ++(-1.15pt,-1.15pt) rectangle ++(2.3pt,2.3pt);
\fill[ochre] (4.636,1.93) ++(-1.15pt,-1.15pt) rectangle ++(2.3pt,2.3pt);
\fill[ochre] (5.246,0.5) ++(-1.15pt,-1.15pt) rectangle ++(2.3pt,2.3pt);
\fill[ochre] (5.246,1.859) ++(-1.15pt,-1.15pt) rectangle ++(2.3pt,2.3pt);
\fill[ochre] (5.734,0.572) ++(-1.15pt,-1.15pt) rectangle ++(2.3pt,2.3pt);
\fill[ochre] (5.734,2.002) ++(-1.15pt,-1.15pt) rectangle ++(2.3pt,2.3pt);
\fill[ochre] (6.71,0.071) ++(-1.15pt,-1.15pt) rectangle ++(2.3pt,2.3pt);
\fill[ochre] (7.076,2.502) ++(-1.15pt,-1.15pt) rectangle ++(2.3pt,2.3pt);
\fill[ochre] (7.442,0.715) ++(-1.15pt,-1.15pt) rectangle ++(2.3pt,2.3pt);
\fill[ochre] (7.564,2.645) ++(-1.15pt,-1.15pt) rectangle ++(2.3pt,2.3pt);
\fill[ochre] (7.93,0.786) ++(-1.15pt,-1.15pt) rectangle ++(2.3pt,2.3pt);
\fill[ochre] (8.906,3.146) ++(-1.15pt,-1.15pt) rectangle ++(2.3pt,2.3pt);
\fill[ochre] (9.394,3.289) ++(-1.15pt,-1.15pt) rectangle ++(2.3pt,2.3pt);
\fill[ochre] (9.638,0.929) ++(-1.15pt,-1.15pt) rectangle ++(2.3pt,2.3pt);
\fill[ochre] (10.126,1.001) ++(-1.15pt,-1.15pt) rectangle ++(2.3pt,2.3pt);
\fill[ochre] (10.736,3.789) ++(-1.15pt,-1.15pt) rectangle ++(2.3pt,2.3pt);
\fill[ochre] (11.224,3.932) ++(-1.15pt,-1.15pt) rectangle ++(2.3pt,2.3pt);
\fill[wine] (0.61,0.286) ++(0,1.9pt) -- ++(1.9pt,-1.9pt) -- ++(-1.9pt,-1.9pt) -- ++(-1.9pt,1.9pt) -- cycle;
\fill[wine] (0.854,0.071) ++(0,1.9pt) -- ++(1.9pt,-1.9pt) -- ++(-1.9pt,-1.9pt) -- ++(-1.9pt,1.9pt) -- cycle;
\fill[wine] (0.854,0.286) ++(0,1.9pt) -- ++(1.9pt,-1.9pt) -- ++(-1.9pt,-1.9pt) -- ++(-1.9pt,1.9pt) -- cycle;
\fill[wine] (0.854,0.357) ++(0,1.9pt) -- ++(1.9pt,-1.9pt) -- ++(-1.9pt,-1.9pt) -- ++(-1.9pt,1.9pt) -- cycle;
\fill[wine] (0.976,0.5) ++(0,1.9pt) -- ++(1.9pt,-1.9pt) -- ++(-1.9pt,-1.9pt) -- ++(-1.9pt,1.9pt) -- cycle;
\fill[wine] (1.098,0.071) ++(0,1.9pt) -- ++(1.9pt,-1.9pt) -- ++(-1.9pt,-1.9pt) -- ++(-1.9pt,1.9pt) -- cycle;
\fill[wine] (1.098,0.572) ++(0,1.9pt) -- ++(1.9pt,-1.9pt) -- ++(-1.9pt,-1.9pt) -- ++(-1.9pt,1.9pt) -- cycle;
\fill[wine] (1.22,0.071) ++(0,1.9pt) -- ++(1.9pt,-1.9pt) -- ++(-1.9pt,-1.9pt) -- ++(-1.9pt,1.9pt) -- cycle;
\fill[wine] (1.22,0.214) ++(0,1.9pt) -- ++(1.9pt,-1.9pt) -- ++(-1.9pt,-1.9pt) -- ++(-1.9pt,1.9pt) -- cycle;
\fill[wine] (1.342,0.143) ++(0,1.9pt) -- ++(1.9pt,-1.9pt) -- ++(-1.9pt,-1.9pt) -- ++(-1.9pt,1.9pt) -- cycle;
\fill[wine] (1.464,0.071) ++(0,1.9pt) -- ++(1.9pt,-1.9pt) -- ++(-1.9pt,-1.9pt) -- ++(-1.9pt,1.9pt) -- cycle;
\fill[wine] (1.464,0.5) ++(0,1.9pt) -- ++(1.9pt,-1.9pt) -- ++(-1.9pt,-1.9pt) -- ++(-1.9pt,1.9pt) -- cycle;
\fill[wine] (1.586,0.572) ++(0,1.9pt) -- ++(1.9pt,-1.9pt) -- ++(-1.9pt,-1.9pt) -- ++(-1.9pt,1.9pt) -- cycle;
\fill[wine] (2.074,0.143) ++(0,1.9pt) -- ++(1.9pt,-1.9pt) -- ++(-1.9pt,-1.9pt) -- ++(-1.9pt,1.9pt) -- cycle;
\fill[wine] (2.074,0.715) ++(0,1.9pt) -- ++(1.9pt,-1.9pt) -- ++(-1.9pt,-1.9pt) -- ++(-1.9pt,1.9pt) -- cycle;
\fill[wine] (3.416,1.215) ++(0,1.9pt) -- ++(1.9pt,-1.9pt) -- ++(-1.9pt,-1.9pt) -- ++(-1.9pt,1.9pt) -- cycle;
\fill[wine] (3.904,1.358) ++(0,1.9pt) -- ++(1.9pt,-1.9pt) -- ++(-1.9pt,-1.9pt) -- ++(-1.9pt,1.9pt) -- cycle;
\draw[ink, line width=0.5pt] (2.074,0.715) circle (3.2pt);
\draw[ink, line width=0.5pt] (3.904,1.358) circle (3.2pt);
\draw[ink, line width=0.5pt] (5.734,2.002) circle (3.2pt);
\draw[ink, line width=0.5pt] (7.564,2.645) circle (3.2pt);
\draw[ink, line width=0.5pt] (9.394,3.289) circle (3.2pt);
\draw[ink, line width=0.5pt] (11.224,3.932) circle (3.2pt);
\draw[ink, line width=0.5pt] (0.122,0) -- (12.322,0);
\draw[ink, line width=0.5pt] (0.122,0) -- (0.122,7.15);
\draw[ink, line width=0.5pt] (1.22,0) -- (1.22,-0.08);
\node[ink, below] at (1.22,-0.08) {$10$};
\draw[ink, line width=0.5pt] (2.44,0) -- (2.44,-0.08);
\node[ink, below] at (2.44,-0.08) {$20$};
\draw[ink, line width=0.5pt] (3.66,0) -- (3.66,-0.08);
\node[ink, below] at (3.66,-0.08) {$30$};
\draw[ink, line width=0.5pt] (4.88,0) -- (4.88,-0.08);
\node[ink, below] at (4.88,-0.08) {$40$};
\draw[ink, line width=0.5pt] (6.1,0) -- (6.1,-0.08);
\node[ink, below] at (6.1,-0.08) {$50$};
\draw[ink, line width=0.5pt] (7.32,0) -- (7.32,-0.08);
\node[ink, below] at (7.32,-0.08) {$60$};
\draw[ink, line width=0.5pt] (8.54,0) -- (8.54,-0.08);
\node[ink, below] at (8.54,-0.08) {$70$};
\draw[ink, line width=0.5pt] (9.76,0) -- (9.76,-0.08);
\node[ink, below] at (9.76,-0.08) {$80$};
\draw[ink, line width=0.5pt] (10.98,0) -- (10.98,-0.08);
\node[ink, below] at (10.98,-0.08) {$90$};
\draw[ink, line width=0.5pt] (12.2,0) -- (12.2,-0.08);
\node[ink, below] at (12.2,-0.08) {$100$};
\draw[ink, line width=0.5pt] (0.122,0.071) -- (0.042,0.071);
\node[ink, left] at (0.042,0.071) {$1$};
\draw[ink, line width=0.5pt] (0.122,0.929) -- (0.042,0.929);
\node[ink, left] at (0.042,0.929) {$13$};
\draw[ink, line width=0.5pt] (0.122,1.787) -- (0.042,1.787);
\node[ink, left] at (0.042,1.787) {$25$};
\draw[ink, line width=0.5pt] (0.122,3.575) -- (0.042,3.575);
\node[ink, left] at (0.042,3.575) {$50$};
\draw[ink, line width=0.5pt] (0.122,5.362) -- (0.042,5.362);
\node[ink, left] at (0.042,5.362) {$75$};
\draw[ink, line width=0.5pt] (0.122,7.078) -- (0.042,7.078);
\node[ink, left] at (0.042,7.078) {$99$};
\node[ink] at (6.161,-0.75) {denominator $Q$};
\node[ink, rotate=90] at (-0.728,3.575) {numerator $P$};
\node[ink, right] at (0.582,6.998) {margin};
\fill[wine] (0.732,6.578) ++(0,2.4699999999999998pt) -- ++(2.4699999999999998pt,-2.4699999999999998pt) -- ++(-2.4699999999999998pt,-2.4699999999999998pt) -- ++(-2.4699999999999998pt,2.4699999999999998pt) -- cycle;
\node[ink, right] at (0.832,6.578) {1 or 2};
\fill[ochre] (0.732,6.218) ++(-1.4949999999999999pt,-1.4949999999999999pt) rectangle ++(2.9899999999999998pt,2.9899999999999998pt);
\node[ink, right] at (0.832,6.218) {3 to 9};
\fill[teal] (0.732,5.858) circle (1.3650000000000002pt);
\node[ink, right] at (0.832,5.858) {10 to 99};
\fill[navy] (0.732,5.498) circle (0.9750000000000001pt);
\node[ink, right] at (0.832,5.498) {100 or more};
\draw[slate, line width=0.5pt] (0.732,5.138) circle (1.3pt);
\node[ink, right] at (0.832,5.138) {one representation};
\draw[ink, line width=0.5pt] (0.732,4.778) circle (3.2pt);
\node[ink, right] at (0.832,4.778) {$(9k+1)/(15k+2)$};
\end{tikzpicture}

\caption{Computed evidence for \cref{conj:euclidean-minimality}. Each mark is
an occurrence $P/Q$ with $Q\leqslant100$, placed at its denominator and
numerator; its color and shape give the margin, the cost of the cheapest
other primitive representation of the label minus the occurrence cost
$P+Q-2$. Hollow circles are labels with a single representation, which have
no competitor. No margin is zero or negative. The small margins thin out
as the denominator grows; the ringed occurrences $(9k+1)/(15k+2)$,
$k=1,\dots,6$, are the last ones with a margin at most $k$, and
\cref{prop:two-ladders} explains them. The shaded rows $P\leqslant13$ are those
of \cref{thm:thirteen-rows}.}
\label{fig:margins}
\end{figure}

\subsection{How small the margin can be}

The rarefaction of small margins in \cref{fig:margins} follows a precise
law in the computed range. Among the occurrences with $P+Q\leqslant160$, a
margin at most $k$, for $k=1,\dots,6$, occurs only when
$$
 P+Q\leqslant24k+3,\qquad P\leqslant9k+1,\qquad Q\leqslant15k+2 ,
$$
and all three bounds are reached by the single occurrence
$(9k+1)/(15k+2)$, whose margin is exactly $k$ (\cref{fig:thresholds}). These
six fractions are rungs of one Farey ladder, the fractions
$(1+3j)/(2+5j)$ that run from $1/2$ to $3/5$, taken at $j=3k$; their labels
are the neighbors of $m_{3/5}=433$ in the Markov tree. On every third rung
the label is divisible by $5=(1+2i)(1-2i)$, and at the six computed rungs
the cheapest competitor of the occurrence pair is the element of the fiber
in which the Gaussian factor above $5$ is replaced by its conjugate. A companion ladder,
$(2+3j)/(3+5j)$ from $2/3$ to $3/5$, behaves in the same way at $j=3k-1$. The
mechanism can be proved for every rung, and it shows that thresholds of
this kind, if they exist, cannot be smaller.

The proof follows the occurrence vectors $v_j=H_je$ of the rungs, $H_j$ being
the Gram factor of \cref{prop:occurrence-cost}. Two rungs apart, the central
words grow by one fixed word at each end, so that $v_{j+2}$ is obtained from
$v_j$ by a matrix that depends only on the parity of $j$; six rungs apart,
the competitor obtained by conjugating the factor above $5$ is transformed
by the same matrix conjugated by the rotation that multiplies by
$(1-2i)/(1+2i)$.
Cost is additive under nonnegative matrices of determinant one, and the
whole statement reduces to the lengths of four explicit matrices.

\begin{lemma}[Finite facts on two ladders]\label{lem:ladder-facts}
For the ladder $(1+3j)/(2+5j)$ put $u=babab$ and $j_0=3$, and for the ladder
$(2+3j)/(3+5j)$ put $u=abbab$ and $j_0=2$. For each ladder and each parity
of $j\geqslant1$, let $T=H_j\,\mu(u)\,H_j^{-1}$, computed at the rungs $j=1$ and
$j=2$. At each of the four rungs $j_0$ and $j_0+3$ of the two ladders, exactly
one of the two matrices $R$ of multiplication by $(1-2i)/(1+2i)$ and by
$(1+2i)/(1-2i)$ maps $v_j$ to an integer vector, and exactly one diagonal sign
matrix $S$ makes that vector positive. Then:
\begin{enumerate}[label=\textup{(\roman*)},ref=\textup{(\roman*)},leftmargin=*]
\item\label{it:ladder-transfer} On the first ladder $T$ is
 $\left(\begin{smallmatrix}1105&463\\463&194\end{smallmatrix}\right)$ at even $j$
 and $\left(\begin{smallmatrix}1130&437\\437&169\end{smallmatrix}\right)$ at odd
 $j$; on the second ladder the two matrices are exchanged.
\item\label{it:ladder-length} At each of the four rungs, with $T$ the matrix of its parity, the
 matrix $N=SRT^3R^{-1}S^{-1}$ is a nonnegative integer matrix of determinant
 one whose expression as a word in $L_a$ and $L_b$ has $50$ letters, while
 $T^3$ has $48$.
\item\label{it:ladder-base} At each of the four rungs the vector $SRv_j$ is primitive, and its
 Euclidean cost exceeds that of $v_j$ by $1$ at $j_0$ and by $2$ at $j_0+3$.
\end{enumerate}
\end{lemma}

\begin{proof}
These are finitely many assertions about explicit integer matrices and
vectors. They were verified by computer in exact arithmetic, and the
verification was carried out twice independently: once through the Gram
factors and the matrices of \ref{it:ladder-length}, and once by Gaussian division of the
occurrence pairs, with the labels taken as continuants of the continued
fraction $[2,\sigma(z),2]$.
\end{proof}

\begin{proposition}[Margins along two ladders]\label{prop:two-ladders}
For every integer $k\geqslant1$, the labels of the reduced fractions
$(9k+1)/(15k+2)$ and $(9k-1)/(15k-2)$ are divisible by $5$. In the occurrence
pair of each, written as a Gaussian integer, replacing the factor above $5$
by its conjugate gives a primitive representation of the same label whose
Euclidean cost exceeds the occurrence cost by exactly $k$. Consequently the
margins of these two occurrences are at most $k$.
\end{proposition}

\begin{proof}
\emph{The words.} Two fractions $a/b<c/d$ with $bc-ad=1$ are Farey
neighbors, and the lower Christoffel word of their mediant is the
product $\gamma_{a/b}\gamma_{c/d}$ \cite{Reutenauer2009}. The fraction
$(1+3j)/(2+5j)$ is a neighbor of $3/5$ lying below it, since
$3(2+5j)-5(1+3j)=1$, and its mediant with $3/5$ is the next rung; starting
from $\gamma_{1/2}=ab$ and $\gamma_{3/5}=ababb$, induction gives
$\gamma_{(1+3j)/(2+5j)}=ab\,(ababb)^j$, with central word
$z_j=b\,(ababb)^{j-1}abab$. Concatenation gives
$z_{j+2}=babab\cdot z_j\cdot babab=u\,z_j\,u$. In the same way
$5(2+3j)-3(3+5j)=1$ places $(2+3j)/(3+5j)$ above $3/5$, the words are
$(ababb)^jabb$ with central words $z_j=babb\,(ababb)^{j-1}ab$, and
$z_{j+2}=babba\cdot z_j\cdot abbab=\widetilde u\,z_j\,u$.

\emph{The transfer.} In both cases the right half of the central word gains
the word $u$ at its far end while its middle letter is unchanged, so that
the Gram factor of \cref{prop:occurrence-cost} satisfies
$H_{j+2}=H_j\,\mu(u)$. Put $T_j=H_j\,\mu(u)\,H_j^{-1}$. Then
$T_jv_j=H_j\,\mu(u)\,e=H_{j+2}e=v_{j+2}$, and
$T_{j+2}=H_j\,\mu(u)\,\mu(u)\,\mu(u)^{-1}H_j^{-1}=T_j$, so that $T_j$ depends only
on the parity of $j$: it is the matrix $T$ of \cref{lem:ladder-facts}\ref{it:ladder-transfer}, and
$v_{j+6}=T^3v_j$.

\emph{The competitor.} Write $v_j=(s,r)^{\mathsf T}$, so that the label is
$m=s^2+r^2$. Fix one of the four rungs $j_0$, $j_0+3$ of
\cref{lem:ladder-facts}, with its matrices $R$, $S$ and $T$, and put
$c_j=SR\,v_j$ for every rung $j$ congruent to it modulo six. Since
$v_{j+6}=T^3v_j$,
$$
 c_{j+6}=SR\,T^3v_j=N\,c_j,\qquad N=SRT^3R^{-1}S^{-1}.
$$
The matrix $N$ has nonnegative integer entries and determinant one
(\cref{lem:ladder-facts}\ref{it:ladder-length}); it maps positive vectors to positive
vectors, since no row of it vanishes, and primitive integer vectors to
primitive integer vectors, since its inverse is integral. By
\cref{lem:ladder-facts}\ref{it:ladder-base} and induction on $j$ in steps of six, $c_j$ is
therefore a positive primitive integer vector at every rung $j=3k$ of the
first ladder and $j=3k-1$ of the second. The matrices $R$ and $S$ preserve
lengths, so $c_j$ has norm $m$: up to the order of its entries it is an
element of $\mathcal F(m)$. As a Gaussian integer, $c_j$ is $s+ir$ multiplied
by $(1\mp2i)/(1\pm2i)$, up to a sign and a conjugation; its integrality means
that $1\pm2i$ divides $s+ir$, so that $5$ divides $m$, and its primitivity that
$25$ does not, by \cref{prop:gaussian-fibre}. The element $c_j$ is thus the
occurrence pair with the factor above $5$ replaced by its conjugate.

\emph{Cost.} By \cref{lem:euclidean-bijection}, every positive primitive
column is $x=L_w{\bf1}$ for a unique word $w$, so that for every word $w'$,
$\ell(L_{w'}x)=\ell(L_{w'w}{\bf1})=\lvert w'\rvert+\ell(x)$. Every nonnegative integer
matrix $M=\left(\begin{smallmatrix}\alpha&\beta\\\gamma&\delta\end{smallmatrix}\right)\neq I$
of determinant one is such an $L_{w'}$. Indeed one of its rows dominates the
other entrywise: if $\alpha>\gamma$ and $\beta<\delta$, then
$\alpha\delta-\beta\gamma\geqslant(\gamma+1)(\beta+1)-\beta\gamma=\beta+\gamma+1$, which equals one
only for $M=I$; and if $\alpha<\gamma$ and $\beta>\delta$, then
$\alpha\delta-\beta\gamma\leqslant(\gamma-1)(\beta-1)-\beta\gamma<0$. Subtracting the dominated
row from the dominating one removes a first letter $L_a$ or $L_b$, keeps the
entries nonnegative and the determinant one, and lowers the sum of the
entries, so that the process ends at $I$. By \cref{lem:ladder-facts}\ref{it:ladder-length},
$$
 \ell(c_{j+6})=\ell(c_j)+50,\qquad \ell(v_{j+6})=\ell(v_j)+48,
$$
so that the excess $\ell(c_j)-\ell(v_j)$ grows by $2$ every six rungs. It is
$1$ and $2$ at the first two rungs of each family
(\cref{lem:ladder-facts}\ref{it:ladder-base}), hence $k$ at the $k$th, which is the rung
$(9k+1)/(15k+2)$ of the first ladder and $(9k-1)/(15k-2)$ of the second.
The margin is at most the excess of this particular competitor.
\end{proof}

The proposition explains the six ringed occurrences of \cref{fig:margins}:
together with the computations, which find no cheaper competitor at those
rungs, it gives their margins exactly. It also shows that the thresholds
observed above cannot be lowered: for every $k$ there is an occurrence with
$P+Q=24k+3$ and margin at most $k$.

\begin{figure}[tb]
\centering
% Authors: Dr. Denys Dutykh (Mathematics Department, Khalifa University of Science and
% Technology, Abu Dhabi, UAE) and Prof. Laurent Vuillon (Univ. Savoie Mont Blanc, CNRS, LAMA)
% Generated by supplement/code/make_figures.py from exact data; do not edit by hand.
\begin{tikzpicture}[font=\footnotesize, >={Stealth[length=4pt]}]
\fill[sand] (0.225,0) -- (12,3.271) -- (12,0) -- cycle;
\draw[slate, line width=0.7pt, dash pattern=on 5pt off 2pt] (0.225,0) -- (12,3.271);
\filldraw[navy, fill=white, line width=0.6pt] (0.6,0.5) circle (1.5pt);
\filldraw[navy, fill=white, line width=0.6pt] (0.675,0.5) circle (1.5pt);
\filldraw[navy, fill=white, line width=0.6pt] (0.675,1.5) circle (1.5pt);
\filldraw[navy, fill=white, line width=0.6pt] (0.75,1) circle (1.5pt);
\filldraw[navy, fill=white, line width=0.6pt] (0.825,1) circle (1.5pt);
\filldraw[navy, fill=white, line width=0.6pt] (0.825,3) circle (1.5pt);
\filldraw[navy, fill=white, line width=0.6pt] (0.9,1) circle (1.5pt);
\filldraw[navy, fill=white, line width=0.6pt] (0.975,0.5) circle (1.5pt);
\filldraw[navy, fill=white, line width=0.6pt] (0.975,1) circle (1.5pt);
\filldraw[navy, fill=white, line width=0.6pt] (1.05,1.5) circle (1.5pt);
\filldraw[navy, fill=white, line width=0.6pt] (1.125,0.5) circle (1.5pt);
\filldraw[navy, fill=white, line width=0.6pt] (1.2,3) circle (1.5pt);
\filldraw[navy, fill=white, line width=0.6pt] (1.2,3.5) circle (1.5pt);
\filldraw[navy, fill=white, line width=0.6pt] (1.2,4) circle (1.5pt);
\filldraw[navy, fill=white, line width=0.6pt] (1.275,1) circle (1.5pt);
\filldraw[navy, fill=white, line width=0.6pt] (1.275,1.5) circle (1.5pt);
\filldraw[navy, fill=white, line width=0.6pt] (1.275,2) circle (1.5pt);
\filldraw[navy, fill=white, line width=0.6pt] (1.35,1.5) circle (1.5pt);
\filldraw[navy, fill=white, line width=0.6pt] (1.35,2.5) circle (1.5pt);
\filldraw[navy, fill=white, line width=0.6pt] (1.425,0.5) circle (1.5pt);
\filldraw[navy, fill=white, line width=0.6pt] (1.425,1) circle (1.5pt);
\filldraw[navy, fill=white, line width=0.6pt] (1.425,1.5) circle (1.5pt);
\filldraw[navy, fill=white, line width=0.6pt] (1.425,2.5) circle (1.5pt);
\filldraw[navy, fill=white, line width=0.6pt] (1.425,3) circle (1.5pt);
\filldraw[navy, fill=white, line width=0.6pt] (1.425,4.5) circle (1.5pt);
\filldraw[navy, fill=white, line width=0.6pt] (1.5,2.5) circle (1.5pt);
\filldraw[navy, fill=white, line width=0.6pt] (1.5,3) circle (1.5pt);
\filldraw[navy, fill=white, line width=0.6pt] (1.575,0.5) circle (1.5pt);
\filldraw[navy, fill=white, line width=0.6pt] (1.575,2) circle (1.5pt);
\filldraw[navy, fill=white, line width=0.6pt] (1.575,2.5) circle (1.5pt);
\filldraw[navy, fill=white, line width=0.6pt] (1.575,3) circle (1.5pt);
\filldraw[navy, fill=white, line width=0.6pt] (1.65,4.5) circle (1.5pt);
\filldraw[navy, fill=white, line width=0.6pt] (1.725,1.5) circle (1.5pt);
\filldraw[navy, fill=white, line width=0.6pt] (1.725,2.5) circle (1.5pt);
\filldraw[navy, fill=white, line width=0.6pt] (1.725,3) circle (1.5pt);
\filldraw[navy, fill=white, line width=0.6pt] (1.725,3.5) circle (1.5pt);
\filldraw[navy, fill=white, line width=0.6pt] (1.8,2) circle (1.5pt);
\filldraw[navy, fill=white, line width=0.6pt] (1.875,2) circle (1.5pt);
\filldraw[navy, fill=white, line width=0.6pt] (1.875,2.5) circle (1.5pt);
\filldraw[navy, fill=white, line width=0.6pt] (1.875,3) circle (1.5pt);
\filldraw[navy, fill=white, line width=0.6pt] (1.875,5) circle (1.5pt);
\filldraw[navy, fill=white, line width=0.6pt] (1.95,3) circle (1.5pt);
\filldraw[navy, fill=white, line width=0.6pt] (1.95,3.5) circle (1.5pt);
\filldraw[navy, fill=white, line width=0.6pt] (1.95,4.5) circle (1.5pt);
\filldraw[navy, fill=white, line width=0.6pt] (2.025,0.5) circle (1.5pt);
\filldraw[navy, fill=white, line width=0.6pt] (2.025,2.5) circle (1.5pt);
\filldraw[navy, fill=white, line width=0.6pt] (2.025,4) circle (1.5pt);
\filldraw[navy, fill=white, line width=0.6pt] (2.1,2.5) circle (1.5pt);
\filldraw[navy, fill=white, line width=0.6pt] (2.1,3.5) circle (1.5pt);
\filldraw[navy, fill=white, line width=0.6pt] (2.175,1.5) circle (1.5pt);
\filldraw[navy, fill=white, line width=0.6pt] (2.175,4) circle (1.5pt);
\filldraw[navy, fill=white, line width=0.6pt] (2.175,4.5) circle (1.5pt);
\filldraw[navy, fill=white, line width=0.6pt] (2.175,5) circle (1.5pt);
\filldraw[navy, fill=white, line width=0.6pt] (2.325,3) circle (1.5pt);
\filldraw[navy, fill=white, line width=0.6pt] (2.325,4) circle (1.5pt);
\filldraw[navy, fill=white, line width=0.6pt] (2.325,5) circle (1.5pt);
\filldraw[navy, fill=white, line width=0.6pt] (2.4,1.5) circle (1.5pt);
\filldraw[navy, fill=white, line width=0.6pt] (2.4,2.5) circle (1.5pt);
\filldraw[navy, fill=white, line width=0.6pt] (2.4,3.5) circle (1.5pt);
\filldraw[navy, fill=white, line width=0.6pt] (2.4,4.5) circle (1.5pt);
\filldraw[navy, fill=white, line width=0.6pt] (2.475,4.5) circle (1.5pt);
\filldraw[navy, fill=white, line width=0.6pt] (2.475,5) circle (1.5pt);
\filldraw[navy, fill=white, line width=0.6pt] (2.55,1.5) circle (1.5pt);
\filldraw[navy, fill=white, line width=0.6pt] (2.625,2) circle (1.5pt);
\filldraw[navy, fill=white, line width=0.6pt] (2.625,5) circle (1.5pt);
\filldraw[navy, fill=white, line width=0.6pt] (2.775,4.5) circle (1.5pt);
\filldraw[navy, fill=white, line width=0.6pt] (2.85,3) circle (1.5pt);
\filldraw[navy, fill=white, line width=0.6pt] (2.925,4) circle (1.5pt);
\filldraw[navy, fill=white, line width=0.6pt] (2.925,5) circle (1.5pt);
\filldraw[navy, fill=white, line width=0.6pt] (3,3.5) circle (1.5pt);
\filldraw[navy, fill=white, line width=0.6pt] (3,4.5) circle (1.5pt);
\filldraw[navy, fill=white, line width=0.6pt] (3.075,4) circle (1.5pt);
\filldraw[navy, fill=white, line width=0.6pt] (3.15,5) circle (1.5pt);
\filldraw[navy, fill=white, line width=0.6pt] (3.225,3.5) circle (1.5pt);
\filldraw[navy, fill=white, line width=0.6pt] (3.225,5) circle (1.5pt);
\filldraw[navy, fill=white, line width=0.6pt] (3.3,3) circle (1.5pt);
\filldraw[navy, fill=white, line width=0.6pt] (3.3,5) circle (1.5pt);
\filldraw[navy, fill=white, line width=0.6pt] (3.375,1) circle (1.5pt);
\filldraw[navy, fill=white, line width=0.6pt] (3.375,3) circle (1.5pt);
\filldraw[navy, fill=white, line width=0.6pt] (3.375,5) circle (1.5pt);
\filldraw[navy, fill=white, line width=0.6pt] (3.45,3) circle (1.5pt);
\filldraw[navy, fill=white, line width=0.6pt] (3.45,4) circle (1.5pt);
\filldraw[navy, fill=white, line width=0.6pt] (3.675,4) circle (1.5pt);
\filldraw[navy, fill=white, line width=0.6pt] (3.675,4.5) circle (1.5pt);
\filldraw[navy, fill=white, line width=0.6pt] (3.75,2.5) circle (1.5pt);
\filldraw[navy, fill=white, line width=0.6pt] (3.825,1) circle (1.5pt);
\filldraw[navy, fill=white, line width=0.6pt] (3.9,4) circle (1.5pt);
\filldraw[navy, fill=white, line width=0.6pt] (3.975,5) circle (1.5pt);
\filldraw[navy, fill=white, line width=0.6pt] (4.125,2.5) circle (1.5pt);
\filldraw[navy, fill=white, line width=0.6pt] (4.125,3.5) circle (1.5pt);
\filldraw[navy, fill=white, line width=0.6pt] (4.125,5) circle (1.5pt);
\filldraw[navy, fill=white, line width=0.6pt] (4.2,4.5) circle (1.5pt);
\filldraw[navy, fill=white, line width=0.6pt] (4.5,4) circle (1.5pt);
\filldraw[navy, fill=white, line width=0.6pt] (4.875,3.5) circle (1.5pt);
\filldraw[navy, fill=white, line width=0.6pt] (5.175,1.5) circle (1.5pt);
\filldraw[navy, fill=white, line width=0.6pt] (5.325,3.5) circle (1.5pt);
\filldraw[navy, fill=white, line width=0.6pt] (5.625,1.5) circle (1.5pt);
\filldraw[navy, fill=white, line width=0.6pt] (5.7,3.5) circle (1.5pt);
\filldraw[navy, fill=white, line width=0.6pt] (6.9,4.5) circle (1.5pt);
\filldraw[navy, fill=white, line width=0.6pt] (6.975,2) circle (1.5pt);
\filldraw[navy, fill=white, line width=0.6pt] (7.275,4.5) circle (1.5pt);
\filldraw[navy, fill=white, line width=0.6pt] (7.425,2) circle (1.5pt);
\filldraw[navy, fill=white, line width=0.6pt] (8.775,2.5) circle (1.5pt);
\filldraw[navy, fill=white, line width=0.6pt] (9.225,2.5) circle (1.5pt);
\filldraw[navy, fill=white, line width=0.6pt] (10.575,3) circle (1.5pt);
\filldraw[navy, fill=white, line width=0.6pt] (11.025,3) circle (1.5pt);
\fill[wine] (2.025,0.5) circle (1.6pt);
\draw[wine, line width=0.5pt] (2.025,0.5) circle (3.2pt);
\fill[wine] (3.825,1) circle (1.6pt);
\draw[wine, line width=0.5pt] (3.825,1) circle (3.2pt);
\fill[wine] (5.625,1.5) circle (1.6pt);
\draw[wine, line width=0.5pt] (5.625,1.5) circle (3.2pt);
\fill[wine] (7.425,2) circle (1.6pt);
\draw[wine, line width=0.5pt] (7.425,2) circle (3.2pt);
\fill[wine] (9.225,2.5) circle (1.6pt);
\draw[wine, line width=0.5pt] (9.225,2.5) circle (3.2pt);
\fill[wine] (11.025,3) circle (1.6pt);
\draw[wine, line width=0.5pt] (11.025,3) circle (3.2pt);
\draw[ink, line width=0.5pt] (0,0) -- (12,0);
\draw[ink, line width=0.5pt] (0,0) -- (0,5);
\draw[ink, line width=0.5pt] (1.5,0) -- (1.5,-0.08);
\node[ink, below] at (1.5,-0.08) {$20$};
\draw[ink, line width=0.5pt] (3,0) -- (3,-0.08);
\node[ink, below] at (3,-0.08) {$40$};
\draw[ink, line width=0.5pt] (4.5,0) -- (4.5,-0.08);
\node[ink, below] at (4.5,-0.08) {$60$};
\draw[ink, line width=0.5pt] (6,0) -- (6,-0.08);
\node[ink, below] at (6,-0.08) {$80$};
\draw[ink, line width=0.5pt] (7.5,0) -- (7.5,-0.08);
\node[ink, below] at (7.5,-0.08) {$100$};
\draw[ink, line width=0.5pt] (9,0) -- (9,-0.08);
\node[ink, below] at (9,-0.08) {$120$};
\draw[ink, line width=0.5pt] (10.5,0) -- (10.5,-0.08);
\node[ink, below] at (10.5,-0.08) {$140$};
\draw[ink, line width=0.5pt] (12,0) -- (12,-0.08);
\node[ink, below] at (12,-0.08) {$160$};
\draw[ink, line width=0.5pt] (0,0) -- (-0.08,0);
\node[ink, left] at (-0.08,0) {$0$};
\draw[ink, line width=0.5pt] (0,1) -- (-0.08,1);
\node[ink, left] at (-0.08,1) {$2$};
\draw[ink, line width=0.5pt] (0,2) -- (-0.08,2);
\node[ink, left] at (-0.08,2) {$4$};
\draw[ink, line width=0.5pt] (0,3) -- (-0.08,3);
\node[ink, left] at (-0.08,3) {$6$};
\draw[ink, line width=0.5pt] (0,4) -- (-0.08,4);
\node[ink, left] at (-0.08,4) {$8$};
\draw[ink, line width=0.5pt] (0,5) -- (-0.08,5);
\node[ink, left] at (-0.08,5) {$10$};
\node[ink] at (6,-0.75) {$P+Q$};
\node[ink, rotate=90] at (-0.75,2.5) {margin};
\node[slate, font=\scriptsize, rotate=15.524, anchor=north] at (10.275,2.492) {$(P+Q-3)/24$};
\end{tikzpicture}

\caption{The margin against $P+Q$ for the $120$ occurrences with
$P+Q\leqslant160$ and margin at most $10$ (hollow circles; several
occurrences may share a mark). The dashed line is $(P+Q-3)/24$, and the
shaded region below it contains no occurrence. The filled wine circles are
the occurrences $(9k+1)/(15k+2)$, $k=1,\dots,6$, which lie on the line and
whose margin is exactly $k$; \cref{prop:two-ladders} shows that the margin
of the next ones on the line is at most $k$ for every $k$.}
\label{fig:thresholds}
\end{figure}

\begin{conjecture}[Thresholds of the margin]\label{conj:thresholds}
For every integer $k\geqslant0$, every occurrence $P/Q$ with $P+Q>24k+3$ has
margin larger than $k$, the margin of a label with a single primitive
representation being taken as infinite.
\end{conjecture}

The case $k=0$ asserts that every occurrence is strictly minimal, which is
stronger than \cref{conj:euclidean-minimality} and therefore implies the
Frobenius--Markov uniqueness conjecture. The conjecture holds for every
occurrence with $P+Q\leqslant160$, and \cref{prop:two-ladders} shows that it
would be false with any smaller threshold.

The minimality is not a consequence of the digit pattern alone. The
representation of a Markov label by its occurrence pair has a half continued
fraction with digits $1$ and $2$ only \cite[Corollary~4.4]{CanakciSchiffler2020},
but such digits do not force minimality for other integers. The integer
$135328989953=37\cdot113\cdot761\cdot42533$, which is not a Markov number,
has eight primitive representations; two of them,
$196223^2+311168^2$ and $131468^2+343577^2$, have continued fractions with
digits $1$ and $2$ only, and their costs are $28$ and $27$. The first is
therefore not minimal. Whatever makes the conjecture true must use the
Markov structure of the label, as the path comparison of this paper does on
thirteen rows.

\subsection{Directives of the Christoffel words of a given length}\label{sec:directives}

The comparison of this paper fixes a label and varies its representation as
a sum of two squares. One may instead fix an integer $m$ and vary the
Christoffel word of length $m$. For each $n$ prime to $m$ with $1\leqslant n<m/2$,
the lower Christoffel word $\gamma_{n/m}$ has a central word, and the central
word has a directive, in the sense of the iterated palindromic closure of de
Luca. This subsection shows that these directives are the Euclidean
directives of \cref{lem:euclidean-bijection}, characterizes those that are
images of the Thue--Morse morphism, and states a conjecture on them that
implies the uniqueness conjecture. It also recasts
\cref{conj:euclidean-minimality} in terms of the square roots of $-1$ modulo
the label.

The palindromic closure $w^{(+)}$ of a word $w$ is the shortest palindrome
having $w$ as a prefix, and the iterated palindromic closure is defined by
$\operatorname{Pal}(\varepsilon)=\varepsilon$ and
$\operatorname{Pal}(vx)=(\operatorname{Pal}(v)x)^{(+)}$ for a word $v$ and a letter $x$
\cite{DeLuca1997}. The map $\operatorname{Pal}$ is injective, and $d$ is called the
\emph{directive} of the central word $\operatorname{Pal}(d)$; the map
$d\mapsto a\operatorname{Pal}(d)b$ is a bijection from the binary words onto the
proper lower Christoffel words
\cite[Section~2 and the proof of Theorem~1]{ReutenauerVuillon2017}, a consequence of
\cite[Corollary~3.1]{BertheDeLucaReutenauer2008}. The Thue--Morse morphism
$\theta$ sends $a$ to $ab$ and $b$ to $ba$; its image is the set $\{ab,ba\}^*$
of the products of the words $ab$ and $ba$.

\begin{lemma}[Directives of Christoffel words]\label{lem:directives}
Let $1\leqslant n<m/2$ be coprime integers, write $\gamma_{n/m}=azb$, and let $d$
be the directive of $z$. Then $L_d{\bf1}=(m-n,n)^{\mathsf T}$: the word $d$ is the
directive of the pair $(m-n,n)$ in the sense of \cref{lem:euclidean-bijection},
and $\abs d=\ell(n,m)-1$. If $n=1$, then $d=a^{m-2}$. If $n\geqslant2$ and
$m/n=[c_0;c_1,\ldots,c_k]$ is the canonical continued fraction, the runs of $d$
have the lengths $c_0-1,c_1,\ldots,c_{k-1},c_k-1$, alternately of the letters
$a$ and $b$, beginning with $a$. In particular, $d$ lies in $\{ab,ba\}^*$ if and
only if
\begin{equation}\label{eq:block-form}
 \frac mn=[2;B_1,\ldots,B_j,2],\qquad j\geqslant0,\quad B_i\in\{(2),\,(1,1)\},
\end{equation}
that is, if and only if the first and the last partial quotients equal $2$
and those in between split into consecutive blocks $(2)$ and $(1,1)$.
\end{lemma}

\begin{proof}
\emph{The pair.} Berth\'e, de Luca and Reutenauer attach to a word $d$ the
incidence matrix $M_d$ of a composition of morphisms, with $M_{dd'}=M_dM_{d'}$,
$M_a=\left(\begin{smallmatrix}1&1\\0&1\end{smallmatrix}\right)=L_a$ and
$M_b=\left(\begin{smallmatrix}1&0\\1&1\end{smallmatrix}\right)=L_b$
\cite[proofs of Proposition~3.1 and Corollary~3.2]{BertheDeLucaReutenauer2008},
so that $M_d=L_d$. By \cite[Corollary~3.1]{BertheDeLucaReutenauer2008}, the
Christoffel word $a\operatorname{Pal}(d)b$ has as many letters $a$ as the sum
of the entries of the first row of $L_d$, and as many letters $b$ as the sum
of the entries of its second row; that is, its two letter counts are the
entries of $L_d{\bf1}$. For $z=\operatorname{Pal}(d)$, \eqref{eq:farey-counts}
therefore gives $L_d{\bf1}=(m-n,n)^{\mathsf T}$, and \cref{lem:euclidean-bijection}
shows that $d$ is the directive of $(m-n,n)$, of length $\ell(m-n,n)$. Since
$m-n>n$, the first subtraction carries the pair $(n,m)$ to $(n,m-n)$, so that
$\ell(n,m)=1+\ell(n,m-n)$.

\emph{The runs.} If $n=1$, the pair $(m-1,1)$ reaches $(1,1)$ after $m-2$
subtractions of its second coordinate, and its directive is $a^{m-2}$. If
$n\geqslant2$, then $m/n$ is not an integer, so $k\geqslant1$; moreover $c_0\geqslant2$
because $m>2n$, and $(m-n)/n=m/n-1=[c_0-1;c_1,\ldots,c_k]$ is a canonical
expansion. The description of directives after \cref{lem:euclidean-bijection}
gives the runs $c_0-1,c_1,\ldots,c_{k-1},c_k-1$, alternately of $a$ and $b$,
beginning with $a$.

\emph{The image of $\theta$.} A word lies in $\{ab,ba\}^*$ exactly when its
length is even and its letters in the positions $2i-1$ and $2i$ differ for
every $i$, that is, when a run ends after every odd position. If $n=1$, the
word $a^{m-2}$, with $m\geqslant3$, has no such run end, and \eqref{eq:block-form}
fails since $m/1$ has a single partial quotient. Let $n\geqslant2$, and write
$r_0=c_0-1$, $r_i=c_i$ for $0<i<k$ and $r_k=c_k-1$, all positive; the runs end
after the positions $r_0$, $r_0+r_1$, \dots, $r_0+\cdots+r_{k-1}$. Suppose that
$d$ lies in $\{ab,ba\}^*$, of length $2h$. A run ending after the position $1$
forces $r_0=1$, that is $c_0=2$. A run ending after the last odd position
$2h-1$ leaves a last run of length $1$, so $r_k=1$, that is $c_k=2$. Between
two consecutive odd positions $2i-1$ and $2i+1$, there is either no other run
end, which makes a run of length $2$, or one run end after the position $2i$,
which makes two runs of length $1$; so the runs $r_1,\ldots,r_{k-1}$, which are
the quotients $c_1,\ldots,c_{k-1}$, split into consecutive blocks $(2)$ and
$(1,1)$. Conversely, when \eqref{eq:block-form} holds, the run lengths are
$1$, then blocks $(2)$ and $(1,1)$, then $1$, so the word has even length and
a run ends after every odd position.
\end{proof}

Every occurrence supplies an integer of the form \eqref{eq:block-form}. Let
$P/Q$ be an occurrence of the label $m$, with central word $z$, and write
$\mu(z)e=(x,y)^{\mathsf T}$. Since $\mu(a)=P_1^2$, $\mu(b)=P_2^2$ and
$P_2(1,0)^{\mathsf T}=e$, the column
$P_2\,\mu(z)\,P_2(1,0)^{\mathsf T}=P_2\,\mu(z)\,e=(2x+y,x)^{\mathsf T}$ is the image
of $(1,0)^{\mathsf T}$ under the product of the matrices $P_c$, $c$ running
through the sequence $(2,\sigma(z),2)$ of \cref{sec:occurrences}, and its
first entry $2x+y=e^{\mathsf T}\mu(z)e$ is $m$ by \eqref{eq:square-count-identity}.
By \cref{lem:cost-descriptions}, the integer $u_{P/Q}=x$, which we call the
\emph{root} of the occurrence, satisfies
\begin{equation}\label{eq:root-cf}
 \frac m{u_{P/Q}}=[2;\sigma(z),2].
\end{equation}
The sequence $\sigma(z)$ consists of the blocks $(1,1)$ and $(2,2)$, so the
root is of the form \eqref{eq:block-form}, and $m/3<u_{P/Q}<m/2$. The matrix
$P_2\,\mu(z)\,P_2$ is symmetric of determinant one, with first column
$(m,u_{P/Q})^{\mathsf T}$, so that $u_{P/Q}^2\equiv-1\pmod m$. Since $\mu(a)e_1=e$,
$\mu(b)e_1=2e+e_1$ and $\mu(z)$ is symmetric, the upper left entry of
$\mu(\gamma_{P/Q})=\mu(a)\mu(z)\mu(b)$ is
$e^{\mathsf T}\mu(z)(2e+e_1)=2m+u_{P/Q}$. By the first step of the proof of
\cref{prop:occurrence-in-cube}, $u_{P/Q}\,m_R\equiv-m_L\pmod m$, and since
$0<u_{P/Q}<m/2$, the root is the characteristic number of the Markov triple
$(m,m_L,m_R)$ in the normalization of \cite[Section~2.2]{Aigner2013}, the
smaller of the two solutions in $(0,m)$ of $m_Rx\equiv\pm m_L\pmod m$.

\begin{conjecture}[Directives in the image of the Thue--Morse morphism]\label{conj:directives}
For every Markov number $m\geqslant5$, exactly one integer $n$ with
$1\leqslant n<m/2$, prime to $m$, has the property that the directive of the
central word of $\gamma_{n/m}$ lies in $\{ab,ba\}^*$; equivalently, by
\cref{lem:directives}, exactly one such $n$ satisfies \eqref{eq:block-form}.
\end{conjecture}

\begin{proposition}[The conjecture implies uniqueness]\label{prop:directives-uniqueness}
\Cref{conj:directives} implies the Frobenius--Markov uniqueness conjecture.
In terms of the Thue--Morse morphism, it states that for every Markov number
$m\geqslant5$ exactly one word $w$ beginning with the letter $a$ satisfies
$\abs{\operatorname{Pal}(\theta(w))}=m-2$.
\end{proposition}

\begin{proof}
\emph{Uniqueness.} Every occurrence $P/Q$ of a Markov number $m$ supplies the
integer $n=u_{P/Q}$, which satisfies \eqref{eq:block-form} by
\eqref{eq:root-cf}. Two distinct occurrences supply distinct integers. Indeed,
the canonical continued fraction \eqref{eq:root-cf} of $m/u_{P/Q}$ determines
the sequence $\sigma(z)$; this sequence determines $z$, because $\sigma$ replaces
each letter by two equal entries that identify it; and $z$ determines the word
$\gamma_{P/Q}=azb$ and the fraction $P/Q$, by \eqref{eq:farey-counts}. If the
integer of the conjecture is unique, so is the occurrence.

\emph{The second form.} Let $n$ be an integer as in the conjecture, and $d$ the
directive of the central word of $\gamma_{n/m}$. By \cref{lem:directives}, $d$
begins with $a$, and it lies in $\{ab,ba\}^*$ exactly when $d=\theta(w)$ for a
word $w$, which then begins with $a$ and is unique, the morphism $\theta$ being
injective. Conversely, let $w$ begin with $a$. Then $d=\theta(w)$ begins with
$a$, and $a\operatorname{Pal}(d)b$ is a proper lower Christoffel word whose numbers
of letters $a$ and $b$ are the entries of $L_d{\bf1}$, by the first step of the
proof of \cref{lem:directives}; the first entry is the larger, since
$L_d{\bf1}=L_aL_{d'}{\bf1}$ with $d=ad'$. This word is therefore
$\gamma_{n/m}$ for coprime integers with $n<m/2$ and
$m=\abs{\operatorname{Pal}(\theta(w))}+2$, and the directive of its central word is
$\theta(w)$. The integers $n$ of the conjecture for the label $m$ thus
correspond one to one to the words $w$ beginning with $a$ such that
$\abs{\operatorname{Pal}(\theta(w))}=m-2$.
\end{proof}

Reutenauer and Vuillon proved that, for every word $v$, the number
$\abs{\operatorname{Pal}(\theta(\operatorname{Pal}(av)))}+2$ is a Markov number, that every
Markov number $m\geqslant5$ arises in this way, and that the injectivity of the
map $v\mapsto\abs{\operatorname{Pal}(\theta(\operatorname{Pal}(av)))}$ is equivalent to the
uniqueness conjecture \cite[Theorem~1]{ReutenauerVuillon2017}. The words
$\operatorname{Pal}(av)$ are the central words beginning with $a$. In the second
form of \cref{prop:directives-uniqueness}, the conjecture contains this
injectivity, and asserts in addition that no word $w$ beginning with $a$ that
is not a central word gives a Markov number. The next lemma supports it for
every Markov number up to $10^{12}$.

\begin{lemma}[Finite facts on directives]\label{lem:directive-facts}
\begin{enumerate}[label=\textup{(\roman*)},leftmargin=*]
\item There are $150$ Markov numbers $m$ with $5\leqslant m\leqslant10^{12}$. Each
 is the label of exactly one occurrence $P/Q$, and its root $u_{P/Q}$ is the
 only integer $n$ prime to $m$ with $1\leqslant n<m/2$ satisfying
 \eqref{eq:block-form}.
\item The integers $31$ and $1130$, which are not Markov numbers, have exactly
 two such integers prime to them each: $31/12=[2;1,1,2,2]$ and $31/13=[2;2,1,1,2]$, and
 $1130/437=[2;\sigma(abba),2]$ and $1130/467=[2;\sigma(baab),2]$.
\item For $m=m_{9/10}=6\,625\,109$, the root is $u_{9/10}=2\,744\,210$, with
 $m/u_{9/10}=[2;2,\ldots,2]$, eighteen partial quotients equal to $2$, so that
 $\ell(u_{9/10},m)=35$. The integer $n=2\,561\,343$ is prime to $m$ and smaller
 than $m/2$, and
 $$
 \frac mn=[2;1,1,2,2,1,1,2,1,1,1,2,1,1,1,2,1,2,2,1,1,1,1,2,2],
 $$
 so that $\ell(n,m)=34$.
\end{enumerate}
\end{lemma}

\begin{proof}
These are finitely many assertions about integers. They were verified by
computer in exact arithmetic, and the verification was carried out twice
independently. The first verification generates the labels from the Markov
triples of \cref{lem:farey-triangle}, takes the roots from the residues of
$-m_L/m_R$ modulo $m$ given by the first step of the proof of
\cref{prop:occurrence-in-cube}, and searches, for each label, the sequences of
blocks from the left. If a prefix $[2;B_1,\ldots,B_i]$ has $p$ and $p'$ as the
numerators of its last two convergents, every completion has a tail column
$(X,Y)^{\mathsf T}$ with coprime entries $0<Y<X$ and $m=pX+p'Y$; hence
$Y<m/(p+p')$ and $Y$ lies in one residue class modulo $p$, so that once
$p(p+p')>m$ at most one candidate remains, and it is tested directly. The
second verification enumerates all the continued fractions
\eqref{eq:block-form} with numerator at most $10^{12}$ and compares their
numerators and denominators with the labels and roots computed from the
Cohn matrices \eqref{eq:cohn-source}. The continued fractions of the last two
items can also be checked by hand.
\end{proof}

The integers $31$ and $1130$ of \cref{lem:directive-facts} show that
\cref{conj:directives} depends on the Markov structure of $m$. Whenever the block sequence $B_1,\ldots,B_j$ of \eqref{eq:block-form}
is not a palindrome, the reversed expansion has the same numerator, since
the product of the symmetric matrices $P_c$ in reverse order is the
transpose, and a different denominator; this is the case of $31$. For
$1130$, the two expansions come from the palindromes $abba$ and $baab$, which
are not central words and have the same label $e^{\mathsf T}\mu(z)e$. The
conjecture therefore asserts that a Markov number is the numerator of no
expansion \eqref{eq:block-form} with a non-palindromic block sequence, and of
no palindromic one other than that of its root.

The roots of the occurrences are square roots of $-1$, and the square roots
of $-1$ carry the whole two-square fiber, as the next lemma shows. This gives
\cref{conj:euclidean-minimality} a form in which only the label and its square
roots of $-1$ appear. In substance the lemma is classical. Smith derived from
two identities of Euler that the numerator of a palindromic continued fraction
with an even number of partial quotients is a sum of two coprime squares, given
by its first half. He proved in this way that every prime $p\equiv1\pmod4$ is a
sum of two squares, since a parity argument shows that one of the expansions of
$p/u$ with $1<u<p/2$ is a palindrome, and he observed that the same argument
applies to every divisor of an integer $1+x^2$ \cite{Smith1855}. Brillhart read
the two squares off the remainders of the Euclidean algorithm on $p$ and the
square root $u$ of $-1$, refining a method of Hermite \cite{Brillhart1972}.
\c{C}anak\c{c}{\i} and Schiffler relate the primitive representations of an
integer to its even palindromic continued fractions \cite[Theorem~3.10 and
Corollary~3.14]{CanakciSchiffler2020}. What the lemma adds is a complete proof,
for every label, of the correspondence between the square roots of $-1$ in
$(1,m/2)$ and the fiber $\mathcal F(m)$, and the identity \eqref{eq:root-cost}
between the costs, on which the reformulation of the conjecture below rests.

\begin{lemma}[Square roots of $-1$]\label{lem:square-roots}
Let $m\geqslant5$ be an integer with a primitive representation as a sum of two
squares, and let $u$ be an integer with $1<u<m/2$ and $u^2\equiv-1\pmod m$.
Then the canonical continued fraction of $m/u$ is a palindrome
$[c_0;c_1,\ldots,c_{2h-1}]$ of even length. If $s/r=[c_{h-1};c_{h-2},\ldots,c_0]$
in lowest terms, then $m=r^2+s^2$, $0<r<s$, and
\begin{equation}\label{eq:root-cost}
 \ell(u,m)=2\,\ell(r,s)+1 .
\end{equation}
The map $u\mapsto(r,s)$ is a bijection from these integers $u$ onto
$\mathcal F(m)$, and it sends the root of an occurrence to its occurrence pair.
\end{lemma}

\begin{proof}
\emph{The palindrome.} Write $m/u=[c_0;c_1,\ldots,c_k]$ canonically. Here
$k\geqslant1$, because $u>1$ is prime to $m$; $c_0\geqslant2$, because $u<m/2$; and
$c_k\geqslant2$. By \cref{lem:cost-descriptions}, the product
$\Pi=P_{c_0}\cdots P_{c_k}$ has first column $(m,u)^{\mathsf T}$; write
$\Pi=\left(\begin{smallmatrix}m&m'\\u&u'\end{smallmatrix}\right)$. Each $P_c$ is
symmetric, so $\Pi^{\mathsf T}=P_{c_k}\cdots P_{c_0}$, whose first column is
$(m,m')^{\mathsf T}$, and \cref{lem:cost-descriptions} gives
$m/m'=[c_k;\ldots,c_0]$, a canonical expansion since $c_0\geqslant2$. As $c_k\geqslant2$
and $k\geqslant1$, $m/m'>2$, so that $0<m'<m/2$. The determinant of $\Pi$ is
$(-1)^{k+1}$, so $mu'-m'u=(-1)^{k+1}$, and $m'u\equiv(-1)^k\pmod m$. Since
$u^2\equiv-1$, the inverse of $u$ modulo $m$ is $-u$, whence
$m'\equiv(-1)^{k+1}u\pmod m$. If $k$ were even, then $m'\equiv-u$ and
$0<m'<m$ would give $m'=m-u>m/2$, a contradiction. Hence $k$ is odd,
$m'\equiv u$, and $m'=u$, both lying in $(0,m)$. The two canonical expansions
$[c_0;\ldots,c_k]$ and $[c_k;\ldots,c_0]$ of $m/u$ then coincide, so that
$c_i=c_{k-i}$ for every $i$, and $k+1=2h$ is even.

\emph{The two squares.} By the symmetry of the expansion, the second half of
the product is $P_{c_{h-1}}\cdots P_{c_0}=\Pi_0^{\mathsf T}$, where
$\Pi_0=P_{c_0}\cdots P_{c_{h-1}}$, so that $\Pi=\Pi_0\Pi_0^{\mathsf T}$. If
$(s',r')$ is the first row of $\Pi_0$, the upper left entry of $\Pi$ is
$m=s'^2+r'^2$. The first column of $\Pi_0^{\mathsf T}=P_{c_{h-1}}\cdots P_{c_0}$ is
$(s',r')^{\mathsf T}$, so \cref{lem:cost-descriptions} shows that $s'$ and $r'$ are
coprime and positive, with $s'/r'=[c_{h-1};\ldots,c_0]$, a canonical
expansion since $c_0\geqslant2$; hence $(s',r')=(s,r)$, and $r<s$ because this
expansion is not $[1]$. Finally \eqref{eq:quotient-cost} gives
$\ell(u,m)=2(c_0+\cdots+c_{h-1})-1$ and $\ell(r,s)=c_0+\cdots+c_{h-1}-1$, which
is \eqref{eq:root-cost}.

\emph{The bijection.} The pair $(r,s)$ determines $u$, because the canonical
expansion of $s/r$ is the reversed first half of the palindrome, which gives
the whole palindrome and hence $m/u$. The map is therefore injective. By \cref{prop:gaussian-fibre},
$\abs{\mathcal F(m)}=2^{t-1}$, where $t$ is the number of distinct odd primes
dividing $m$, as in \cref{prop:gaussian-fibre}. By the Chinese remainder theorem and Hensel's lemma, $-1$ has
exactly $2^{t}$ square roots modulo $m$, two modulo each odd prime power
dividing $m$ and one modulo $2$ when $m$ is even. They pair off as $u$ and
$m-u$, and none of them is $1$ or $m/2$, because $m\geqslant5$; so exactly
$2^{t-1}$ of them lie in $(1,m/2)$. An injective map between two finite sets
of the same size is a bijection. Finally, for the root of an occurrence, the
expansion is \eqref{eq:root-cf}, a palindrome because $z$ is one, and the
reversal of its first half is the right half $[\sigma(w),2]$ or
$[j,\sigma(w),2]$ of the paragraph after \cref{prop:occurrence-cost}, which is
the continued fraction of the occurrence pair.
\end{proof}

By \cref{lem:square-roots} and \cref{prop:occurrence-cost},
\cref{conj:euclidean-minimality} is equivalent to the following statement:
for every occurrence $P/Q$ of a label $m$, no square root of $-1$ modulo $m$
in $(1,m/2)$ has a shorter Euclidean algorithm with $m$ than the root, whose
cost is $\ell(u_{P/Q},m)=2(P+Q-2)+1$. \Cref{fig:directives} shows the label
$m_{4/5}=985$, whose two representations differ in cost by one. The
restriction to square roots of $-1$ cannot be dropped, even on the rows of
\cref{thm:thirteen-rows}: by the last item of \cref{lem:directive-facts}, the
integer $n=2\,561\,343$, which is not a square root of $-1$ modulo $m_{9/10}$
since its cost is even, has a shorter Euclidean algorithm with $m_{9/10}$
than the root. Among all the integers $n$ prime to $m$, the condition
\eqref{eq:block-form} isolates the root where the minimum of the cost does
not.

\begin{figure}[tb]
\centering
% Authors: Dr. Denys Dutykh (Mathematics Department, Khalifa University of Science and
% Technology, Abu Dhabi, UAE) and Prof. Laurent Vuillon (Univ. Savoie Mont Blanc, CNRS, LAMA)
% Generated by supplement/code/make_figures.py from exact data; do not edit by hand.
\begin{tikzpicture}[font=\footnotesize]
\draw[slate, line width=0.6pt, dash pattern=on 4pt off 2pt] (-0.2,0.458) -- (12,0.458);
\filldraw[navy, fill=white, line width=0.45pt] (0.024,5.182) circle (1.15pt);
\filldraw[navy, fill=white, line width=0.45pt] (0.048,4.401) circle (1.15pt);
\filldraw[navy, fill=white, line width=0.45pt] (0.072,3.948) circle (1.15pt);
\filldraw[navy, fill=white, line width=0.45pt] (0.096,3.63) circle (1.15pt);
\filldraw[navy, fill=white, line width=0.45pt] (0.144,3.193) circle (1.15pt);
\filldraw[navy, fill=white, line width=0.45pt] (0.168,3.012) circle (1.15pt);
\filldraw[navy, fill=white, line width=0.45pt] (0.192,2.896) circle (1.15pt);
\filldraw[navy, fill=white, line width=0.45pt] (0.216,2.748) circle (1.15pt);
\filldraw[navy, fill=white, line width=0.45pt] (0.264,2.542) circle (1.15pt);
\filldraw[navy, fill=white, line width=0.45pt] (0.288,2.518) circle (1.15pt);
\filldraw[navy, fill=white, line width=0.45pt] (0.312,2.362) circle (1.15pt);
\filldraw[navy, fill=white, line width=0.45pt] (0.336,2.29) circle (1.15pt);
\filldraw[navy, fill=white, line width=0.45pt] (0.384,2.148) circle (1.15pt);
\filldraw[navy, fill=white, line width=0.45pt] (0.408,2.245) circle (1.15pt);
\filldraw[navy, fill=white, line width=0.45pt] (0.432,2.023) circle (1.15pt);
\filldraw[navy, fill=white, line width=0.45pt] (0.456,2.004) circle (1.15pt);
\filldraw[navy, fill=white, line width=0.45pt] (0.504,1.965) circle (1.15pt);
\filldraw[navy, fill=white, line width=0.45pt] (0.528,1.84) circle (1.15pt);
\filldraw[navy, fill=white, line width=0.45pt] (0.552,1.817) circle (1.15pt);
\filldraw[navy, fill=white, line width=0.45pt] (0.576,2.096) circle (1.15pt);
\filldraw[navy, fill=white, line width=0.45pt] (0.624,1.747) circle (1.15pt);
\filldraw[navy, fill=white, line width=0.45pt] (0.648,1.817) circle (1.15pt);
\filldraw[navy, fill=white, line width=0.45pt] (0.672,1.647) circle (1.15pt);
\filldraw[navy, fill=white, line width=0.45pt] (0.696,2.042) circle (1.15pt);
\filldraw[navy, fill=white, line width=0.45pt] (0.744,1.537) circle (1.15pt);
\filldraw[navy, fill=white, line width=0.45pt] (0.768,1.507) circle (1.15pt);
\filldraw[navy, fill=white, line width=0.45pt] (0.792,1.507) circle (1.15pt);
\filldraw[navy, fill=white, line width=0.45pt] (0.816,2.042) circle (1.15pt);
\filldraw[navy, fill=white, line width=0.45pt] (0.864,1.415) circle (1.15pt);
\filldraw[navy, fill=white, line width=0.45pt] (0.888,1.382) circle (1.15pt);
\filldraw[navy, fill=white, line width=0.45pt] (0.912,1.537) circle (1.15pt);
\filldraw[navy, fill=white, line width=0.45pt] (0.936,1.477) circle (1.15pt);
\filldraw[navy, fill=white, line width=0.45pt] (0.984,2.096) circle (1.15pt);
\filldraw[navy, fill=white, line width=0.45pt] (1.008,1.313) circle (1.15pt);
\filldraw[navy, fill=white, line width=0.45pt] (1.032,1.415) circle (1.15pt);
\filldraw[navy, fill=white, line width=0.45pt] (1.056,1.241) circle (1.15pt);
\filldraw[navy, fill=white, line width=0.45pt] (1.104,1.202) circle (1.15pt);
\filldraw[navy, fill=white, line width=0.45pt] (1.128,1.673) circle (1.15pt);
\filldraw[navy, fill=white, line width=0.45pt] (1.152,1.382) circle (1.15pt);
\filldraw[navy, fill=white, line width=0.45pt] (1.176,1.348) circle (1.15pt);
\filldraw[navy, fill=white, line width=0.45pt] (1.224,1.202) circle (1.15pt);
\filldraw[navy, fill=white, line width=0.45pt] (1.248,1.477) circle (1.15pt);
\filldraw[navy, fill=white, line width=0.45pt] (1.272,1.122) circle (1.15pt);
\filldraw[navy, fill=white, line width=0.45pt] (1.296,1.202) circle (1.15pt);
\filldraw[navy, fill=white, line width=0.45pt] (1.344,1.079) circle (1.15pt);
\filldraw[navy, fill=white, line width=0.45pt] (1.368,1.079) circle (1.15pt);
\filldraw[navy, fill=white, line width=0.45pt] (1.392,2.245) circle (1.15pt);
\filldraw[navy, fill=white, line width=0.45pt] (1.416,1.035) circle (1.15pt);
\filldraw[navy, fill=white, line width=0.45pt] (1.464,1.122) circle (1.15pt);
\filldraw[navy, fill=white, line width=0.45pt] (1.488,1.202) circle (1.15pt);
\filldraw[navy, fill=white, line width=0.45pt] (1.512,1.035) circle (1.15pt);
\filldraw[navy, fill=white, line width=0.45pt] (1.536,0.989) circle (1.15pt);
\filldraw[navy, fill=white, line width=0.45pt] (1.584,1.278) circle (1.15pt);
\filldraw[navy, fill=white, line width=0.45pt] (1.608,1.035) circle (1.15pt);
\filldraw[navy, fill=white, line width=0.45pt] (1.632,1.348) circle (1.15pt);
\filldraw[navy, fill=white, line width=0.45pt] (1.656,0.94) circle (1.15pt);
\filldraw[navy, fill=white, line width=0.45pt] (1.704,1.163) circle (1.15pt);
\filldraw[navy, fill=white, line width=0.45pt] (1.728,1.035) circle (1.15pt);
\filldraw[navy, fill=white, line width=0.45pt] (1.752,1.817) circle (1.15pt);
\filldraw[navy, fill=white, line width=0.45pt] (1.776,0.94) circle (1.15pt);
\filldraw[navy, fill=white, line width=0.45pt] (1.824,1.537) circle (1.15pt);
\filldraw[navy, fill=white, line width=0.45pt] (1.848,0.94) circle (1.15pt);
\filldraw[navy, fill=white, line width=0.45pt] (1.872,0.89) circle (1.15pt);
\filldraw[navy, fill=white, line width=0.45pt] (1.896,0.989) circle (1.15pt);
\filldraw[navy, fill=white, line width=0.45pt] (1.944,0.989) circle (1.15pt);
\filldraw[navy, fill=white, line width=0.45pt] (1.968,2.518) circle (1.15pt);
\filldraw[navy, fill=white, line width=0.45pt] (1.992,0.989) circle (1.15pt);
\filldraw[navy, fill=white, line width=0.45pt] (2.016,0.94) circle (1.15pt);
\filldraw[navy, fill=white, line width=0.45pt] (2.064,0.989) circle (1.15pt);
\filldraw[navy, fill=white, line width=0.45pt] (2.088,1.035) circle (1.15pt);
\filldraw[navy, fill=white, line width=0.45pt] (2.112,0.94) circle (1.15pt);
\filldraw[navy, fill=white, line width=0.45pt] (2.136,1.241) circle (1.15pt);
\filldraw[navy, fill=white, line width=0.45pt] (2.184,0.89) circle (1.15pt);
\filldraw[navy, fill=white, line width=0.45pt] (2.208,0.838) circle (1.15pt);
\filldraw[navy, fill=white, line width=0.45pt] (2.232,0.838) circle (1.15pt);
\filldraw[navy, fill=white, line width=0.45pt] (2.256,1.079) circle (1.15pt);
\filldraw[navy, fill=white, line width=0.45pt] (2.304,0.89) circle (1.15pt);
\filldraw[navy, fill=white, line width=0.45pt] (2.328,0.989) circle (1.15pt);
\filldraw[navy, fill=white, line width=0.45pt] (2.352,1.313) circle (1.15pt);
\filldraw[navy, fill=white, line width=0.45pt] (2.376,1.313) circle (1.15pt);
\filldraw[navy, fill=white, line width=0.45pt] (2.424,1.477) circle (1.15pt);
\filldraw[navy, fill=white, line width=0.45pt] (2.448,0.989) circle (1.15pt);
\filldraw[navy, fill=white, line width=0.45pt] (2.472,0.838) circle (1.15pt);
\filldraw[navy, fill=white, line width=0.45pt] (2.496,0.989) circle (1.15pt);
\filldraw[navy, fill=white, line width=0.45pt] (2.544,0.725) circle (1.15pt);
\filldraw[navy, fill=white, line width=0.45pt] (2.568,0.89) circle (1.15pt);
\filldraw[navy, fill=white, line width=0.45pt] (2.592,0.94) circle (1.15pt);
\filldraw[navy, fill=white, line width=0.45pt] (2.616,1.537) circle (1.15pt);
\filldraw[navy, fill=white, line width=0.45pt] (2.664,1.163) circle (1.15pt);
\filldraw[navy, fill=white, line width=0.45pt] (2.688,0.838) circle (1.15pt);
\filldraw[navy, fill=white, line width=0.45pt] (2.712,0.838) circle (1.15pt);
\filldraw[navy, fill=white, line width=0.45pt] (2.736,0.725) circle (1.15pt);
\filldraw[navy, fill=white, line width=0.45pt] (2.784,1.537) circle (1.15pt);
\filldraw[navy, fill=white, line width=0.45pt] (2.808,0.664) circle (1.15pt);
\filldraw[navy, fill=white, line width=0.45pt] (2.832,0.89) circle (1.15pt);
\filldraw[navy, fill=white, line width=0.45pt] (2.856,0.783) circle (1.15pt);
\filldraw[navy, fill=white, line width=0.45pt] (2.904,0.989) circle (1.15pt);
\filldraw[navy, fill=white, line width=0.45pt] (2.928,1.079) circle (1.15pt);
\filldraw[navy, fill=white, line width=0.45pt] (2.952,2.896) circle (1.15pt);
\filldraw[navy, fill=white, line width=0.45pt] (2.976,1.163) circle (1.15pt);
\filldraw[navy, fill=white, line width=0.45pt] (3.024,0.989) circle (1.15pt);
\filldraw[navy, fill=white, line width=0.45pt] (3.048,0.94) circle (1.15pt);
\filldraw[navy, fill=white, line width=0.45pt] (3.072,0.725) circle (1.15pt);
\filldraw[navy, fill=white, line width=0.45pt] (3.096,0.94) circle (1.15pt);
\filldraw[navy, fill=white, line width=0.45pt] (3.144,0.989) circle (1.15pt);
\filldraw[navy, fill=white, line width=0.45pt] (3.168,0.989) circle (1.15pt);
\filldraw[navy, fill=white, line width=0.45pt] (3.192,0.783) circle (1.15pt);
\filldraw[navy, fill=white, line width=0.45pt] (3.216,0.725) circle (1.15pt);
\filldraw[navy, fill=white, line width=0.45pt] (3.264,0.89) circle (1.15pt);
\filldraw[navy, fill=white, line width=0.45pt] (3.288,0.725) circle (1.15pt);
\filldraw[navy, fill=white, line width=0.45pt] (3.312,0.838) circle (1.15pt);
\filldraw[navy, fill=white, line width=0.45pt] (3.336,0.94) circle (1.15pt);
\filldraw[navy, fill=white, line width=0.45pt] (3.384,2.305) circle (1.15pt);
\filldraw[navy, fill=white, line width=0.45pt] (3.408,1.079) circle (1.15pt);
\filldraw[navy, fill=white, line width=0.45pt] (3.432,1.202) circle (1.15pt);
\filldraw[navy, fill=white, line width=0.45pt] (3.456,0.783) circle (1.15pt);
\filldraw[navy, fill=white, line width=0.45pt] (3.504,1.163) circle (1.15pt);
\filldraw[navy, fill=white, line width=0.45pt] (3.528,1.035) circle (1.15pt);
\filldraw[navy, fill=white, line width=0.45pt] (3.552,0.89) circle (1.15pt);
\filldraw[navy, fill=white, line width=0.45pt] (3.576,0.783) circle (1.15pt);
\filldraw[navy, fill=white, line width=0.45pt] (3.624,0.89) circle (1.15pt);
\filldraw[navy, fill=white, line width=0.45pt] (3.648,1.035) circle (1.15pt);
\filldraw[navy, fill=white, line width=0.45pt] (3.672,0.838) circle (1.15pt);
\filldraw[navy, fill=white, line width=0.45pt] (3.696,0.838) circle (1.15pt);
\filldraw[navy, fill=white, line width=0.45pt] (3.744,0.725) circle (1.15pt);
\filldraw[navy, fill=white, line width=0.45pt] (3.768,0.725) circle (1.15pt);
\filldraw[navy, fill=white, line width=0.45pt] (3.792,0.664) circle (1.15pt);
\filldraw[navy, fill=white, line width=0.45pt] (3.816,0.838) circle (1.15pt);
\filldraw[navy, fill=white, line width=0.45pt] (3.864,0.989) circle (1.15pt);
\filldraw[navy, fill=white, line width=0.45pt] (3.888,1.035) circle (1.15pt);
\filldraw[navy, fill=white, line width=0.45pt] (3.912,1.348) circle (1.15pt);
\filldraw[navy, fill=white, line width=0.45pt] (3.936,3.193) circle (1.15pt);
\filldraw[navy, fill=white, line width=0.45pt] (3.984,1.241) circle (1.15pt);
\filldraw[navy, fill=white, line width=0.45pt] (4.008,0.838) circle (1.15pt);
\filldraw[navy, fill=white, line width=0.45pt] (4.032,0.725) circle (1.15pt);
\filldraw[navy, fill=white, line width=0.45pt] (4.056,0.725) circle (1.15pt);
\filldraw[navy, fill=white, line width=0.45pt] (4.104,0.783) circle (1.15pt);
\filldraw[navy, fill=white, line width=0.45pt] (4.128,1.035) circle (1.15pt);
\filldraw[navy, fill=white, line width=0.45pt] (4.152,0.599) circle (1.15pt);
\filldraw[navy, fill=white, line width=0.45pt] (4.176,1.163) circle (1.15pt);
\filldraw[navy, fill=white, line width=0.45pt] (4.224,0.89) circle (1.15pt);
\filldraw[navy, fill=white, line width=0.45pt] (4.248,0.725) circle (1.15pt);
\filldraw[navy, fill=white, line width=0.45pt] (4.272,0.989) circle (1.15pt);
\filldraw[navy, fill=white, line width=0.45pt] (4.296,2.542) circle (1.15pt);
\filldraw[navy, fill=white, line width=0.45pt] (4.344,0.664) circle (1.15pt);
\filldraw[navy, fill=white, line width=0.45pt] (4.368,0.838) circle (1.15pt);
\filldraw[navy, fill=white, line width=0.45pt] (4.416,0.89) circle (1.15pt);
\filldraw[navy, fill=white, line width=0.45pt] (4.464,0.531) circle (1.15pt);
\filldraw[navy, fill=white, line width=0.45pt] (4.488,0.664) circle (1.15pt);
\filldraw[navy, fill=white, line width=0.45pt] (4.512,0.664) circle (1.15pt);
\filldraw[navy, fill=white, line width=0.45pt] (4.536,0.599) circle (1.15pt);
\filldraw[navy, fill=white, line width=0.45pt] (4.584,0.664) circle (1.15pt);
\filldraw[navy, fill=white, line width=0.45pt] (4.608,0.89) circle (1.15pt);
\filldraw[navy, fill=white, line width=0.45pt] (4.632,0.89) circle (1.15pt);
\filldraw[navy, fill=white, line width=0.45pt] (4.656,1.278) circle (1.15pt);
\filldraw[navy, fill=white, line width=0.45pt] (4.704,1.771) circle (1.15pt);
\filldraw[navy, fill=white, line width=0.45pt] (4.752,1.723) circle (1.15pt);
\filldraw[navy, fill=white, line width=0.45pt] (4.776,1.313) circle (1.15pt);
\filldraw[navy, fill=white, line width=0.45pt] (4.824,1.348) circle (1.15pt);
\filldraw[navy, fill=white, line width=0.45pt] (4.848,1.035) circle (1.15pt);
\filldraw[navy, fill=white, line width=0.45pt] (4.872,0.664) circle (1.15pt);
\filldraw[navy, fill=white, line width=0.45pt] (4.896,0.725) circle (1.15pt);
\filldraw[navy, fill=white, line width=0.45pt] (4.944,0.531) circle (1.15pt);
\filldraw[navy, fill=white, line width=0.45pt] (4.968,0.838) circle (1.15pt);
\filldraw[navy, fill=white, line width=0.45pt] (4.992,0.531) circle (1.15pt);
\filldraw[navy, fill=white, line width=0.45pt] (5.016,0.989) circle (1.15pt);
\filldraw[navy, fill=white, line width=0.45pt] (5.064,2.29) circle (1.15pt);
\filldraw[navy, fill=white, line width=0.45pt] (5.088,0.599) circle (1.15pt);
\filldraw[navy, fill=white, line width=0.45pt] (5.112,1.382) circle (1.15pt);
\filldraw[navy, fill=white, line width=0.45pt] (5.136,0.94) circle (1.15pt);
\filldraw[navy, fill=white, line width=0.45pt] (5.184,0.725) circle (1.15pt);
\filldraw[navy, fill=white, line width=0.45pt] (5.208,0.783) circle (1.15pt);
\filldraw[navy, fill=white, line width=0.45pt] (5.232,1.079) circle (1.15pt);
\filldraw[navy, fill=white, line width=0.45pt] (5.256,2.748) circle (1.15pt);
\filldraw[navy, fill=white, line width=0.45pt] (5.304,0.725) circle (1.15pt);
\filldraw[navy, fill=white, line width=0.45pt] (5.328,0.664) circle (1.15pt);
\filldraw[navy, fill=white, line width=0.45pt] (5.352,1.122) circle (1.15pt);
\filldraw[navy, fill=white, line width=0.45pt] (5.376,0.989) circle (1.15pt);
\filldraw[navy, fill=white, line width=0.45pt] (5.424,0.599) circle (1.15pt);
\filldraw[navy, fill=white, line width=0.45pt] (5.448,1.163) circle (1.15pt);
\filldraw[navy, fill=white, line width=0.45pt] (5.472,0.94) circle (1.15pt);
\filldraw[navy, fill=white, line width=0.45pt] (5.496,0.725) circle (1.15pt);
\filldraw[navy, fill=white, line width=0.45pt] (5.544,0.599) circle (1.15pt);
\filldraw[navy, fill=white, line width=0.45pt] (5.568,1.035) circle (1.15pt);
\filldraw[navy, fill=white, line width=0.45pt] (5.592,0.838) circle (1.15pt);
\filldraw[navy, fill=white, line width=0.45pt] (5.616,0.664) circle (1.15pt);
\filldraw[navy, fill=white, line width=0.45pt] (5.664,0.89) circle (1.15pt);
\filldraw[navy, fill=white, line width=0.45pt] (5.688,0.783) circle (1.15pt);
\filldraw[navy, fill=white, line width=0.45pt] (5.712,0.725) circle (1.15pt);
\filldraw[navy, fill=white, line width=0.45pt] (5.736,0.89) circle (1.15pt);
\filldraw[navy, fill=white, line width=0.45pt] (5.784,1.079) circle (1.15pt);
\filldraw[navy, fill=white, line width=0.45pt] (5.808,1.035) circle (1.15pt);
\filldraw[navy, fill=white, line width=0.45pt] (5.832,1.163) circle (1.15pt);
\filldraw[navy, fill=white, line width=0.45pt] (5.856,1.537) circle (1.15pt);
\filldraw[navy, fill=white, line width=0.45pt] (5.904,3.63) circle (1.15pt);
\filldraw[navy, fill=white, line width=0.45pt] (5.928,2.443) circle (1.15pt);
\filldraw[navy, fill=white, line width=0.45pt] (5.952,1.62) circle (1.15pt);
\filldraw[navy, fill=white, line width=0.45pt] (5.976,1.241) circle (1.15pt);
\filldraw[navy, fill=white, line width=0.45pt] (6.024,0.94) circle (1.15pt);
\filldraw[navy, fill=white, line width=0.45pt] (6.048,1.415) circle (1.15pt);
\filldraw[navy, fill=white, line width=0.45pt] (6.072,0.725) circle (1.15pt);
\filldraw[navy, fill=white, line width=0.45pt] (6.096,0.838) circle (1.15pt);
\filldraw[navy, fill=white, line width=0.45pt] (6.144,0.599) circle (1.15pt);
\filldraw[navy, fill=white, line width=0.45pt] (6.168,1.817) circle (1.15pt);
\filldraw[navy, fill=white, line width=0.45pt] (6.192,1.313) circle (1.15pt);
\filldraw[navy, fill=white, line width=0.45pt] (6.216,0.94) circle (1.15pt);
\filldraw[navy, fill=white, line width=0.45pt] (6.264,0.531) circle (1.15pt);
\filldraw[navy, fill=white, line width=0.45pt] (6.288,0.664) circle (1.15pt);
\filldraw[navy, fill=white, line width=0.45pt] (6.312,0.89) circle (1.15pt);
\filldraw[navy, fill=white, line width=0.45pt] (6.336,0.531) circle (1.15pt);
\filldraw[navy, fill=white, line width=0.45pt] (6.384,0.458) circle (1.15pt);
\filldraw[navy, fill=white, line width=0.45pt] (6.408,0.531) circle (1.15pt);
\filldraw[navy, fill=white, line width=0.45pt] (6.432,0.89) circle (1.15pt);
\filldraw[navy, fill=white, line width=0.45pt] (6.456,1.241) circle (1.15pt);
\filldraw[navy, fill=white, line width=0.45pt] (6.504,0.599) circle (1.15pt);
\filldraw[navy, fill=white, line width=0.45pt] (6.528,0.458) circle (1.15pt);
\filldraw[navy, fill=white, line width=0.45pt] (6.552,0.531) circle (1.15pt);
\filldraw[navy, fill=white, line width=0.45pt] (6.576,0.664) circle (1.15pt);
\filldraw[navy, fill=white, line width=0.45pt] (6.624,0.664) circle (1.15pt);
\filldraw[navy, fill=white, line width=0.45pt] (6.648,1.507) circle (1.15pt);
\filldraw[navy, fill=white, line width=0.45pt] (6.672,0.599) circle (1.15pt);
\filldraw[navy, fill=white, line width=0.45pt] (6.696,0.664) circle (1.15pt);
\filldraw[navy, fill=white, line width=0.45pt] (6.744,1.883) circle (1.15pt);
\filldraw[navy, fill=white, line width=0.45pt] (6.768,1.62) circle (1.15pt);
\filldraw[navy, fill=white, line width=0.45pt] (6.792,0.89) circle (1.15pt);
\filldraw[navy, fill=white, line width=0.45pt] (6.816,0.664) circle (1.15pt);
\filldraw[navy, fill=white, line width=0.45pt] (6.864,1.537) circle (1.15pt);
\filldraw[navy, fill=white, line width=0.45pt] (6.888,0.599) circle (1.15pt);
\filldraw[navy, fill=white, line width=0.45pt] (6.912,0.458) circle (1.15pt);
\filldraw[navy, fill=white, line width=0.45pt] (6.936,0.531) circle (1.15pt);
\filldraw[navy, fill=white, line width=0.45pt] (6.984,1.241) circle (1.15pt);
\filldraw[navy, fill=white, line width=0.45pt] (7.008,0.725) circle (1.15pt);
\filldraw[navy, fill=white, line width=0.45pt] (7.032,0.664) circle (1.15pt);
\filldraw[navy, fill=white, line width=0.45pt] (7.056,0.725) circle (1.15pt);
\filldraw[navy, fill=white, line width=0.45pt] (7.104,1.163) circle (1.15pt);
\filldraw[navy, fill=white, line width=0.45pt] (7.128,0.664) circle (1.15pt);
\filldraw[navy, fill=white, line width=0.45pt] (7.152,0.531) circle (1.15pt);
\filldraw[navy, fill=white, line width=0.45pt] (7.176,1.079) circle (1.15pt);
\filldraw[navy, fill=white, line width=0.45pt] (7.224,1.415) circle (1.15pt);
\filldraw[navy, fill=white, line width=0.45pt] (7.248,0.599) circle (1.15pt);
\filldraw[navy, fill=white, line width=0.45pt] (7.272,2.362) circle (1.15pt);
\filldraw[navy, fill=white, line width=0.45pt] (7.296,0.989) circle (1.15pt);
\filldraw[navy, fill=white, line width=0.45pt] (7.344,0.599) circle (1.15pt);
\filldraw[navy, fill=white, line width=0.45pt] (7.368,0.94) circle (1.15pt);
\filldraw[navy, fill=white, line width=0.45pt] (7.392,1.241) circle (1.15pt);
\filldraw[navy, fill=white, line width=0.45pt] (7.416,1.202) circle (1.15pt);
\filldraw[navy, fill=white, line width=0.45pt] (7.464,2.004) circle (1.15pt);
\filldraw[navy, fill=white, line width=0.45pt] (7.488,0.599) circle (1.15pt);
\filldraw[navy, fill=white, line width=0.45pt] (7.512,0.89) circle (1.15pt);
\filldraw[navy, fill=white, line width=0.45pt] (7.536,0.725) circle (1.15pt);
\filldraw[navy, fill=white, line width=0.45pt] (7.584,0.783) circle (1.15pt);
\filldraw[navy, fill=white, line width=0.45pt] (7.608,1.035) circle (1.15pt);
\filldraw[navy, fill=white, line width=0.45pt] (7.632,0.94) circle (1.15pt);
\filldraw[navy, fill=white, line width=0.45pt] (7.656,0.89) circle (1.15pt);
\filldraw[navy, fill=white, line width=0.45pt] (7.704,0.989) circle (1.15pt);
\filldraw[navy, fill=white, line width=0.45pt] (7.728,1.477) circle (1.15pt);
\filldraw[navy, fill=white, line width=0.45pt] (7.752,1.241) circle (1.15pt);
\filldraw[navy, fill=white, line width=0.45pt] (7.776,1.537) circle (1.15pt);
\filldraw[navy, fill=white, line width=0.45pt] (7.824,1.883) circle (1.15pt);
\filldraw[navy, fill=white, line width=0.45pt] (7.848,2.443) circle (1.15pt);
\filldraw[navy, fill=white, line width=0.45pt] (7.872,3.948) circle (1.15pt);
\filldraw[navy, fill=white, line width=0.45pt] (7.896,3.179) circle (1.15pt);
\filldraw[navy, fill=white, line width=0.45pt] (7.944,1.747) circle (1.15pt);
\filldraw[navy, fill=white, line width=0.45pt] (7.968,1.507) circle (1.15pt);
\filldraw[navy, fill=white, line width=0.45pt] (7.992,1.278) circle (1.15pt);
\filldraw[navy, fill=white, line width=0.45pt] (8.016,1.163) circle (1.15pt);
\filldraw[navy, fill=white, line width=0.45pt] (8.064,0.94) circle (1.15pt);
\filldraw[navy, fill=white, line width=0.45pt] (8.088,1.537) circle (1.15pt);
\filldraw[navy, fill=white, line width=0.45pt] (8.112,0.989) circle (1.15pt);
\filldraw[navy, fill=white, line width=0.45pt] (8.136,0.838) circle (1.15pt);
\filldraw[navy, fill=white, line width=0.45pt] (8.184,1.747) circle (1.15pt);
\filldraw[navy, fill=white, line width=0.45pt] (8.208,1.035) circle (1.15pt);
\filldraw[navy, fill=white, line width=0.45pt] (8.232,0.838) circle (1.15pt);
\filldraw[navy, fill=white, line width=0.45pt] (8.256,0.783) circle (1.15pt);
\filldraw[navy, fill=white, line width=0.45pt] (8.304,0.783) circle (1.15pt);
\filldraw[navy, fill=white, line width=0.45pt] (8.328,0.94) circle (1.15pt);
\filldraw[navy, fill=white, line width=0.45pt] (8.352,0.838) circle (1.15pt);
\filldraw[navy, fill=white, line width=0.45pt] (8.376,0.531) circle (1.15pt);
\filldraw[navy, fill=white, line width=0.45pt] (8.424,0.725) circle (1.15pt);
\filldraw[navy, fill=white, line width=0.45pt] (8.448,1.278) circle (1.15pt);
\filldraw[navy, fill=white, line width=0.45pt] (8.472,0.599) circle (1.15pt);
\filldraw[navy, fill=white, line width=0.45pt] (8.496,0.989) circle (1.15pt);
\filldraw[navy, fill=white, line width=0.45pt] (8.544,0.599) circle (1.15pt);
\filldraw[navy, fill=white, line width=0.45pt] (8.568,0.599) circle (1.15pt);
\filldraw[navy, fill=white, line width=0.45pt] (8.592,1.84) circle (1.15pt);
\filldraw[navy, fill=white, line width=0.45pt] (8.616,0.783) circle (1.15pt);
\filldraw[navy, fill=white, line width=0.45pt] (8.664,0.664) circle (1.15pt);
\filldraw[navy, fill=white, line width=0.45pt] (8.688,0.664) circle (1.15pt);
\filldraw[navy, fill=white, line width=0.45pt] (8.712,1.348) circle (1.15pt);
\filldraw[navy, fill=white, line width=0.45pt] (8.736,1.202) circle (1.15pt);
\filldraw[navy, fill=white, line width=0.45pt] (8.784,0.599) circle (1.15pt);
\filldraw[navy, fill=white, line width=0.45pt] (8.808,0.838) circle (1.15pt);
\filldraw[navy, fill=white, line width=0.45pt] (8.832,0.838) circle (1.15pt);
\filldraw[navy, fill=white, line width=0.45pt] (8.856,1.747) circle (1.15pt);
\filldraw[navy, fill=white, line width=0.45pt] (8.904,0.783) circle (1.15pt);
\filldraw[navy, fill=white, line width=0.45pt] (8.928,0.664) circle (1.15pt);
\filldraw[navy, fill=white, line width=0.45pt] (8.952,0.458) circle (1.15pt);
\filldraw[navy, fill=white, line width=0.45pt] (8.976,0.458) circle (1.15pt);
\filldraw[navy, fill=white, line width=0.45pt] (9.024,0.531) circle (1.15pt);
\filldraw[navy, fill=white, line width=0.45pt] (9.048,0.531) circle (1.15pt);
\filldraw[navy, fill=white, line width=0.45pt] (9.072,0.599) circle (1.15pt);
\filldraw[navy, fill=white, line width=0.45pt] (9.096,1.673) circle (1.15pt);
\filldraw[navy, fill=white, line width=0.45pt] (9.144,0.725) circle (1.15pt);
\filldraw[navy, fill=white, line width=0.45pt] (9.168,0.599) circle (1.15pt);
\filldraw[navy, fill=white, line width=0.45pt] (9.192,2.023) circle (1.15pt);
\filldraw[navy, fill=white, line width=0.45pt] (9.216,1.035) circle (1.15pt);
\filldraw[navy, fill=white, line width=0.45pt] (9.264,0.94) circle (1.15pt);
\filldraw[navy, fill=white, line width=0.45pt] (9.288,1.647) circle (1.15pt);
\filldraw[navy, fill=white, line width=0.45pt] (9.312,1.507) circle (1.15pt);
\filldraw[navy, fill=white, line width=0.45pt] (9.336,0.664) circle (1.15pt);
\filldraw[navy, fill=white, line width=0.45pt] (9.384,0.989) circle (1.15pt);
\filldraw[navy, fill=white, line width=0.45pt] (9.408,1.313) circle (1.15pt);
\filldraw[navy, fill=white, line width=0.45pt] (9.432,1.723) circle (1.15pt);
\filldraw[navy, fill=white, line width=0.45pt] (9.504,1.163) circle (1.15pt);
\filldraw[navy, fill=white, line width=0.45pt] (9.528,0.89) circle (1.15pt);
\filldraw[navy, fill=white, line width=0.45pt] (9.552,0.783) circle (1.15pt);
\filldraw[navy, fill=white, line width=0.45pt] (9.576,0.989) circle (1.15pt);
\filldraw[navy, fill=white, line width=0.45pt] (9.624,0.599) circle (1.15pt);
\filldraw[navy, fill=white, line width=0.45pt] (9.648,0.783) circle (1.15pt);
\filldraw[navy, fill=white, line width=0.45pt] (9.672,1.84) circle (1.15pt);
\filldraw[navy, fill=white, line width=0.45pt] (9.696,0.599) circle (1.15pt);
\filldraw[navy, fill=white, line width=0.45pt] (9.744,0.664) circle (1.15pt);
\filldraw[navy, fill=white, line width=0.45pt] (9.768,0.458) circle (1.15pt);
\filldraw[navy, fill=white, line width=0.45pt] (9.816,0.531) circle (1.15pt);
\filldraw[navy, fill=white, line width=0.45pt] (9.864,0.94) circle (1.15pt);
\filldraw[navy, fill=white, line width=0.45pt] (9.888,0.531) circle (1.15pt);
\filldraw[navy, fill=white, line width=0.45pt] (9.912,0.989) circle (1.15pt);
\filldraw[navy, fill=white, line width=0.45pt] (9.936,0.94) circle (1.15pt);
\filldraw[navy, fill=white, line width=0.45pt] (9.984,0.531) circle (1.15pt);
\filldraw[navy, fill=white, line width=0.45pt] (10.008,0.725) circle (1.15pt);
\filldraw[navy, fill=white, line width=0.45pt] (10.032,0.725) circle (1.15pt);
\filldraw[navy, fill=white, line width=0.45pt] (10.056,0.725) circle (1.15pt);
\filldraw[navy, fill=white, line width=0.45pt] (10.104,1.163) circle (1.15pt);
\filldraw[navy, fill=white, line width=0.45pt] (10.128,3.012) circle (1.15pt);
\filldraw[navy, fill=white, line width=0.45pt] (10.152,1.348) circle (1.15pt);
\filldraw[navy, fill=white, line width=0.45pt] (10.176,0.89) circle (1.15pt);
\filldraw[navy, fill=white, line width=0.45pt] (10.224,0.94) circle (1.15pt);
\filldraw[navy, fill=white, line width=0.45pt] (10.248,0.664) circle (1.15pt);
\filldraw[navy, fill=white, line width=0.45pt] (10.272,0.725) circle (1.15pt);
\filldraw[navy, fill=white, line width=0.45pt] (10.296,0.531) circle (1.15pt);
\filldraw[navy, fill=white, line width=0.45pt] (10.344,2.148) circle (1.15pt);
\filldraw[navy, fill=white, line width=0.45pt] (10.368,1.079) circle (1.15pt);
\filldraw[navy, fill=white, line width=0.45pt] (10.392,0.531) circle (1.15pt);
\filldraw[navy, fill=white, line width=0.45pt] (10.416,0.531) circle (1.15pt);
\filldraw[navy, fill=white, line width=0.45pt] (10.464,1.122) circle (1.15pt);
\filldraw[navy, fill=white, line width=0.45pt] (10.488,1.079) circle (1.15pt);
\filldraw[navy, fill=white, line width=0.45pt] (10.512,2.042) circle (1.15pt);
\filldraw[navy, fill=white, line width=0.45pt] (10.536,0.783) circle (1.15pt);
\filldraw[navy, fill=white, line width=0.45pt] (10.584,0.599) circle (1.15pt);
\filldraw[navy, fill=white, line width=0.45pt] (10.608,0.89) circle (1.15pt);
\filldraw[navy, fill=white, line width=0.45pt] (10.632,0.783) circle (1.15pt);
\filldraw[navy, fill=white, line width=0.45pt] (10.656,0.531) circle (1.15pt);
\filldraw[navy, fill=white, line width=0.45pt] (10.704,0.783) circle (1.15pt);
\filldraw[navy, fill=white, line width=0.45pt] (10.728,0.838) circle (1.15pt);
\filldraw[navy, fill=white, line width=0.45pt] (10.752,1.507) circle (1.15pt);
\filldraw[navy, fill=white, line width=0.45pt] (10.776,0.725) circle (1.15pt);
\filldraw[navy, fill=white, line width=0.45pt] (10.824,0.599) circle (1.15pt);
\filldraw[navy, fill=white, line width=0.45pt] (10.848,0.531) circle (1.15pt);
\filldraw[navy, fill=white, line width=0.45pt] (10.872,0.664) circle (1.15pt);
\filldraw[navy, fill=white, line width=0.45pt] (10.896,0.783) circle (1.15pt);
\filldraw[navy, fill=white, line width=0.45pt] (10.944,1.202) circle (1.15pt);
\filldraw[navy, fill=white, line width=0.45pt] (10.968,0.664) circle (1.15pt);
\filldraw[navy, fill=white, line width=0.45pt] (10.992,0.725) circle (1.15pt);
\filldraw[navy, fill=white, line width=0.45pt] (11.016,0.725) circle (1.15pt);
\filldraw[navy, fill=white, line width=0.45pt] (11.064,0.783) circle (1.15pt);
\filldraw[navy, fill=white, line width=0.45pt] (11.088,0.838) circle (1.15pt);
\filldraw[navy, fill=white, line width=0.45pt] (11.112,0.783) circle (1.15pt);
\filldraw[navy, fill=white, line width=0.45pt] (11.136,0.989) circle (1.15pt);
\filldraw[navy, fill=white, line width=0.45pt] (11.184,0.783) circle (1.15pt);
\filldraw[navy, fill=white, line width=0.45pt] (11.208,0.838) circle (1.15pt);
\filldraw[navy, fill=white, line width=0.45pt] (11.232,0.783) circle (1.15pt);
\filldraw[navy, fill=white, line width=0.45pt] (11.256,1.965) circle (1.15pt);
\filldraw[navy, fill=white, line width=0.45pt] (11.304,1.382) circle (1.15pt);
\filldraw[navy, fill=white, line width=0.45pt] (11.328,1.382) circle (1.15pt);
\filldraw[navy, fill=white, line width=0.45pt] (11.352,1.035) circle (1.15pt);
\filldraw[navy, fill=white, line width=0.45pt] (11.376,0.94) circle (1.15pt);
\filldraw[navy, fill=white, line width=0.45pt] (11.424,0.989) circle (1.15pt);
\filldraw[navy, fill=white, line width=0.45pt] (11.448,0.989) circle (1.15pt);
\filldraw[navy, fill=white, line width=0.45pt] (11.472,1.348) circle (1.15pt);
\filldraw[navy, fill=white, line width=0.45pt] (11.496,1.163) circle (1.15pt);
\filldraw[navy, fill=white, line width=0.45pt] (11.544,1.382) circle (1.15pt);
\filldraw[navy, fill=white, line width=0.45pt] (11.568,1.673) circle (1.15pt);
\filldraw[navy, fill=white, line width=0.45pt] (11.592,1.348) circle (1.15pt);
\filldraw[navy, fill=white, line width=0.45pt] (11.616,1.537) circle (1.15pt);
\filldraw[navy, fill=white, line width=0.45pt] (11.664,1.673) circle (1.15pt);
\filldraw[navy, fill=white, line width=0.45pt] (11.688,1.84) circle (1.15pt);
\filldraw[navy, fill=white, line width=0.45pt] (11.712,2.042) circle (1.15pt);
\filldraw[navy, fill=white, line width=0.45pt] (11.736,2.305) circle (1.15pt);
\filldraw[navy, fill=white, line width=0.45pt] (11.784,3.179) circle (1.15pt);
\filldraw[navy, fill=white, line width=0.45pt] (11.808,4.401) circle (1.15pt);
\fill[ochre] (4.392,0.599) circle (1.9pt);
\node[ochre, below] at (4.392,0.368) {$183$};
\fill[wine] (9.792,0.458) circle (1.9pt);
\node[wine, below] at (9.792,0.368) {$408$};
\draw[wine, line width=0.6pt, fill=none] (9.792,0.458) circle (3.2pt);
\draw[ink, line width=0.5pt] (-0.2,0) -- (12,0);
\draw[ink, line width=0.5pt] (-0.2,0) -- (-0.2,5.2);
\draw[ink, line width=0.5pt] (0,0) -- (0,-0.08);
\node[ink, below] at (0,-0.08) {$0$};
\draw[ink, line width=0.5pt] (2.4,0) -- (2.4,-0.08);
\node[ink, below] at (2.4,-0.08) {$100$};
\draw[ink, line width=0.5pt] (4.8,0) -- (4.8,-0.08);
\node[ink, below] at (4.8,-0.08) {$200$};
\draw[ink, line width=0.5pt] (7.2,0) -- (7.2,-0.08);
\node[ink, below] at (7.2,-0.08) {$300$};
\draw[ink, line width=0.5pt] (9.6,0) -- (9.6,-0.08);
\node[ink, below] at (9.6,-0.08) {$400$};
\draw[ink, line width=0.5pt] (12,0) -- (12,-0.08);
\node[ink, below] at (12,-0.08) {$500$};
\draw[ink, line width=0.5pt] (-0.2,0) -- (-0.28,0);
\node[ink, left] at (-0.28,0) {$10$};
\draw[ink, line width=0.5pt] (-0.2,0.458) -- (-0.28,0.458);
\node[ink, left] at (-0.28,0.458) {$15$};
\draw[ink, line width=0.5pt] (-0.2,0.783) -- (-0.28,0.783);
\node[ink, left] at (-0.28,0.783) {$20$};
\draw[ink, line width=0.5pt] (-0.2,1.817) -- (-0.28,1.817);
\node[ink, left] at (-0.28,1.817) {$50$};
\draw[ink, line width=0.5pt] (-0.2,2.6) -- (-0.28,2.6);
\node[ink, left] at (-0.28,2.6) {$100$};
\draw[ink, line width=0.5pt] (-0.2,3.383) -- (-0.28,3.383);
\node[ink, left] at (-0.28,3.383) {$200$};
\draw[ink, line width=0.5pt] (-0.2,4.417) -- (-0.28,4.417);
\node[ink, left] at (-0.28,4.417) {$500$};
\draw[ink, line width=0.5pt] (-0.2,5.2) -- (-0.28,5.2);
\node[ink, left] at (-0.28,5.2) {$1000$};
\node[ink] at (6,-0.75) {$n$};
\node[ink, rotate=90] at (-1.2,2.6) {$\ell(n,985)$};
\end{tikzpicture}

\caption{The Euclidean cost $\ell(n,m)$ for $m=m_{4/5}=985$ and the $392$
integers $n$ prime to $m$ with $1\leqslant n<m/2$ (hollow circles), on a
logarithmic scale; the directive of the central word of $\gamma_{n/m}$ has
$\ell(n,m)-1$ letters. The two square roots of $-1$ modulo $m$ in $(1,m/2)$ are
filled: the root $408$ of the occurrence $4/5$, of cost $15=2\cdot7+1$, and
$183$, of cost $17=2\cdot8+1$, which correspond by \cref{lem:square-roots} to the
representations $12^2+29^2$ and $16^2+27^2$. Six other integers, $266$,
$272$, $288$, $373$, $374$ and $407$, share the least cost $15$ with the root
(dashed line), but only the root, ringed, satisfies \eqref{eq:block-form}.}
\label{fig:directives}
\end{figure}

\Cref{conj:directives} is a rule of the kind discussed in
\cref{rem:pair-determines}: it designates the root from the label alone, and
with it the occurrence pair, which \cref{lem:square-roots} reads off the
first half of the continued fraction of $m/u_{P/Q}$. Like
\cref{conj:euclidean-minimality}, it can hold only because of the Markov
structure of the label, as the integers $31$ and $1130$ of
\cref{lem:directive-facts} show, and it
reaches the uniqueness conjecture by a different road: through the
palindromic closure and the Thue--Morse morphism rather than through the
comparison of costs.

\section*{Supplementary material}
The supplementary material accompanying this article contains the exact
computations behind \cref{prop:finite-verifications,lem:ladder-facts}: the
factorizations and complete Gaussian fibers of the finite windows, the
covers of the unbounded tails, the source comparisons, the matrices of the
two ladders, two independent verifications of each, and instructions for
repeating them. It also contains the complete factorizations of the labels
behind the margins of \cref{sec:outlook}, verified a second time
independently, and the two independent computations behind
\cref{lem:directive-facts}.

\section*{Code and data availability}
{\clubpenalty=0 \widowpenalty=0
The source code and the data supporting this article, with the instructions
that repeat every calculation, are contained in the supplementary material of
this article, which is available, together with that of the other articles of
this series, at
\par}
\begin{center}
  \url{https://github.com/dutykh/fibonacci-pell-nearest-gaps/}
\end{center}

\section*{Funding}
This research did not receive any specific grant from funding agencies in the public, commercial,
or not-for-profit sectors.

\section*{Declaration of generative AI and AI-assisted technologies in the manuscript preparation process}
Generative AI (Claude Code by Anthropic) was used exclusively for grammatical polishing, LaTeX typesetting,
and generating TikZ figure templates. No AI tools were used in the formulation, derivation, or verification
of the mathematical results.

\section*{Conflict of interest}
The authors declare that they have no conflict of interest.

\bibliographystyle{amsplain}
\bibliography{references}
\end{document}
